% 本文件由 LaTeX 排版 Agent 根据有效 translation.json 自动生成。
% author=Clement of Rome; work_id=0804; title=克莱孟的认亲记

\chapter{克莱孟的认亲记}

% structure: p0001 source=h1 target=omitted reason=page-title-already-rendered-by-parent

\Emph{序言}\newline{}\Emph{第一卷}\newline{}\Emph{第二卷}\newline{}\Emph{第三卷}\newline{}\Emph{第四卷}\newline{}\Emph{第五卷}\newline{}\Emph{第六卷}\newline{}\Emph{第七卷}\newline{}\Emph{第八卷}\newline{}\Emph{第九卷}\newline{}\Emph{第十卷}\par

\SourceRecord{Translated by Thomas Smith.；Ante-Nicene Fathers；Vol. 8.；Edited by Alexander Roberts, James Donaldson, and A. Cleveland Coxe.；Buffalo, NY: Christian Literature Publishing Co.,；1886.；<http://www.newadvent.org/fathers/0804.htm>.}

% structure: p0001 source=h1 target=section reason=page-title

\section{认亲记（序言）}

\Emph{由\Person{阿奎莱亚的鲁菲努斯}撰写，致\Person{加乌登提乌斯}主教}\par

的确，\Person{加乌登提乌斯}啊，你是我们导师中的精英佼佼，拥有如此敏锐的心智，是的，如此丰沛的圣神恩宠，以致凡你在日常讲道中所说的，凡你在教会中所宣讲的，都理当载入书籍，传于后世，以启发后人。然而我们这些人，因才智贫弱而显得不够敏捷，如今又因年迈而变得迟缓、懈怠，虽然几经耽搁，最终还是把这份作品呈献给你——这份作品是当年可敬的贞女\Person{西尔维亚}托付我们，要求我们将\Person{克莱门特}的作品译成我们的语言，而后又由你按继承权向我们索要的。因此，我们算是从希腊人的书库里夺来了不小的战利品，以此滋养我们的人民，用外来的养分喂养那些我们无法用自己的产物养育的人。因为外来之物通常似乎更为悦人，有时也更为有益。简言之，几乎一切能医治我们身体、对抗疾病、中和毒物之物都是外来的。因为\Place{犹太地}给我们送来\Emph{泪香脂}，\Place{克里特}送来\Emph{白鲜枝}，\Place{阿拉伯}献上她的香料之花，\Place{印度}收获她的甘松香。它们虽在抵达我们时比离开原产地时略为破碎，但仍完整保留其芳香的甜美和治疗的功效。因此，我的灵魂啊，请迎接\Person{克莱门特}回到你身边；请迎接他，如今他身披罗马的衣装。倘若他往昔那华丽的辞藻或许显得不如往常，请不必惊讶。只要内容与过去相同，便无关紧要。因此，我们费尽辛苦将外来的货物运回本国。我不知道我的同胞们以多么欢喜的表情欢迎我，因我为他们带来了希腊丰厚的战利品，用我们语言的钥匙打开了隐藏的智慧宝库。愿\Person{天主}应允你的祈祷，使不祥的目光或任何嫉妒的面容不要临到我们，以免出现极端的怪异之事：拿走他的那些人并不嫉妒，而得到他的那些人反而抱怨。确实，有必要向你说明我们的翻译计划，因你已读过这些作品的希腊文本，以免你在某些部分可能认为翻译的顺序没有保持。我想你知晓，\Person{克莱门特}的这部作品有两个希腊文版本——即\OriginalTerm{认亲记}{\GreekText{᾿Αναγνώσεις}}，也就是\WorkTitle{认亲记}；并且有两个书卷集，在某些点上相异，但在许多方面\Emph{包含}同样的叙述。简言之，这部作品的最后部分，即关于\Person{西门}变形的记述，包含在一个集子中，但在另一个集子中完全没有。此外，在两个集子中都有一些关于\OriginalTerm{不生之天主与受生之天主}{Unbegotten God and the Begotten}以及其他题目的论述，这些超过我们的理解（不再赘述）。因此，由于超出我们的能力，我选择将其留给他人，而非以不完善的状态呈现。但在其余部分，我们已尽力做到既不离弃原意，甚至也不背离语言和表达方式；这虽然使叙述风格不那么华丽，但却更为忠实。\Person{克莱门特}写给主的兄弟\Person{雅各}的书信，在其中告知他\Person{伯多禄}的去世，以及\Person{伯多禄}已留下他作为自己座椅和教导的继任者，并且该书信也处理了教会秩序的全部主题，我没有将其置于这部作品之前，既因为其年代较晚，也因为我已翻译并出版了它。但我不认为在此解释那封信中可能令一些人感到矛盾之处是不合适的。因为有人会问：既然\Person{理诺}和\Person{克雷特}在\Person{克莱门特}之前已是\Place{罗马}城的主教，那么\Person{克莱门特}本人写信给\Person{雅各}时，怎能说教导的座椅是由\Person{伯多禄}交给他的呢？关于此事，我们听到这样的解释：\Person{理诺}和\Person{克雷特}确实在\Person{克莱门特}之前是\Place{罗马}城的主教，但这是在\Person{伯多禄}生前：也就是说，他们承担主教职的职责，而\Person{伯多禄}则完成宗徒职务；正如在\Place{凯撒勒雅}也曾发现的情况，当\Person{伯多禄}亲自在场时，他仍立了由他祝圣的\Person{匝凯}为主教。按此方式，两种说法都将被视为真实的：这些主教在\Person{克莱门特}之前被计算，而\Person{克莱门特}在\Person{伯多禄}去世时接受了导师的座椅。现在，让我们看看\Person{克莱门特}写给主的兄弟\Person{雅各}的信如何开始叙述。\par

\SourceRecord{Translated by Thomas Smith.；Ante-Nicene Fathers；Vol. 8.；Edited by Alexander Roberts, James Donaldson, and A. Cleveland Coxe.；Buffalo, NY: Christian Literature Publishing Co.,；1886.；<http://www.newadvent.org/fathers/080400.htm>.}

% structure: p0001 source=h1 target=section reason=page-title

\section{《认主集》（卷一）}

% structure: p0002 source=h2 target=subsection reason=relative-heading-rank; base=h2

\subsection{第一章 \Person{克莱孟}的早期经历；疑惑}

我\Person{克莱孟}，生于\Place{罗马}城，自幼便热爱贞洁；而我的心志如同被焦虑与悲伤的锁链牢牢束缚。因为我心中有一种想法——不知从何而起——不断引导我思考自己必死的境况，并探讨这类问题：死后是否还有生命，还是我将彻底消亡？我是否在出生之前不存在，死后是否不再有今生的记忆，以至于无垠的时间将所有事物交付遗忘与沉默？不仅我们不再存在，也没有人记得我们曾经存在过。我也在心中思量：世界何时被创造，或创造之前有什么，或它是否从永恒就已存在？因为似乎确定，如果世界是被创造的，它必然注定要消亡；如果它消亡了，之后会是什么？——除非，或许一切都会沉入遗忘与沉默，或者会有某种人心如今无法想象的事物。\par

% structure: p0004 source=h2 target=subsection reason=relative-heading-rank; base=h2

\subsection{第二章 他的痛苦}

我不断在心中反复思考这些以及诸如此类的问题——不知从何而来——因过度悲伤而异常憔悴；更糟的是，若我有时想抛开这些并无多大用处的忧虑，焦虑的波涛反而更高地涌起。因为我内在有那位最优秀的伴侣，不肯让我安歇——即对不朽的渴望；因为正如后来的结果所显示，并全能的\Person{天主}的\OriginalTerm{恩宠}{grace}所引导，这种心志引领我寻求真理，承认真光；因此后来我很快就怜悯那些我在无知时曾以为有福的人。\par

% structure: p0006 source=h2 target=subsection reason=relative-heading-rank; base=h2

\subsection{第三章 他对哲学家学派的失望}

因此，我自幼便有如此心志，求知的渴望使我常去哲学家们的学府。在那里我看到，除了无休止地主张和辩驳学说、进行争论、探讨三段论的艺术和结论的微妙之外，别无他事。若有时灵魂不朽的学说占上风，我便感激；若有时它被反驳，我便悲伤地离去。然而，两种学说都未能在我的心中取得真理的力量。我只明白一件事：事物的意见和定义之所以被视为真或假，并非按照其本性和论据的真实性，而是按照支持者才能的大小。而在我内心深处，我更加痛苦，因为我既不能抓住那些被坚定宣称的言论中的任何一项，也无法放弃探究的渴望；反而我越是努力忽视和轻视它们，这种渴望就越是——如我所说——秘密地带着某种愉悦潜入我内心，占据我的心智。\par

% structure: p0008 source=h2 target=subsection reason=relative-heading-rank; base=h2

\subsection{第四章 他日益增加的不安}

因此，在探究事物中感到困窘，我对自己说：既然事物结局是显然的，我们为何徒然劳苦？因为若死后我不再存在，我目前的折磨便是无用的；但若死后我将有生命，那就让我们为那生命保存它应有的激励，免得某些比我现在所受的更悲惨的事临到我——除非我敬虔而清醒地生活，并且按照某些哲学家的意见，我被投入昏暗流淌的\Place{弗勒革同}之河，或\Place{塔尔塔罗斯}，如同\Person{西西弗斯}和\Person{提堤俄斯}，并受到地狱中永恒的惩罚，如同\Person{伊克西翁}和\Person{坦塔罗斯}。而我又会回答自己：但这些是神话；或者若果真如此，既然事情不确定，最好还是敬虔地生活。但我又暗自思量：当我对正义的赏报不确定时，我该如何克制自己犯罪的私欲？——尤其当我无法确定什么是正义，什么是\Person{天主}所喜悦的；也无法确知灵魂是否不朽，是否具有可盼望的事物；更不知道未来确实会怎样。然而，我仍无法停止这类思虑。\par

% structure: p0010 source=h2 target=subsection reason=relative-heading-rank; base=h2

\subsection{第五章 他设计检验灵魂不朽}

那么，我该怎么办？我要这样做：我要前往\Place{埃及}，在那里结交那些管理神庙的\OriginalTerm{圣师}{hierophants}或先知。然后我用金钱收买一位魔法师，恳求他通过所谓\OriginalTerm{招魂术}{necromantic art}，从冥界召来一个灵魂，仿佛我想就某事咨询它。但这将是我的咨询：灵魂是否不朽。现在，灵魂不朽的证据将无可置疑地确定，不是通过它所说的，或我所听到的，而是通过我所看见的：因为亲眼看见它，我以后将持有最确定的信念，相信它的不朽；言辞的谬误或听说的不确定性将永远无法动摇由视觉所产生的信念。然而，我把这个计划告诉了一位与我亲密的哲学家，他劝我不要冒险；他说：\Quote{因为，若灵魂不听从魔法师的召唤，你以后将更绝望地生活，认为死后一无所有，并且还尝试了不法之事。但若你似乎看见了什么，从非法和不敬的事中能产生什么宗教或虔敬呢？因为人们说这类行为是\Person{天主}所憎恶的，\Person{天主}反对那些在灵魂脱离身体后搅扰它们的人。}我听到这话，确实在意图上动摇了；然而我无论如何既不能放下我的渴望，也不能摆脱那痛苦的思虑。\par

% structure: p0012 source=h2 target=subsection reason=relative-heading-rank; base=h2

\subsection{第六章 听闻基督}

简而言之，当我在这思潮的波涛中颠簸时，有一则传言，发源于\Person{提庇留·凯撒}统治时期的东方地区，逐渐传到我们这里；它每经过一处，就愈发有力，如同来自天主的好消息，充满整个世界，不愿让天主的旨意被沉默掩盖。这传言传遍各处，宣告在\Place{犹太}有一个人，从春季开始，向犹太人宣讲天主的国，并说那些遵守祂诫命和教导的人将获得它。为使祂的话值得相信，充满神性，据说祂仅凭言语就施行了许多大能事迹、奇事和奇迹；如此，作为有权柄来自天主者，祂使聋子听见，瞎子看见，瘸子站立，驱逐一切疾病和一切魔鬼；甚至使带到祂面前的死人复活；远远地一看就治愈麻风病人；根本没有什么是祂不能的。这些和类似的事，随着时间的推移，不仅通过频繁的传闻，还通过从那些地区来的人明确陈述得到证实；这事日日更显真切。\par

% structure: p0014 source=h2 target=subsection reason=relative-heading-rank; base=h2

\subsection{第七章 \Person{巴尔纳伯}抵达\Place{罗马}}

最后，城中各处开始举行集会，此事成为谈话讨论的主题，令人惊奇的是，这位出现者究竟是谁，祂从天主那里给人类带来了什么消息；直到大约同年，有一个人站在城中一个最拥挤的地方，向百姓宣告说：\Quote{\Place{罗马}的公民们，请听我说。天主之子如今在\Place{犹太}地区，向凡愿意听从祂的人应许永生，但条件是：他要按差遣祂的那位，即天主圣父的旨意，来规范自己的行为。所以，你们要从邪恶转向良善，从暂时转向永恒。要承认有一位天主，天地的统治者，在祂公义的眼中，你们这些不义的人居住在祂的世界里。但如果你们悔改，按祂的旨意行事，那么，在来世到来时，成为不朽的，你们将享受祂不可言喻的福乐和赏报。}现在，向百姓说这些话的人来自东方地区，按民族是希伯来人，名叫\Person{巴尔纳伯}，他说自己就是祂的门徒之一，并且为此目的而被差遣——向那些愿意听的人宣告这些事。我听到这些事，便开始与众人一起跟随他，听他说话。我确实看出这人没有任何辩证技巧，只是简朴地、毫无诡辩地讲述他从天主之子那里听见或看见的事。因为他不用论证的力量来证实自己的主张，而是从周围站着的百姓中提出许多见证人，来证明他所叙述的言词和奇事。\par

% structure: p0016 source=h2 target=subsection reason=relative-heading-rank; base=h2

\subsection{第八章 他的宣讲}

既然百姓开始欣然赞同这些真诚的话，接受他单纯的言论，那些自以为有学问或哲学的人就开始嘲笑他，讥讽他，向他抛出三段论的铁钩，如同强壮的臂膀。但他毫不畏惧，把他们的诡辩视为狂言，甚至不认为值得回答，而是勇敢地继续他先前设定的话题。最后，当他说话时，有人向他提出这个问题：为什么蚊子如此被造——虽是小生物，有六只脚，却又添了翅膀；而大象，虽然是巨兽，没有翅膀，却只有四只脚？他不理会这个问题，继续他那被不合时宜的挑战打断的讲话，只是在每次打断时加上这个告诫：\Quote{我们受命向你们宣告差遣我们的那位的言语和奇事，并证实我们所说之事的真理，不是靠巧辩的论证，而是靠你们中间产生的见证人。因为我认出你们中间站着许多人，我记得他们曾同我们一起听过我们所听见的，见过我们所看见的。但你们可以选择接受或拒绝我们带给你们的消息。因为我们不能隐瞒我们知道对你们有益的事，因为如果我们沉默，我们有祸了；但对你们，如果你们不接受我们所说的话，就是毁灭。我本可以轻易回答你们愚昧的挑战，如果你们是为学习真理而问——我是指关于蚊子和大象的区别；但现在，当万物的创造者和造主你们都不认识时，谈论这些受造物是荒谬的。}\par

% structure: p0018 source=h2 target=subsection reason=relative-heading-rank; base=h2

\subsection{第九章 \Person{克莱孟}为\Person{巴尔纳伯}说情}

他说完这话后，众人仿佛同心合意，粗野地高声嘲笑，要使他蒙羞并住口，喊着说他是蛮族和疯子。我看见事情如此进展，不知何处涌起一股热忱，被宗教热诚点燃，我无法保持沉默，便大胆喊道：\Quote{全能的天主最公义地隐藏他的旨意，不让你知道，因为他预见到你们不配认识他，这在真正有智慧的人看来，从你们现在的行为就显明了。因为当你们看到天主圣意的宣讲者来到你们中间，他们的言论并不显示文法知识的技巧，而是以简单朴素的语言向你们陈明天主诫命，使所有听者都能追随并理解所说的话时，你们竟嘲笑你们得救的仆人和使者，不知道正是你们这些自以为有技巧和口才的人，反倒要受责备：粗野蛮荒的人拥有真理的知识；而真理来到你们这里，却连客人的待遇都得不到——倘若你们的放纵和情欲不加阻挡，它本应成为你们本地的公民和土著。这样，你们被证明不是真理和哲学的朋友，而是夸夸其谈者的追随者。你们以为真理不在于简单的话语，而在于巧妙和精微的言论，于是制造出无数话语，却不及一句话的价值。那么，你们这些希腊人群，如果有如他所说的天主的审判，你们以为自己会怎样？但现在，不要为了自取灭亡而嘲笑这个人；请你们中谁愿意就来回答我。因为你们的狂吠确实打扰了那些渴望得救的人的耳朵，你们吵闹的声音使那些预备接受信德的心灵转向不信的堕落。当真理的使者向你们提供天主的知识时，你们嘲笑并施以暴力，还能有什么宽恕呢？即使他并没有带给你们真理，但仅凭他对你们的善意，你们也该怀着感激和欢迎之心接受。}\par

% structure: p0020 source=h2 target=subsection reason=relative-heading-rank; base=h2

\subsection{第十章 与巴尔纳伯的往来}

当我力陈这些和类似论点时，旁观者中掀起了巨大骚动，一些人出于对外来人的怜悯而感动，并赞同我的发言，认为合乎这种情感；另一些人则固执轻率，不仅对巴尔纳伯，也对我，激起了他们未受约束心灵中的怒气。但天色渐晚，我握住巴尔纳伯的右手，尽管他不愿意，还是把他领到我的住所；我让他留下，恐怕粗鲁人群中有谁对他动手。我们这样相处了几天，我愉快地听他宣讲真理之言；但他急于启程，说他必须到犹太去庆祝他宗教即将到来的节庆日，并说他今后要与同胞和弟兄们留在那里——显然表明他对自己所受的伤害感到惊骇。\par

% structure: p0022 source=h2 target=subsection reason=relative-heading-rank; base=h2

\subsection{第十一章 巴尔纳伯的离开}

最后我对他说：\Quote{你只需向我讲解你所说那位显现者的道理，我将用我的语言编排你的言论，并宣讲全能天主的国度与正义；之后，你若愿意，我甚至可以与你一同航行，因为我极渴望看到犹太，也许我会永远与你在一起。}他回答说：\Quote{如果你真想看看我们的家乡，并学习你所渴望的事，现在就与我一同起航；或者，如果现在有什么事留你，我会留给你我住所的指引，以便你何时想来，都能轻易找到我；因为明天我就要上路了。}我见他心意已决，就跟他到港口，仔细收好他给我的指引，以找到他的住所。我告诉他，若非必须收取一笔欠款，我绝不会延迟，但我会很快跟随他。说完这话，我把他托付给船上的管事，悲伤地返回；因为我心中充满了与一位优秀客人和挚友交往的记忆。\par

% structure: p0024 source=h2 target=subsection reason=relative-heading-rank; base=h2

\subsection{第十二章 克莱孟抵达凯撒勒雅并引见给伯多禄}

于是停留了几天，大致处理完收债事务（因为我为了目的而急于赶路，忽略了许多事），我便直接起航往犹太去，十五天后在凯撒勒雅·斯特拉托尼登陆，那是巴勒斯坦最大的城市。登陆后，我寻找客栈，从人们的谈话中得知，一位名叫伯多禄的人，是那位在犹太显现、并在人间行了诸多神圣奇迹的至可信赖的门徒，将于次日与一位撒玛黎雅人西满举行一场言语与问题的辩论。听后，我请求指示他的住处；找到后，我站在门前，向守门人说明我是谁，从何而来。看，巴尔纳伯出来，一看见我就跑入我怀中，欢喜流泪，抓住我的手领我到伯多禄那里。他远远指给我说：\Quote{这就是我向你提到的伯多禄，他在天主的智慧中最为伟大，我也时常向他提到你。进去吧，你已是他的熟识。因为他深知你的一切好处，并仔细了解了你宗教上的意向，因此他极其渴望见你。今天我将你作为一件大礼物献给他。}同时，他介绍说：\Quote{啊，伯多禄，这位就是克莱孟。}\par

% structure: p0026 source=h2 target=subsection reason=relative-heading-rank; base=h2

\subsection{第十三章 伯多禄对他的热诚接待}

但\Person{伯多禄}最和蔼地，听到我的名字后，立刻跑向我并亲吻了我。然后，使我坐下后，他说：\Quote{你接待了真理的宣讲者\Person{巴尔纳伯}为客人，做得好，毫不惧怕那些疯狂之人的愤怒。你将有福。因为正如你使真理的使者堪受一切尊荣，同样，真理本身将接待你这漂泊者和陌生人，并将你登记为她本城的公民；那时你将有大喜乐，因为施予一个小小的恩惠，你将被称为永恒福乐的继承者。因此，现在不要麻烦向我解释你的心意；因为\Person{巴尔纳伯}已以忠实的言谈告知我关于你和你的意向的一切，几乎每日不辍，追忆你的美德。并且，简略地向你指出，如同对一位已与我们同心合意的朋友，什么是对你最好的道路；若没有什么妨碍你，就与我们同行，聆听真理之言，我们将要在各处宣讲，直到我们抵达罗马城；现在，你若有所愿，请说。}\par

% structure: p0028 source=h2 target=subsection reason=relative-heading-rank; base=h2

\subsection{第十四章 他的自述}

我向他详细说明了从起初我所怀有的意图，以及我如何被虚妄的探究所困扰，以及起初我曾向您——我的主\Person{雅各伯}——所提及的一切，为免再重复同样的事。我欣然同意与他同行；\Quote{因为那，}我说，\Quote{正是我最渴望的。但我首先希望将真理的体系向我解释，使我知道灵魂是有死的还是不朽的；若是不朽的，它是否要因它在此所行的事而受审判。此外，我渴望知道什么是天主所喜悦的义；然后，世界是否受造，为何受造，是否将要解散，是否将被更新变得更好，或者此后根本不再有世界；并且，不必列举一切，我希望被告知关于这些以及类似的事情是如何的。}对此，\Person{伯多禄}回答说：\Quote{我将简略地把这些事的知识传授给你，啊，\Person{克莱孟}：因此请听。}\par

% structure: p0030 source=h2 target=subsection reason=relative-heading-rank; base=h2

\subsection{第十五章 伯多禄的首次教导：无知的原因}

\Quote{天主的旨意和计划为许多理由向人隐藏了；首先，确实，通过坏的教导、邪恶的交往、不良的习惯、无益的交谈、以及不义的猜度。由于所有这些，我说，首先错误，然后轻蔑，然后不信与恶意，贪婪也，以及虚妄的夸耀，以及其他类似的邪恶，充满了这整个世界的房屋，如同某种浓烟，并阻止那些住在其中的人正确看见其创造者，并觉察什么事是祂所喜悦的。那么，对于里面的人，除了从他们内心深处发出呼喊来恳求祂的帮助，是合适的吗？因为只有祂没有被关在那充满烟尘的房屋里，愿祂前来打开房屋的门，使里面的烟被驱散，使外面照耀的太阳光得以进入。}\par

% structure: p0032 source=h2 target=subsection reason=relative-heading-rank; base=h2

\subsection{第十六章 教导继续：真先知}

\Quote{因此，那位其援助是为充满无知黑暗与恶习烟雾的房屋所需要的，我们所说，就是被称为真先知的那一位，祂唯独能够光照人的灵魂，使他们用眼睛清楚地看见得救的道路。因为否则，除非从那位真先知学习，就不可能获得关于神圣和永恒事物的知识；因为，正如你刚才自己所说，对事物的信念和对原因的意见，是根据其倡导者的才能而被评估的：因此，同一件原因有时被认为公正，有时不公正；那现在看来是真的，很快因另一人的断言而成为假的。为此，宗教和虔敬的信用要求真先知的临在，以便祂自己可以告诉我们关于每个细节，真理是如何的，并教导我们关于每一件事我们应当如何相信。因此，在一切之前，必须先谨慎地检查先知本人的凭据；一旦你确认他是一位先知，你就应当从此相信他的一切，不再讨论他所教导的细节，而要把他所说的话当作确定和神圣的；这些事情，虽然看似由信德接受，却是在先前建立的检验基础上被相信的。因为一旦在开端通过检验确立了先知之真，其余部分就应当基于已经确立他是真理导师的信德来聆听和持守。并且，正如所有属于神圣知识的事物都应当依照真理的准则来持守是确定的，同样，除了祂自己，从别处不可能知道什么是真实的，这也是无疑的。}\par

% structure: p0034 source=h2 target=subsection reason=relative-heading-rank; base=h2

\subsection{第十七章 伯多禄请求他作自己的随从}

他这样说了之后，便极其开诚布公地向我说明那位先知是谁，以及如何可以找到他，以致我仿佛亲眼看见、亲手触摸了他关于先知的真理所提出的证据。我深感惊异：为什么没有人看见那些众人都追求的事物，尽管它们就在眼前。于是，遵从他的命令，我将他对我说的话加以整理，编成了一本关于真先知的书，并遵从他的命令从\Place{凯撒勒雅}寄给了你。因为他说他从你那里得到命令，要每年向你报告他的言行。与此同时，在他第一天讲论的开端，当他非常详细地向我讲解了关于真先知以及许多其他事情后，他又加上了这句话：\Quote{看，}他说，\Quote{从今以后，你要出席讨论会，每逢有需要时，我将与那些反对者进行辩论。与他们辩论时，即使我似乎处于劣势，我也不怕你会对我所告诉你的那些事产生怀疑；因为即使我似乎被打败了，但真先知所传授给我们的那些事也不因此显得不确定。然而，我也希望我们在辩论中也不会被打败，只要我们的听众是理性的、热爱真理的，能辨别言语的力量和意义，并能识别出哪种言论出自诡辩之术，不包含真理，而只是真理的影像；哪种言论是质朴而不加诡诈地说出，其全部力量不在于炫耀和修饰，而在于真理和理性。}\par

% structure: p0036 source=h2 target=subsection reason=relative-heading-rank; base=h2

\subsection{第十八章　他因\Person{伯多禄}的教导获益}

\Quote{我感谢全能的天主，因为我已如我所愿、所渴望的那样受到了教导。无论如何，你可以信赖我到如此地步：我永远不可能对你所教给我的那些事产生怀疑；所以，即使你本人有时想使我的信德离开真先知，你也不可能做到，因为我已全心全意地吸收了你的话。而且，为了使你不认为我在说大话——当我说我无法被从这信德移开时——我确信，凡接受了这真先知论述的人，此后绝不可能再怀疑它的真理。因此，我深信这来自天上的教导，在其中一切邪恶的伎俩都被克服。因为与\Emph{这个}预言相对抗，没有任何技巧能够站立，也没有诡辩和推论的精巧能够胜过；每一个听到真先知的人，必会立刻渴望真理本身，并且此后不会在寻求真理的借口下忍受各种谬误。所以，我的主上\Person{伯多禄}，不要再为我担心，好像我是一个不知道自己所领受的是什么、不知道赐予他的是何等恩惠的人。请确信，您已将恩惠赐给了一个知道并理解其价值的人；因此，我不会轻易受骗，因为我似乎已经迅速得到了长久渴望的东西；因为一个人可能迅速得到他所渴望的，而另一个人即使缓慢地也不能得到他所渴望的东西。}\par

% structure: p0038 source=h2 target=subsection reason=relative-heading-rank; base=h2

\subsection{第十九章　\Person{伯多禄}的满意}

然后\Person{伯多禄}听到我这样说，就说：\Quote{我感谢我的天主，既为你的救恩，也为我的平安；因我非常高兴地看到你已明白先知的德能何其伟大，并且因为你所说，即使我自己愿意（但愿天主不许！）也不能使你转向另一种信仰。现在，从今以后，你要与我们在一起，明天要出席我们的讨论，因为我将与行邪术的\Person{西满}进行辩论。}他这样说了之后，便与朋友们一同去进餐；但他命令我独自进食；饭后，他赞美了天主并谢了恩，便向我汇报了这件事，又补充说：\Quote{愿主恩赐你，使你在一切事上得以相似我们，使你领受圣洗后，能与我们同席。}说了这话，他便命令我去休息，因为此时疲倦和时辰都催人入眠。\par

% structure: p0040 source=h2 target=subsection reason=relative-heading-rank; base=h2

\subsection{第二十章　推迟与行邪术的\Person{西满}的讨论}

第二天清早，匝凯进来见我们，问候后对伯多禄说：\Quote{西满将辩论推迟到本月十一日，也就是七天后，因为他说那时他会有更多闲暇来辩论。但在我看来，他的推迟对我们也有利，好让更多人聚集，他们可以作我们辩论的听众或评判者。然而，如果你觉得合适，让我们利用这段间隔，在我们中间讨论我们认为可能进入争论的问题；这样我们每个人知道要提出什么，要给出什么答案，就可以自己考虑是否一切都正确，或者对手是否能找到任何反对的理由，或推翻我们对他提出的东西。但如果我们要说的话显然在各方面无懈可击，我们就会有信心进入审查。确实，这是我的意见：首先应当探究万物的起源，或者什么是可以直接被称为万有之因的东西；然后，关于所有存在的事物，它们是否被造，被谁，通过谁，为了谁；它们是从一个、两个还是多个得到实存；它们是否是从没有先前存在的材料中取来并塑造的，还是从某些材料中；然后，在最高的事物中是否有任何德能，或在较低的事物中；是否有任何东西比一切都好，或任何东西比一切都低；是否有任何运动，还是没有；那些看得见的事物是否一直存在且将永远存在；它们是否在没有创造者的情况下出现，并在没有毁灭者的情况下消失。我说，如果讨论从这些开始，我认为那些将被探究的事物经过勤奋的审查讨论，将很容易确定。当这些确定后，后续的知识将容易找到。我已陈述了我的意见；请告知你对此事的看法。}\par

% structure: p0042 source=h2 target=subsection reason=relative-heading-rank; base=h2

\subsection{第二十一章 延迟的好处。}

对此，伯多禄回答说：\Quote{告诉西满，他在这期间可以随意行事，并放心，只要天主眷顾，他将永远发现我们已准备好。}然后匝凯出去向西满转达了所听到的话。但伯多禄看着我们，发现我因辩论推迟而感伤，就说：\Quote{相信世界是由至高天主的眷顾所治理的人，啊，我的朋友克莱孟，不应因任何具体事件的发生方式而见怪，要确信天主的公义将那些在事务中看似多余或相反的事，导向有利且合适的结果，尤其是对那些更亲密敬拜祂的人；因此，确信这些事的人，如我所说，如果有什么与他预期相反的事发生，他知道如何因此驱散心中的忧伤，在他更好的判断中视为毫无疑问，即由于善良天主的治理，甚至看似相反的事也可以转为善。所以，啊，克莱孟，现在不要让术士西满的延迟使你感伤：因为我深信这是出于天主的眷顾，为你的益处；使我能在这七天的间隔中，毫无干扰地向你阐述我们信德的方式和连续的顺序，按照真先知传统，祂独自知道过去如何，现在如何，将来如何；这些事确实由祂明确说出，但并未明确写出；以至于当它们被阅读时，没有注释者就无法理解，因为罪与人一同成长，如我之前所说。因此，我将向你解释一切，以便在那些写下的东西中，你能清楚地看出立法者的心意。}\par

% structure: p0044 source=h2 target=subsection reason=relative-heading-rank; base=h2

\subsection{第二十二章 重复教导。}

他说完这话后，便开始逐点向我讲解那些似乎有争议的法律章节，从创世之初一直到我来到凯撒勒雅遇见他的时候，告诉我西满的延迟有助于我按顺序学习一切。\Quote{在其他时候，}他说道，\Quote{我们将根据我们谈话的场合，更全面地讨论我们现在简短说过的各个要点；这样，按照我的承诺，你可以获得对一切圆满和完美的认识。既然因着这次延迟，我们今天有空闲，我想再次向你重复已经说过的话，以便你能更好地记住。}然后他这样开始唤起我的记忆：\Quote{啊，朋友克莱孟，你还记得我向你讲述的那个没有终结的永恒世代吗？}于是我说：\Quote{啊，伯多禄，如果我能失去或忘记那事，我就什么也记不住了。}\par

% structure: p0046 source=h2 target=subsection reason=relative-heading-rank; base=h2

\subsection{第二十三章 重复继续。}

伯多禄听了我的回答，甚为喜悦，说：\Quote{我祝贺你，因为你这样回答，不因你能轻易谈论这些事，而是因你承认你记得它们；因为最崇高的真理藉沉默得到最佳的尊崇。然而，为了你记得那些关于不可言说之事的正确性，请告诉我你记下的、我们其次谈论的那些可以轻易言说之事，好让我看出你记忆的牢固，从而更乐意地向你指明，并自由地展开我愿谈论的事。}然后我，当我觉察到他因听众的良好记忆而喜乐时，说：\Quote{我不仅记得你的定义，也记得那置于定义前的序言；你阐释的几乎所有内容，我都完整保留了含义，尽管并非全部言词；因为你所说的，已成为我灵魂的、天生的本性。因为你在我极度干渴时，递给了我极其甘甜的杯。而且，免得你以为我用言语拖延，却忘却事物，我现在便要忆起所说的话，你的讨论次序极大地帮助了我；因为你所说之事依次递进、均衡排列，使之藉其次序的线索易于回忆。因为言论的次序对记忆有益：当你开始逐点依次跟进时，若有欠缺，理智立即寻求；找到后便保留，或者，若无法发现，也无需犹豫向老师询问。但为不耽搁成全你对我的要求，我将简述你关于真理定义所传授给我的。}\par

% structure: p0048 source=h2 target=subsection reason=relative-heading-rank; base=h2

\subsection{第二十四章 重复（续）}

始终有、现在有、且永远会有那使首生自永恒的\OriginalTerm{第一次意志}{first Will}得以持存者；从第一次意志发出\OriginalTerm{第二次意志}{second Will}。之后来了世界；从世界来了时间；从时间来了众多的人；从众多的人来了蒙拣选的爱者，从他们的一心一德建立起和平的天主之国。至于其余应随之而来的，你应允另时告诉我。此后，当你解释了受造的世界，你暗示了天主的谕令，\Quote{祂凭自己的善意，在全体最初的众天使面前所宣布的，}并作为永恒的法律订立给一切；祂如何建立了两个王国——即今世与来世——并为二者指定了时期，且规定要期待审判之日，祂确定此日，届时事物与灵魂将被分开：恶者必因其罪被投入永火；但那按造物主天主的旨意生活者，因善工获得祝福，以最明亮的光辉闪耀，被引入永恒居所，常存于不朽，领受不可言喻福乐之永恒恩赐。\par

% structure: p0050 source=h2 target=subsection reason=relative-heading-rank; base=h2

\subsection{第二十五章 重复（续）}

我正这样继续时，伯多禄满心喜乐，为我挂虑，仿佛我是他的儿子，惟恐我忘记其余而当着在场的人蒙羞，说：\Quote{够了，克莱孟；因为你讲得比我本人阐释的还要清楚。}于是我说：\Quote{博雅学问赋予我有序叙述的能力，并清晰陈述需要之事的能力。若我们运用学问来主张古时的谬误，便会因言语的优雅与流畅而自毁；但若将学问与言语的优美用于主张真理，我以为这不无裨益。尽管如此，我的老师伯多禄，你只能想象我是何等感恩地沉浸于你的一切教训，尤其是你所传授的教理陈述：只有一位天主，世界是祂的工程，因为祂全然公义，必按各人的行为报应各人。此后你补充道：为论证这一教条，将有无数言词被提出；但在那些蒙赐\OriginalTerm{真先知}{true Prophet}知识的人，这言语的丛林全被砍伐。因此，既然你授予了我关于真先知的论述，你便用你论断的一切自信坚固了我。}然后，当我领悟到所有宗教与虔敬的总纲即在此，我立即回答：\Quote{你进行得极为出色，伯多禄；因此，从今以后，毫不踌躇地向那已知道信德与虔敬基础的人，阐释真先知的传承；因唯有他，正如已清楚证明的，当被相信。但那需要论断与论证的阐述，请保留给不信者，你尚未认为适于将先知恩宠的无疑信德托付给他们。}我说了这话，又补充道：\Quote{你应允在适当时候给我两样：首先是这样简洁且完全无误的阐述；其次是对各点随不同问题演进时的阐述。此后你从世界的起初到现今，按次序解释了事物的延续；如果你愿意，我可将全部从记忆中复述。}\par

% structure: p0052 source=h2 target=subsection reason=relative-heading-rank; base=h2

\subsection{第二十六章 天主的友谊；如何获得}

伯多禄回答说：我极其欣喜，哦，\Person{克莱孟}，因我将我的话交托给如此稳妥的心；因为记住所说的话，表示已预备好行为的信德。但那邪恶的魔鬼从他偷走救恩之语、并从他记忆中夺去的人，即使他愿意，也不能得救；因为他失去了通往生命的道路。因此，我们宁可重复已经说过的话，并在你心中加以确认，即世界是如何被造的，或被谁造的，好使我们能进到与造物主的友谊。但祂的友谊是通过善生和服从祂的旨意而获得的；那旨意是所有活物的法律。因此，我们将向您简要地阐述这些事，以便它们能被更确凿地记住。\par

% structure: p0054 source=h2 target=subsection reason=relative-heading-rank; base=h2

\subsection{第二十七章 创造记述。}

\Quote{在起初，天主创造了天和地，\BibleRef{创 1:1}如同一座房屋，由世俗物体投射的阴影使其中所包含的事物陷入黑暗。但当\OriginalTerm{天主的旨意}{the will of God}引入了光时，那由物体阴影所导致的黑暗便立即消散：于是光被定为白天，黑暗被定为夜晚。那时，在世界之内的水，在起初的天和地中间的空间，像霜一样凝结，坚固如水晶，被展开，天地之间的中间空间被这样的一座穹苍分开；那穹苍，造物主称之为天，与先前所造的天同名：这样，祂将宇宙的构造分成了两部分，尽管它只是同一座房屋。划分的理由是：上部可以为天使提供居所，下部则为人。之后，海洋和混沌的位置遵永恒旨意之命，接受了留在下面的那一部分水；这些水流到低洼之处，干地显露出来；水的聚集成为海洋。此后，显露出来的大地产生了各种草木和灌木。它也涌出泉源和河流，不仅在平原上，也在山上。万事都已预备好，使将要居住其上的人可以按自己的意愿使用这一切，即或善或恶。}\par

% structure: p0056 source=h2 target=subsection reason=relative-heading-rank; base=h2

\subsection{第二十八章 创造记述（续）。}

\Quote{此后，祂用星辰点缀那可见的天。祂也在其中安放了太阳和月亮，使白昼享有太阳的光，黑夜享有月亮的光；同时，它们可以指示过去、现在和未来的事物。因为它们被造作为节令和日子的标记，这些标记虽然人人都能看见，却只有有学问和聪明的人才能理解。此后，祂命令从地和水中生出活物，祂造了乐园，并称之为喜悦之地。但在这一切之后，祂造了人，祂为了人而预备了万物，人的内在形像更为古老，并且万有都因他而被造，交付给他服务，并指定为他居住之用。}\par

% structure: p0058 source=h2 target=subsection reason=relative-heading-rank; base=h2

\subsection{第二十九章 巨人：洪水。}

\Quote{因此，当天地和水中一切事物都完成后，人类也繁衍增多，在第八代，曾过着天使般生活的义人，被妇女的美貌所引诱，与她们发生了混乱和不合法的结合；从那时起，他们在一切事上毫无节制、混乱行事，改变了人类事务的状态和天主所规定的生命秩序，以致通过劝说或强迫，他们迫使所有人犯罪得罪他们的造物主天主。在第九代，降生了称为古时巨人的，并非如希腊神话所传的有龙脚，而是身躯庞大的人，他们的巨型骨头在一些地方仍被展示作为证实。但针对这些，天主公义的眷顾在世界上降下洪水，以使大地从他们的玷污中被净化，并借由毁灭恶者使每一处都变为海洋。然而那时仍有一个义人，名叫诺厄，他与三个儿子和儿媳在一艘方舟中被救，成为洪水退去后世界的开拓者，带着与他一同封闭在内的动物和种子。}\par

% structure: p0060 source=h2 target=subsection reason=relative-heading-rank; base=h2

\subsection{第三十章 诺厄的儿子们。}

\Quote{在第十二代，天主祝福了人类，他们开始繁衍，\BibleRef{创 9:1}便领受了一个命令：不可尝血，因为洪水也正是为此而发。在第十三代，诺厄三子中的次子对父亲做了伤害，并被他诅咒，他便使他的后代沦为奴隶。他的长兄获得了世界中间区域的居住地份，其中是犹太地；幼子获得了东方区域，而次子获得了西方区域。在第十四代，受诅咒的后裔中的一个首先为魔鬼筑了一座祭台，目的是行魔法技艺，并在那里献上血腥的祭献。在第十五代，人们首次立了一个偶像并崇拜它。直到那时，天主赐予人的希伯来语独享统治地位。在第十六代，人们从东方迁移，来到分配给他们祖先的土地，每个人以自己的名字标记他所得之地的位置。在第十七代，宁录一世在巴比伦统治，建了一座城，并从那里迁移到波斯人那里，教他们崇拜火。}\par

% structure: p0062 source=h2 target=subsection reason=relative-heading-rank; base=h2

\subsection{第三十一章 洪水后的世界。}

\Quote{在第十八代，建起了有城墙的城邑，组织了军队并配备了武器，设立了审判官和法律，建造了庙宇，且万民的首领被尊为神明。在第十九代，那在洪水后受诅咒者的后代越过了他们在西方地区抽签所得的地界，驱赶那些分得了世界中部地区的人进入东方之地，并将他们追至波斯，而他们自己则强占了被驱逐者的土地。在第二十代，首次有儿子因乱伦之罪先于父亲去世，\BibleRef{创 11:28}。}\par

% structure: p0064 source=h2 target=subsection reason=relative-heading-rank; base=h2

\subsection{第三十二章 亚巴辣黑。}

在第二十一代，有一位智者，出自被驱逐者的种族，即诺厄长子的家族，名叫亚巴辣黑，我们的希伯来民族即源自于他。当全世界再次被谬误笼罩，且因其罪行的可憎而即将面临毁灭——这次不是洪水，而是烈火——当惩罚已悬于全地之上，始于索多玛时，此人因与乐纳他的天主为友，从天主那里获得恩惠，使全世界不至同归于尽。起初，此人身为占星家，能从星辰的排列与秩序中认识造物主，而其余人皆在迷误中，并明白万事均由祂的上智安排所管理。因此，一位天使在异象中显现于他，更充分地教导他开始觉察的那些事。天使也指示他关于他的种族与后代的事，并应许他说，那些地区应归还而非赐予他们。\par

% structure: p0066 source=h2 target=subsection reason=relative-heading-rank; base=h2

\subsection{第三十三章 亚巴辣黑：他的后裔。}

\Quote{因此，当亚巴辣黑渴望学习事物的原因，并深思所告知于他的事时，真先知——他独知人的心与意向——显现给他，向他启示了他所渴望的一切。真先知教给他神性的知识；启示了世界的起源，同样也启示了世界的终结；向他显示了灵魂的不朽，以及中悦天主的生活方式；也宣告了死者的复活、未来的审判、善人的赏报、恶人的惩罚——这一切都按公义的审判来管理。在明确而充分地传授所有这些知识后，他又回到那不可见的居所。然而，正如我们之前对你们所说的，当亚巴辣黑尚在无知中时，他生了两个儿子：一个名叫依市玛耳，另一个名叫赫利埃斯德罗斯。从前者生出了野蛮的民族，从后者生出了波斯人，其中有些采用了邻近民族婆罗门的生活方式和制度。另一些定居在阿拉伯，他们的后代中有些也散布到埃及。从他们中，有些印度人和埃及人学会了受割礼，并比他人守持更纯洁的礼规，尽管随着时间的推移，他们中大多数人将那纯洁的证明与标记变成了不敬虔。}\par

% structure: p0068 source=h2 target=subsection reason=relative-heading-rank; base=h2

\subsection{第三十四章 在埃及的以色列子民。}

然而，由于他在仍处于对事物无知的时候得了这两个儿子，在获得对天主的认识后，他向那位公义者祈求，使他堪当由合法妻子撒辣得后裔，尽管她不生育。她生了一个儿子，他起名叫依撒格；从他生了雅各伯，从雅各伯生了十二位圣祖，从这十二位生了七十二人。当饥荒发生时，这些人带着全家来到埃及。在四百年间，因天主的祝福与应许而繁衍增多，他们受埃及人压迫。当他们受压迫时，真先知显现给梅瑟，\EditorialNote{出谷纪第三章}，并以十灾打击埃及人，因为他们拒绝让希伯来人民离开他们，回返故土；他带领天主的子民出了埃及。但那些在灾祸中幸存的埃及人，受他们国王敌意的感染，追赶希伯来人。他们在海边追上了他们，并以为要毁灭消灭他们全部时，梅瑟向天主倾心祈祷，将海分为两半，使水停在左右两边，如同冻结了一般，天主的子民如行干地般过去；而那些追赶他们的埃及人，鲁莽地进入，被淹死了。因为当最后一个希伯来人出来时，最后一个埃及人下了海；立刻，那由祂命令如同霜冻般被束缚的海水，因那束缚者的命令而解开了，恢复其自然的自由，惩罚了那邪恶的民族。\par

% structure: p0070 source=h2 target=subsection reason=relative-heading-rank; base=h2

\subsection{第三十五章 出谷。}

\Quote{此后，\Person{梅瑟}奉天主的命令——天主的眷顾临于万有——率领希伯来人出到旷野；他舍去了从埃及直通犹太的最短道路，带领百姓绕行旷野的漫长迂回，好藉四十年之久的训练，生活方式的革新能根除因长期与埃及人的习俗亲近而缠住他们的恶习。其间他们来到西乃山，法律便藉天上的声音和异象赐给了他们，写在十条诫命中，其中第一条且最大的是：他们应只敬拜天主本身，不可为自己制造任何形象或样式来敬拜。但当\Person{梅瑟}上山，在那里停留了四十天时，百姓虽然见过埃及被十灾击打，海水分开让他们步行穿过，玛纳从天上赐给他们作食物，饮料从跟随他们的磐石中涌出\BibleRef{格前 10:4}，这种食物能变成任何人所愿的味道；并且虽然他们处在酷热的天下，白天有云彩遮蔽他们，免得被炎热灼伤，夜间有火柱照亮，免得黑暗的恐怖加于旷野的荒凉之上；——我所说的这些百姓，当\Person{梅瑟}留在山上时，却造了一个金牛犊的头，仿效他们在埃及所见的阿庇斯样式，并敬拜它；在见过如此多、如此大的奇迹之后，仍不能洗净并洗除旧习惯的污秽。因此，\Person{梅瑟}舍去从埃及直通犹太的短道，带领他们绕行旷野的巨大迂回，但愿他能如前所述，藉新教育的变化摇脱旧习惯的恶。}\par

% structure: p0072 source=h2 target=subsection reason=relative-heading-rank; base=h2

\subsection{第三十六章 暂时允许祭献}

其间，当那忠信而聪明的管家\Person{梅瑟}察觉到，百姓因与埃及人的交往，拜偶像的恶习已深植于心，这邪恶的根无法从他们中拔除时，他确实允许他们献祭，但只允许向天主献祭，好藉此断绝根深蒂固的邪恶的一半，留下另一半由另一位在将来时候来纠正；即由那一位，关于他\Person{梅瑟}自己曾说：\Quote{主，你的天主必为你兴起一位先知，你要听从祂，如同听从我一样，凡祂对你们所说的一切。凡不听从那位先知的，他的灵魂必从百姓中剪除。}\par

% structure: p0074 source=h2 target=subsection reason=relative-heading-rank; base=h2

\subsection{第三十七章 圣所}

\Quote{此外，他还指定了一个地方，只有在那里他们才可以向天主献祭。这一切的安排都是为了：当适当时机来到，他们藉先知得知天主喜爱仁慈胜过祭献时，他们能看见那一位，祂要教导他们：天主所拣选的、适合向天主奉献牺牲的地方，就是祂的智慧；而另一方面，他们也要听见：这个似乎暂时被拣选的地方，虽屡遭敌人入侵和掠夺，终将完全被毁灭。为了将此铭记于他们心中，甚至在真先知到来——祂将同时弃绝祭献和地方——之前，这地方就常被敌人掠夺并焚烧，百姓被掳到外邦中，而当他们投靠天主的仁慈时又被带回；好藉这些事教导他们：献祭的百姓被驱逐，交在敌人手中；而行仁慈和正义的人，虽无祭献，却脱离被掳，归回本土。但结果很少有人明白这事；因为大多数人，虽然能察觉并观察这些事，却被庸众的非理性意见所束缚：因为正当的意见加上自由，乃是少数人的特权。}\par

% structure: p0076 source=h2 target=subsection reason=relative-heading-rank; base=h2

\subsection{第三十八章 以色列人的罪}

\Quote{\Person{梅瑟}安排了这些事，并指派一个名叫\Person{若苏厄}的人带领百姓进入祖地，他自己则遵永生天主的命令上了一座山，并在那里去世。但他的死法是如此，以至直到今日无人找到他的埋葬之处。因此，当百姓到达祖地时，藉天主的眷顾，第一次进攻中邪恶种族的居民就被击溃，他们进入祖传产业，由抽签分配给他们。此后一段时间，他们不是由君王统治，而是由民长治理，并处于较为和平的状态。但当他们为自己寻求暴君而非君王时，也随着王室的野心，在为他们指定祈祷的地方建造了一座圣殿；这样，通过一连串邪恶的君王，百姓堕落至更大更深的邪恶。}\par

% structure: p0078 source=h2 target=subsection reason=relative-heading-rank; base=h2

\subsection{第三十九章 以圣洗取代祭献}

\Quote{但当开始临近时，正如我们所说，\Person{梅瑟}制度中所缺乏的应被补足，那位他预先宣告的先知应出现，他要藉天主的仁慈警告他们停止献祭；免得他们可能认为献祭停止后就没有罪的赦免，他便在他们中设立了圣洗，使他们能因呼求祂的名而从一切罪中获得赦免，并从此追随完美的生活，住在不朽之中，不是被兽血净化，而是被天主智慧的净化所洁净。随后，这伟大奥秘的一个明显证据也\Emph{事实上}提供出来，即凡信这\Person{梅瑟}所预言的先知、并因祂的名受洗的人，必从不信民族及其地方本身行将临到的战争毁灭中安然无恙；而不信的人则将从他们的地方和王国中被逐出，使他们甚至违背自己的意愿也明白并服从天主的旨意。}\par

% structure: p0080 source=h2 target=subsection reason=relative-heading-rank; base=h2

\subsection{第四十章 真先知的来临}

\Quote{这些事情既已预先安排，那所期待的一位便来了，带着迹象和奇迹作为祂的凭据，好使祂被显明。但即便这样，百姓仍不信，尽管他们历代以来已被训练去相信这些事。不仅不信，他们还在不信上加添了亵渎，说祂是个贪食好酒的人，并被魔鬼驱使——那正是为他们救恩而来的祂。邪恶借着恶者之手竟如此猖獗，以至于若非天主的智慧帮助那爱真理的人，几乎所有人都将陷入不虔的迷惑中。因此，祂拣选了我们十二人\EditorialNote{玛窦福音第十章}，就是最先信祂的人，并称我们为宗徒；后来又拣选了其他七十二位最被认可的门徒\EditorialNote{路加福音第十章}，好叫群众至少由此认出梅瑟的样式\BibleRef{户 11:16}，从而相信这位就是梅瑟所预言将要来的先知。}\BibleRef{申 18:15}\par

% structure: p0082 source=h2 target=subsection reason=relative-heading-rank; base=h2

\subsection{第四十一章 对真先知的拒绝}

\Quote{但或许有人说，任何人都能模仿一个数字；然而我们对于祂所行的迹象和奇迹又能说什么呢？因为梅瑟曾在埃及行过奇迹和治愈。祂——就是梅瑟所预言要兴起的像他自己的那位先知——虽然治愈了百姓中各样的疾病和软弱，行了无数的奇迹，并宣讲永生，却被恶人匆忙地推向十字架；然而，这举动却因祂的大能转变为善。总之，当祂受苦时，全世界都与祂一同受苦：因为太阳变暗了，山岭崩裂了，坟墓打开了，圣殿的帐幔裂开了，如同为那即将临到此地的毁灭而哀悼。然而，尽管整个世界都震动了，他们自己却至今仍未被这些如此重大的事所打动而悔悟。}\par

% structure: p0084 source=h2 target=subsection reason=relative-heading-rank; base=h2

\subsection{第四十二章 外邦人的蒙召}

\Quote{但是，既然外邦人必须被召来填补那些仍不信者留下的空缺，以便那曾向亚巴郎所启示的数目得以满足，那么，有福的天主之国的宣讲便被遣送到普世。为此，世俗的诸灵受到搅动，他们常与寻求自由的人为敌，并用错误的器械来摧毁天主的工程；而那些奋力奔向安全与自由之荣耀的人，因抵抗这些神灵而变得更加勇敢，并在与它们激烈斗争的努力中，获得安全的冠冕，并伴随着胜利的棕榈。与此同时，当祂受苦，黑暗从第六时辰直到第九时辰笼罩了世界，\BibleRef{玛 27:45}一旦太阳重新照耀，万物恢复正常，连恶人也回到自己身上，回到他们先前的行径，因为他们的恐惧已经消退。他们中的一些人，极其小心地看守那地方，当无法阻止祂复活时，便说祂是个术士；另一些人则谎称祂被偷走了。}\BibleRef{玛 28:13}\par

% structure: p0086 source=h2 target=subsection reason=relative-heading-rank; base=h2

\subsection{第四十三章 福音的成功}

\Quote{然而，真理在处处得胜；因为，为证明这些事是凭天主的大能而成就，我们这些原本极少数的人，在几天之内，借着天主的帮助，变得比他们多得多了。以致司祭们一度害怕，恐怕因天主的眷顾，到他们的困惑，全体百姓都投奔我们的信仰。因此他们常派人到我们这里来，请我们向他们讲论关于耶稣的事，祂是否是梅瑟所预言的那位先知，即永久的基督。\BibleRef{若 12:34}因为只在这一点上，我们这些相信耶稣的人与不信的犹太人之间似乎存在分歧。但尽管他们屡次向我们提出这样的请求，我们也寻找合适的时机，从主受难起一周的年份已经完成，主在耶路撒冷建立的教会大大地繁衍增长，由\Person{雅各伯}以最公义的规章治理，他是由主祝圣为该教会的主教。}\par

% structure: p0088 source=h2 target=subsection reason=relative-heading-rank; base=h2

\subsection{第四十四章 盖法的挑战}

\Quote{但是，当我们十二宗徒在逾越节那一天，与一大群人聚集，并进入弟兄们的教会时，我们每人在\Person{雅各伯}的要求下，简要地向会众叙述了我们在各处所行的事。正在此时，大司祭\Person{盖法}派司祭到我们这里来，请我们到他那里去，要我们向他证明耶稣就是永久的基督，或者由他向我们证明祂不是，好使全体百姓在这样或那样的信仰上达成一致；他屡次恳求我们这样做。但我们常常拖延，总是等待更合适的时机。}然后我，\Person{克莱孟}，对此回答说：\Quote{我认为，这个关于祂是否是基督的问题，对于奠定信仰至关重要；否则大司祭不会如此频繁地询问，好使自己能就基督有所学习或教导。}于是\Person{伯多禄}说：\Quote{你回答得对，\Person{克莱孟}；因为正如没有人能没有眼睛而看见，没有耳朵而听见，没有鼻孔而闻气味，没有舌头而尝味，没有手而触摸任何东西，同样，没有真先知，就不可能知道什么是中悦天主的。}我又回答说：\Quote{我已从你的教导中得知，这真先知就是\Emph{基督}；但我希望了解\Emph{基督}是什么意思，或者祂为何如此称呼，免得这件如此重要的事在我心中模糊不清。}\par

% structure: p0090 source=h2 target=subsection reason=relative-heading-rank; base=h2

\subsection{第四十五章 真先知：为何称为基督}

那时伯多禄开始这样教导我：\Quote{当天主创造了世界，作为宇宙的主宰，祂为各种受造物设立了首领：树木、山岳、泉源、河流，以及祂所造的一切，正如我们告诉你的；因为逐一细述太长了。因此，祂设立一位天使作天使的首领，一位神灵作神灵的首领，一颗星辰作星辰的首领，一个魔鬼作魔鬼的首领，一只鸟作鸟的首领，一只野兽作野兽的首领，一条蛇作蛇的首领，一条鱼作鱼的首领，一个人作人的首领，他就是耶稣基督。但祂之所以称为\Emph{基督}，乃是由于一种卓越的宗教礼仪；因为正如国王有某些共同的称号，如波斯人中的阿尔萨息，罗马人中的凯撒，埃及人中的法老，同样在犹太人中间，国王被称为\Emph{基督}。这个称号的缘由是：虽然祂确实是天主子，万物的起始，祂却成了人；天主首先用取自生命树的油傅了祂；因此，从那次傅油起，祂被称为\Emph{基督}。此外，祂自己按照圣父的安排，也用同样的油傅每一位来到祂国度的虔诚人，作为他们劳苦后的安息，因为他们已经克服了路途的艰难；好使他们的光可以照耀，并充满圣神，得享不朽。但我想我已经充分向你解释了那支从中取得香膏的树枝的全部性质。}\par

% structure: p0092 source=h2 target=subsection reason=relative-heading-rank; base=h2

\subsection{第四十六章 傅油}

\Quote{但现在我也要用一个非常简短的叙述，使你想起这一切。在今世，亚郎，第一位大司祭，是用一种由香膏调制的傅油受傅的，这傅油是按照我们先前所说过的属灵香膏的样式制成的。他是民众的首领，如同君王一样，从民众那里一一收取初果和贡物；他承担了审判民众的职务，判断洁净与不洁之物。但如果别人也接受了同一种傅油，就从中获得德能，成为君王、先知或司铎。如果这由人调制的时间性的恩宠都有如此效力，那么你想，由天主从生命树的树枝中提取的香膏该有多大能力？因为人调制的傅油都能在人中赋予如此卓越的尊位。在今世，还有什么比先知更光荣，比司铎更显赫，比君王更崇高呢？}\par

% structure: p0094 source=h2 target=subsection reason=relative-heading-rank; base=h2

\subsection{第四十七章 亚当受傅为先知}

对此，我回答说：\Quote{伯多禄，我记得你告诉过我，第一个人是先知；但你没有说他曾受傅。如果先知不能没有傅油，那么第一个人既是先知，又怎可能未受傅呢？}于是伯多禄微笑着说：\Quote{如果第一个人预言过，那他肯定也受过傅。因为虽然记录法律的人在他的书页中对他的傅油保持沉默，但他显然让我们明白这些事。因为，如果他说过他受了傅，那么他同时也是先知，虽然法律中没有记载，也是无可置疑的；同样，既然他肯定是先知，那么他也肯定受过傅，因为没有傅油他不能成为先知。但你更应该这样说：如果香膏是由亚郎用香料调制而成的，那么第一个人怎么能在亚郎时代之前受傅呢？那时调香的技术尚未被发现。}于是我回答说：\Quote{伯多禄，请不要误解我；因为我所说的不是那种调制的香膏和时间性的油，而是那种简单而永恒的\Emph{香膏}，你告诉我是由天主所造，而你说另一种是效仿它由人调制的。}\par

% structure: p0096 source=h2 target=subsection reason=relative-heading-rank; base=h2

\subsection{第四十八章 真先知、司铎}

于是伯多禄面带愠色回答说：怎么！你以为，克莱孟，我们大家都能在时机未到之前知道一切事吗？但现在不要偏离我们提出的论述，等以后你进步更明显时，我们再更清楚地解释这些事。\par

\Quote{然而，一位司铎或先知，用调制的香膏受傅后，在天主的祭台上点火，便在全世界上享有盛誉。但在亚郎（他是司铎）之后，另有一位从水中取出。我说的不是梅瑟，而是祂——在洗礼的水中被天主称为儿子。\BibleRef{玛 3:17}因为正是耶稣，借着洗礼的恩宠，熄灭了司铎为罪所点的火；因为自从祂显现以来，那赋予司铎、先知或王权的香膏就停止了。}\par

% structure: p0099 source=h2 target=subsection reason=relative-heading-rank; base=h2

\subsection{第四十九章 基督的两次来临}

\Quote{因此，祂的来临曾由梅瑟预言——他将天主的法律传给世人；但在梅瑟之前，还有另一个人，正如我已经告诉过你的。所以他暗示，祂要来，第一次来临是谦卑的，第二次却是荣耀的。第一次已经成就了，因为祂已经来了并教导了，祂——万民的审判者——曾被审判并被杀。但第二次来临时，祂要来审判，确实要定恶人的罪，却要将虔诚人纳入祂国度中的一份和交往。现在，对祂第二次来临的信德取决于第一次。因为先知们——尤其是雅各伯和梅瑟——曾论及第一次，有些也论及第二次。但预言卓越之处主要在于此：先知们不是按照事物的顺序预言未来的事；否则他们似乎只不过是聪明人猜测了事物的顺序所指出的事。}\par

% structure: p0101 source=h2 target=subsection reason=relative-heading-rank; base=h2

\subsection{第五十章 祂被犹太人所弃绝}

\Quote{但我要说的是：按常理，基督应当被祂所来的犹太人接纳，他们理应相信这位按照祖传传统为百姓救恩所期待的那一位；但外邦人却应抵制祂，因为既没有关于祂的应许或宣告传给他们，甚至他们连祂的名字都未曾听过。然而，先知们却违反事物的次序和顺序，说祂应成为外邦人的期待，而非犹太人的期待。\BibleRef{创 49:10}事实正是如此。因为当祂来临的时候，那些似乎因祖先传统而期待祂的人，竟完全不认识祂；而那些对祂一无所闻的人，却既相信祂已来到，又盼望祂将要来临。如此，预言在一切事上都显得信实可靠，它曾说祂是外邦人的期待。因此，犹太人对主的第一次来临犯了错误；我们与他们之间只有这一点分歧。因为他们自己知道并期待基督要来；但他们不知道祂已经谦卑地来了——就是那位被称为\Person{耶稣}的。这正是一个强有力的证据，证明祂的来临，即并非所有人都相信祂。}\par

% structure: p0103 source=h2 target=subsection reason=relative-heading-rank; base=h2

\subsection{第五十一章 唯一的救主}

\Quote{因此，天主在世界的终结时立定了祂；因为除非人类的本性保持完整，即保存意志的自由，否则人的邪恶无法被任何其他方式除去。故此，这条件被不可侵犯地保存下来，祂来邀请所有义人和那些渴望中悦祂的人进入祂的王国。为这些人，祂预备了不可言说的美善，以及天上的城\Place{耶路撒冷}，它将闪耀超过太阳的光辉，作为圣人们的居所。但不义的人、邪恶的人、藐视天主的人、将赐予他们的生命用于各种邪恶的人、把本应用于正义工作的时间用来行恶的人，祂将把他们交付给合适且应得的惩罚。至于那时将要发生的其余事情，无论是天使还是人都无法述说或描述。我们只需知道这一点就够了：天主将为善人赐予永恒的美善产业。}\par

% structure: p0105 source=h2 target=subsection reason=relative-heading-rank; base=h2

\subsection{第五十二章 基督来临前的圣人们}

他说了这话之后，我回答说：\Quote{如果那些在基督来临时被祂发现是义人的人将享受基督的王国，那么那些在祂来临之前已经死去的人，就完全被剥夺了王国吗？}\newline{} 于是\Person{伯多禄}说：\Quote{哦，\Person{克莱孟}，你强迫我触及一些不可言说的事。但既然允许宣告，我就不退缩。那么你要知道，基督从起初就存在，且永远存在，祂一直隐秘地与历代虔诚的人同在，尤其是那些等候祂的人，祂常常向他们显现。但时候还没有到，使已分解的身体复活；相反，这似乎是来自天主的赏报，即凡被看为义的人，应更长久地留在身体里；或者，至少正如法律书关于某位义人清楚记载的，天主将他提升了。\BibleRef{创 5:24}同样，其他中悦祂旨意的人也受到如此对待，他们被提升到乐园，为王国而保存。至于那些未能完全满足正义的规范，但在肉身上仍有些许残余邪恶的人，他们的身体固然分解了，但灵魂被保存在善良而蒙福的居所，以便在死人复活时，当他们重新获得自己因分解而甚至被洁净的身体时，可以按他们的善行获得永恒的产业。因此，所有那些将到达基督王国的人是有福的；因为他们不仅将逃脱地狱的惩罚，还将保持不朽，并最先看到天主圣父，在天主面前获得首批的尊荣之位。}\par

% structure: p0107 source=h2 target=subsection reason=relative-heading-rank; base=h2

\subsection{第五十三章 犹太人的敌意}

\Quote{因此，关于基督没有丝毫怀疑；所有不信的犹太人都被无边的愤怒攻击我们，唯恐他们犯罪得罪的那位正是祂。他们的恐惧越来越大，因为他们知道，当他们把祂钉在十字架上时，全世界都向祂表示同情；并且他们的身体虽然严密看守，却无处可寻；无数群众正在归附祂的信仰。因此，他们同大司祭\Person{盖法}一起，被迫一次又一次地派人到我们这里来，要求对祂名字的真相进行调查。当他们不断恳求，要么学习要么教导关于\Person{耶稣}是否是基督时，我们决定上到圣殿去，在全人民面前为祂作证，同时也指控犹太人许多他们所做的愚妄事情。因为自\Person{若翰}的日子以来，民众已分裂成许多派别。}\par

% structure: p0109 source=h2 target=subsection reason=relative-heading-rank; base=h2

\subsection{第五十四章 犹太教派}

\Quote{因为当\Person{基督}的兴起临近，要废除祭献并赐予圣洗的恩宠时，仇敌从预言中得知那时刻临近，就在民众中制造各种分裂，以便如果可能的话，可以废除先前的罪，而后来的过错却不可纠正。所以，第一个分裂就是那些被称为撒杜塞人的分裂，他们几乎在\Person{若翰}时代兴起。这些人自以为比别人更正义，就开始脱离民众的集会，否认死者的复活，\BibleRef{玛 22:23}并借着不信的论证断言，说不值得在应许赏报下来敬拜天主。这个意见的首创者是\Person{多西特奥斯}，其次是\Person{西满}。另一个分裂是撒玛黎雅人的分裂；因为他们否认死者的复活，并断言天主不应在\Place{耶路撒冷}受敬拜，而应在\Place{革黎斤山}。他们确实根据\Person{梅瑟}的预言正确期待那位唯一的真先知；但因\Person{多西特奥斯}的邪恶，他们被阻碍相信\Person{耶稣}就是他们所期待的那一位。经师和法利塞人也迷失进入另一个分裂；但这些人，既由\Person{若翰}施洗，又持守从\Person{梅瑟}传统所领受的真理之言作为天国的钥匙，却将它隐藏起来，不使民众听见。\BibleRef{路 11:52}是的，甚至\Person{若翰}的某些门徒，似乎是重要人物，也脱离民众，宣称自己的老师是\Person{基督}。但所有这些分裂都已预备好，为的是借此阻碍对\Person{基督}的信德和圣洗。}\par

% structure: p0111 source=h2 target=subsection reason=relative-heading-rank; base=h2

\subsection{第五十五章 公开讨论}

\Quote{然而，正如我们正要说的，当大司祭多次派司祭来要求我们彼此讨论关于\Person{耶稣}的事时；当合适的时机似乎来到，并且全体教会都同意时，我们上到\Place{圣殿}，与忠信的弟兄们一起站在台阶上，民众保持肃静；首先大司祭开始劝勉民众要耐心安静地听，同时见证和判断那些将要说出的事。然后，接着，他以许多赞美之词颂扬天主赐予人类用以赦罪的礼仪或祭献，却指责我们的\Person{耶稣}的圣洗是最近才引进而与祭献对立的。但\Person{玛窦}回应了他的论点，清楚地表明，凡没有得到\Person{耶稣}的圣洗的人，不但要被剥夺天国，而且在死者复活时也难免危险，即使他因良好的生活和正直的品行而得到特权。他作了这些和类似陈述后，就停了下来。}\par

% structure: p0113 source=h2 target=subsection reason=relative-heading-rank; base=h2

\subsection{第五十六章 撒杜塞人被驳斥}

\Quote{但是撒杜塞党人，他们否认死者的复活，非常愤怒，以至他们中有一个人从民众中喊叫说，那些以为死者会复活的人大大错了。我的兄弟\Person{安德肋}回答他说，那不是错误，而是最确定的信德之事，即死者复活，符合那一位的教导，\Person{梅瑟}曾预言他要来作真先知。\InnerQuote{或者，他说，如果你们不认为这就是\Person{梅瑟}所预言的那一位，那么先查究这个问题，当此法清楚证明他就是那一位时，关于他所教导的事就再也没有怀疑了。} \Person{安德肋}宣告了这些和许多类似的话，然后停了下来。}\par

% structure: p0115 source=h2 target=subsection reason=relative-heading-rank; base=h2

\subsection{第五十七章 撒玛黎雅人被驳斥}

\Quote{但有一个撒玛黎雅人，说话反对民众和反对天主，断言死者不会复活，也不应维持\Place{耶路撒冷}的敬拜，而应尊敬\Place{革黎斤山}，他还加上了这个反对我们的话，说我们的\Person{耶稣}不是\Person{梅瑟}预言要来到世上的那位先知。\Person{载伯德}的儿子\Person{雅各伯}和\Person{若望}，与他以及另一个支持他言论的人激烈争辩；虽然他们得到命令不要进入他们的城市，\BibleRef{玛 10:5}也不要向他们宣讲真道，但唯恐他们的言论若不加以限制会伤害别人的信德，他们回答得如此谨慎而有力，以至使他们永远沉默。因为\Person{雅各伯}作了一篇关于死者复活的演说，获得全体民众的赞同；而\Person{若望}则表明，如果他们抛弃\Place{革黎斤山}的错误，就会因此承认\Person{耶稣}确实是按照\Person{梅瑟}的预言所期待的那一位；因为，正如\Person{梅瑟}行了奇迹，\Person{耶稣}也行过奇事。而且无疑，奇迹的相似性证明他就是梅瑟所说要来的那位先知，\InnerQuote{像他自己一样}。他们宣告了这些和更多类似的话，就停了下来。}\par

% structure: p0117 source=h2 target=subsection reason=relative-heading-rank; base=h2

\subsection{第五十八章 经师被驳斥}

\Quote{并且，看，有一个经师从民众中间喊叫说：\InnerQuote{你们的\Person{耶稣}所行的奇迹，他不是作为先知而是作为术士行的。} \Person{斐理伯}急切地反驳他，表明用这个论点他也指控了\Person{梅瑟}。因为当\Person{梅瑟}在\Place{埃及}行奇迹时，正如\Person{耶稣}在\Place{犹太}行过一样，毫无疑问，对耶稣所说的同样也可以对\Person{梅瑟}说。他作了这些和类似的抗议后，就沉默了。}\par

% structure: p0119 source=h2 target=subsection reason=relative-heading-rank; base=h2

\subsection{第五十九章 法利塞人被驳斥}

\Quote{当时，一个法利塞人听见这话，就责备\Person{斐理伯}，因为他将\Person{耶稣}与\Person{梅瑟}相提并论。\Person{巴尔多禄茂}回答说，大胆宣称：我们不仅说\Person{耶稣}与\Person{梅瑟}平等，而且说他比\Person{梅瑟}更大，因为\Person{梅瑟}确实是一位先知，\Person{耶稣}也是如此，但\Person{梅瑟}并非基督，而\Person{耶稣}是基督，因此，那位既是先知又是基督的，无疑大于那位仅仅是先知的。在论述完这一连串论证后，他停了下来。在他之后，\Person{阿尔斐}的儿子\Person{雅各伯}对民众讲论，意图表明：我们相信\Person{耶稣}，并非基于先知们曾预言过他，而是我们相信先知们，他们确实是先知，因为基督为他们作证；因为基督的临在与来临才表明他们是真先知：因为作证必须由更高者向较低者作，而非由较低者向较高者作。在说了这些和许多类似的话后，\Person{雅各伯}也沉默了。在他之后，\Person{雷巴尤}开始激烈地指责民众，因为他们不相信\Person{耶稣}，而\Person{耶稣}曾以教导他们属于天主的事、安慰受苦者、治愈病人、周济穷人等方式，为他们行了如此多的善事；然而对于这一切恩惠，他们的回报却是仇恨与死亡。当他向民众宣告了这些和更多这样的事后，他便停止了。}\par

% structure: p0121 source=h2 target=subsection reason=relative-heading-rank; base=h2

\subsection{第60章 被驳斥的若翰门徒}

\Quote{且看，一个若翰的门徒断言若翰是基督，而不是耶稣，因为耶稣亲自宣称若翰比所有人、所有先知都大\BibleRef{玛 11:9}。然后他说：\InnerQuote{如果这样，他比众人更大，就必须被认为比梅瑟和耶稣本人更大。但如果他是最大的，那么他必然是基督。}对此，客纳罕人西满回答说，若翰确实比所有先知和所有由女人所生的人都大，但他并不比人子更大。因此耶稣是基督，而若翰只是先知；他们之间的差别，就如先驱与所预备的主之间的差别，或如赐律法者与守律法者之间的差别。在说了这些和类似的话后，客纳罕人也沉默了。在他之后，巴尔纳伯，又称玛弟亚，被选为宗徒代替犹大，开始劝勉民众不要以仇恨看待耶稣，也不要诽谤他。因为即使有人对耶稣无知或怀疑，爱他也远比恨他更合宜。因为天主为爱设定了赏报，为恨设定了惩罚。他说：\InnerQuote{正是因为他取了犹太人的身体，生在犹太人当中，这岂不激励我们所有人都爱他吗？}当他讲了这些和更多同样的话后，便停止了。}\par

% structure: p0123 source=h2 target=subsection reason=relative-heading-rank; base=h2

\subsection{第61章 盖法回答}

\Quote{于是盖法试图诋毁耶稣的教导，说耶稣讲的是虚妄之事，因为他称穷人为有福的；应许属地的赏报；将首要的恩赐置于属地的产业中；应许持守正义的人要饱足饮食；耶稣还被指控教导许多这类的事。多默在回答中证明他的指控是轻率的；指出盖法所相信的先知们更大量地教导这些事，却没有指明这些事要如何实现或如何理解；而耶稣则指出了该如何领会它们。当多默说了这些和类似的话后，也沉默了。}\par

% structure: p0125 source=h2 target=subsection reason=relative-heading-rank; base=h2

\subsection{第62章 宣讲的愚妄}

\Quote{于是盖法又看着我，时而警告，时而指控，说我将来应当停止宣讲基督耶稣，以免自取灭亡，也免得我被欺骗又欺骗别人。此外，他还指责我狂妄，因为尽管我未受过教育，是个渔夫、乡下人，却敢于担任教师的职务。当他讲了这些和许多类似的话后，我回答说，如果如他所说，这位耶稣不是基督，那么我所冒的风险较小，因为我接受了他作为律法的教师；但如果这位确实是基督——他确实是——那么他就处于可怕的危险中：因为我相信了这位已经显现的；而他却将为那从未显现的保留他的信仰？但是，如果像我这样未受过教育、没有学问的人，如你所说，一个渔夫、乡下人，却比有智慧的老年人更有见识，那么，我说，这更应该令你惊恐。因为如果我是凭学识辩倒你们这些有智慧、有学问的人，那就会显得我是靠长期学习获得这种能力，而非靠天主德能的恩宠；但现在，正如我所说，我们这些没有技巧的人却说服并战胜了你们这些有智慧的人，有理智的人谁看不出这不是人的巧妙，而是天主的旨意和恩赐呢？}\par

% structure: p0127 source=h2 target=subsection reason=relative-heading-rank; base=h2

\subsection{第63章 向犹太人的呼吁}

\Quote{我们就这样辩论并作证；我们这些未受过教育的渔夫，教导司铎们关于唯一的天上之主；教导撒杜塞人关于死人的复活；教导撒玛黎雅人关于耶路撒冷的神圣（我们并未进入他们的城，而是在公开场合与他们辩论）；教导经师和法利塞人关于天国；教导若翰的门徒，不要使若翰成为他们的绊脚石；教导全体民众，耶稣是永恒的基督。最后，我警告他们，在我们出发前往外邦人那里，向他们宣讲天主父的知识之前，他们应当与天主和好，接受他的圣子；因为我向他们表明，他们无法以其他方式得救，除非藉着圣神的恩宠，他们急忙受洗，接受三重呼号的洗礼，并领受主基督的圣体圣事；他们只应相信他所教导的那些事，如此才能配得获得永恒的救恩；否则，他们绝不可能与天主和好，即使为天主点燃一千座祭台和一千座高祭台。}\par

% structure: p0129 source=h2 target=subsection reason=relative-heading-rank; base=h2

\subsection{第64章 圣殿将被毁灭}

\Quote{因为我们已确凿无疑地确认，天主远更不悦于你们所奉献的祭献，因为祭献的时期现已过去；并且因为你们不肯承认奉献牺牲的时期已经过去，因此圣殿将被毁灭，\OriginalTerm{荒凉的可憎之物}{abomination of desolation}将立于圣地；那时福音将传于外邦人，作为反对你们的证据，使你们的无信因他们的信德而受审判。因为整个世界在不同时期遭受各种病患，或普遍蔓延于众人，或特别影响某些人。因此它需要一位医生来访问它，为它的救恩。因此我们向你们作证，并向你们宣布对你们每个人隐藏的事。由你们考虑什么对你们有利。}\par

% structure: p0131 source=h2 target=subsection reason=relative-heading-rank; base=h2

\subsection{第65章 \Person{加玛里耳}平息骚动}

\Quote{当我这样说了以后，全体司铎都大为恼怒，因为我预言了圣殿的倾覆。当民众的首领\Person{加玛里耳}看见这个情况——他暗中是我们信仰中的弟兄，但根据我们的建议仍留在他们中间——因为他们对我们极其愤怒，动了强烈的怒，他站起来说：\BibleRef{宗 5:35} \InnerQuote{你们暂且安静，以色列人啊，因为你们没有察觉悬在你们头上的考验。因此，不要干涉这些人；如果他们所从事的是出于人的计谋，很快就会结束；但如果是出于天主，你们何必无故犯罪，一无所成呢？因为谁能胜过天主的旨意？现在，既然白日渐近黄昏，我明天将在同一地点，当着你们的面，与这些人辩论，使我公开反对并清楚驳斥每一个错误。}通过这番话，他们的怒气稍减，尤其希望明天我们会被公开定罪；于是他平和地遣散了百姓。}\par

% structure: p0133 source=h2 target=subsection reason=relative-heading-rank; base=h2

\subsection{第66章 讨论恢复}

现在当我们来到我们的\Person{雅各伯}那里，向他详细陈述所说所行的一切后，我们用了晚餐，并与他同住，整夜向全能天主恳求，使即将进行的辩论能显示我们信仰的无可置疑的真理。因此，次日，\Person{雅各伯}主教与我们和整个教会一同上到圣殿。在那里我们发现了大群百姓，他们从半夜就开始等候我们。于是我们站在先前同一地方，站在高处，使众人能看到我们。然后，当获得深沉的寂静时，\Person{加玛里耳}，如前所述，是信我们信仰的，但因天主的安排仍留在他们中间，以便倘若他们任何时候企图对我们行不义或邪恶之事，他或是通过巧妙地采纳建议来阻止他们，或是提醒我们，使我们或能戒备，或能避开——于是他，仿佛在反对我们，首先望着\Person{雅各伯}主教，这样对他说：\par

% structure: p0135 source=h2 target=subsection reason=relative-heading-rank; base=h2

\subsection{第67章 \Person{加玛里耳}的讲话}

\Quote{如果我，\Person{加玛里耳}，认为那对我的学识或我的高龄都不是耻辱，从婴孩和未曾受教育者那里学习一些东西，倘若有什么有益或安全的东西可以获取（因为合理生活的人知道没有什么比灵魂更宝贵），这岂不应成为众人爱慕和渴望的对象，即学习他们所不知道的，教导他们所学习到的？因为最确实的是，友谊、亲情、或崇高权力，都不应比真理对人更宝贵。因此，诸位弟兄，如果你们知道更多，不要退缩，不在在场的天主之民和你们的弟兄面前提出；而全体百姓将乐意且极其安静地听你们所说。因为百姓为何不这样做呢？当他们看见连我也与他们一样愿意从你们学习，倘若天主向你们启示了更多？但如果你们在某些方面有欠缺，你们也不要同样以向我们学习为耻，好使天主补足双方所欠缺的。但如果现在有恐惧激动你们，因着我们中一些人心存偏见反对你们，且因害怕他们的暴力，你们不敢公开说话，那么为使你们免于这恐惧，我向你们公开宣誓，凭着永生的全能天主，我决不容许任何人下手攻击你们。既然你们有全体百姓作我这誓言的见证，且你们以我们圣事盟约为合宜的保证，愿你们每个人毫不迟疑地宣布他所学到的；让我们，弟兄们，热切且安静地聆听。}\par

% structure: p0137 source=h2 target=subsection reason=relative-heading-rank; base=h2

\subsection{第68章 信德的准则}

\Quote{这些话并不太使\Person{盖法}喜悦；他怀疑\Person{加玛里耳}，便开始狡猾地介入讨论：因为，对\Person{加玛里耳}所说的话报以微笑，大司祭问主教长\Person{雅各伯}，关于基督的讨论不应从圣经之外引用；他说，好使我们知道，他说，耶稣是否就是那基督。然后\Person{雅各伯}说：我们必须首先查考我们尤其该从哪些圣经进行讨论。然后他经过困难，终于被理性说服，回答说必须从法律中引证；之后他也提到了先知。}\par

% structure: p0139 source=h2 target=subsection reason=relative-heading-rank; base=h2

\subsection{第69章 基督的两次来临}

\newline{}\Quote{我们的雅各伯开始向他们说明：先知们所说的一切话都是从法律中取得的，他们所说的话符合法律。他也讲了关于\BibleBook{列王}{纪}的一些事：这些书是怎样、何时、由谁写的，以及应如何使用。当他最详尽地讨论了法律，并以最清晰的解释阐明了其中有关基督的一切后，他用最丰富的证据证明耶稣就是基督，并且所有关于祂谦卑来临的预言都在祂身上应验了。因为他说明，祂的两次来临已被预言：一次在谦卑中，已经实现；另一次在荣耀中，希望将来实现，那时祂要来将王国赐给那些相信祂并遵守祂所命令的一切的人。当他清楚地教导了人民关于这些事后，他又补充了这一点：除非一个人在水里，以\OriginalTerm{三重福分}{the threefold blessedness}的名受洗，正如\OriginalTerm{真先知}{the true Prophet}所教导的，他既不能获得罪的赦免，也不能进入天国；他宣布这是\OriginalTerm{无生之天主}{the unbegotten God}的规定。对此他又补充说：\InnerQuote{不要以为我们谈论两位无生之神，或者一位被分成两位，或者同一位变成男性和女性。但我们谈论的是天主独生子，不是从其他根源发出的，而是不可言喻地自有的；同样，我们也谈论\OriginalTerm{护慰者}{the Paraclete}。}但当他也讲了一些关于洗礼的话后，连续七天他说服了全体人民和大司祭，要他们立即去领受洗礼。}\par

% structure: p0141 source=h2 target=subsection reason=relative-heading-rank; base=h2

\subsection{第七十章　扫罗引起的骚动}

\newline{}\Quote{当事情发展到他们即将前来受洗的地步时，我们敌人中的一个人带着几个人进入圣殿，开始喊叫说：\InnerQuote{以色列人哪，你们什么意思？你们为何如此轻易被冲动？为何被这些最可怜的人所迷惑，被术士西蒙欺骗？}他这样说着，又说了更多同样的话，而雅各伯主教正在反驳他，他就开始煽动群众，制造骚乱，使人们无法听到所说的话。于是他开始用喊叫扰乱一切，破坏苦心安排好的事，同时辱骂司祭们，用斥责和辱骂激怒他们，像一个疯子一样煽动每个人去杀人，说：\InnerQuote{你们在做什么？为何犹豫？哦，迟钝而懒惰的人，我们为何不抓住他们，把所有这些家伙撕碎？}他说完这话，首先从祭台抓起一根粗大的木柴，带头殴打。于是其他人看见他，也带着同样的冲动被带动。然后双方打了起来，打人的人和被打的人一片混乱。流了很多血，混乱的逃窜中，那个敌人攻击了雅各伯，把他从台阶顶上扔了下去；以为他死了，就不在乎再对他施加暴力。}\par

% structure: p0143 source=h2 target=subsection reason=relative-heading-rank; base=h2

\subsection{第七十一章　逃往耶里哥}

\newline{}\Quote{但我们的朋友们把他扶起来，因为他们人数更多、力量更强；但因他们对天主的敬畏，他们宁愿被较弱的一方杀死，也不愿杀死别人。到了晚上，司祭们关闭了圣殿，我们回到雅各伯的家，在那里整夜祈祷。然后天亮前我们下到耶里哥，大约五千人。三天后，一位弟兄从我们之前提到的加玛里耳那里来到我们这里，秘密地告诉我们，那个敌人已经从大司祭盖法那里得到委任，要逮捕所有相信耶稣的人，并带着他的信去大马士革，在那里也要借助不信之人的帮助，对信友们进行破坏；他之所以急于去大马士革，主要是因为他相信伯多禄已经逃到那里。大约三十天后，他在经过耶里哥去大马士革的路上停留。那时我们不在，出去探望两位弟兄的坟墓，这两座坟墓每年都会自动变白，这个奇迹抑制了许多人对我们的愤怒，因为他们看到我们的弟兄在天主面前被纪念。}\par

% structure: p0145 source=h2 target=subsection reason=relative-heading-rank; base=h2

\subsection{第七十二章　伯多禄被派往凯撒勒雅}

\newline{}\Quote{因此，当我们住在耶里哥，专心祈祷和禁食时，雅各伯主教派人来找我，把我派到这里凯撒勒雅，说匝凯从凯撒勒雅写信给他，说有一个撒玛黎雅术士西蒙正在使我们的许多人动摇，声称他自己就是\OriginalTerm{那站立者}{Stans}——也就是说，就是基督，以及至高天主的伟大权能，这权能超越世界的创造者；同时他行了许多奇迹，使一些人怀疑，另一些人则投向他。他向我详细报告了从那些从前是他的同伙或门徒、后来归向匝凯的人那里得来的关于这个人的一切情况。\InnerQuote{因此，伯多禄，}雅各伯说，\InnerQuote{有许多人为了他们的安全，你应该去驳斥那个术士，教导真理之言。所以不要迟延；也不要因你独自出发而难过，要知道天主藉着耶稣会与你同行，会帮助你，并且很快，因祂的恩宠，你会有许多同伴和同情者。现在你要确保每年给我寄一份你言行事迹的书面报告，特别是每七年年底。}说完这些话，他打发我走了，我在六天内到达了凯撒勒雅。}\par

% structure: p0147 source=h2 target=subsection reason=relative-heading-rank; base=h2

\subsection{第七十三章　匝凯的欢迎}

\Quote{当我进城时，我们至爱的弟兄匝凯遇见了我；他拥抱了我，带我来到这住处，他自己就住在这里，他询问我各位弟兄的情况，特别是关于我们可敬的弟兄雅各伯。我告诉他，他仍是一足跛行，他立刻问起原因，我就向他讲述了刚才向你们详细叙述的一切：我们如何被司祭们和大司祭盖法叫到圣殿；总主教雅各伯如何站在台阶顶端，一连七天向全体百姓从主的圣经指出耶稣就是基督；当众人都同意应由他因耶稣的名给他们付洗时，一个仇敌如何行了所有我已然提及、无需重复之事。}\par

% structure: p0149 source=h2 target=subsection reason=relative-heading-rank; base=h2

\subsection{第74章 西满术士挑战伯多禄}

\Quote{当匝凯听了这些事后，他也向我讲述了西满的行径；与此同时，西满本人——我不知道他如何得知我到来——派人捎信给我说：\InnerQuote{我们明天当着百姓的面辩论。}我回答说：\InnerQuote{就按你的意思吧。}我的这个承诺传遍了全城，因此就连你——你在那天到达——也得知我第二天要与西满辩论，并按照你从巴尔纳伯那里得到的指示，找到我的住处，来到我这里。但我对你到来如此欢喜，以致我的心——不知怎的受推动——急忙向你迅速阐述一切，尤其是我们信德的核心要点，即关于真先知的道理；我毫不怀疑，仅这一点就足以作为我们全部教义的基础。然后，其次，我向你阐明了成文法律更隐秘的含义，按其各个要点，当时有机会阐明；我也没有向你隐瞒传统中的美好事物。但其余部分，从明天起，你将逐日听到，它们与同西满的讨论中所提出的问题相关，直到赖天主的恩宠，我们抵达那座我们认为旅程所向的罗马城。}\par

我随即表示，我对他所告知的一切心怀感激，并允诺我将欣然遵从他的一切命令。然后，用过餐后，他吩咐我休息，他自己也去休息了。\par

\SourceRecord{Translated by Thomas Smith.；Ante-Nicene Fathers；Vol. 8.；Edited by Alexander Roberts, James Donaldson, and A. Cleveland Coxe.；Buffalo, NY: Christian Literature Publishing Co.,；1886.；<http://www.newadvent.org/fathers/080401.htm>.}

% structure: p0001 source=h1 target=section reason=page-title

\section{\WorkTitle{众认录（卷二）}}

% structure: p0002 source=h2 target=subsection reason=relative-heading-rank; base=h2

\subsection{第一章 习惯的力量}

当与\Person{西满}辩论所定的那天破晓时，\Person{伯多禄}在第一次鸡叫时起身，也唤醒了我们：因为我们都睡在同一个房间里，共十三人；次于\Person{伯多禄}，首先是\Person{匝凯}，然后是\Person{索福尼}、\Person{若瑟}和\Person{米该亚}、\Person{厄里厄则尔}、\Person{丕乃哈斯}、\Person{拉匝禄}和\Person{厄里叟}；此后是我（\Person{克肋孟}）和\Person{尼苛德摩}；然后是\Person{尼塞塔}和\Person{阿桂拉}，他们曾是\Person{西满}的门徒，并在\Person{匝凯}的教导下归依了基督信仰。妇女中没有人在场。当时夕光仍存，我们全都坐下；\Person{伯多禄}见我们醒着，并且注意听他，向我们致意后，立刻开始讲话如下：——\par

\Quote{弟兄们，我承认，我惊叹人性的力量，我看到它适应并适合每一种呼唤。然而，我想说的这一点是我凭经验发现的：当半夜过后，我自动醒来，睡眠不再临到我。这发生在我身上，原因是我养成了回忆我主话语的习惯，这些话语是亲耳听祂说的；并且因为我渴望它们，我强迫自己的心思和意念被唤醒，以便在醒着时，逐一回忆并整理它们，好能牢记在心。因此，当我渴望以心中所有的喜悦珍视主的话时，醒着的习惯就临到了我，即使没有什么我想思考的事。就这样，以某种不可解释的方式，一旦任何习惯建立起来，旧习惯就会被改变，只要你不过分勉强，而是按照本性所容许的程度。因为完全无眠是不可能的；否则夜晚就不会被造来休息了。}\par

% structure: p0005 source=h2 target=subsection reason=relative-heading-rank; base=h2

\subsection{第二章 减少睡眠}

\Quote{那时，我一听见这话，就说：\Person{伯多禄}，您说得很好；因为一个习惯被另一个习惯所取代。当我在海上时，起初我苦恼，整个系统紊乱，感觉像被打了一样，无法忍受海的颠簸和喧嚣；但几天后，当我习惯了，我开始能忍受了，以致我乐意在早晨立刻与船员们一起进食，而以前我在第七时辰之前从不吃东西。现在，仅仅由于我当时养成的习惯，饥饿在约莫我惯常与船员们进食的时候提醒我；然而，我希望一旦另一个习惯形成，就能摆脱它。因此，我相信您也已经获得了醒着的习惯，如您所述；并且您愿意在合适的时候向我们解释这一点，使我们也不吝惜抛弃并舍去一部分睡眠，以便能接受活生生的教理。因为当食物消化了，心灵在夜晚的寂静影响下，适时教导的那些事物就会存留在其中。}\par

% structure: p0007 source=h2 target=subsection reason=relative-heading-rank; base=h2

\subsection{第三章 需要谨慎}

于是\Person{伯多禄}，很高兴听到我理解了他序言的用意，即他是为了我们的益处而讲的；并称赞我，无疑是为了鼓励和激励我，开始发表以下讲话：在我看来，似乎适时且有必要对近在眼前的事情进行一些讨论，即关于\Person{西满}的事。因为我想知道他是怎样的人，品行如何。所以，如果你们中有人了解他，请务必告诉我；因为事先知道这些事是很重要的。因为如果我们受命，当我们进入一座城时，要先查问谁是可堪当的，\BibleRef{玛 10:11} 好能与他同吃，那么把永生的话语托付给谁，或者怎样的人，岂不更应当确定吗？因为我们应当小心，极其小心，免得把珍珠丢在猪面前。\BibleRef{玛 7:6}\par

% structure: p0009 source=h2 target=subsection reason=relative-heading-rank; base=h2

\subsection{第四章 对待反对者的明智}

\Quote{但还有其他原因，我有必要了解这个人。因为如果我知道，在那些无可置疑是好事的事上，他是无可指摘、无可非议的——也就是说，如果他节制、怜悯、正直、温和、有人性，这些无人怀疑是优良品质——那么，似乎就适合将信德和知识的欠缺赐予那位拥有这些美德的人；如此，他原本值得赞许的生命，应在那些显得不完美的方面得到修正。但如果他仍然身陷于那些明显是罪的过犯中并被玷污，我就不应向他说论关于神圣知识的更深奥、更神圣的事，反而应当抗议并当面反对他，使他停止犯罪，净化其行为，脱离恶习。但如果他潜入我们的谈话，诱使我们说出他——虽然他行为不当——不应听的话，我们就要谨慎地避开他。因为完全不回答他似乎不合适，为了听众的缘故，以免他们可能以为我们因缺乏回答能力而回避辩论，从而因误解我们的意图而损害他们的信德。}\par

% structure: p0011 source=h2 target=subsection reason=relative-heading-rank; base=h2

\subsection{第五章 \Person{西满}术士，一个可畏的对手}

当\Person{伯多禄}这样对我们说完后，\Person{尼塞塔}请求允许他说话；\Person{伯多禄}允许后，他说：\Quote{请原谅，我恳求您，我主\Person{伯多禄}，请听我说，我非常为您焦急，我担心在与\Person{西满}的辩论中，您会显得处于劣势。因为那捍卫真理的人常常不能得胜，原因是听众要么抱有偏见，要么对更好的事业没有多大兴趣。但除此之外，\Person{西满}本人是个极其激烈的演说家，受过\OriginalTerm{辩证术}{dialectic art}和\OriginalTerm{三段论}{syllogisms}的训练；而最糟糕的是，他精通\OriginalTerm{魔术}{magic art}。因此我担心，万一他如此全副武装，在不认识他的人面前，他会被认为是捍卫真理，而实际上他是在散播谎言。因为就连我们自己，若不是作为他的助手和与他同犯错误的人，得以查明他是个欺骗者和魔术师，我们也不可能逃脱他而归向主。}\par

% structure: p0013 source=h2 target=subsection reason=relative-heading-rank; base=h2

\subsection{第六章 西满术士：他的邪恶}

当\Person{尼塞塔}这样说完后，\Person{阿奎拉}也请求允许他说话，然后继续如下：\Quote{请接受，最卓越的\Person{伯多禄}，我对您的爱戴的保证；因为我也确实极为您焦虑。不要为此责备我们，因为关心一个人是出于爱；而冷淡无异于仇恨。但我呼求天主作证，我为您感到担心，并非因为知道您在辩论中较弱——因为我从未参加过您参与的辩论——而是因为我深知此人的不虔敬，我关心您的声誉，同时关心听众的灵魂，尤其关心真理本身的利益。因为这位\OriginalTerm{魔术师}{magician}对他所希望的一切都极为激烈，而且邪恶无度。因为我们从童年起就是他邪恶的助手和仆役，对这一切都很了解；若不是天主的爱将我们从他那里解救出来，我们甚至现在还会与他一同行恶。但一种内在的对天主之爱使他的邪恶在我们眼中变得可憎，使天主的敬拜对我们有吸引力。因此，我认为也正是\OriginalTerm{天主眷顾}{Divine Providence}的作为，让我们先成为他的同伙，从而得以了解他以何种方式或技艺行出那些看似神奇的事。因为谁不会对他所做的奇妙之事感到惊奇呢？谁不会认为他是从天降下为了人类的救恩的天主呢？我自己承认，如果不是与他亲密相识，并参与了他的所作所为，我很容易被他带走。因此，我们与他的交往分离并不困难，因为我们知道他所依赖的是\OriginalTerm{魔术技艺}{magic arts}和邪恶的计谋。但如果您自己也希望了解他的全部——他是谁，是什么，从哪里来，以及他如何策划他所做的事——那么请听。}\par

% structure: p0015 source=h2 target=subsection reason=relative-heading-rank; base=h2

\subsection{第七章 西满术士：他的历史}

\Quote{这个\Person{西满}的父亲是\Person{安东尼}，母亲是\Person{辣黑耳}。按民族他是\OriginalTerm{撒玛黎雅人}{Samaritan}，来自\Place{该通}村；按职业是\OriginalTerm{魔术师}{magician}，但极其精通\OriginalTerm{希腊文学}{Greek literature}；爱好虚荣，自夸超过全人类，以至于他希望被相信是超越造物主天主之\OriginalTerm{崇高的能力}{exalted power}，并被认为就是基督，被称为\Emph{站立的}。他使用这个名字暗示自己永不消解，声称他的肉体被其神性的能力如此凝聚，以至于可以永恒存续。因此，他被叫做\Emph{站立的}，仿佛他不会因任何朽坏而倒下。}\par

% structure: p0017 source=h2 target=subsection reason=relative-heading-rank; base=h2

\subsection{第八章 西满术士：他的历史}

\Quote{因为当\Person{若翰洗者}被杀后，正如您自己所知，当\Person{多西特乌斯}开创了他的\OriginalTerm{异端}{heresy}，连同三十个主要门徒，还有一个女人名叫\Emph{路娜}——由此看出这三十人似乎是按照月亮的运行，以日数而设立的——这个\Person{西满}，如我们所说，野心勃勃，图谋邪恶荣耀，去到\Person{多西特乌斯}那里，假装友善，恳求他，如果那三十人中有人死去，他就立即取代死者的位置：因为违反他们的规矩，要么超过固定数目，要么接纳一个未知或尚未考验的人；因此，其余的人渴望配得上这地位和数目，都急切地取悦，按照他们教派的制度，每个渴望被接纳进入数目的人，希望当有人死去时，他能被认为配得上接替死者的位置。因此，\Person{多西特乌斯}被此人极力催促，当数目中出现了空缺时，就接纳了\Person{西满}。}\par

% structure: p0019 source=h2 target=subsection reason=relative-heading-rank; base=h2

\subsection{第九章 西满术士：他的职业}

\Quote{但不久之后，他爱上了那个他们称为\Person{路娜}的女人；他把一切都向我们这些朋友吐露：他是如何作一个术士，如何爱\Person{路娜}，如何因渴求荣耀而不愿不光彩地享有她，而是耐心等待，直到能体面地享有她；但条件是，我们也必须与他合谋，协助他实现欲望。他还答应，作为这项服务的回报，他将使我们获得最高的尊荣，使人们相信我们是神；\InnerQuote{惟有一个条件，}他说，\InnerQuote{你们要把首席让给我，我\Person{西满}，凭魔术能显现许多征兆和奇事，藉以建立我的荣耀或我们的派别。因为我能使想要捉拿我的人看不见我，又能在愿意被看见时显现。若我想逃跑，我能穿山而过，穿过岩石如同穿过泥土。若我从高山顶上一头栽下，我能安然落地，仿佛被托住；被捆绑时，我能自解，并能捆绑捆绑我的人；被关在监牢里，我能使栅栏自动打开；我能使雕像活起来，以致看见的人以为它们是真人。我能使新树突然长出，并立刻发芽。我能投身火中而不被烧；我能改变容貌，以致无人能认出我；我还能向人显示我有两张面孔。我将把自己变成绵羊或山羊；我将使小男孩长出胡须；我将腾空飞行；我将展示大量黄金，并扶立和废黜君王。我将被当作天主来敬拜；我将被公开授予神圣的尊荣，以致我的雕像将被设立，我将被当作天主来敬拜和朝拜。还需要多说什么？凡我所愿，我都能成就。因为我已通过试验成就了许多事。总而言之，}他说，\InnerQuote{有一次，我母亲\Person{辣黑耳}吩咐我去田里收割，我看见一把镰刀放在那里，我命令它去收割；它收割的比其他人多十倍。最近，我使许多新芽从地里长出，并使它们瞬间长叶结果；我还成功钻穿了最近的一座山。}}\par

% structure: p0021 source=h2 target=subsection reason=relative-heading-rank; base=h2

\subsection{第十章 西满术士：他的欺骗}

\Quote{但是，当他说到发芽和穿山的事时，我对此感到困惑，因为他想欺骗我们这些他似乎信任的人；因为我们知道，那些事情从我们祖先的时代就已存在，他却说成是他最近做的。我们虽然从他那里听到这些残暴之事，甚至更甚，但我们却追随他的罪行，任由其他人被他欺骗，还为他编造许多谎言；而且这发生在他兑现任何承诺之前，以致他尚未做成任何事，就已被一些人当作天主。}\par

% structure: p0023 source=h2 target=subsection reason=relative-heading-rank; base=h2

\subsection{第十一章 西满术士：多西特乌斯派的首领}

\Quote{与此同时，一开始，当他被算作\Person{多西特乌斯}的三十门徒之一时，他开始贬损\Person{多西特乌斯}本人，说他教导得不纯正、不完美，而且这不是出于恶意，而是出于无知。但\Person{多西特乌斯}觉察到\Person{西满}在贬低他，害怕自己在人中的声誉受损（因为他自己被认为是\OriginalTerm{站立者}{Standing One}），勃然大怒，当他们在学堂照常聚会时，他抓起一根棍子，开始殴打\Person{西满}；但突然那根棍子似乎穿过了他的身体，如同穿过烟雾。\Person{多西特乌斯}见此大为惊异，对他说：\InnerQuote{告诉我，你是不是那\OriginalTerm{站立者}{Standing One}，好让我朝拜你。}而\Person{西满}回答说他是；于是\Person{多西特乌斯}意识到自己不是那\OriginalTerm{站立者}{Standing One}，便俯伏朝拜他，并把自己的首席地位让给\Person{西满}，命令那三十人的等级都服从他；他自己则占据了\Person{西满}原先的较低位置。此后不久，他就死了。}\par

% structure: p0025 source=h2 target=subsection reason=relative-heading-rank; base=h2

\subsection{第十二章 西满术士与路娜}

\Quote{因此，在\Person{多西特乌斯}死后，\Person{西满}将\Person{路娜}据为己有；他带着她四处游走，如你所见，欺骗众人，声称他自己是高于造物主天主的某种权能，而与他一起的\Person{路娜}是从高天被带下来的，她就是\OriginalTerm{智慧}{Wisdom}，万物的母亲，他说，希腊人和异邦人为此争辩，却只能在一定程度上看到她的影像；但至于她本身，作为与唯一且首要的天主同住者，他们全然无知。他宣扬这些以及诸如此类的其他说法，欺骗了许多人。但我也应当提及我亲眼所见的一件事。有一次，他的这个\Person{路娜}在某座塔中，一大群人聚集来看她，围站在塔的四周；但所有人都看见她探身向前，从塔的各个窗户向外张望。他还做了许多其他奇妙的事，至今仍在做；以致人们对此感到惊异，以为他自己就是伟大的天主。}\par

% structure: p0027 source=h2 target=subsection reason=relative-heading-rank; base=h2

\subsection{第十三章 西满术士：他魔术的秘密}

\Quote{有一次，\Person{尼塞塔}和我请他向我们解释，这些事如何能藉魔术完成，以及那东西的性质是什么，\Person{西满}便开始这样向我们他的同伴解释。\InnerQuote{我}，他说，\InnerQuote{取得了一个纯洁且被强暴杀害的男孩的灵魂，用不可言喻的咒语召来协助我；藉着它，我所命令的一切都得以成就。}\InnerQuote{但是}，我说，\InnerQuote{一个灵魂可能做这些事吗？}他回答说：\InnerQuote{我要你们知道，人的灵魂一旦从身体的黑暗中释放出来，就占据仅次于天主的地位。它立即获得预知：因此它被召来行招魂术。}于是我回答说：\InnerQuote{那么，为什么被杀者的灵魂不报复杀害他们的人呢？}\InnerQuote{你不记得吗}，他说，\InnerQuote{我告诉过你，当它离开身体时，它获得对未来的知识？}\InnerQuote{我记得}，我说。\InnerQuote{那么}，他说，\InnerQuote{一旦它离开身体，它立即知道将来会有审判，每个人都要为自己所做的恶事受罚；因此它们不愿报复杀害它们的人，因为它们自己正因在这里所做的恶事而遭受痛苦，并且知道在审判中更严厉的惩罚等着它们。此外，管理它们的天使不准它们出去或做任何事。}\InnerQuote{那么}，我回答，\InnerQuote{如果天使不准它们到这里来，或做它们所乐意的事，灵魂怎能服从那召唤它们的魔术士呢？}\InnerQuote{并非}，他说，\InnerQuote{它们纵容那些愿意来的灵魂：但是当管辖天使被一位比它们更大的恳求时，它们便以我们这些恳求者的暴力为借口，允许我们所召唤的灵魂出去：因为受暴力者并不犯罪，而是我们，将必要性加于它们。}于是\Person{尼塞塔}无法再克制，急忙回答说——其实我也正要这样做，只是我想先从他那里得到关于几个问题的信息——但如我所说，\Person{尼塞塔}抢先说道：\InnerQuote{你们难道不怕审判的日子吗？你们对天使施暴、召唤灵魂、欺骗人、为自己从人那里换取神圣的尊荣。你们怎能说服我们，说没有审判，如同有些犹太人承认的那样，而且灵魂不是不死的，如同许多人以为的那样，尽管你们亲眼看见它们，并从它们那里得到天主审判的确证？}}\par

% structure: p0029 source=h2 target=subsection reason=relative-heading-rank; base=h2

\subsection{第十四章 西满魔术士，宣称自己是天主。}

\Quote{听着他的话，\Person{西满}面色苍白；但过了一会儿，他镇定下来，这样回答说：\InnerQuote{不要以为我是你们族类的人。我既不是魔术士，也不是路娜的情人，也不是\Person{安东尼乌斯}的儿子。因为在我母亲\Person{辣黑耳}和他结合以前，她仍是一个童贞女时，就怀了我，而我能够随意成为微小或伟大，并在人中间显现为一个人。因此我首先选择了你们作为我的朋友，目的是要考验你们，好使我在考验你们之后，将你们安置在我天上不可言喻的地方。因此我假装是一个人，以便我更清楚地查明你们是否对我怀有完全的爱。}但我听到这些话时，认为他确实是个可怜虫，却惊异于他的无耻；我为他脸红，同时担心他会设法加害我们，便示意\Person{尼塞塔}与我一起暂时假装一下，并对他说：\InnerQuote{不要对我们这些可朽坏的人生气，啊，不朽坏的天主，反而接受我们的爱和愿意认识天主的意愿；因为我们直到现在还不知道你是谁，也不曾察觉你就是我们所寻找的那一位。}}\par

% structure: p0031 source=h2 target=subsection reason=relative-heading-rank; base=h2

\subsection{第十五章 西满魔术士，宣称用空气造了一个男孩。}

\Quote{当我们以适合场合的表情说着这些和类似的话时，这个最虚妄的人相信我们被骗了；因而更加得意，他又加了这话：\InnerQuote{我现在要善待你们，因为你们对我怀有如同对天主的爱；你们在不知道我的时候就爱了我，在无知中寻找我。但我不愿你们怀疑，当一个人能随意变大变小时，这真正是成为天主；因为我能以我所喜悦的任何方式向人显现。现在，我要开始向你们揭示真相。有一次，我藉我的能力，将空气变成水，再将水变成血，并将其凝固成肉，形成了一个新的人造物——一个男孩——产生了一个比天主创造者更崇高的作品。因为祂用泥土造人，而我用空气——一个困难得多的材料；我又将他消散，恢复为空气，但在此之前，我将他的画像和肖像放在我的卧室里，作为我工作的证据和纪念。}于是我们明白，他指的是那个男孩，他在被强暴杀害后，其灵魂被他用来做他所需要的事情。}\par

% structure: p0033 source=h2 target=subsection reason=relative-heading-rank; base=h2

\subsection{第十六章 西满魔术士：他的情形无望。}

伯多禄听了这些话，含泪说道：\Quote{我极为惊叹天主的无限忍耐，而另一方面，又惊叹某些人莽撞的胆大妄为。因为还有什么理由能说服西满，让他相信天主会审判不义的人呢？他既然说服自己，认为他可以运用灵魂的服从服务于他的罪行。但事实上，他被魔鬼欺骗了。然而，尽管他通过这些事情确知灵魂是不朽的，并且会为他们所行的恶事受审判，尽管他以为他真的看见了那些我们因着信德所相信的事；尽管，如我所说，他被魔鬼欺骗了，他却以为他看见了灵魂的本质。这样的人，我说，他既占据如此邪恶的立场，又怎能承认他作恶，或者承认他要为他所做的事受审判呢？他明知天主的审判，却藐视它，显明自己是天主的仇敌，还敢犯下如此可怕的事。因此，我的弟兄们，可以肯定，有些人反对天主的真理和宗教，并非因为他们认为理性与信德绝不相容，而是因为他们要么被过度的邪恶所缠累，要么被自身的恶行所阻碍，要么被内心的骄傲所抬举，以致他们甚至不相信那些他们以为亲眼看见的事。}\par

% structure: p0035 source=h2 target=subsection reason=relative-heading-rank; base=h2

\subsection{第十七章 人是天主的仇敌}

\Quote{然而，既然对造物主天主天生的爱似乎足以使那些爱祂的人获得救恩，仇敌就设法扭曲人里面的这种爱，使他们变得敌对和忘恩负义于他们的造物主。因为我呼天唤地作证，如果天主允许仇敌随意肆虐，所有人在很久以前就灭亡了；但为了祂的仁慈，天主不容许他这样做。但如果人们把他们的爱转向天主，所有人无疑都会获得救恩，即使为了某些过失，他们可能看来要为了正义而被纠正。但现在大多数人成了天主的仇敌，邪恶者进入了他们的心，把造物主天主原本植入他们里面、好让他们对祂怀有的爱转向了他自己。至于其余那些一时看似警醒的人，仇敌以荣耀和辉煌的幻象显现，向他们许下某些伟大有力的事，导致他们的心思和心远离天主；然而，他被准许成就这些事，是出于某种公正的理由。}\par

% structure: p0037 source=h2 target=subsection reason=relative-heading-rank; base=h2

\subsection{第十八章 人的责任}

对此，阿桂拉回答说：\Quote{那么，如果邪恶者把自己装作光明的天使，\BibleRef{格后 11:14}向人应许比造物主自己更大的东西，人又怎么会是有过错的呢？}伯多禄回答说：\Quote{我认为，没有比这更不义的了；现在请听我告诉你这是多么不义。如果你的儿子，你悉心教养、精心培育，把他养大成人，他却对你忘恩负义，离开你去找另一个人，或许他看见那人更富有，就把本应给你的尊敬献给他，并且因为希望得到更大的利益而否认自己的出身，拒绝你作父亲的权利——这对你来说是正义的还是邪恶的呢？}阿桂拉回答说：\Quote{众人皆知，这是邪恶的。}伯多禄说：\Quote{如果你说这在人中间是邪恶的，那么在神身上岂不更是如此？祂超乎万有，配得人尊敬；我们不仅享受祂的恩惠，而且当我们还不存在时，是藉着祂的手段和能力我们才开始存在，如果我们愿意，我们将从祂那里得到永远在福乐中的恩赐！因此，为了使不信的人与信的人分开，虔诚的人与不虔诚的人分开，邪恶者被允许使用那些手段来考验每个人对真正父亲的感情。但如果确实有某个陌生的神，我们离开那创造我们、是我们的父亲和造物主的天主，而转投另一个，这是正当的吗？}阿桂拉说：\Quote{决不。}伯多禄说：\Quote{那么，我们怎能说邪恶者是我们犯罪的原因呢？因为这是天主容许的，为要叫那些被更大应许所诱惑、背弃了对真正父亲和造物主的本分的人，在审判之日被显明并定罪；而那些即使贫穷、遭患难仍持守信德和对自己的父亲的爱的人，将在祂的国度里享受天上的恩赐和不朽的尊荣。但我们将在其他时候更详细地阐述这些。此刻，我想知道西满之后做了什么。}\par

% structure: p0039 source=h2 target=subsection reason=relative-heading-rank; base=h2

\subsection{第十九章 辩论开始}

尼刻达回答说：\Quote{当他察觉我们已经识破他，我们彼此谈论了他的罪行之后，便离开了他，来到匝凯那里，把现在告诉你的那些事同样告诉了他。他非常和蔼地接待了我们，教导我们有关主耶稣基督的信德，并将我们登记在信友的行列中。}尼刻达说完后，匝凯（他之前出去了一会儿）进来，说：\Quote{伯多禄啊，时候到了，你该去进行辩论了；因为一大群人聚集在屋外的庭院里等候你，西满站在他们中间，由许多随从簇拥着。}伯多禄听了这话，吩咐我为祈祷而退下（因为我还没有从无知所犯的罪中洗净），然后对其他人说：\Quote{弟兄们，让我们祈祷，祈求天主因祂藉着基督不可言喻的怜悯，帮助我出去为那些由祂所造的人的救恩工作。}说完这话，他祈祷后，便走到屋外的庭院里，那里聚集了一大群人。当他看到众人都肃静地注视着他，而术士西满像一个旗手站在他们中间时，便这样开始了。\par

% structure: p0041 source=h2 target=subsection reason=relative-heading-rank; base=h2

\subsection{第二十章 天主的国和祂的义德}

愿平安归于你们众人，就是那些准备伸出右手接纳真理的人。因为凡是服从真理的人，似乎确实是在向天主施恩；但他们自己却从天主那里得到了祂最大恩赐的礼物，行走于祂正义的途径。因此，众人的首要职责是探究天主的正义和祂的王国\BibleRef{玛 6:33}：祂的正义，好叫我们被教导行事正直；祂的王国，好叫我们知道为劳苦和忍耐所定的赏报是什么；在这个王国里，确实有永生的美物赐予善人，但对于那些行事违背天主旨意的人，则根据每人的行为施以相应的惩罚。因此，你们当趁着还在这里——即在此生——的时候，确实知道天主的旨意，同时也还有机会去实行它。因为若有人在改正自己的行为之前，就想去探究他无法发现的事，这样的探究是愚蠢且无益的。因为时日短促，天主的审判将专注于行为，而非问题。所以，在一切之前，让我们探究这件事：我们必须做什么，或怎样做，才能配得获得永生。\par

% structure: p0043 source=h2 target=subsection reason=relative-heading-rank; base=h2

\subsection{正义是通往王国的道路。}

\Quote{因为如果我们用虚浮无用的问题占据此生的短暂时光，我们必将空虚无善工地走到天主面前，那时，如我所说，我们的行为将被带到审判中。因为万事各有其时间和地点。这里是行为的地点和时间；来世是赏报的地点和时间。因此，为避免因颠倒地点和时间的顺序而陷入困惑，让我们首先探究什么是天主的正义；好使我们像准备出发旅行的人一样，充满善工如同带足丰富的行粮，从而能够来到天主的王国，如同来到一座大城。因为对于思想正确的人来说，天主甚至透过祂所创造的世界的运作而显明，藉着祂造物的证据\BibleRef{罗 1:20}；因此，既然对天主不应有任何怀疑，我们现在只需探究祂的正义和祂的王国。但是，如果我们的心念建议我们在探究正义的行为之前，去探究隐密和隐藏的事，我们应当给自己一个理由：因为如果我们行为良好，我们将配得获得救恩；然后，纯洁干净地走向天主，我们将充满圣神，并且将知道一切隐密和隐藏的事，无需任何吹毛求疵的问题；然而现在，即使有人耗尽一生去探究这些事，他不仅不能发现它们，反而会陷入更大的错误，因为他没有首先通过正义的道路进入，并努力抵达生命的港湾。}\par

% structure: p0045 source=h2 target=subsection reason=relative-heading-rank; base=h2

\subsection{正义；它是什么？}

\Quote{因此我建议，应先探究祂的正义，使我们沿着正义前行，置身于真理之道，能够找到真先知，不是靠脚快，而是靠善工；并享有祂的引导，就无迷路的危险。因为如果在祂的引导下，我们配得进入那所渴望到达的城，则我们现在所探究的一切，都将亲眼看见，几乎成为万有的承继者。所以，应当明白：道路就是我们此生的历程；旅客是那些行善工的人；大门就是我们所谈论的真先知；城就是全能圣父所居住的王国，只有心地纯洁的人才能看见祂\BibleRef{玛 5:8}。因此，不要认为这旅程的辛劳是艰难的，因为它的终点将是安息。因为真先知本身也从太初、经过时间的历程，急忙趋向安息。因为祂时刻与我们同在；每当需要时，祂就显现并纠正我们，为把服从祂的人带入永生。因此，这是我的判断，也是真先知的意愿：应当首先探究正义，特别是对于那些自认认识天主的人。所以，如果有人提出什么更好的建议，让他发言；他发言后，也当倾听，但要耐心和安静：因为为了这个目的，起初在问候时，我曾为你们众人祈求平安。}\par

% structure: p0047 source=h2 target=subsection reason=relative-heading-rank; base=h2

\subsection{息孟拒绝和平。}

对此，\Person{息孟}回答说：\Quote{我们不需要你的和平；因为如果有了和平与和谐，我们就无法在发现真理方面取得任何进展。因为强盗和荡妇彼此有和平，每种邪恶都与其自身相合；如果我们以这样的观点聚在一起，即为了和平而赞同所说的一切，我们将对听众毫无益处；相反，我们将欺骗他们，并友好地分手。所以，不要祈求和平，而要祈求战争，因为战争是和平之母；如果你能，就消除谬误。不要寻求通过不公正的让步而获得的友谊；因为我希望你特别知道，当两人相争时，只有当一方被击败并倒下时，才会有和平。因此，尽你所能战斗，不要期望没有战争的和平，这是不可能的；如果你能实现和平，请告诉我们如何做到。}\par

% structure: p0049 source=h2 target=subsection reason=relative-heading-rank; base=h2

\subsection{伯多禄的解释。}

伯多禄回答说：\Quote{请你们全神贯注地听，诸位，我们所说的话。我们假设这个世界是一个大平原，来自两个彼此敌对的王国的两位将军被派去作战；并假设好国王的将军提出这样的忠告：双方军队应不流血地服从于更善国王的权威，这样所有人都不受危险而安全；但敌对的将军却说：不，我们必须战斗；这样统治的就不是那有德者，而是那更强壮者，与那些幸存的人；——请问你们会选择哪一个？我毫不怀疑你们会向更好的国王伸出援手，以保全所有人。我现在并不希望，正如西满所说，为了和平而赞同那些错误的话，而是希望安静而有秩序地寻求真理。}\par

% structure: p0051 source=h2 target=subsection reason=relative-heading-rank; base=h2

\subsection{第二十五章 讨论应遵循的原则}

\Quote{有些人在辩论中，当他们意识到自己的错误被驳倒时，为了自圆其说而立即开始制造混乱、挑起争端，以免众人看穿他们的败局；因此我常常恳求，要以极大的耐心和安静进行讨论，使得如果偶然有什么说得不妥，可以允许回头再讲，更清楚地解释。因为有时一件事可能这样说却那样听，或者说得太晦涩，或者听得不够仔细；为此，我希望我们的谈话能够耐心进行，这样既不会一人夺过另一人的话，也不会因反驳者的不合时宜的插话而打断另一人的发言；我们不应怀有吹毛求疵的欲望，而应允许像我说过的，对讲得不够清楚的地方重新讲述，以便通过最公正的审查使真理的认识更加清晰。因为我们应当知道，如果一个人被真理征服，被征服的不是他，而是他里面的无知，那是最坏的恶魔；因此，能将无知驱逐出去的人，就赢得了救赎的棕榈枝。因为我们的目的是使听者受益，不是为了一味地获胜，而是为了承认真理而甘愿被良好地征服。因为如果我们的言谈是由寻求真理的渴望所驱动，即使因人性软弱而说得不够完美，天主也会以其不可言喻的良善，在听者的心中暗暗地补足那些欠缺。因为祂是公义的；祂按照每人的意向，使一些人容易找到他们所寻求的，而使另一些人连眼前的也看不明白。既然天主的道路是和平的道路，就让我们以和平寻求属于天主的事物。如果有人对此有什么可反驳的，请他说出来；但如果没有谁想回答，我就开始发言，由我自己来提出别人可能反对我的论点，并加以驳斥。}\par

% structure: p0053 source=h2 target=subsection reason=relative-heading-rank; base=h2

\subsection{第二十六章 西满的插话}

于是，当伯多禄开始继续他的讲论时，西满打断他说：\Quote{你为什么急于想说你想说的话？我明白你的花招。你希望提出那些你已经精心研究过解释的问题，以便在无知的群众面前显得说得很好；但我不会让你得逞。现在既然你勇敢地承诺要回答任何人提出的任何问题，那就请你首先回答我。}随后伯多禄说：\Quote{我准备好了，只要我们的讨论能在和平中进行。}西满又说：\Quote{你这傻瓜，难道你没有看到吗？你为和平辩护，是在违背你的师傅，你所提议的并不适合那承诺推翻无知的人？或者，如果你向听众求和平是对的，那么你的师傅说\BibleQuote{我来不是给地上带来和平，而是带来刀剑。}\BibleRef{玛 10:34}就是错的。因为要么你说得对，他不对；要么如果你师傅说得对，那你就完全不对：因为你不明白你的话与他的话相违背，而他正是你自称其门徒的那位。}\par

% structure: p0055 source=h2 target=subsection reason=relative-heading-rank; base=h2

\subsection{第二十七章 问答}

伯多禄说：\Quote{那派遣我的并没有错，错在于把刀剑送到地上，而我在向听众求和平也并没有违背祂。但是你们既笨拙又轻率地指责你们所不明白的事：因为你们听说了师傅来到地上不是送和平；但祂也说：\BibleQuote{缔造和平的人是有福的，因为他们要被称为天主的子女。}\BibleRef{玛 5:9}你们却没有听到。因此，当我推荐和平时，我的意见与师傅的并无不同，祂曾把真福赐给守和平的人。}西满说：\Quote{伯多禄啊，你急于为你的师傅辩护，却给他带来了更严重的指控：如果他自己没有带来和平，却又命令别人要持守和平。那么，他另一句话\BibleQuote{门徒能像师傅一样就够了}\BibleRef{玛 10:25}又怎么一致呢？}\par

% structure: p0057 source=h2 target=subsection reason=relative-heading-rank; base=h2

\subsection{第二十八章 基督教导的一致性}

\Quote{我们的主，祂是真正的先知，始终记念自己，既不与自己矛盾，也不命令我们做任何与祂自己所行不同的事。因为祂说：\InnerQuote{我来不是为地上送和平，而是送刀剑；从此以后，你们将要看见父亲与儿子分离，儿子与父亲分离，丈夫与妻子、妻子与丈夫分离，母亲与女儿、女儿与母亲分离，兄弟与兄弟分离，公公与儿媳分离，朋友与朋友分离。} 这一切都包含和平的教导；我将告诉你们如何理解。在祂传道之初，如同愿意邀请并引导众人走向救恩，促使他们耐心忍受劳苦和考验，祂祝福了穷人，并应许他们将因忍受贫穷而获得天国，为的是在这种希望的感召下，他们能平静地承受贫穷的重担，摒弃贪吝；因为贪吝是最大、最有害的罪恶之一。祂也应许饥渴的人将得享永生公义的饱足，为的是他们能耐心忍受贫穷，不致因此而从事任何不义的工作。同样，祂也说心里洁净的人是有福的，他们将因此看见天主，为的是每个渴望如此大善的人都能远离邪恶和污秽的思念。}\par

% structure: p0059 source=h2 target=subsection reason=relative-heading-rank; base=h2

\subsection{第二十九章 和平与纷争}

\Quote{因此，我们的主邀请祂的门徒忍耐，并教导他们要以忍耐的劳苦来保存和平的祝福。但另一方面，祂为那些生活在财富和奢华之中、不周济穷人的人而哀恸，证明他们必须交账，因为他们不怜悯邻舍，即使邻舍处于贫困之中，他们本应爱邻如己。通过这样的言论，祂使一些人听从祂，而另一些人则成为敌对者。因此，祂嘱咐那些相信并服从的人彼此保持和平，对他们说：\BibleQuote{缔造和平的人是有福的，因为他们要称为天主的子女。}\BibleRef{玛 5:9}但对于那些不仅不信、反而反对祂教导的人，祂宣布了话语和辩驳的战争，并说：\InnerQuote{从此以后，你们将要看见儿子与父亲分离，丈夫与妻子分离，女儿与母亲分离，兄弟与兄弟分离，儿媳与婆婆分离，人的仇敌就是自己家里的人。}因为每当一个家庭中开始出现信徒与不信者之间的分歧时，就必然有争斗：不信者对抗信仰，而信徒则驳斥他们旧的错误和罪恶的恶习。}\par

% structure: p0061 source=h2 target=subsection reason=relative-heading-rank; base=h2

\subsection{第三十章 和平归于和平之子}

\Quote{同样，在祂教导的最后阶段，祂向经师和法利塞人宣战，指责他们行为邪恶、教理不正，并隐藏了他们从梅瑟传下来的知识的钥匙，这把钥匙本可开启天国之门。但是，当我们主派遣我们出去传道时，祂命令我们，无论进入哪座城或哪家，都要说：\InnerQuote{愿这平安归于这家。}\InnerQuote{如果那里有和平之子，你们的平安就会临到他；如果没有，你们的平安就会回归你们。}并且，从那个家或城出来时，要连沾在你们脚上的尘土也抖落掉。\InnerQuote{在审判之日，索多玛和哈摩辣地所受的惩罚，比那城或那家还要容易忍受。}祂确实命令这样做，但首先要在那城或那家宣讲真理之言，好让接受真理信仰的人成为和平之子和天主的子女；而那些不愿接受的人，则被证实为和平与天主的仇敌。}\par

% structure: p0063 source=h2 target=subsection reason=relative-heading-rank; base=h2

\subsection{第三十一章 和平与战争}

\Quote{因此，我们遵守我们主的命令，首先向听众提供平安，使救恩之道得以不受干扰地显明。但如果有人不接受和平的话语，也不服从真理，我们就知道如何用言语的战争对付他，通过驳斥他的无知、严厉责备他的罪过。因此，我们必须先提供和平，这样如果有谁是和平之子，我们的平安就会临到他；但如果有人自绝于和平，我们的平安就会回归我们自己。我们不像你所说的，与恶人达成和平协议，因为如果那样，我们本会立即与你握手；而是为了通过我们安静而耐心的讨论，使听众更容易辨别哪一个是真实的言语。但如果你与自己观点相左、不一致，你怎能站立得住？一个分裂的人必然跌倒；\BibleQuote{凡自相分裂的国，必不能存立。}\BibleRef{玛 12:25}如果你对此有话要说，请说。}\par

% structure: p0065 source=h2 target=subsection reason=relative-heading-rank; base=h2

\subsection{第三十二章 西满的挑战}

西满说：\Quote{我对你们的愚昧感到惊讶。因为你们这样提出你们主的话，仿佛已确定祂是先知；而我很容易证明祂常常自相矛盾。简言之，我要用你们自己提出的话来驳斥你们。因为你们说，祂说凡自相分裂的国或城必不能存立；而其他地方你们说，祂说祂将送来刀剑，为要分开同一家中的人，使儿子与父亲分裂，女儿与母亲分裂，兄弟与兄弟分裂；以至如果一家有五个人，三人要与两人分裂，两人要与三人分裂。\BibleRef{路 12:51}如果你能回答，就回答。}\par

% structure: p0067 source=h2 target=subsection reason=relative-heading-rank; base=h2

\subsection{第三十三章 权威}

接着\Person{伯多禄}说：\Quote{西满，不要轻率地对你所不明白的事提出异议。首先，我要回答你的断言，即我陈述了我主的言语，并从中解决尚有疑问的事。我们的主，当祂派遣我们宗徒去宣讲时，嘱咐我们要教导万民\BibleRef{玛 28:19}那些交托给我们的事。因此，我们不能照祂亲自所说的那样说出那些话。因为我们的使命不是要\Emph{说}，而是要\Emph{教导}那些事，并从它们中表明每件事是如何建立在真理上的。再者，我们也不允许说任何自己的话。因为我们是被派遣的；被派遣的人必然按照所吩咐的传递信息，并陈述派遣者的旨意。因为如果我所说的与派遣我的主所吩咐的任何不同，我就成了一个假宗徒，没有说我所受命要说的，而是说自己认为好的。这样做的人显然想表明自己比派遣他的那一位更好，无疑是个叛徒。相反，如果他遵守所吩咐的事，并提出最清楚的论证，那么他就显得是在完成宗徒的工作；我正是努力履行这一点而令你不悦。因此，不要因为我提出派遣我之主的言语而责备我。但如果其中有任何说得不恰当的地方，你有自由反驳我；但这绝不可能，因为祂是一位先知，不能自相矛盾。但如果你不认为祂是先知，那么首先要查考此事。}\par

% structure: p0069 source=h2 target=subsection reason=relative-heading-rank; base=h2

\subsection{第三十四章 论证的顺序}

接着\Person{西满}说：\Quote{我不需要从你这里学习这一点，而是这些事如何彼此一致。因为如果他被证明是不一致的，那么他同时也被证明不是先知。} 接着\Person{伯多禄}说：\Quote{但如果我先证明祂是先知，那么看似不一致的事就不是不一致。因为没有人能仅仅通过一致性被证明是先知，因为许多人可以达到这一点；但如果一致性不能证明是先知，那么不一致更不能。因为，有很多事在有些人看来不一致，但在更深入的研究中却有一致性；同样，另一些事看似一致，但经过更仔细的讨论，却发现是不一致的；因此，我认为没有比首先确定祂是否是那位说了看似不一致之事的先知更好的判断方法了。因为显然，如果祂被证实是先知，那么那些看似矛盾的事必然有一致性，只是被误解了。因此，关于这些事，将适当地要求证据。因为我们宗徒被派遣来解释那位派遣我们之主的言论并确认祂的判断；但我们没有被授权说任何自己的话，而是要如我所说的，阐明祂话语的真理。}\par

% structure: p0071 source=h2 target=subsection reason=relative-heading-rank; base=h2

\subsection{第三十五章 错误如何不能与真理并存}

接着\Person{西满}说：\Quote{因此，请教导我们，那引起分裂的人，而这些分裂使那些分裂的人跌倒，如何可能是善的，或是为了人的救恩而来的。} 接着\Person{伯多禄}说：\Quote{我要告诉你，我们的主说过，凡自相分争的国和家必不能站立；而祂亲自做了这事，看它如何有益于救恩。祂以真理之言确实分裂了建立在错误之上的世界之国及其中的每一个家，好使错误跌倒，真理掌权。但如果发生在某个家里，错误被某人引入而分裂了真理，那么，在错误立足之处，真理必然不能站立。} 接着\Person{西满}说：\Quote{但你的主是分裂错误还是真理，却是不确定的。} 接着\Person{伯多禄}说：\Quote{那属于另一个问题；但如果你同意凡分裂的必然跌倒，那么我所要做的就是——只要你愿意和平地听——表明我们的\Person{耶稣}如何藉教导真理而分裂并驱散了错误。}\par

% structure: p0073 source=h2 target=subsection reason=relative-heading-rank; base=h2

\subsection{第三十六章 争论}

接着\Person{西满}说：\Quote{不要一再重复你的和平之言，而是简要地阐述你所想或所信的。} \Person{伯多禄}回答说：\Quote{你为什么害怕常常听到和平？难道你不知道和平是法律的圆满吗？因为战争和纷争源自罪；没有罪的地方，就有灵魂的和平；而有和平的地方，在辩论中就找到真理，在行为中就找到正义。} 接着\Person{西满}说：\Quote{在我看来，你似乎无法表达你的想法。} 接着\Person{伯多禄}说：\Quote{我将发言，但按照我自己的判断，不受你伎俩的约束。因为我希望那有益和有利的事能让所有人知道，因此我将尽快地陈述。只有一位天主；祂是世界的创造者，公义的审判者，在各样的时候按各人的行为报应各人。但现在为了确认这些事，我知道可以说出无数的话。}\par

% structure: p0075 source=h2 target=subsection reason=relative-heading-rank; base=h2

\subsection{第三十七章 \Person{西满}的诡诈}

于是西满说：\Quote{我确实佩服你思维的敏捷，但我并不接受你信仰的错谬。因为你已经明智地预见到自己可能被反驳；你甚至礼貌地承认，为了证明这些事，将需要数不清的言辞，因为没有人同意你信仰的宣认。简而言之，至于有一位天主，世界是他的造化，谁能接受这个道理？我认为，无论是外邦人，即使是无学问之人，还是哲学家，肯定没有一人；甚至连最粗鲁、最悲惨的犹太人，以及我自己——我深谙他们的法律——也不会接受。}于是伯多祿说：\Quote{撇开那些不在这里之人的意见，面对面告诉我们你自己是怎么想的。}于是西满说：\Quote{我能说出我的真实想法；但我之所以不愿这样做，是因为考虑到：如果我说出既不合你心意、也不被这无知的乌合之众视为正确的话，那么你，作为理屈词穷者，会立刻堵住耳朵，以免被亵渎所玷污（真的），并因无法回答而逃走；而无知的大众会赞同你，把你当作教导他们之间普遍接受之事的人来拥抱；他们却会咒骂我，因为我说出新颖闻所未闻的事，并将我的谬误灌输给别人。}\par

% structure: p0077 source=h2 target=subsection reason=relative-heading-rank; base=h2

\subsection{第三十八章 西满的信仰。}

于是伯多祿：\Quote{难道你不是像你指责我们那样在长篇大论地作序，因为你没有真理可提出吗？因为如果你有真理，就开门见山地说吧，如果你如此有信心。而且，如果你的话让任何一位听众不悦，他可以离开；而留下的人将被你的主张迫使赞同真理。那么，开始阐述你认为正确的事吧。}于是西满说：\Quote{我说有许多神；但有一位天主是不可理解的，不为众人所知，他是这些众神的天主。}于是伯多祿答道：\Quote{你所断言的那位不可理解、不为众人所知的这位天主，你能从被认作权威的犹太人的圣经中，或从我们全然不知的其他经书中，或从希腊作者的作品中，或从你自己的著作中证明他的存在吗？当然，你可以随意引用任何著作，但必须首先表明它们是先知性的；如此它们的权威才被毫无疑问地接受。}\par

% structure: p0079 source=h2 target=subsection reason=relative-heading-rank; base=h2

\subsection{第三十九章 多神论的论证。}

于是西满说：\Quote{我将只引用犹太人的法律中的主张。因为凡关心宗教的人都知道，这部法律具有普世权威，但每个人都是按照自己的判断来理解这部法律。因为它是由创造世界的那位如此书写，以至于事物的信仰依赖于此。因此，无论是谁想提出真理，或是谁想提出谬误，没有这部法律，任何主张都不会被接受。因此，既然我的知识完全符合法律，我正确地宣告有许多神，其中一位比其余更崇高，且不可理解，他即是万神的天主。但有多位神，法律本身告诉了我。首先，经上有一处，以蛇的形象对第一个女人厄娃说话：\BibleQuote{你们将如同神，}\BibleRef{创 3:5} 即如同造人的那一位；他们尝了果子之后，天主亲自作证，向其余的神说：\BibleQuote{看，亚当成了我们中的一个；}\BibleRef{创 3:22} 因此，显然有许多神参与了造人。此外，起初天主对别的神说：\BibleQuote{让我们照我们的肖像和样式造人；}\BibleRef{创 1:26} 他又说：\BibleQuote{让我们驱逐他；}\BibleRef{创 3:22} 又说：\BibleQuote{来，我们下去，混乱他们的语言；}\BibleRef{创 11:7} 这一切都表明有许多神。而经上也写着：\BibleQuote{不可诅咒神，也不可诅咒你百姓的首领；}\BibleRef{出 22:28} 这段经文又写着：\BibleQuote{上主独自领导他们，没有外神与他们同在，}\BibleRef{申 32:12} 这显明有许多神。从法律中还可以举出许多其他见证，不仅晦涩，而且清晰，教导有许多神。其中一位被抽签选出，作犹太人的神。但我说的不是他，而是那位也是他的天主的那位，连犹太人自己也不知道他。因为他不是他们的天主，而是认识他之人的天主。}\par

% structure: p0081 source=h2 target=subsection reason=relative-heading-rank; base=h2

\subsection{第四十章 伯多祿的回答。}

伯多祿听了这话，回答说：\Quote{不要害怕，西满；因为，看，我们既没有堵住耳朵，也没有逃走；而是用真理之言回答你所说的那些虚假之事，首先肯定：只有一位天主，就是犹太人的天主，他是唯一的天主，天地万物的创造者，也是你称为诸神的那一切的天主。因此，如果我向你证明，没有比他更崇高的，他自身超越一切，你就会承认你的谬误超越一切。}于是西满说：\Quote{确实，即使我不愿意承认，在场的听众难道不会指责我不愿意承认真实的事吗？}\par

% structure: p0083 source=h2 target=subsection reason=relative-heading-rank; base=h2

\subsection{第四十一章 回答（续）。}

\Quote{那麼，請聽，}伯多祿說，\Quote{首先，你們要知道，即使有許多神明，正如你們所說的，他們也要服從猶太人的天主，沒有人能與祂相等，沒有人能比祂更大；因為經上記載，梅瑟先知曾這樣對猶太人說：\InnerQuote{上主你們的天主是萬神之神，萬主之主，偉大的天主。}\BibleRef{申 10:17} 因此，雖然有許多被稱為神的，但猶太人的天主卻唯獨被稱為萬神之神。因為並非每一個被稱為神的必然是神。事實上，連梅瑟也被稱為法郎的神，\BibleRef{出 7:1} 而我們確知他只不過是人；審判官也被稱為神，而他們顯然是有死的。外邦人的偶像也被稱為神，而我們都知道它們不是；但這乃是加於惡人的刑罰，因為他們不願承認真天主，便將自己所遇到的任何形像或偶像當作神。因為他們拒絕接受那一位——如我所說的——是萬有的天主的知識，故此准許他們以那些不能為其敬拜者做任何事者為神。因為無論是死的偶像或是活的受造物，能給人什麼呢？既然萬有的權能都在於那一位。}\par

% structure: p0085 source=h2 target=subsection reason=relative-heading-rank; base=h2

\subsection{第四十二章　護守天使}

因此，\Emph{神}這個名稱有三種應用方式：或者接受此名者是真神，或者他是那真神的僕人；為了差遣者的榮耀，使祂的權威得以完全，被差遣者被稱為差遣者的名字，正如天使們常被對待那樣：因為當他們向人顯現時，如果那人是一個有智慧的人，他會問顯現者的名字，以便同時認識被差遣者的尊榮和差遣者的權威。因為每一個民族都有一位天使，天主將治理該民族的責任交託給他；當這樣的一位顯現時，雖然被他所管理的人認為並稱為神，然而若問他，他卻不這樣自稱。因為至高無上的天主，唯獨掌握萬有的權能，將地上的萬國分為七十二份，並派天使為首領治理它們。但在總領天使中最偉大的那一位，被委託治理那些在所有民族之先接受了至高天主的敬拜和知識的民族。然而，聖善的人，如我們所說，對惡人而言也被視為神，因為他們獲得了掌管他們生死之權，正如我們上面提到梅瑟和審判官時所說的。因此，經上也有關乎他們的記載：\Quote{不可褻瀆神，也不可詛咒你百姓的首領。}\BibleRef{出 22:28} 這樣，各個民族的首領被稱為神。但基督是首領之神，祂是眾人的審判者。因此，無論天使、人，或任何受造物，都不能真正是神，因為他們都處於權下，是受造的、可變的：天使，因為他們先前不存在而後存在；人，因為他們是有死的；一切受造物，因為它們能解體，只要造它們的那一位願意解體。所以，唯獨祂是真天主，祂不僅自己活著，也賜生命給他人，並能隨意收回生命。\par

% structure: p0087 source=h2 target=subsection reason=relative-heading-rank; base=h2

\subsection{第四十三章　除雅威外沒有別的神}

因此，聖經以猶太人的天主之名呼喊說：\Quote{看，看，我是天主，除我以外沒有別的神；我殺死，也使我生存；我打擊，也醫治；誰也不能從我手中救出。}\BibleRef{申 32:39} 請看，聖經以一種不可言喻的力量，反對那些將來會肯定除了猶太人的天主外在天地間還有別的神的人的錯誤，如此斷定：\Quote{上主你們的天主是唯一的天主，在天上地下；除祂以外沒有別的神。}\BibleRef{申 4:39} 那麼，你們怎敢說除了猶太人的天主以外還有別的神呢？聖經又說：\Quote{看，上主你們的天主擁有天，及天上的天，地和地上的一切；然而我揀選了你們的祖先，為愛他們，並在他們之後愛你們。}\BibleRef{申 10:14} 這樣，聖經從各方面支持了這個判斷：創造世界的天主是真實而唯一的天主。\par

% structure: p0089 source=h2 target=subsection reason=relative-heading-rank; base=h2

\subsection{第四十四章　蛇是多神論的始作俑者}

但是即使像我們所說的，還有其他被稱為神的，他們都在猶太人的天主的權下；因為聖經這樣對猶太人說：\Quote{上主我們的天主，祂是萬神之神，萬主之主。}\BibleRef{申 10:17} 聖經也命令唯獨敬拜祂，說：\Quote{你要敬拜上主你的天主，惟獨事奉祂；}又說：\Quote{以色列，請聽：上主你的天主是唯一的天主。}\BibleRef{申 6:4} 是的，連那些充滿天主聖神、並蒙祂仁慈的露珠滋潤的聖人們也呼喊說：\Quote{在眾神中誰能像祢？上主啊，誰能像祢？}又說：\Quote{誰是神，除了上主？誰是神，除了我們的上主？}因此梅瑟看見百姓逐漸前進，便逐步引導他們理解獨一真主的王權與信仰，正如他以下的話所說：\Quote{你們不要提別神的名字；}無疑地，他記住了那首先說出\Emph{神}這個名字的蛇所受的懲罰。因為蛇被判以塵土為食，並被認為配得那樣的食物，因為它首先將\Emph{神}的名稱引入世界。但如果你們也想引入許多神，請當心不要分受蛇的厄運。\par

% structure: p0091 source=h2 target=subsection reason=relative-heading-rank; base=h2

\subsection{第四十五章　多神論無可推諉}

因为这事实你务要确信：你必不能使我们成为这尝试的同谋；我们也不容许自己被你所欺骗。因为在审判时，如果我们说被你欺骗了，这对我们毫无助益；就如同第一个女人不能因为不幸相信了蛇而得赦免，反而因坏的信德而被判处死亡。为此，梅瑟也向百姓赞美独一天主的信德，说：\BibleQuote{你要小心，不要被诱惑离开上主你的天主。}\BibleRef{申 8:11}请注意，他使用了与第一个女人在为自己辩护时相同的词，她说她被诱惑了；但这对她毫无益处。此外，即使有一位真先知兴起，能行神迹奇事，却想说服我们敬拜犹太人的天主以外的其他神明，我们绝不能相信他。因为神圣的法律这样教导我们，藉着传统更纯洁地传递一条隐秘的诫命，因为如此说：\BibleQuote{如果你们中有先知，或做梦的人兴起，给你神迹或奇事，这些神迹奇事应验了，他对你说：我们去敬拜你们不认识的外邦神吧。你不要听那先知的话，也不要听那做梦者的梦，因为上主你们的天主是在考验你们，要知道你们是否爱慕上主你们的天主。}\BibleRef{申 13:1}\par

% structure: p0093 source=h2 target=subsection reason=relative-heading-rank; base=h2

\subsection{第四十六章 基督承认犹太人的天主。}

\Quote{因此，我们的主也行神迹奇事，宣讲犹太人的天主；所以我们相信他所宣讲的是正确的。至于你，即使你真是一位先知，并如你所应许的那样行神迹奇事，但如果你宣扬那位真天主以外的其他神明，那就显然你是被兴起作为天主百姓的考验；因此你绝不能得到相信。因为只有祂是真天主，祂是犹太人的天主；为此，我们的主耶稣基督并没有教导他们必须寻求天主，因为他们已经很认识祂，而是要他们寻求祂的国和祂的义德，\BibleRef{玛 6:33}而这义德，经师和法利塞人领受了知识的钥匙，却没有关闭，反而关闭了外面。\BibleRef{路 11:52}因为如果他们不认识真天主，祂绝不会离开这最重要之事的认识，反而责备他们那些微小的事，如放宽衣繸、在宴席上争坐首位、在十字路口站着祈祷等类似的事；这些事与不认识天主这重大罪过相比，无疑显得微小而不重要。}\par

% structure: p0095 source=h2 target=subsection reason=relative-heading-rank; base=h2

\subsection{第四十七章 西满的刁难。}

\Quote{我要用你师父的话来驳斥你，因为他也向众人介绍了一位已知的天主。因为虽然亚当认识创造他的天主，即世界的造物主；哈诺客认识祂，因为他被祂提升；诺厄认识祂，因为他受命建造方舟；虽然亚巴郎、依撒格、雅各伯、梅瑟以及所有民族、万国都认识世界的造物主，并承认祂是天主，但你的耶稣在众圣祖之后出现，却说：\BibleQuote{除了父，没有人认识子；除了子和子所愿意启示的人，没有人认识父。}因此，即使你的耶稣也承认另有位不可理解、无人知晓的天主。}\par

% structure: p0097 source=h2 target=subsection reason=relative-heading-rank; base=h2

\subsection{第四十八章 伯多禄的回答。}

\Quote{你没有察觉到你在自相矛盾。因为如果我们的耶稣也认识你所谓的未知天主，那么祂就不是唯独被你认识。是的，如果我们的耶稣认识祂，那么那位预言耶稣将要到来的梅瑟，当然也不会不认识祂。因为梅瑟是先知；预言子的那位，一定认识父。因为如果子有自由将父启示给祂所愿意的人，那么从起初就与父同在且贯穿万代的子，既已启示父给梅瑟，也同样启示给其他先知；但如果是这样，显然父对他们都不是未知的。然而，父又怎能启示给你这不信子的人呢？因为除了子和子所愿意启示的人，没有人认识父。但子启示父给那些尊重子如同尊重父的人。}\BibleRef{若 5:23}\par

% structure: p0099 source=h2 target=subsection reason=relative-heading-rank; base=h2

\subsection{第四十九章 至高之光。}

\Quote{你记得你说过天主有一个儿子，这是对天主的冒犯；因为祂怎么可能有儿子呢？除非祂像人或动物一样受情欲支配。但现在没有时间指出你极深的愚蠢，因为我急于要阐述关于至高之光的无限性，所以现在请听。我的意见是：有一种无限且不可言喻的光的力量，其伟大被认为不可理解，甚至世界的造物主、立法者梅瑟和你的师父耶稣都不认识这种力量。}\par

% structure: p0101 source=h2 target=subsection reason=relative-heading-rank; base=h2

\subsection{第五十章 西满的狂妄。}

然后\Person{伯多禄}说：\Quote{难道你们不认为这是疯狂吗，竟有人胆敢断言存在另一位天主，而不是万有的天主；并且说他假设有一种权能，还在他自己对自己所说的话没有把握之前，就冒昧向别人肯定这一点？有谁如此鲁莽，竟相信你的话，他看出你自己都怀疑呢，并承认有一种权能，不为创造主天主、梅瑟、先知、法律，甚至我们的主耶稣所知，这权能如此良善，以至于它不愿向任何人启示自己，只向一个人，而且是这样一个人——像你一样的人！再者，如果那是一个新的权能，为什么它不赐给我们某种新的感官，以补充我们拥有的那五种，好使我们借着它赐给我们的这新感官，能够接受并理解这新的权能本身？或者，如果它不能赐给我们这样的感官，又怎能赐给你呢？或者，如果它向你启示了自己，为什么不同样启示给我们呢？但如果你独自明白连先知都无法看清或理解的事，那么来吧，告诉我们我们每个人现在在想什么；因为如果你里面有这样的灵，使你知道那些高于诸天、无人知晓、无人能理解的事，你就更容易知道地上人的思想。但如果你不能知道我们这些站在这里的人的思想，你怎么能说你知道那些你断言无人知道的事呢？}\par

% structure: p0103 source=h2 target=subsection reason=relative-heading-rank; base=h2

\subsection{第五十一章 第六种感官}

\Quote{但请相信我，你永远无法知道什么是光，除非你从光本身接受了视觉和理解；其他事物也是如此。因此，你接受了理解力，便在想象中建构出某种更伟大、更崇高的东西，仿佛在做梦，但你所有的提示都来源于那五种感官，而你对它们的赐予者却忘恩负义。但请确信这一点：除非你找到一种超越我们大家都享有的那五种感官的新感官，否则你不能断言一位新天主的存在。} 然后\Person{西满}回答说：\Quote{既然所有存在的事物都与那五种感官相符，那比一切都更卓越的权能不能增添任何新事物。} 然后\Person{伯多禄}说：\Quote{这是假的；因为还有第六种感官，即预知：因为那五种感官能获取知识，但第六种是预知：先知们就拥有它。那么，你怎么能认识一位无人所知的天主呢？你并不认识预知的感官，即预知的能力。} 然后\Person{西满}开始说：\Quote{我所说的这个权能，不可理解且比一切都更卓越，是的，甚至比创造世界的天主更卓越，任何天使、魔鬼、犹太人，甚至任何依靠创造主天主而存在的受造物都不曾认识它。那么，创造主的法律怎能教导创造主自己都不知道的事呢？既然法律本身也不知道它，又怎能教导它呢？}\par

% structure: p0105 source=h2 target=subsection reason=relative-heading-rank; base=h2

\subsection{第五十二章 \OriginalTerm{归谬法}{Reductio Ad Absurdum}}

然后\Person{伯多禄}说：\Quote{我很奇怪，你怎么能比法律本身所能知道或教导的，从法律中学习到更多；并且你说你从法律中为你所乐于断言的那些事提出证明，而你又宣称法律本身和那赐下法律者——即世界的创造主——都不知道你所说的那些事！但我也对此感到奇怪，你独自知道这些事，怎么现在和我们大家站在一起，被这小小庭院的界限所限制？} 然后\Person{西满}，看见\Person{伯多禄}和众人在笑，说：\Quote{你在笑吗，伯多禄，当如此重大崇高的问题正在讨论之时？} 然后\Person{伯多禄}说：\Quote{不要发怒，西满，因为我们不过是在遵守承诺：因为我们既没有像你说的那样塞住耳朵，也没有一听到你提出你那些不可言说之事就逃跑；我们甚至没有离开这个地方。的确，你甚至没有提出任何与真理相似的事，这些事可能会在一定程度上吓唬我们。不过，无论如何，请向我们解释这句话的含义：你怎么从法律中认识了一位法律本身所不知道、而且那赐下法律者也对其无知的天主？} 然后\Person{西满}说：\Quote{如果你们笑够了，我将用清晰的断言来证明。} 然后\Person{伯多禄}说：\Quote{我当然会停止，好让我从你那里学习，你怎么从法律中知道了既非法律本身，也非法律的天主自己所知道的事。}\par

% structure: p0107 source=h2 target=subsection reason=relative-heading-rank; base=h2

\subsection{第五十三章 西满的亵渎}

然后\Person{西满}说：\Quote{请听：这是众人皆知的，并且以无法解释的方式确定，有一位天主，他比万有更美善，万有都从他开始；因此，所有在他之后的事物都必然隶属于他，作为万有的首领和最卓越者。因此，当我根据法律所教导的，确知那创造世界的天主在许多方面是软弱的，而软弱与一位完美的天主完全不相容，并且我看见他并不完美时，我就必然断定有另一位完美的天主。因为这位天主，正如我所说，根据法律的文字所教导的，在许多事上显示出软弱。首先，因为他所塑造的人不能保持他本意要让他成为的那样；而且，那给第一个人颁布法律，说他可以吃乐园中所有树上的果子，但不可触碰知识树，如果他吃了就必定死的人，不能是良善的。因为他为什么禁止他吃，禁止他知道善恶，好使他知道了可以避恶择善呢？但这一点他却不允许；因为他违反命令吃了，发现了什么是善，并为了尊荣而学会了遮盖自己的裸体（因为他觉得赤身站在创造主面前是不体面的），他就将死亡判处给那学会了向天主致敬的人，并且诅咒那向他显示这些事的蛇。但确实，如果人将因此受到伤害，他为什么要把伤害的起因放在乐园里呢？但如果他放在乐园里的是好的，那么一位良善者就不应该阻止别人得到这善。}\par

% structure: p0109 source=h2 target=subsection reason=relative-heading-rank; base=h2

\subsection{第五十四章 西蒙如何从法律中学到法律所不教导的}

\Quote{因此，既然按法律所载，那位造了人和世界的是不完美的，我们无疑就明白了，还有另一位是完美的。因为必然有一位超越万有的，万有因祂而各安其位。所以我也知道，必须有一位比那位赐下法律的不完美的天主更仁慈、更有能力，通过比较不完美而理解完美的，我甚至从圣经中明白了那在那里没有提及的天主。这样，哦，伯多禄，我就能从法律中学到法律所不知道的。但即使法律没有给出暗示，让人推断出造世界的天主是不完美的，我仍然可以从这世界上所行而未改正的恶事中推断出：要么造物主无能为力（若他不能改正做错了的事），要么若他不愿除去恶事，那他自己是恶的；若他既不能也不愿，那他既非有能力，也非善。由此不得不结论说，另有一位天主更加卓越，更有能力。如果你对此有话要说，请说。}\par

% structure: p0111 source=h2 target=subsection reason=relative-heading-rank; base=h2

\subsection{第五十五章 西蒙的异议被用来反对他自己}

伯多禄回答说：\Quote{哦，西蒙，那些没有在有学识的导师指导下阅读法律，却自以为是导师，以为能理解法律而那位从导师学习的人没有给他们解释的人，惯常对天主生出如此荒谬的想法。不过现在，为了让我们也似乎按你的理解来遵循法律书：既然你说创造世界的天主被证明既是无能又是邪恶的，你怎么看不到你那自称超越万有的能力也失败并遭受同样的指控呢？因为可以对它说同样的话：要么它无能为力，因为它不改正这里做错的事；要么它能而愿，则它是恶的；要么它既不能也不愿，则它既无能又不完美。因此，你那新的能力不仅同样受指控，甚至更糟，若此外它还被相信存在而实际上并不存在。因为创造世界的那位，祂存在由祂创造世界的行为本身彰显，正如你本人也承认。但你说只有你知道的这个能力，却没有给出任何迹象，让我们至少能感知它是存在并存留的。}\par

% structure: p0113 source=h2 target=subsection reason=relative-heading-rank; base=h2

\subsection{第五十六章 没有天主在创造者之上}

\Quote{那么，我们离弃天主——在祂的世界中我们生活并享受生命所需的一切——而去追随一个我们不仅得不到好处，甚至不能知道祂是否存在的神，这是何等行为？的确祂并不存在。因为无论你称祂为光，比我们所见的光更亮，你还是从创造世界的那位借用了那名；或者你说祂是超越万有的实体，你从祂那里借用了这个观念并加以夸大。无论你提到心智、良善、生命或其他任何东西，你都是从祂那里借用了词汇。既然你关于所说的那个能力，不仅在理解上，甚至在命名上都没有任何新的东西，你怎能引入一位新的天主，你甚至不能为祂找到一个新的名字？因为不仅创造世界的那位被称为能力，连祂光荣的仆役和所有天军也是如此。那么，你不认为更好的做法是跟随我们的创造者天主，如同一位父亲，按祂所知的方式训练我们并赐予我们恩赐吗？但若如你所说，有一位比所有人都更仁慈的天主，那么肯定祂不会对我们发怒；或者若祂发怒，祂就是恶的。因为若我们的天主发怒并惩罚，祂不是恶的，而是公义的，因为祂纠正和改正祂自己的儿子。但那位与我们无关的天主，若祂惩罚我们，祂怎么能是善的呢？因为我们没有因虚妄的想象而被引诱去离弃我们自己的父亲并跟随祂，就给我们惩罚，你怎能宣称祂是如此良善，以至于祂甚至不能被看作公义的呢？}\par

% structure: p0115 source=h2 target=subsection reason=relative-heading-rank; base=h2

\subsection{第五十七章 西蒙的矛盾}

然后西蒙说：\Quote{伯多禄，你如此错误，难道不知道我们的灵魂是由那位良善的天主——最卓越的一位——造的，但被俘掳到这个世界来了吗？}伯多禄回答：\Quote{那么祂并非如你刚才所说为所有人所不知；而且，若祂是全能者，良善的天主又怎会允许祂的灵魂被俘呢？}然后西蒙说：\Quote{祂派遣创造者天主来造世界；而他造了以后，就宣称自己是天主。}伯多禄说：\Quote{那么祂并非如你所说，不为造世界的那位所知；灵魂也并非不认识祂，若他们确实是从祂那里被偷来的。那么，祂能对谁隐藏呢？如果创造世界的那位知道祂，因为受祂派遣；所有灵魂也知道祂，因为被强行从祂那里夺走。此外，我希望你告诉我们：派遣创造世界的那位，难道不知道祂不会守信吗？若祂不知道，则祂没有预知；若祂预知并容忍了这事，祂自己就是此事的罪人，因为祂没有阻止；但若祂不能阻止，则祂不是全能的。但若祂明知是善而没禁止，那么那个大胆做了差遣祂的那位不知道是善的事的，反而显得更好。}\par

% structure: p0117 source=h2 target=subsection reason=relative-heading-rank; base=h2

\subsection{第五十八章 西蒙的天主是不义的}

然后西满说：\Quote{他接纳那些前来投靠他的人，并善待他们。} 伯多禄回答说：\Quote{但这并不是什么新事；因为你也承认是世界的创造者的那一位也这样做。} 然后西满说：\Quote{但善天主只要被人认识就赐予救恩；而世界的创造者却要求遵守法律。} 于是伯多禄说：\Quote{他拯救奸夫和杀人犯，只要他们认识他；但善良、节制、仁慈的人，如果因为无人告知而不认识他，他却不拯救！你所宣扬的这位真是伟大而良善，他与其说是恶人的救主，不如说他对善人毫不怜悯。} 然后西满说：\Quote{人处于肉身之内时，要认识他确实非常困难；因为包围灵魂的这身体比一切黑暗更黑，比一切泥土更重。} 于是伯多禄说：\Quote{你的那位善天主要求的是艰难的事；但那真正的天主寻求更容易的事。那么，既然他是如此良善，就让他把我们留给我们的父和创造者吧；一旦我们离开身体，离开你所说的那种黑暗，我们就会更容易认识他；那时灵魂会更清楚地明白天主是它的创造者，并与他同在，不再被各种幻想扰乱；也不会想逃往另一个权能，那个权能除了西满无人知晓，而且如此良善，以致除非人首先对自己父亲行不敬，否则无人能到达它！我不知道这个权能如何能被称作良善或正义，因为只有通过对自己创造者行不敬才能取悦它！}\par

% structure: p0119 source=h2 target=subsection reason=relative-heading-rank; base=h2

\subsection{第五十九章 创造者我们的父}

然后西满说：\Quote{为了更大的利益和好处，逃向一位更有荣耀者，并不是不敬。} 伯多禄说：\Quote{如果正如你所说的，逃向一位陌生者并非不敬，那么无论如何，留在我们自己的父亲身边，即使他贫穷，也更为虔诚。但是，如果你认为离开我们的父亲、逃向另一位，以为他更好，并不是不敬；并且你相信我们的创造者不会因此介意；那么善天主就更加不会生气，因为当我们还是他的陌生者时，我们没有逃向他，而是留在我们自己的创造者身边。是的，我认为他反而会更称赞我们，因为我们忠于天主我们的创造者；因为他会想到，如果我们是他造的，我们就绝不会被任何其他人的引诱所迷惑而离弃他。因为如果有人被更丰富的应许所引诱，离开自己的父亲，投靠一位陌生者，那么他可能也会在日后离开他，再去投靠另一位许诺更大事物的人，这尤其因为他不是那人的儿子，因为他甚至可以离开那位按本性是他的父亲的人。} 然后西满说：\Quote{但如果灵魂是从他而来，却不认识他，而他确实是他们的父亲呢？}\par

% structure: p0121 source=h2 target=subsection reason=relative-heading-rank; base=h2

\subsection{第六十章 创造者至高的天主}

于是伯多禄说：\Quote{你把他描述得够软弱的。因为，如果如你所说，他比万物都更有能力，那么绝不可能相信较弱的能从较强的夺走战利品。\BibleRef{路 11:22} 或者说，如果创造者天主能够以暴力将灵魂带进这个世界，那么当灵魂脱离身体、挣脱囚禁的锁链之后，善天主又怎能召唤他们去承受惩罚，借口他们（因他的疏忽或软弱）被拖到这里，并卷入身体之中，如同在无知的黑暗里？在我看来，你似乎不知道什么是父亲和天主；但我可以告诉你灵魂的来源，以及它们何时和如何被造；但现在不允许我将这些事透露给你，你关于天主的认知竟如此错误。} 然后西满说：\Quote{终有一天你会后悔没有理解我所说的关于那不可言说的权能的事。} 伯多禄说：\Quote{那么就请你，正如我常说的，既然你是一位新天主，或是从他那降下的，就赐给我们一种新的官能，藉以认识你所说的那位新天主；因为天主我们的创造者赐给我们的这五种官能，忠于自己的创造者，无法察觉有其他天主存在，因为它们的本性必然如此。}\par

% structure: p0123 source=h2 target=subsection reason=relative-heading-rank; base=h2

\subsection{第六十一章 想象}

为此西满回答说：\Quote{请你把心思用于我将要说的事，并使它行在和平的道路上，以达至我将要证明的事物。现在请听。你是否曾在思想中将你的心智延伸到遥远的地区或岛屿，并如此专注于其中，以至于甚至看不见眼前的人，也不知道自己坐在何处，只因你注视的那些事物令人愉悦？} 伯多禄说：\Quote{确实如此，西满，我常遇到这种情况。} 然后西满说：\Quote{现在请以同样方式将你的感官延伸到天上，甚至超越上天，看那必定存在一个超越世界或在世界之外的地方，那里既无天也无地，也没有任何这些东西的阴影产生黑暗；因此，由于其中既无物体，也无物体引起的黑暗，必然有无限的光明；并想一想那种光是什么样的，它永远不被黑暗接替。因为如果这太阳的光充满了整个世界，那么你想那无物体且无限的光会有多大？毫无疑问，如此之大，以至于这太阳的光相比之下会显得是黑暗而非光明。}\par

% structure: p0125 source=h2 target=subsection reason=relative-heading-rank; base=h2

\subsection{第六十二章 伯多禄关于想象的经验}

当西满这样说了以后，伯多禄回答说：\Quote{如今请耐心听关于这两件事，就是关于伸展感官的例子和关于光的无限。我知道我自己，西满，有时在思想中曾如你所说，将我的感官延伸到远处的地区和岛屿，并以我的心灵看见它们，不亚于用眼睛看见。当我在葛法翁，忙于捕鱼，坐在一块岩石上，手里拿着一个连在线上、适于诱骗鱼的钩子时，\Emph{我是如此地全神贯注}，以致没有感觉到有鱼上钩，而我的心灵急切地穿过我所爱的耶路撒冷，我曾多次醒着上去那里为献祭和祈祷。但我也习惯羡慕这座凯撒勒雅，从别人那里听说，并渴望见到它；我似乎自己看到了它，尽管从未去过；我根据在其他城市所见，想着它作为一座大城市应有的东西：城门、城墙、浴场、街道、小巷、市场等等；我如此陶醉于这样审视的专注，以致如你所说，我没有看见站在我身边、在我旁边的人，也不知道自己坐在哪里。}然后西满说：\Quote{现在你说得很好。}\par

% structure: p0127 source=h2 target=subsection reason=relative-heading-rank; base=h2

\subsection{伯多禄的遐想}

然后伯多禄说：\Quote{简而言之，当我由于心灵的占据而没有察觉到我已经钓到一条非常大的鱼附着在钩上，尽管它正从我手中拖走钓线，我兄弟安德肋坐在我旁边，看见我陷入遐想几乎要跌倒，用他的手肘戳我的肋旁，好像要把我从睡梦中唤醒，说：\InnerQuote{你没有看到吗，伯多禄，你钓到一条多大的鱼？你神志不清了吗，为何这样惊呆？告诉我，你怎么了？}但我对他有点生气，因为他把我从沉思那些事的喜悦中拉回来；然后我回答说我没有遭受任何病痛，而是我心中正注视着可爱的耶路撒冷，同时也注视着凯撒勒雅；虽然身体与他同在，心灵却完全被带往那里。但他，不知从何处得到灵感，说出了一个隐藏而奥秘的真理之言。}\par

% structure: p0129 source=h2 target=subsection reason=relative-heading-rank; base=h2

\subsection{安德肋的责备}

\Quote{放弃吧，他说，啊，伯多禄。你在做什么？因为那些开始被魔鬼附身或精神错乱的人，就这样开始。他们首先被幻想带向一些愉快和令人高兴的事物，然后他们向着不存在的事物徒然无益地倾泻。这是由于某种心灵的疾病，因为他们看不见存在的事物，却渴望把不存在的事物带到眼前。但那些患热狂的人也这样，他们似乎看见许多形象，因为他们的灵魂，由于过冷或过热而脱离了其位置，遭受着自然服务的失败。那些因口渴而痛苦的人，当入睡时，似乎看见河流和泉水，并喝水；但这发生在他们因身体未湿润而干燥受苦。因此，可以肯定这发生在灵魂或身体的某种疾病。}\par

% structure: p0131 source=h2 target=subsection reason=relative-heading-rank; base=h2

\subsection{想像的谬误}

\Quote{简而言之，为了使你相信这事：关于耶路撒冷，我经常见过，我告诉了我的兄弟我曾以为自己看见哪些地方和哪些聚集的人群。但关于凯撒勒雅，我从未见过，我却坚持认为它就像我在心中和思想中所设想的那样。但当我来到这里，看见完全不像我在幻象中所见的一切时，我责备自己，并清楚地观察到，我给它分配了从我所见的其他城市中看到的城门、城墙和建筑，实际是从其他城市借取相似性。确实，没有人能想象出什么新东西，或从未存在过的形式。因为即使有人凭想象造出五头公牛，他也只是从他所见过的一头牛造出五头牛。因此，你，现在，如果你真的认为自己用思想感知到什么，并看在天之上，毫无疑问你是在根据你在地上所见到的事物来想象它们。但如果你认为你的心灵容易升到天上，并且你能够理解那里的事物，并领会那无限光的知识，我认为对于一个能理解这些事物的人来说，更容易把他的感官（知道如何升到那里的感官）投入到我们旁边某人的心里和胸中，并说出他在胸中怀有什么思想。因此，如果你能说出我们当中任何一个人的内心思想——这个人并非事先偏向于你——我们也许就能相信你，认为你能知道天上那些事物，虽然这些事物要高得多。}\par

% structure: p0133 source=h2 target=subsection reason=relative-heading-rank; base=h2

\subsection{存在与概念}

西满回答说：\Quote{啊，你编织了许多轻浮之网的人，现在请听。任何进入人思想的东西，不可能不在真理和现实中存在。因为不存在的事物没有表象；没有表象的事物，不能呈现于我们的思想。}然后伯多禄说：\Quote{如果一切能进入我们思想的事物都存在，那么，关于你所说的世界之外的那无限地方，如果一个人心里认为它是光，另一个人认为它是黑暗，那么同一个地方如何能同时是光和黑暗，根据他们不同的想法呢？}然后西满说：\Quote{暂且放过我所说的；告诉我们你认为在天之上是什么。}\par

% structure: p0135 source=h2 target=subsection reason=relative-heading-rank; base=h2

\subsection{法律教导无限}

然后伯多禄说：\Quote{如果你相信真正的光明之泉，我就可以教导你什么是无限以及无限是怎样的，并会提出一个连贯而必然的真理叙述，不是使用诡辩的主张，而是使用法律和本性的见证，好让你知道法律特别包含了我们对于无限所应相信的。但如果无限的教义并非法律所不知，那么法律肯定也不可能有其他未知之事；因此，你的假设是错误的，即认为法律有所不知。何况，赐予法律的那位更不会有所不知。然而，关于无限和那些没有界限的事物，我无法对你说什么，除非你先接受我们对那些被某个界限所界定的天的解释，或者提出你自己的解释。但如果你不能理解那些被固定界限所包含的事物，那么你更无法知道或学习那些没有界限的事物。}\par

% structure: p0137 source=h2 target=subsection reason=relative-heading-rank; base=h2

\subsection{第六十八章 可见与不可见的天。}

西满回答说：\Quote{在我看来，最好简单地相信天主存在，并且我们所见到的这个天是整个宇宙中唯一的天。}但伯多禄说：\Quote{并非如此；而应当宣认唯一真实存在之天主；但有诸天，是由祂所造，正如法律所说，其中一个是更高的，在其中也包含了可见的穹苍；那个更高的天是永恒而永久的，连同住在其中的那些人；但这个可见的天在世界终结时将溶解和消逝，以便那个更古老更高的天在审判后显现给圣者和堪当者。}西满回答说：\Quote{这些事如你所说，可能对那些相信的人显得如此；但对于寻求这些事理由的人，不可能从法律中得出证明，尤其是关于光明的无限。}\par

% structure: p0139 source=h2 target=subsection reason=relative-heading-rank; base=h2

\subsection{第六十九章 信德与理性。}

伯多禄说：\Quote{不要以为我们说这些事仅仅凭信德接受，而且也要凭理性断言。因为确实，没有理性就把这些事托付给单纯的信德是不安全的，因为真理不可能没有理性。因此，凡凭理性巩固而接受这些事的人，永远不会失去它们；而那些没有证明就接受它们的人，仅仅同意一个简单的陈述，既不能安全地保存它们，也不确定它们是否为真；因为容易相信的人，也容易放弃。但那些为他所相信和接受的事寻求理由的人，如同被理性本身的锁链束缚，永远不能从他相信的事中被撕开或分离。因此，一个人越是急切地要求理由，他就越坚定地保持他的信德。}\par

% structure: p0141 source=h2 target=subsection reason=relative-heading-rank; base=h2

\subsection{第七十章 休会。}

西满回答：\Quote{你承诺了一件大事，即无边光明的永恒可以从法律中证明。}伯多禄说：\Quote{只要你愿意，我就证明。}西满回答说：\Quote{如今时候已晚，我明天再与你对辩；如果你能证明这个世界是被造的，并且灵魂是不死的，我就协助你宣讲。}他说完这话，就离去了，跟随他的是同来的所有人中的三分之一，约一千人。其余的人双膝跪下，俯伏在伯多禄面前；伯多禄呼求天主之名，治愈了一些附魔的人，医治了其他患病的人，然后欢欢喜喜地遣散了群众，吩咐他们第二天早点来。群众散去后，伯多禄命令在院子里，即辩论举行的露天地方，把桌子摆在地上，与那十一人一同坐下；但我与一些其他刚开始聆听天主圣言、且深受爱戴的人斜倚着吃饭。\par

% structure: p0143 source=h2 target=subsection reason=relative-heading-rank; base=h2

\subsection{第七十一章 与不洁者分离。}

伯多禄极其仁慈地看着我，惟恐那种分离可能使我忧伤，便对我说：\Quote{克莱孟，我并非出于傲慢而不与尚未净化的人一起吃饭；而是恐怕我可能伤害自己，对他们也无益。因为我希望你们确知，凡曾崇拜偶像、敬拜外教人所称的神明、或吃过祭献给他们的东西的人，都不是没有不洁之神的；因为他已成为魔鬼的宾客，并分享了那在他心中形成形象的魔鬼，无论是出于恐惧还是爱。\BibleRef{格前 10:20} 借此，他并非没有不洁之神，因此需要圣洗的净化，使那不洁之神离开他，这神已潜伏在他灵魂深处的情感中，更糟的是，它不表示它潜伏在内，恐怕被揭露和驱逐。}\par

% structure: p0145 source=h2 target=subsection reason=relative-heading-rank; base=h2

\subsection{第七十二章 补救法。}

\Quote{因为这些邪魔喜爱居住在人的身体里，好借人的服事来满足自己的欲望；它们牵引人的灵魂动向，使之趋向它们自己所渴望的事物，迫使人们顺从它们的私欲，使人完全成为魔鬼的器皿。其中一个便是这个\Person{西满}，他被这种病患所困，如今已无法治愈，因为他的意志和意图都生了病。魔鬼并非违背他的意愿住在他内；因此，若有人要把它从他身上赶出去，由于它已与他不可分离，并且可以说已成了他的灵魂本身，那倒更像是杀了他，并担上杀人罪责。因此，你们中谁也不要因不能与我们一同用餐而悲伤，因为每个人都该明白，分别的时间长短全凭他乐意。因为愿意尽快领洗的人，只分别很短的时间；而愿意稍后领洗的人，则分别较长的时间。因此，每个人都可以自行决定要求较短或较长的悔改时间；所以，何时愿意来与我们共桌，在于你们，而不在于我们，因为我们不被允许与任何尚未受洗的人一同进食。因此，如果你们在净化的道路上设置拖延，延迟你们的领洗，那反倒是你们妨碍了我们与你们一同用餐。}说完这话，祝福后，他进食了。之后，他向天主谢恩，便进入屋内上床就寝；我们也都如此行，因为已是夜晚。\par

\SourceRecord{Translated by Thomas Smith.；Ante-Nicene Fathers；Vol. 8.；Edited by Alexander Roberts, James Donaldson, and A. Cleveland Coxe.；Buffalo, NY: Christian Literature Publishing Co.,；1886.；<http://www.newadvent.org/fathers/080402.htm>.}

% structure: p0001 source=h1 target=section reason=page-title

\section{认识（卷三）}

% structure: p0002 source=h2 target=subsection reason=relative-heading-rank; base=h2

\subsection{第一章 珍珠投于猪前}

与此同时，伯多禄在鸡叫时起身，想要叫醒我们，发现我们已醒，晚灯还亮着。他照例问候我们后，大家都坐下，他便这样开始了。\Quote{我的弟兄们，没有什么比在一群混杂的人群中论及真理更困难的了。因为那真实的事，不能因着那些怀着恶意和诡诈倾听的人而照实说出；但为了那些真诚渴望听见真理的人，也不该欺骗。那么，面对混杂的人群，他该做什么呢？他该隐藏真实吗？那样又如何教导那些配得的人呢？但如果他把纯正的真理摆给那些不愿获得救恩的人，他就伤害了差遣他的那一位，他从那一位领受了命令，不可将祂话语的珍珠投在猪和狗前，\BibleRef{玛 7:6}它们用辩论和诡辩攻击这些话语，把它们滚入属肉体理解的泥淖中，用它们的吠叫和卑劣的回答，使天主圣言的宣讲者疲乏不堪。因此，我大部分时候也使用某种间接的说法，尽力避免向不配的耳朵发表关于至高神性的首要知识。}然后，从圣父、圣子及圣神开始，他简洁明了地向我们讲解，以至我们所有听见的人都惊奇：人们竟抛弃真理，转向虚妄。\par

% structure: p0004 source=h2 target=subsection reason=relative-heading-rank; base=h2

\subsection{第十二章 第二日的讨论}

及至天亮，有人进来说：\Quote{庭院里聚集了极多的人群，西满站在他们中间，企图用最邪恶的劝说占据百姓的耳朵。}于是伯多禄立刻出去，站在前一天辩论的地方，所有百姓都欢喜地转向他，留心听他。但当西满察觉百姓看见伯多禄就欢喜，且被吸引去爱他时，他困惑地说：\Quote{我奇怪人的愚昧，竟称我为魔术士，却爱伯多禄；而他们对我早有认识，反倒应该爱我才是。因此，有理智的人可以从这个迹象明白，伯多禄倒更像是魔术士，因为人们对我不生爱慕——我几乎因熟识而应得这份爱——却对他慷慨地倾注了爱，而他与这份爱毫无亲缘。}\par

% structure: p0006 source=h2 target=subsection reason=relative-heading-rank; base=h2

\subsection{第十三章 西满——诱惑者}

当西满这样讲论时，伯多禄照常问候了百姓，便这样回答：\Quote{西满啊，每人自己的良心足以驳斥他；但如果你对此感到奇怪——那些认识你的人不仅不爱你，甚至恨你——那就从我这里学习缘由吧。既然你是诱惑者，却声称宣讲真理，因此你曾有许多渴望学习真理的朋友。但当他们看见你身上与你所宣讲的相反，他们，正如我说的，既是真理的爱好者，便开始不仅不爱你，甚至恨你。然而他们并没有立刻抛弃你，因为你仍承诺能向他们指明真理。因此，只要没有能向他们指明的人，他们就容忍你；但自从更好的教导的希望向他们显现，他们就轻视你，并寻求那他们认为更好的认识。而你，确实以邪恶的技艺行事，起初以为能逃脱察觉。但你已暴露了。因为你是被逼入角落，出乎意料地变得声名狼藉，不仅不懂真理，而且不愿从认识真理的人那里听取。因为假若你愿意听，那一位所说的这句话就应验在你身上，祂说：\InnerQuote{没有隐藏的事，不会被揭露；也没有遮盖的事，不会被彰显。}}\BibleRef{玛 10:26}\par

% structure: p0008 source=h2 target=subsection reason=relative-heading-rank; base=h2

\subsection{第十四章 西满要求伯多禄履行承诺}

当伯多禄说这些话以及类似的话时，西满回答说：\Quote{我不愿你用长篇大论拖延我，伯多禄；我要求你履行昨天承诺的。你当时说，你能表明法律教导关于永恒之光的无限性，只有两个天，且是受造的，较高者是那光的居所，不可言喻的圣父独自永远居住在那里；而可见的这个天是仿照那个天造的，你说它将要过去。因此，你说万有的父是独一的，因为不能有两个无限者；否则两者都不是无限的，因为在一个存在之处，它就构成了另一个存在的界限。既然你不仅承诺了这一点，而且能从法律中表明，那就放下其他事，着手这个吧。}接着伯多禄说：\Quote{如果我只是为了你的缘故——你来的唯一目的是反驳——才被要求谈论这些事，那你永远也不该从我这里听到一篇讲论；但既然农夫在愿意播种好土壤时，也必须在石头地或被人践踏的地方，或长满荆棘和蒺藜的地方撒下一些种子（正如我们的师傅也曾设喻，指出这些对应不同灵魂意向的多样性），我就不会拖延。}\par

% structure: p0010 source=h2 target=subsection reason=relative-heading-rank; base=h2

\subsection{第十五章 西满的傲慢}

那时西满说：\Quote{你似乎对我发怒；但若是如此，就不必进入争辩了。}伯多禄说：\Quote{我看得出你意识到自己将受谴责，便想彬彬有礼地逃避辩论；因为你究竟看到了什么，使我对你发怒？你是一个想要欺骗如此众多群众的人，却无话可说，假装节制，甚至还凭你的权威命令，辩论要按你的意思进行，而不是按秩序的要求？}西满说：\Quote{我要强迫自己耐心忍受你的无知，以便显明你确实想引诱百姓，而我教导真理。但现在我不讨论那无限的光明。回答我所问你的：既然照你所说，天主造了万物，邪恶从何而来？}伯多禄说：\Quote{以这种方式提问不是对手的姿态，而是学习者的姿态。因此，如果你愿意学习，就承认这一点；我会先教你该如何学习，等你学会了聆听，我就立即开始教导你。但如果你不愿学习，好像你什么都知道，我会先陈述我所宣讲的信仰，你也陈述你认为正确的事；当我们各自的表白都公开后，让听众判断谁的言论有真理支持。}西满回答说：\Quote{这真是个笑话：瞧这个人居然要教导我！不过我还是容忍你，忍受你的无知和傲慢。我承认，我确实愿意学习；让我们看看你如何教导我。}\par

% structure: p0012 source=h2 target=subsection reason=relative-heading-rank; base=h2

\subsection{第十六章 邪恶的存在}

伯多禄说：\Quote{如果你真正愿意学习，那么就首先学习这一点：你的问题提得多么拙劣；因为你说：既然天主造了万物，邪恶从何而来？但在你问这个问题之前，应当先提出三类问题：\Emph{第一}，邪恶是否存在？\Emph{第二}，邪恶是什么？\Emph{第三}，它对谁存在，从何而来？}西满回答说：\Quote{哦，你这个最愚拙无知的人，难道有谁不承认此生中有邪恶吗？因此，我也以为你具有众人共有的常识，才问邪恶从何而来；并非由于我愿意学习——因为我通晓一切，更不是向你学习，你一无所知——而是为了向你显明你一无所知。而且，为了不让你认为我因发怒而说话严厉些，要知道我是出于对在场之人的怜悯，你正企图欺骗他们。}伯多禄说：\Quote{你更加邪恶了，如果你能不生气就做这样不义的事；但有火的地方必有烟。不过我还是告诉你，免得显得我只抓住你的话，而不回答你那些语无伦次的话。你说人人都承认邪恶的存在，这其实是假的；因为首先，整个希伯来民族否认其存在。}\par

% structure: p0014 source=h2 target=subsection reason=relative-heading-rank; base=h2

\subsection{第十七章 并非人人承认}

西满打断他的话说：\Quote{说没有邪恶的人做得对。}伯多禄回答说：\Quote{我们现在不打算讨论这一点，只是陈述一个事实：邪恶的存在并非普遍承认。但你应该问的第二个问题是：邪恶是什么？——是实体，是属性，还是行为？以及其他许多类似的问题。然后，它是针对什么、如何存在、或对谁是邪恶的？——对天主、天使、人类、义人、恶人、所有人、一部分人、自己、还是无人？然后你应当追问：它从何而来？——来自天主还是来自虚无？是永恒存在还是在时间中开始？是有用还是无用？以及这类命题所要求的许多其他问题。}西满回答说：\Quote{请原谅，我在第一个问题上犯错了；但假设我现在先问：邪恶是否存在？}\par

% structure: p0016 source=h2 target=subsection reason=relative-heading-rank; base=h2

\subsection{第十八章 进行讨论的方式}

伯多禄说：\Quote{你以何种方式提问？是愿意学习，还是教导，或是为了质疑？如果确实是愿意学习，我首先要教导你，使你通过合理的顺序和正确的教义，从自身理解邪恶是什么。但如果你是作为教导者提问，我不需要受你教导，因为我有一位师傅，从他那里我学到一切。但如果你仅仅为了质疑和争论而问，那就让我们各自先陈明意见，再辩论此事。因为你不应当既像愿意学习的人那样提问，又像教导的人那样反驳，以至于在我回答之后，你可以随意说我说得好或不好。因此，你既不能站在反对者的位置，也不能评判我们所说的话。所以，正如我所说，如果进行讨论，让每个人陈述自己的观点；当我们处于争论中时，这些虔诚的听众将是公正的法官。}\par

% structure: p0018 source=h2 target=subsection reason=relative-heading-rank; base=h2

\subsection{第十九章 求知欲}

然后西满说：\Quote{你们不认为让一群无知的百姓来评判我们的话是荒谬的吗？}伯多禄答道：\Quote{并非如此；因为一个人可能不太清楚的事，可以由多人共同探究，常常连民间的传言也具有预言的性质。但除此之外，所有这些人都站在这里，被对天主的爱和对认识真理的渴望所约束，因此他们所有人都应被视为一体，因为他们对真理的情感是同一的；反之，两个人若彼此不合，就是众多且相异的。但是，若你愿意得到一个迹象，说明这些站在我们面前的人如何如同一人，你可以从他们的沉默和安静中看出，他们以何等的忍耐，如你所见，甚至在尚未认识天主的真理之前就尊敬它；因为他们尚未学会他们所欠的更大敬意。因此，我盼望，藉着天主的仁慈，祂会接纳他们心中对祂的虔诚意向，并将胜利的棕榈枝赐给宣讲真理的人，好叫祂向他们彰显真理的宣告者。}\par

% structure: p0020 source=h2 target=subsection reason=relative-heading-rank; base=h2

\subsection{第二十章　共同原则。}

然后西满说：\Quote{你希望就什么题目进行讨论？请告诉我，我也好阐明我的想法，这样探讨就可以开始了。}伯多禄答道：\Quote{你若真按我所认为正确的去做，我愿依照我师傅的训诫来办，祂首先命令希伯来民族——祂知道他们认识天主，且知道是祂创造了世界——不是要他们探究他们所认识的那一位，而是要在认识祂的基础上，探究祂的旨意和祂的正义；因为人有能力去探究这些事，从而发现、实行并遵守那些他们将受审判的事。因此，祂命令我们不要追问邪恶从何而来，像你刚才问的，而要寻求善的天主的正义和祂的国，并且这一切，祂说，都要加给你们。}\BibleRef{玛 6:33}然后西满说：\Quote{既然这些事是命令给希伯来人的，因为他们对天主有正确的认识，并且认为每个人都有能力去做那些他将受审判的事——但我的意见与他们不同——你想让我从哪里开始？}\par

% structure: p0022 source=h2 target=subsection reason=relative-heading-rank; base=h2

\subsection{第二十一章　意志的自由。}

然后伯多禄说：\Quote{我建议先探究：我们是否有能力知道我们将因什么受审判。}但西满说：\Quote{并非如此；而是关于天主，所有在座的人都渴望听。}伯多禄说：\Quote{那么，你承认某件事是在意志的能力之下的：只要承认这一点，如果真是如此，我们就按你所说，来探究关于天主的事。}西满答道：\Quote{绝不。}伯多禄说：\Quote{那么，如果什么都不在我们能力之内，我们探究任何关于天主的事就是徒劳的，因为寻求的人没有能力找到；因此我说得好，这应当是首要的探究：意志的能力之下是否有任何东西。}西满说：\Quote{连你所说的这一点——意志的能力之下是否有任何东西——我们也无法理解。}伯多禄见他转向争辩，并且因害怕被驳倒而将一切事物都搅得模糊不清，就答道：\Quote{那么，你怎么知道人没有能力知道任何事呢？因为至少你知道这一点。}\par

% structure: p0024 source=h2 target=subsection reason=relative-heading-rank; base=h2

\subsection{第二十二章　责任。}

然后西满说：\Quote{我不知道我是否连这一点也知道；因为每个人按照命运给他所定的，去行事、理解或受苦。}伯多禄说：\Quote{看哪，我的弟兄们，西满陷入了何等荒谬的境地！他在我来之前还教导说，人有能力变得聪明并做他们所愿意的事，但现在被我的论据逼到角落，竟否认人有任何知觉或行动的能力，却还胆敢自称为师傅！但请告诉我，如果人没有能力做任何事，天主怎会按真理审判每个人的行为呢？若持有这意见，一切都会被连根拔起：追求良善的渴望便是徒然；甚至世上的审判官施行法律、惩罚行恶的人也是徒然，因为他们没有能力不犯罪；列国的法律对恶行定下刑罚也是徒然。那些辛苦持守正义的人将是可怜的；而那些生活在享乐中、施行暴政、过着奢华和邪恶生活的人则有福了。照此说来，便没有正义、没有良善、没有美德，也没有你所说的天主。但是，西满啊，我知道你为何这样说：实在是因为你想避免探究，免得被公开驳倒；因此你说人没有能力知觉或分辨任何事。但若这真是你的意见，你在我来之前就不会在百姓面前自称师傅了。因此，我说，人是在自己的控制之下的。}西满说：\Quote{在自己控制之下是什么意思？请告诉我们。}伯多禄答道：\Quote{若什么都不能学，你为什么还想听？}西满说：\Quote{你对这一点无话可答。}\par

% structure: p0026 source=h2 target=subsection reason=relative-heading-rank; base=h2

\subsection{第二十三章　罪恶的起源。}

伯多禄接着说：\Quote{我并非受你强迫，而是应听者的请求才发言。选择的能力是灵魂的知觉，具有一种性质，能倾向于它愿意采取的行动。} 西满对伯多禄的讲话表示赞许，说道：\Quote{你解释得真是宏伟无比，无可比拟，因为我有责任为你的善言作证。现在，如果你能向我解释我此刻提出的问题，那么在其他一切事上，我都将服从你。我所希望知道的便是：凡天主愿意有的，便存在；凡祂不愿意有的，便不存在。请回答我。} 伯多禄说：\Quote{如果你不知道你问的是一个荒谬而不合逻辑的问题，我会原谅你并加以解释；但如果你明知你在不合逻辑地提问，那就不对了。} 西满说：\Quote{我以那至高的神性（无论祂是谁）起誓，祂审判并惩罚犯罪者，我确实不知道我所说的话有何不合逻辑之处，也不知道我的话——即我刚才所说的话——有何荒谬之处。}\par

% structure: p0028 source=h2 target=subsection reason=relative-heading-rank; base=h2

\subsection{第二十四章　天主是善的作者，而非恶的作者。}

伯多禄回答说：既然你承认你无知，那么现在学习吧。你的问题要求我们解决两个互相矛盾的事项。因为每个运动都分为两部分：一部分受必然性推动，另一部分受意志推动；受必然性推动的事物总是在运动中，而受意志推动的事物则不总是如此。例如，太阳的运行是必然地完成其指定的轨道，而天界的每一状态和服务都依赖于必然的运动。但人支配着自己行为的自愿行动。因此，有些事物被造的目的就是在它们服务中受必然性支配，不能做任何其他被指派给它们的事；当它们完成这项服务时，万物的创造者——祂按自己的旨意如此安排——保存它们。但另一些事物则具有意志的力量，并能自由选择它们愿意做的事。这些事物，如我所说，并不总是停留在它们被造时的状态中；而是依照它们的意志引导和心智的判断倾向，它们产生或善或恶；因此祂为行善者设定了赏报，为作恶者设定了惩罚。\par

% structure: p0030 source=h2 target=subsection reason=relative-heading-rank; base=h2

\subsection{第二十五章　\Quote{谁曾抗拒过祂的旨意？}}

因此，你说：如果天主愿意什么存在，它就存在；如果祂不愿意，它就不存在。但如果我回答：祂愿意的便存在，祂不愿意的便不存在，那么你会说：既然如此，祂就愿意世界上所行的恶事存在，因为凡祂愿意的便存在，凡祂不愿意的便不存在。但如果我回答：并非天主愿意的便存在、祂不愿意的便不存在，那么你会反诘我：天主一定是无能的，因为祂不能做祂愿意做的事；你甚至会更加无礼，以为你已获胜，尽管你所说的完全不着边际。因此，西满啊，你实在无知，非常无知，不懂得天主的旨意如何在每个个别情况下运作。因为有些事物，如我们所说，祂已如此愿意它们存在，以至于它们不能与祂所规定的方式不同；对于这些，祂既未指定赏报也未指定惩罚；但那些祂愿意它们具备行使其意志能力的事物，祂根据它们的行为和意志指定它们应得赏报或惩罚。因此，正如我告诉你的，一切运动的事物按照我先前陈述的区分分为两部分，凡天主愿意的便存在，凡祂不愿意的便不存在。\par

% structure: p0032 source=h2 target=subsection reason=relative-heading-rank; base=h2

\subsection{第二十六章　没有自由就没有善。}

对此西满回答说：\Quote{难道祂不能把我们造成都是善的，且不能有别的样子吗？} 伯多禄回答说：\Quote{这也是一个荒谬的问题。因为如果祂把我们造成一种不变的本性，不能从善移开，那么我们就不是真正的善，因为我们不能是别的样子；我们行善也不是出于本意；我们所做的也不是我们自己的，而是我们本性的必然。但怎么能把并非出于本意的事称为善呢？因此，世界需要漫长的时期，直到预定充满它的灵魂数目完成为止，然后那可见的天要像卷轴一样被卷起，那更高的天将显现，有福者的灵魂恢复身体后被引入光明；而恶者的灵魂因他们不洁的行为被火灵环绕，将被投入不灭之火的深渊，承受永罚。这些事确实如此，那真先知曾向我们作证；关于祂，如果你想知道祂是一位先知，我可以用无数的声明来教导你。因为凡祂所说的话，甚至现在祂所说的一切都在应验；而祂关于未来所说的话，人们相信将要应验，因为信德是从已发生的事给予未来的。}\par

% structure: p0034 source=h2 target=subsection reason=relative-heading-rank; base=h2

\subsection{第二十七章　可见的天：为何被造。}

但\Person{西满}察觉\Person{伯多禄}明显是在从预言的根本上提出理由，整个问题因此得以解决，便不愿让谈话朝向这个方向发展，于是回答说：\Quote{请回答我提出的问题，告诉我，如果这可见的天如你所说将要解散，那最初为何要造它？}\Person{伯多禄}回答说：\Quote{它是为了人类现世的生命而造的，好有某种屏障与分离，免得任何不配之人看见天上居所和\Person{天主}本身的住所；这些住所预备好只为心地纯洁的人所见。\BibleRef{玛 5:8}但现在，在争战之时，祂乐意使那些作为胜利者奖赏的事物成为不可见的。}于是\Person{西满}说：\Quote{如果造物主是善的，世界也是善的，那善的祂怎会摧毁善的事物呢？但如果祂要摧毁善的事物，祂自己又怎能被认为是善的呢？但如果祂要将它作为恶的解散并摧毁，那造了恶之物的祂，又怎能不被认为是恶的呢？}\par

% structure: p0036 source=h2 target=subsection reason=relative-heading-rank; base=h2

\subsection{第28章 为何要解散}

对此\Person{伯多禄}回答说：\Quote{既然我们承诺不逃避你的亵渎之语，我们便耐心忍受，因为你将为自己所说的话交账。现在请听。如果那可见且可逝的天是为它自己而造的，你说它不应被解散还有几分道理。但如果它不是为了自己，而是为了别的事物而造的，那它必须被解散，好使它所为之造的那事物显现出来。打比方说，鸡蛋的壳无论看似多么匀称精巧，仍必须破开，才能使小鸡从中出来，整个蛋壳形式所为之塑造的那东西得以显现。同样，这世界的状况也必须要过去，好使那更高的天上国度的状况得以显现。}\par

% structure: p0038 source=h2 target=subsection reason=relative-heading-rank; base=h2

\subsection{第29章 可朽坏与暂时之物由不朽与永恒者所造}

于是\Person{西满}说：\Quote{在我看来，由\Person{天主}所造的天不能被解散。因为由永恒者所造的是永恒的，而由可朽者所造的是暂时的和腐朽的。}\Person{伯多禄}说：\Quote{并非如此。事实上，各种可朽坏和暂时之物都由必死之物所造；但永恒者并非总造可朽之物，也并非总造不朽之物；而是按照\Person{天主}造物主的旨意，祂所造之物将如何。因为\Person{天主}的权能不受规律约束，祂的旨意就是祂造物的规律。}于是\Person{西满}回答说：\Quote{我唤你回到最初的问题。你方才说\Person{天主}无人能看见；但当那天解散、天上国度的更高境界显现时，心地纯洁的人\BibleRef{玛 5:8}将看见\Person{天主}；这话与法律相悖，因为法律上写着：\Person{天主}说：\BibleQuote{无人能见我的面而存活。}\BibleRef{出 33:20}}\par

% structure: p0040 source=h2 target=subsection reason=relative-heading-rank; base=h2

\subsection{第30章 心地纯洁者如何看见\Person{天主}}

然后\Person{伯多禄}回答说：\Quote{对于那些不按\Person{梅瑟}传统读法律的人，我的话似乎与之相悖；但我要向你们表明如何并不矛盾。\Person{天主}是以理智而非身体看见的，以精神而非肉体。因此，作为精神的众\Person{天使}看见\Person{天主}；所以人，只要是人，就不能看见祂。但在死人复活之后，当他们被造得像\Person{天使}一样\BibleRef{玛 22:30}，他们就能看见\Person{天主}。因此，我的话并不违反法律；我们的主所说的\BibleQuote{心地纯洁的人是有福的，因为他们要看见\Person{天主}}\BibleRef{玛 5:8}也不违反法律。因为祂表明，将来会有一个时期，人将变为\Person{天使}，他们在心灵的精神中看见\Person{天主}。}在这些话和许多类似的话之后，\Person{西满}开始以许多誓言断言说：\Quote{只就一件事给我一个理由，告诉我灵魂是否不死，我便在一切事上顺服你的意愿。但明天再说吧，因为今天已经晚了。}因此，当\Person{伯多禄}开始说话时，\Person{西满}出去了，跟着他出去的同伙很少；那是出于羞惭。但其余所有人都转向\Person{伯多禄}，屈膝跪在他面前；其中一些患各种疾病或被魔鬼侵扰的人，因\Person{伯多禄}的祈祷而获得痊愈，高兴地离去，仿佛同时获得了真\Person{天主}的道理和祂的仁慈。当人群退去，只有我们这些侍从留下与他在一起时，我们坐在铺在地上的榻上，每人认出自己习惯的位置，吃了食物，感谢了\Person{天主}，便睡觉了。\par

% structure: p0042 source=h2 target=subsection reason=relative-heading-rank; base=h2

\subsection{第31章 勤于研习}

但次日，\Person{伯多禄}照常黎明前起身，发现我们已经醒来并准备好听讲，便这样开始：\Quote{我恳求你们，我的弟兄和同仆，如果你们中有人无法保持清醒，不要因尊重我在场而折磨自己，因为突然的改变是困难的；但如果一个人长期逐渐习惯，那因习惯而来的就不会令人痛苦。因为我们并非都有同样的训练；虽然随着时间的推移，我们将能塑造成同一习惯，因为人们说习惯占据第二天性的地位。但我请\Person{天主}作证，如果有人无法保持清醒，我并不生气；反而，如果有人在整夜沉睡之后，白天没有补足夜间所缺，我才生气。因为必须全神贯注、不间断地研习道理，使我们的心智充满唯独对\Person{天主}的思念；因为充满对\Person{天主}思念的心智中，不会给恶者留出位置。}\par

% structure: p0044 source=h2 target=subsection reason=relative-heading-rank; base=h2

\subsection{第32章 \Person{伯多禄}的私下教导}

当\Person{伯多禄}这样对我们说话时，我们每个人都热切地向他保证，那时我们早已醒着，只睡了很短的时间就满足了，但我们不敢叫醒他，因为门徒指挥师傅是不合适的；\Quote{而且，\Person{伯多禄}啊，就连这件事我们也几乎冒险承担起来了，因为我们的心因渴望你的话语而激动，完全驱散了睡意。但我们对你的爱又阻止了我们，不让我们粗暴地叫醒你。}然后\Person{伯多禄}说：\Quote{既然你们声称因为渴望听道而情愿醒着，我愿意更仔细地重复给你们，并有序地解释昨天随意讲论的事。我打算在这些每日的辩论中这样做，以便在夜晚，当时间和地点都私密时，我将按正确的顺序和直接的解释，阐明在争论中尚未充分陈述的任何事。}然后他开始向我们指出，昨天的讨论应该怎样进行，以及由于对方的争强好辩或缺乏技巧，它未能如此进行；因此他仅使用了断言，只推翻了对手所说的，却没有完全或清楚地阐述自己的教义。然后他把几件事重复给我们，按正常顺序并充分理由地讨论它们。\par

% structure: p0046 source=h2 target=subsection reason=relative-heading-rank; base=h2

\subsection{第三十三章 学习者和好辩者}

但当天开始破晓时，他在祈祷后出去到人群那里，站在他惯常的地方准备讨论；看见\Person{西满}站在人群中间，他照常问候了人们，并对他们说：我承认，我对某些人感到难过，他们来到我们这里，本为了学习一些东西，但当我们开始教导他们时，他们却自称是师傅；当他们作为无知者提问时，却又作为知者来反驳。但也许有人会说，提问的人确实是为了学习而提问，但当他听到的似乎不对时，他必须回答，这看似反驳，其实不是反驳，而是进一步的询问。\par

% structure: p0048 source=h2 target=subsection reason=relative-heading-rank; base=h2

\subsection{第三十四章 违反秩序就是违反理性}

\Quote{那么让这样的人听这个：所有教义的教导都有一定的秩序，有些东西必须先讲，另一些其次，还有一些再次，如此按序进行；如果这些按序传授，就会变得清晰；但如果打乱顺序提出，就会显得违背理性。因此，如果我们寻求的目的是找到我们所寻求的，那么秩序是最应被遵守的。因为正确上路的人，会按顺序注意第二步，从第二步更容易找到第三步；他走得越远，知识之路就对他越敞开，直到他到达真理之城——他奔向并渴望到达的地方。但那些缺乏技巧、不知探索之道的人，就像在异国他乡的旅行者，无知而迷途，若不愿意雇用当地人为向导——毫无疑问，当他偏离真理之路时，将留在生命之门以外，如此，被黑夜的黑暗笼罩，行走在丧亡之路上。因此，既然那些当被寻求的事若有序地寻求最容易找到，而无技巧者不知探索的秩序，那么无知者应当服从有知者，先学习探索的秩序，以便最终找到提问和回答的方法。}\par

% structure: p0050 source=h2 target=subsection reason=relative-heading-rank; base=h2

\subsection{第三十五章 先学后教}

对此\Person{西满}回答说：\Quote{那么真理并非属于众人，而只属于那些知道辩论技巧的人，这是荒谬的；因为既然祂同样是众人的天主，众人就不能同样地知道祂的旨意。}然后\Person{伯多禄}说：\Quote{众人都是由祂所造平等，祂平等地赐给众人领受真理的能力。但没有人会怀疑，出生的人中没有人生来就有教育，教育是在出生之后才有的。因此，人的出生在这方面是平等的，即所有人都同样有能力接受训练，差别不在于本性，而在于教育。谁不知道，一个人所学的东西，在学之前是不知道的呢？}然后\Person{西满}说：\Quote{你说得对。}然后\Person{伯多禄}说：\Quote{如果在那些常用的技艺中，一个人先学习然后教导，那么那些自称是灵魂教育者的人，岂不更应该先学习然后教导，以免当他们自己缺乏技巧却承诺向别人提供知识时，使自己成为笑柄？}然后\Person{西满}说：\Quote{这在那些常用的技艺方面是真实的；但在知识的言语上，人一听就学会了。}\par

% structure: p0052 source=h2 target=subsection reason=relative-heading-rank; base=h2

\subsection{第三十六章 真理的明证}

那时伯多禄说：\Quote{如果有人以有序且规律的方式聆听，他就能认识真理；但若有人拒绝接受改过迁善的生活和纯洁品行的规诫——这实在是认识真理的正常结果——他就不会承认他所知道的。因为我们确实看到一些人，放弃了年轻时学到的技艺，转去从事其他行业，并借口懒惰，开始挑剔那技艺无益。} 于是西满说：\Quote{难道所有听到的人，都应当相信他们所听到的一切都是真的吗？} 伯多禄说：\Quote{凡听到真理的有序陈述，绝不能否认它，反而知道所说的真实，只要他也自愿服从生活的规诫。但那些听到却不愿行善的人，被作恶的欲望所阻止，不肯赞同他们判断为正确的事。由此可见，听众有能力选择他们所偏爱的。但如果所有听到的人都服从，那将是本性的必然，引领所有人走同一条路。因为正如没有人能被说服变矮或变高，因为本性力量不允许；同样，如果所有人都因一句话而皈依真理，或者全都不皈依，那将是本性力量在一个情况下迫使所有人，在另一个情况下无人皈依。}\par

% structure: p0054 source=h2 target=subsection reason=relative-heading-rank; base=h2

\subsection{第三十七章　天主既公义又良善}

西满说：\Quote{因此，请教我们，想要认识真理的人必须首先学习什么。} 伯多禄说：首要是必须探究人能够发现什么。因为天主的审判必然关乎这一点：一个人能否行善却未行。因此，人必须省察自己是否有能力通过寻求找到善，并在找到后行善；因为这是他们将要受审判的原因。但除此之外，除了先知，没有人需要知道更多：因为人需要知道世界如何被造呢？如果我们必须从事类似的建造，这确实需要学习。但现在，为了敬拜天主，知道祂创造了世界就足够了；但祂如何创造世界，不是我们要探究的课题，因为正如我所说，我们不必获得那种技艺的知识，好像我们要制造类似的东西。但我们也不会因未学习世界如何被造而受审判，而只会因我们是否不认识造物主而受审判。因为如果我们沿着正义的道路寻求祂，我们就会知道世界的造物主是公义而良善的天主。因为如果我们只知道祂是良善的，这种知识并不足以得救。因为在现世生活中，不仅配得的人，连不配得的人也享受祂的良善和恩惠。但如果我们相信祂不仅是良善的，也是公义的，并且按照我们对天主的信仰，在整个人生中持守正义，我们就将永远享受祂的良善。总之，对于希伯来人，他们关于天主的意见是祂只是良善的，我们的主教导他们也应寻求祂的义德，\BibleRef{玛 6:33}也就是说，他们应当知道祂在现世确实良善，使所有人都能在祂的良善中生活，但祂在审判之日将是公义的，将永恒的赏报赐给配得的人，不配得的人将被排除在外。\par

% structure: p0056 source=h2 target=subsection reason=relative-heading-rank; base=h2

\subsection{第三十八章　天主的公义在审判之日彰显}

西满说：\Quote{同一位存在怎能既良善又公义？} 伯多禄回答说：\Quote{因为没有公义，良善就会成为不义；因为良善的天主将祂的阳光和雨水均等地赐给义人和不义的人，\BibleRef{玛 5:45}但如果祂总是以同样的命运对待善人和恶人，这似乎不义，除非祂这样做是为了果实，所有生于这世界的人都能平等地享受。然而，正如天主赐下的雨水同样滋养麦子和莠子，但在收割时，庄稼被收入仓里，而糠秕或莠子被火焚烧，\BibleRef{玛 3:12}同样，在审判之日，当义人被引入天国，不义的人被抛出去时，天主的公义也将彰显。因为如果祂永远同样对待恶人和善人，这不仅不是良善，甚至是不义和不公正的：义人和不义的人被祂看作同等功过。}\par

% structure: p0058 source=h2 target=subsection reason=relative-heading-rank; base=h2

\subsection{第三十九章　灵魂的不朽}

西满说：\Quote{我希望得到满足的一点是：灵魂是否不朽？因为除非我首先知道灵魂的不朽，否则我无法承担正义的重负；因为如果灵魂不是不朽的，你们宣讲的道理就无法成立。} 伯多禄说：\Quote{让我们首先探究天主是否公义；因为如果这一点确定了，完美的宗教秩序就会立即建立。} 西满说：\Quote{虽然你夸口对讨论秩序的知识，但在我看来，你现在回答得违反秩序；当我请你表明灵魂是否不朽时，你却说要先探究天主是否公义。} 伯多禄说：\Quote{这是完全正确和有序的。} 西满说：\Quote{我希望知道如何如此。}\par

% structure: p0060 source=h2 target=subsection reason=relative-heading-rank; base=h2

\subsection{第四十章　由现世中恶人的昌盛所证明}

\Person{伯多禄}说：\Quote{请听，有一些亵渎天主的人，他们一生行不义和享乐，却安卧病榻而终，并获得体面的葬礼；另有一些人敬拜天主，以极度的诚实和节制过简朴的生活，却因坚守正义而死在荒野之地，甚至不配得到埋葬。那么，如果将来没有不死的灵魂为不虔敬的行为受苦，或为虔敬和正直而享受赏报，天主的公义何在？}于是\Person{西满}说：\Quote{正是这一点使我不信，因为许多行善的人悲惨地死去，而许多作恶的人却长寿享福。}\par

% structure: p0062 source=h2 target=subsection reason=relative-heading-rank; base=h2

\subsection{第41章 \Person{西满}的诡辩}

于是\Person{伯多禄}说：\Quote{正是这件使你陷入不信的事，反倒给我们提供了确证，证明将有审判。因为既然天主是公义的，就必然有另一个世界，在那里每人按其所应得而领受，从而证明天主的公义。但如果现在人人都已按其所应得而领受，那么当我们说将来有审判时，就真的像是骗子了。因此，正是在今生没有人按行为得报偿这一事实，给那些知道天主是公义的人提供了无可置疑的证据，证明将有审判。}\Person{西满}说：\Quote{那么为什么我仍不相信呢？}\Person{伯多禄}说：\Quote{因为你没有听真先知说过：\InnerQuote{你们先寻求祂的国和祂的义德，这一切自会加给你们。}}\BibleRef{玛 6:33} 于是\Person{西满}说：\Quote{请原谅我，在我知道灵魂是否不死之前，我不愿寻求义德。}\Person{伯多禄}说：\Quote{也请你原谅我这一点：因为我不能违背真理的先知所教导我的。}\Person{西满}说：\Quote{显然你不能断言灵魂是不死的，所以你诡辩，因为你知道如果它被证明是有死的，那么你所试图传播的那种宗教的整个体系就会被连根拔起。因此，我固然赞赏你的谨慎，却不赞同你的说服力；因为你用将来好处的希望说服许多人接受你的宗教，并克制享乐；结果他们失去了现世的享受，又因未来之事的希望而受骗。因为他们一死，灵魂就同时熄灭了。}\par

% structure: p0064 source=h2 target=subsection reason=relative-heading-rank; base=h2

\subsection{第42章 \Quote{满有各样诡诈和恶毒}}

但\Person{伯多禄}听他这样说，就咬牙切齿，用手擦额头，深深叹息说：\Quote{你武装着古蛇的狡猾，站出来欺骗灵魂；因此，正如蛇比任何走兽更狡猾，你自称从起初就是教师。再者，像蛇一样，你曾想引入多神；但现在，在那点上被驳倒后，你竟然声称根本没有神。因为你以某个不知名的神为借口，否认世界的创造者是神，却声称祂要么是恶的，要么有许多平等者，或者如我们所说，祂根本不是神。当你在这立场上被击败后，你现在断言灵魂是有死的，好使人不因将来之事的希望而正直地生活。因为如果将来没有希望，为何不放弃慈悲，纵情于奢侈和享乐？显然一切不义皆源于此。而你竟将如此不虔敬的教义引入人类悲惨的生活中，却自称虔敬，称我为不虔敬，因为我因将来好事的希望，不容许人们拿起武器互相攻击，掠夺和颠覆一切，并企图任何欲望所指示的事。你所引进的那种生活会是什么样子？人们将攻击与被攻击，愤怒与不安，永远活在恐惧中。因为作恶的人必期待同样的恶报。你难道看不出你是骚乱而非和平的领袖，是不义而非正直的领袖吗？但我假装生气，不是因为不能证明灵魂不死，而是因为我怜悯你企图欺骗的灵魂。因此我要说话，但不是因你逼我；因为我知道该如何说话，而在这事上，你所需要的不是说服，而是警示。但对于那些真不知此事的人，我将按适当的方式教导他们。}\par

% structure: p0066 source=h2 target=subsection reason=relative-heading-rank; base=h2

\subsection{第43章 \Person{西满}的遁词}

于是\Person{西满}说：\Quote{你若生气，我就不再问你任何问题，也不愿听你说了。}\Person{伯多禄}说：\Quote{如果你现在找借口逃离，你完全自由，无需用任何特别的借口。因为所有人都听见你说了许多错话，并且看出你什么也不能证明，你只是为反驳而提问——这是任何人都会做的事。因为当最清晰的证据被提出后，回答\InnerQuote{你说的毫无用处}有什么难处呢？但为使你知道，我能用一句话向你证明灵魂不死，我要问你一个众人皆知的问题；请回答我，我就用一句话证明它不死。}\Person{西满}原本以为从\Person{伯多禄}的怒气中找到了离开的借口，但因这显著的承诺而停下，说：\Quote{那么你问我，我会回答你所问的众人皆知之事，好让我如你所应许的，用一句话听到灵魂如何不死。}\par

% structure: p0068 source=h2 target=subsection reason=relative-heading-rank; base=h2

\subsection{第44章 看见还是听见？}

伯多禄说：\Quote{我这样说，是为了在众人面前向你们证明。因此，请回答我，两种方法中哪一种更能说服一个不信的人：看见还是听见？} 西满说：\Quote{看见。} 伯多禄说：\Quote{那么，你为何想从我的话中学习那已由事物本身和亲眼看见向你证明的事呢？} 西满说：\Quote{我不明白你的意思。} 伯多禄说：\Quote{你若不明白，现在就去你家，进入内室，你会看见那里放着一尊像，是一个被杀害的男孩的形像，穿着紫袍；你问他，他会用听觉或视觉告诉你。既然你看见它站在你面前，又何必听他讲灵魂是否不死呢？因为若它不存在，确实不可能被看见。但若你不知道我说的是哪尊像，我们就立刻带着在场的另外十个人去你家。}\par

% structure: p0070 source=h2 target=subsection reason=relative-heading-rank; base=h2

\subsection{第四十五章 直刺要害。}

但西满听了这话，良心受责，面色苍白，毫无血色；因为他害怕，若否认此事，他的房子会被搜查，或者伯多禄愤怒之下会更公开地揭露他，于是所有人都会知道他的真面目。于是他回答说：\Quote{我恳求你，伯多禄，藉着你内那良善的天主，克服我内的邪恶。收纳我悔改，你便会有一个助手协助你宣讲。因为现在我确实明白了你是真天主的先知，因此唯独你知道人隐秘和隐藏的事。} 伯多禄说：\Quote{诸位弟兄，你们看见西满寻求悔改；但片刻之后，你们会看见他又回到他的不信。因为他以为我是先知，既然我揭露了他以为隐秘隐藏的恶行，他就许诺要悔改。但我不可以说谎，也不可欺骗，无论这个不信者是否得救。因为我呼天地作证，我所说的并非藉着先知的神，我所暗示的，凡可能的，是对听讲的群众说的；而是我从一些曾是他同伙、如今已归信我们信仰的人那里，得知他暗中所行的事。因此，我说的我所知道的，并非我预知的。}\par

% structure: p0072 source=h2 target=subsection reason=relative-heading-rank; base=h2

\subsection{第四十六章 西满的愤怒。}

但西满听了这话，就咒骂并责难伯多禄，说：\Quote{你这个最邪恶、最诡诈的人，是命运而非真理把胜利给了你。我寻求悔改，并非因为知识欠缺，而是为了让你以为我会因悔改而成为你的门徒，从而把你的行业的所有秘密都托付给我，这样我最终知道这一切，就可以驳倒你。但你狡猾地明白了为什么我假装忏悔，并且默许，好像你不知道我的计谋，为的是先当着百姓的面揭露我是无能的，然后预见我被这样暴露给群众后，必然愤怒，并承认我并非真心忏悔，你抢先一步，以便你可以说，我忏悔之后会再次回到我的不信，这样你就似乎在各方面都获胜了，无论我继续我声称的忏悔还是不继续；于是你就会被认为有智慧，因为你预见了这些事，而我则显得受骗，因为我没有预见你的诡计。但你预见我的，却用诡计胜了我。但我说过，你的胜利来自命运，而非真理：可我知道为什么我没有预见这一点：因为我站在你身边，以我的良善与你说话，并忍耐了你。但现在我要向你显示我神性的力量，使你立刻俯伏朝拜我。}\par

% structure: p0074 source=h2 target=subsection reason=relative-heading-rank; base=h2

\subsection{第四十七章 西满的夸耀。}

\Quote{我是首要能力\OriginalTerm{首要能力}{first power}，我始终存在，且无始\OriginalTerm{无始}{without beginning}。但我进入了辣黑耳的子宫，由她而生为人，使我可见于世人。我飞越空中；我曾与火混合，成为一体；我使雕像移动；我赋予无生命之物生气；我使石头变成饼；我从山飞到山；我由天使之手扶持，从一处移到另一处，落在地上。我不但做了这些事，而且现在仍能如此行，为用事实向所有人证明，我是天主子，永恒长存，并且我能使相信我的人也永远如此存续。但你的话全是虚空的；你不能行任何真实的工作，\Emph{如我刚刚所提的}，那派遣你的人也是一个术士，他甚至连自己都不能从十字架的苦难中拯救出来。}\par

% structure: p0076 source=h2 target=subsection reason=relative-heading-rank; base=h2

\subsection{第四十八章 企图制造骚乱。}

对于西满的这番话，伯多禄回答说：\Quote{不要插手属于别人的事；因为你是术士，你已经承认，并由你所行的事显明了；但我们的主，祂是天主子和人之子\OriginalTerm{天主子和人之子}{Son of God and of man}，显然是良善的；祂真是天主子，这话已经告诉过那些人，并要告诉那些适合知道的人。但若你不承认你是术士，我们就和这所有群众一起去你的家，那时谁才是术士便会显明。} 伯多禄这样说的时候，西满开始以亵渎和咒骂攻击他，为了制造骚乱，激动众人，使他无法被驳倒，并使伯多禄因这亵渎而退避，似乎被战胜了。但伯多禄站稳了，开始更激烈地指责他。\par

% structure: p0078 source=h2 target=subsection reason=relative-heading-rank; base=h2

\subsection{第四十九章 西满的撤退。}

于是，民众愤慨地将\Person{西满}逐出法庭，把他赶出大门，只有一个人在他被赶出去时跟随了他。然后，静默下来后，\Person{伯多禄}开始这样对民众说话：\Quote{你们应当耐心忍受恶人，因为知道即使天主能剪除他们，却仍容忍他们留存，直到那为所有人施行审判的指定之日。既然如此，我们为何不能容忍天主所容忍的人呢？为何不能勇敢地忍受他们对我们所做的恶事，当那全能者并不向他们报复，好使祂自己的良善和恶人的不虔敬得以显明呢？但如果那恶者没有找到\Person{西满}作他的仆人，他无疑会找到别人，因为在这世上，绊倒人的事是免不了的，\InnerQuote{但是那人的确有祸了，因为绊倒人的事从他而来。}\BibleRef{玛 18:7}因此，\Person{西满}更值得哀悼，因为他已成为恶者的上选器皿，这无疑不会发生，除非他因从前的罪获得了控制他的权柄。我何必再说他曾信了我们的\Person{耶稣}，并相信灵魂是不朽的呢？\BibleRef{宗 8:13}虽然在这事上他被魔鬼欺骗，但他却说服自己，他有一个被杀男孩的灵魂供他随意驱使；在这方面，如我所说，他确实被魔鬼欺骗了，因此我照他自己的观念对他说话，因为他从犹太人那里学到，审判和报复要临到那些反对真信德且不悔改的人。但这里有些人，因作恶满盈，恶者向他们显现，为要欺骗他们，使他们永不回头悔改。}\par

% structure: p0080 source=h2 target=subsection reason=relative-heading-rank; base=h2

\subsection{第五十章 伯多禄的祝福}

\Quote{你们这些藉着悔改转向主的人，要向祂屈膝。}当他说了这话，全体群众都向天主屈膝；\Person{伯多禄}举目望天，含泪为他们祈祷，求天主因祂的良善，恩准接纳那些投奔祂的人。祈祷后，他嘱咐他们次日清早来聚会，便遣散了群众。然后，依照习惯，我们进食后便就寝了。\par

% structure: p0082 source=h2 target=subsection reason=relative-heading-rank; base=h2

\subsection{第五十一章 伯多禄的平易近人}

于是，\Person{伯多禄}在惯常的夜间时分起身，发现我们醒着。他以惯常的方式问候我们，坐下后，\Person{尼刻塔}先说道：\Quote{我的主\Person{伯多禄}，如果您允许，我想问您一件事。}\Person{伯多禄}说：\Quote{我不只允许你，也允许所有的人；不但现在，而且随时，每个人都可以说出心中所感、思虑所痛之处，好能获得医治。因为隐藏在沉默中、不让我们知道的事，如同久病，难以治愈。因此，既然适时而必要的教导之药难以施与沉默的人，每个人都应当宣言自己内心因无知而软弱的方面。但对保持沉默的人，只有天主能给予医治。我们确实也能做到，但需要相当长的时间。因为教义的谈论必须从起初循序渐进，应对每一个问题，揭示一切，解析一切，触及每个人心中所需求的一切；但如我所说，这只能在长期的过程中实现。现在，请随意提问。}\par

% structure: p0084 source=h2 target=subsection reason=relative-heading-rank; base=h2

\subsection{第五十二章 虚假标记与奇迹}

\Person{尼刻塔}说：\Quote{至为仁慈的\Person{伯多禄}，我衷心感谢您；但我想知道的是，\Person{西满}——天主的敌人，如何能行如此伟大之事？因为他所宣称的，确实没有谎言。}对此，真福\Person{伯多禄}回答说：\Quote{唯一真实的天主，决意为祂的首生子预备良善而忠信的朋友。但祂知道，除非人拥有凭以成为良善的感知能力，否则无人能成为良善；人必须出于自己的意愿选择成为良善，否则若因必然而非自愿而被拘于良善，就不能真正良善。因此，祂赋予每个人自主意志的能力，使人可以成为自己所愿成为的。同时，祂预见到这意志的能力会使一些人选择善事，另一些人选择恶事，从而人类必然分为两类，于是祂容许每类选择自己所愿的地方和君王。因为良善的君王喜悦良善，而恶者喜悦邪恶。\Person{克肋孟}啊，虽然我在那篇论预定与终局的讲稿中已经更充分地向你阐述过这些，但现在既然\Person{尼刻塔}问我，我也应当向他澄清，为什么\Person{西满}——心思与天主为敌——能行如此奇事。}\par

% structure: p0086 source=h2 target=subsection reason=relative-heading-rank; base=h2

\subsection{第五十三章 自爱——良善的基础}

首先，在天主的眼中，那不愿探寻对自己有益之事的人便是恶人。因为如果一个人不爱自己，怎能爱别人呢？或者，那不能做自己朋友的人，不会成为所有人的敌人吗？因此，为了在择善与择恶者之间有所区分，天主将有益于人的事，即天国的拥有，隐藏起来，像秘藏的宝藏一样埋藏起来，使人不能凭自己的力量或知识轻易获得。然而，祂通过历代的传承，以各种名称和意见，把它的消息传到所有人的耳中：这样，凡是爱好善良的人，听到它就会探寻并发现对自己有益和得救的事；但他们不应向自己求问，而应向那隐藏它的天主求问，并应祈祷，求赐给他们接近和认识的道路；这条道路只向那些爱它超过世上一切美善的人敞开；否则，无论如何聪明，任何人都不能理解它；而那些忽略探寻对自己有益和得救之事的人，如同自恨者和自敌者，将作为爱恶事者而被剥夺它的美善。\par

% structure: p0088 source=h2 target=subsection reason=relative-heading-rank; base=h2

\subsection{第五十四章  当至爱天主}

\Quote{因此，善人应当爱这\Emph{道路}超过一切，即超过财富、荣耀、休息、父母、亲属、朋友和世上的一切。但那完全爱这天国拥有的人，必将抛弃一切恶习、疏忽、懒惰、恶意、恼怒等情愫。因为如果你将任何这些置于它之上，贪爱自己私欲的恶习超过爱天主，你就不能获得天国的拥有；因为爱任何事物超过爱天主，实在是愚蠢的。因为无论是父母，他们会死去；或是亲属，他们不会持久；或是朋友，他们会改变。唯独天主是永恒的，恒久不变。因此，那不愿寻求对自己有益之事的人，是恶的，甚至他的邪恶超过了不敬之人的首领。因为他滥用天主的仁慈来达到自己的邪恶目的，并自得其乐；而前者却忽略了自己救恩的美善，以自我毁灭来取悦恶者。}\par

% structure: p0090 source=h2 target=subsection reason=relative-heading-rank; base=h2

\subsection{第五十五章  十诫对应于埃及的灾祸}

因此，为了那些因忽略自己救恩而取悦恶者的人，以及那些因寻求自己利益而取悦善者的人，按照降在埃及的十灾之数，规定了十件事作为这现世的考验。因为当梅瑟依照天主的命令向法郎要求释放百姓时，并显了征兆作为他受天主差遣的标记，他把棍杖丢在地上，它就变成了蛇。当法郎因有自由意志而不能因此同意时，由于天主的许可，术士们似乎也行出相似的征兆，以使国王的意图可以从其自由意志中得到证明，看他是否宁愿相信由天主派遣的梅瑟所行的征兆，还是相信那些术士似乎行出、实则并非真正行使的征兆。因为的确，他本应从他们的名称就知道他们不是真理的施行者，因为他们没有被称作天主的使者，而是术士，正如传统所暗示的。此外，他们似乎在某种程度上维持了竞争，但后来他们自己承认并屈服于更强者。\BibleRef{出 8:19}因此，最后一灾降临了，\EditorialNote{出 12}就是长子之死，然后梅瑟受命用洒血来祝圣百姓；这样，献上礼物后，他被再三恳求与百姓一同离开。\par

% structure: p0092 source=h2 target=subsection reason=relative-heading-rank; base=h2

\subsection{第五十六章  西满抗拒伯多禄，如同术士对抗梅瑟}

\Quote{我看出我目前正从事于一桩类似的事。因为正如那时，当梅瑟劝国王相信天主时，术士们用假装行使相似征兆的方式反对他，从而阻止不信者得救；同样，现在当我出来教导万民信仰真天主时，术士西满抗拒我，反对我，正如他们反对梅瑟一样；以便那些在外邦人中不用正确判断的人得以显明，而那正确辨别征兆的人则可得救。}伯多禄这样说时，尼塞塔回答说：\Quote{我恳求你允许我说出我心中所想的一切。}伯多禄见门徒们如此热切，高兴地说：\Quote{你随意说吧。}\par

% structure: p0094 source=h2 target=subsection reason=relative-heading-rank; base=h2

\subsection{第五十七章  术士的奇迹}

尼塞塔于是说：\Quote{埃及人在不相信梅瑟这件事上有什么罪呢？因为术士也行出了相似的征兆，尽管他们行出的是外表而非真实。因为如果我当时在场，看到术士行了与梅瑟相似的作为，我岂不是会认为梅瑟是术士，或者术士是受天主差遣而行征兆的吗？因为我不会认为，即使是外表上，术士也能行出与天主所差遣者所行之同样的作为。那么，现在那些相信西满的人有什么罪呢？因为他们看到他行了如此大的奇事。在天上飞行、与火交融成为一体、使雕像行走、使铜狗吠叫，以及其他类似之事，对于那些不知道如何辨别的人而言，岂不是很奇妙吗？是的，也有人见过他用石头做面包。但如果相信那行征兆的人有罪，那么因我们的主所行的征兆和异能而相信祂的人，又如何能显明是无罪的呢？}\par

% structure: p0096 source=h2 target=subsection reason=relative-heading-rank; base=h2

\subsection{第五十八章  以爱覆蔽的真理}

那时\Person{伯多禄}说：\Quote{我很高兴你把真理引到准则，不让信德的阻碍潜伏在你的灵魂里。因为这样你就能轻易获得医治。你可记得我说过，最坏的事就是有人忽略学习对自己有益的事？}\Person{尼塞塔}回答说：\Quote{我记得。}\Person{伯多禄}又说：\Quote{再者，天主隐藏了祂的真理，为要向那些忠心跟随祂的人启示出来？}\Quote{不，}\Person{尼塞塔}说，\Quote{我也没有忘记这件事。}\Person{伯多禄}于是说：那么你怎么想？天主把祂的真理深埋在地里，并堆上山岭，以便只有那些能深挖到深处的人才能找到它？不是这样；反之，正如祂用天穹环绕山岭和大地，祂也用祂自己爱的帷幔遮蔽真理，以便只有先敲响天主之爱大门的人才能达到它。\par

% structure: p0098 source=h2 target=subsection reason=relative-heading-rank; base=h2

\subsection{第五十九章　善与恶成对出现。}

\Quote{因为，正如我开始说的，天主为这个世界指定了某些对；每一对中首先来的属于恶，其次来的属于善。在此，人人都有机会做出正确的判断，无论他是单纯还是谨慎。因为若他是单纯的，并相信先来的那一位，即使是被征兆和奇事所驱动，他也同样必须相信后来的一位；因为他会被征兆和奇事说服，如先前一样。当他相信这第二位时，他会从后者那里学到他不该相信那先来的、属于恶的一位；于是前者的错误被后者的改正所纠正。但如果他因为已经相信了先来的而不愿接受后来的，他将理所当然地被判为不义；因为不义的是，当他因先来的的征兆而相信了他时，却不愿相信后来的，即使后者带来相同甚至更大的征兆。但如果他不相信先来的，那么他可能会被感动去相信后来的。因为他的心智并没有变得完全迟钝，以至于不能被奇迹的重现所唤醒。但若是他谨慎，他就能分辨征兆。如果他真的相信了先来的，他会因奇迹的增加而被感动去趋向第二位，并通过比较而领会哪一个是更好的；虽然明确检验\Emph{奇迹的}为所有有学识的人所识别，正如我们在讨论的正常顺序中已经指出的。但如果有人，作为健全者而不需要医生，没有被先来的所感动，他会因事情本身的持续性而被吸引到后来的，并以下面这种方式分辨征兆和奇迹——属于恶者的人，他所行的奇迹对人毫无益处；但那善者所行的奇迹却是有益于人的。}\par

% structure: p0100 source=h2 target=subsection reason=relative-heading-rank; base=h2

\subsection{第六十章　假奇迹的无效性。}

请告诉我，你们说西满所做的那些事——雕像行走、铜狗或石狗吠叫、山脉跳舞、在空中飞翔以及诸如此类的东西——有什么用处呢？但那属于善者的\Emph{征兆}，是为了人类的益处，如同我们的主所行的那些：祂使瞎子看见，聋子听见，扶起软弱者和瘸子，驱除疾病和魔鬼，复活死者，并行了其他类似的事，正如你们也看见我所行的。因此，那些有益于人类并赐予他们好处的征兆，恶者不能行，除非在世界末日。因为那时将被允许他将一些好的征兆混入他的征兆中，例如驱逐魔鬼或医治疾病；借此越界，并与自己分裂，与自己相争，他终将被消灭。因此，主曾预言，在末世将有如此的诱惑，以致如果可能，连被选者也要被迷惑；也就是说，由于征兆的标记混乱，连那些似乎善于辨别神灵和区分奇迹的人也会被扰乱。\par

% structure: p0102 source=h2 target=subsection reason=relative-heading-rank; base=h2

\subsection{第六十一章　十对。}

\Quote{因此，从太初开始，我们所说的十对被分配给了这个世界。加音和亚伯尔是一对。第二对是巨人和诺厄；第三对是法老和亚巴郎；第四对是培肋舍特人和依撒格；第五对是厄撒乌和雅各伯；第六对是术士和立法者梅瑟；第七对是诱惑者和人子；第八对是西满和我伯多禄；第九对是万民和那将被派遣在万民中播种圣言者；第十对是假基督和基督。关于这些对，我们将在另一时间给你们更充分的讲解。}\Person{伯多禄}这样说时，\Person{阿桂拉}说：\Quote{确实需要不断的教导，好使人能学习关于一切事物的真理。}\par

% structure: p0104 source=h2 target=subsection reason=relative-heading-rank; base=h2

\subsection{第六十二章　基督徒的生活。}

伯多禄却说：\Quote{谁是这样向往训诲、勤恳探究每一个细节的人呢？岂非那爱惜自己的灵魂以致得救，并弃绝此世一切事务，以便仅有闲暇专注于天主圣言的人？这样的人才是真先知唯独视为有智慧的人，甚至是那变卖所有、买那颗真珍珠的人，\BibleRef{玛 13:46}他明白暂存之物与永恒、微小与伟大、人与天主的差别。因为他懂得在真实而良善的天主面前永生的希望是什么。然而，除了认识天主智慧的人，谁能爱天主？人若不恒心聆听祂的圣言，怎能获得天主的智慧？由此，他生发对祂的爱慕，并以相称的尊荣敬拜祂，向祂献上赞美诗和祈祷，且最愉快地安息其中，若在任何时刻说或做任何其他事，甚至片刻，他都视为最大的损失；因为，事实上，充满天主之爱的灵魂，除了与天主有关的事，不能再顾及其他，而且因着对祂的爱，默思那些它知道取悦祂的事，也无法满足。但那些没有对祂生发情感，也未在心灵中点燃祂之爱的人，仿佛被置于黑暗之中，不能看见光明；因此，甚至在他们开始学习任何关于天主的事之前，就立刻像因劳苦而疲惫一样昏厥；且充满厌倦，他们随即被各自的习惯急忙带到那些令他们愉悦的言语上。因为对于这样的人，听任何关于天主的事都是累赘和烦扰；这正是我所说的原因，因为他们的心灵未曾领受神圣之爱的甘甜。}\par

% structure: p0106 source=h2 target=subsection reason=relative-heading-rank; base=h2

\subsection{第六十三章 来自西满营地的逃兵}

伯多禄正在这样说时，天亮了；看，西满的一个门徒前来呼喊说：\Quote{我恳求你，伯多禄，收纳我这个可怜的人，我被西满术士欺骗了，我曾将他当作一位天上的天主来信从，因我看见他所行的奇迹。但当我听了你的讲论，便开始认为他是一个人，而且是一个邪恶的人；然而，当他离开这里时，我独自跟随了他，因为我还未清楚看出他的不敬。但当他看见我跟从他，就称我为有福的，领我到他家；约在半夜，他对我说：\InnerQuote{我要使你比众人更好，只要你一直与我相处，直到最后。}我答应了他，他便要求我起誓坚持；得到誓愿后，他把一些受污染、受诅咒的秘密之物放在我肩上，让我背负着，并命令我跟随他。但我们到了海边，他上了刚好在那里的一条船，从我颈上取下他命令我背负之物。稍后他出来，两手空空，必定是将其丢入海中。于是他请我与他同去，说他要去罗马，在那里他将如此取悦众人，以至于被尊为神，公开获得神的尊荣。\InnerQuote{那时，}他说，\InnerQuote{若你愿意返回这里，我会送你回去，满载一切财富，并由各种服务扶持。}我听到这话，见在他身上并无任何与这话相符之处，反而看出他是个术士和骗子，便回答：\InnerQuote{求你宽恕我；因我脚痛，不能离开凯撒勒雅。此外，我有妻子和幼小的儿女，决不能留下他们。}他听了这话，责怪我懒惰，就向多辣出发，说：\InnerQuote{你将会后悔，当你听到我在罗马城获得怎样的荣耀。}此后他出发去了罗马，正如他所说的；但我急忙返回这里，恳求你接纳我悔改，因为我被他欺骗了。}\par

% structure: p0108 source=h2 target=subsection reason=relative-heading-rank; base=h2

\subsection{第六十四章 关于西满邪恶的宣告}

当那个从西满那里回来的人这样说了以后，伯多禄吩咐他坐在院子里。他自己出去，看见大批群众，远超前几日，便站在他惯常的位置上，指着那个来的人，开始讲话如下：\Quote{诸位弟兄，我指给你们的这个人，刚刚来到我这里，告诉我关于西满的邪恶行径，以及他如何将邪恶的工具投入海中，这不是出于悔改，而是害怕被发觉而受公法惩处。他请求这个人，据他告诉我，与他同住，并应许巨额的赠礼；当他无法说服他这样做时，就离开了他，责骂他懒惰，于是出发往罗马去了。}伯多禄向群众示意此事后，那个从西满回来的本人就站起来，开始向众人陈述关于西满罪行的一切。当人们听到西满用魔法所做之事而感到震惊时，伯多禄说：\par

% structure: p0110 source=h2 target=subsection reason=relative-heading-rank; base=h2

\subsection{第六十五章 伯多禄决定跟随西满}

\Quote{我的弟兄们，不要因所发生的事而苦恼，而要留意将来；因为过去的事已经结束，但那些威胁的事对于将要遭遇的人却是危险的。罪恶在这世界上永远不会缺乏，只要敌人被允许按他的意愿行事；为使那些精明和了解他诡计的人，在他所兴起的争战中得胜；但那些忽略学习关乎自己灵魂救恩之事的人，可能被他用应得的欺骗所掳获。因此，正如你们所听见的，西满已经出去抢先占据那些被召得救的外邦人的耳朵，我也必须跟随他的踪迹，以便他所挑起的一切争论都能由我们纠正。然而，既然对你们这些已经进入生命墙垣之内的人应该更加忧虑——因为如果已经实际获得的丧亡了，就遭受了实在的损失；而对于尚未获得的，如果能得到，就有那么多收获；但如果得不到，唯一的损失就是没有增益——因此，为使你们在真理中更加坚定，并使被召得救的万民不致因西满的邪恶而受阻碍，我决定委任匝凯作你们的牧者，我自己与你们再留三个月；然后去往外邦人那里，免得我们耽延太久，而西满的恶行在各处横行，以致他们变得不可救药。}\par

% structure: p0112 source=h2 target=subsection reason=relative-heading-rank; base=h2

\subsection{第六十六章 匝凯被立为凯撒勒雅主教；祝圣长老与执事}

听到这宣布，全体民众都哭了，因为他要离开他们；伯多禄同情他们，自己也流泪了。他举目望天，说：\Quote{天主啊，祢创造了天地和其中的万物，我们向祢献上恳求的祈祷，求祢安慰那些在苦难中投靠祢的人。因为他们对祢的爱，使他们爱我，因我向他们宣讲了祢的真理。因此，求祢用祢慈悲的右手护佑他们；因为匝凯或任何其他人都不能作他们足够的守护者。}说了这话及类似的话后，他给匝凯覆手，祈祷他能无可指摘地履行主教职责。然后他祝圣了十二位长老和四位执事，并说：\Quote{我祝圣了这匝凯为主教，因为我知道他敬畏天主，精通圣经。因此，你们应当尊敬他，如尊敬基督的代表，为你们的救恩而服从他，并要知道凡是对他所行的尊敬或伤害，都归于基督，并由基督归于天主。因此，要全神贯注地听他，从他那里接受信德的道理；从长老那里接受生活的训诫；从执事那里接受纪律的秩序。要虔诚地照顾寡妇，积极帮助孤儿，怜悯穷人，教导年轻人谦逊——总之，按情况需要彼此支持。敬拜创造天地的天主，相信基督，彼此相爱，对所有人慈悲为怀，不只以言语，也以行动和行为实践爱德。}\par

% structure: p0114 source=h2 target=subsection reason=relative-heading-rank; base=h2

\subsection{第六十七章 邀请接受圣洗}

当他给了他们这些及类似的训诫后，他向民众宣布说：\Quote{既然我决定与你们再留三个月，若有人愿意，就让他受洗；使他脱去先前的恶，今后因自己的行为成为天上祝福的继承者，作为善行的赏报。因此，凡愿意的，让他到匝凯那里，向他报上名字，从他那里听天国的奥秘。让他常守斋戒，在一切事上证明自己，以便在三个月结束时，在庆节日受洗。你们每人要受洗于活水中，并呼求三倍真福的名号；先被以祈祷祝圣的油傅油，如此，因这些事而被祝圣后，才能获得对神圣事物的感知。}\par

% structure: p0116 source=h2 target=subsection reason=relative-heading-rank; base=h2

\subsection{第六十八章 先他而去的十二人}

当他详细谈论了洗礼后，就遣散了群众，回到他惯常的住所。在那里，十二人站在他周围（即匝凯、索佛尼亚、若瑟、米该亚、厄肋阿匝尔、丕乃哈斯、拉匝禄、厄里叟、我，克肋孟、尼苛德摩、尼塞达、阿桂拉），他向我们说了以下的话：\Quote{我的弟兄们，让我们思考什么是正确的；因为给那些被召得救的万民提供一些帮助是我们的责任。你们自己已经听说西满出发了，想要抢在我们前面。我们本应步步紧跟，使他在任何地方试图颠覆人时，我们能立即驳斥他。但既然在我看来，遗弃那些已经归向天主的人，而将我们的关怀给予那些尚在远方的人，是不公义的；我认为我应当与这城中已归信的人再留三个月，坚固他们；然而我们也不该忽略那些尚在远方的人，免得他们若长期受有害教义的影响，就更难恢复。因此，我希望（但只如果你们也认为合适），为匝凯（我们现在已祝圣为主教）有撒巴之子本雅明代替；为克肋孟（我决定把他常留身边，因为他来自外邦人，极愿听天主的圣言）有撒弗辣之子阿纳尼雅代替；为尼塞达和阿桂拉（他们最近才归信基督信仰）有匝凯的兄弟鲁贝鲁斯和建筑者匝加利亚代替。因此，我愿意用这四个人代替那四个人，完成十二人的数目，使西满感到我藉着他们常与他同在。}\par

% structure: p0118 source=h2 target=subsection reason=relative-heading-rank; base=h2

\subsection{第六十九章 安排得到众弟兄的赞同}

于是，\Person{克莱孟}、\Person{尼塞塔}和\Person{阿奎拉}被分开后，他对那十二人说：\Quote{我希望你们后天前往外邦人那里，跟踪\Person{西满}的足迹，以便将他的所作所为报告给我。你们也要勤问每个人的意见，并告诉他们我将毫不迟延地到他们那里；简言之，在所有地方都要教导外邦人期待我的到来。}他说了这些话以及类似的话之后，又说：\Quote{你们，我的弟兄们，如果对这些事有什么要说的，就说吧，免得只有我认为正确的事可能并非正确。}于是，众人齐声称赞他说：\Quote{我们反倒请求你按自己的判断安排一切，吩咐你认为合适的事；因为我们认为，若我们完成你所命令的，这才算是完美的虔敬之工。}\par

% structure: p0120 source=h2 target=subsection reason=relative-heading-rank; base=h2

\subsection{第七十章 十二人的出发}

因此，到了约定的日子，他们站在\Person{伯多禄}面前说：\Quote{\Person{伯多禄}，不要以为我们因三个月不能听你讲论而失去的特恩是微不足道的；但既然听从你的命令对我们有益，我们将欣然服从。我们心中将永远保存你面容的记忆，就这样我们活跃地出发，正如你命令我们的。}于是，他为主祈祷后，便打发他们走了。那十二位被派往前去的人离开后，\Person{伯多禄}照常进入，站在辩论的地方。一大群人聚集而来，甚至比平时还多；众人都流泪注视着他，因为前一天听他说要因\Person{西满}的缘故而离开。看到他们哭泣，他自己也受了感动，虽然努力掩饰并抑制泪水。但他声音的颤抖和话语的中断，暴露了他同样被情感所困扰。\par

% structure: p0122 source=h2 target=subsection reason=relative-heading-rank; base=h2

\subsection{第七十一章 \Person{伯多禄}为离开准备\Place{凯撒勒雅}人}

然而，他用手指擦了擦额头，说道：\Quote{我的弟兄们，要鼓起勇气，用劝慰安慰你们忧伤的心，将一切交托于天主，在任何事上都只应承行并优先选择祂的旨意。让我们暂且假设，出于对你们的爱，我们违背祂的旨意，与你们同留，难道祂不能藉着使死亡降临于我，安排我与你们更久的分离吗？因此，我们更好按祂的旨意完成这短暂的分离，作为那些被命令在一切事上服从天主的人。同样，你们也应怀着同样的顺服服从祂，因为你们爱我，无非是因为爱祂的缘故。因此，作为天主的朋友，要顺服祂的旨意；但你们也要自行判断何为正确。如果当\Person{西满}欺骗你们时，我因弟兄们在\Place{耶路撒冷}的挽留而未能来到你们那里，尽管你们中间有善良而善于言辞的\Person{匝凯}，这难道不显得邪恶吗？同样，现在也当认为，如果\Person{西满}已经出去攻击完全没有守护者的外邦人，而你们留住我，我不去追赶他，那就是邪恶的。所以，让我们注意，不要因不合理的感情而去成就恶者的意愿。}\par

% structure: p0124 source=h2 target=subsection reason=relative-heading-rank; base=h2

\subsection{第七十二章 超过一万人受洗}

\Quote{在这期间，我要照我所应许的，与你们同留三个月。你们要恒心听道；期满时，若有人能够并愿意跟随我们，只要职责容许，他们可以这样做。我说\Emph{只要职责容许}，意思是，谁也不可因自己的离开而令不应忧伤的人忧伤，比如离开不应离开的父母，或忠信的妻子，或任何其他为了天主的缘故他应该给予安慰的人。}在此期间，他日复一日地辩论和教导，以教导之劳填补了所定的时间；当节期到来时，超过一万人受了洗。\par

% structure: p0126 source=h2 target=subsection reason=relative-heading-rank; base=h2

\subsection{第七十三章 \Person{西满}的消息}

但是，那些日子里，收到了先前出发的弟兄们的一封信，信中详述了\Person{西满}的罪行：他如何从一城到另一城欺骗众人，到处诽谤\Person{伯多禄}，以至于当\Person{伯多禄}到来时，没有人肯听他讲。因为他断言\Person{伯多禄}是术士、不敬神的人、有害、狡猾、无知，并声称不可能的事。\Quote{因为，\InnerQuote{他说死人会复活，这是不可能的。但若有人试图驳斥他，他就会通过随从用秘密的圈套将他除掉。因此，我也，}他说，\InnerQuote{当我击败他并胜过他之后，因害怕他的圈套而逃跑，免得他用咒语毁灭我，或用阴谋置我于死地。}}\par

% structure: p0128 source=h2 target=subsection reason=relative-heading-rank; base=h2

\subsection{第七十四章 告别\Place{凯撒勒雅}}

\Person{伯多禄}于是吩咐把信读给民众听。读完信后，他向他们讲话，就一切事给了他们充分的指示，尤其要他们服从\Person{匝凯}，因为他已祝圣他为他们的主教。他也将长老和执事托付给民众，同样也将民众托付给他们。然后，他宣布要在\Place{的黎波里}过冬，说道：\Quote{我把你们托付给天主的恩宠，明天，若天主愿意，我就要出发了。}但是，在\Place{凯撒勒雅}停留的整整三个月里，为了教导之故，他白天在民众面前所谈论的一切，到了夜里更详尽、更完美地私下向我们解释，因为我们更为可靠，也完全蒙他认可。同时，他吩咐我——因为他知道我仔细地把所听见的记在记忆里——将任何值得记录的事写下来，并送给你——我的主上\Person{雅各伯}，我也遵命照做了。\par

% structure: p0130 source=h2 target=subsection reason=relative-heading-rank; base=h2

\subsection{第七十五章 \Person{克莱孟}致\Person{雅各伯}信函的内容}

因此，我先前寄给您的第一卷书，包含关于真先知以及按照\Person{梅瑟}传统所教导的理解法律的独特性的记述。第二卷包含关于元始的记述，以及元始是一个还是多个，并且希伯来人的法律知道什么是无限。第三卷关于天主以及由祂所规定的事物。第四卷，虽有众神之称，但根据圣经的见证，唯有一位真天主。第五卷，有诸天两层，其一为可见的穹苍，将要消逝，另一为永恒且不可见的。第六卷关于善与恶；万物皆由圣父归服于善；以及恶是什么，为何存在，如何存在，且它怎样与善合作，但非出于善的目的；以及善的标记是什么，恶的标记是什么；以及双重与结合之间的区别是什么。第七卷，十二宗徒在圣殿中当着民众的面所讨论的事。第八卷，关于主的话语，看似矛盾实则不然；以及它们的解释是什么。第九卷，由天主所赐的法律是公义且完美的，并且唯有它能使人洁净。第十卷，关于人类肉身的出生，以及藉着圣洗而生的世代；并且人内肉身种子的传承是什么；以及关于他的灵魂的说明，以及自由意志如何在其中，既然它不是非受生的，而是受造的，就不能从善坚定不移。因此，对于这些不同的主题，凡\Person{伯多禄}在\Place{凯撒勒雅}所讲述的，遵照他的命令，如我所说，我已将其写定成十卷书寄给您。但在次日，按既定计划，我们与一些决心陪同\Person{伯多禄}的忠信之人从\Place{凯撒勒雅}出发了。\par

\SourceRecord{Translated by Thomas Smith.；Ante-Nicene Fathers；Vol. 8.；Edited by Alexander Roberts, James Donaldson, and A. Cleveland Coxe.；Buffalo, NY: Christian Literature Publishing Co.,；1886.；<http://www.newadvent.org/fathers/080403.htm>.}

% structure: p0001 source=h1 target=section reason=page-title

\section{\WorkTitle{认语录（卷四）}}

% structure: p0002 source=h2 target=subsection reason=relative-heading-rank; base=h2

\subsection{第一章 在多辣歇脚}

从凯撒勒雅出发前往特里玻里斯的路上，我们在一个叫做多辣的小镇作了第一次停留，因为它不远；几乎所有因伯多禄的讲道而相信的人，都难以与他分离，而是与我们同行，一次又一次地注视他，一次又一次地拥抱他，一次又一次地与他交谈，直到我们来到旅店。次日我们到了仆托肋买，在那里住了十天；当相当多的人接受了天主的圣言时，我们向其中一些似乎特别留心且希望我们多留几日以便教导的人示意，如果他们愿意，可以跟随我们到特里玻里斯。我们在提洛、漆冬和贝里图斯也做了同样的事，并向那些渴望听更多讲论的人宣告，我们要在特里玻里斯过冬。因此，那些渴望的人从各城跟随伯多禄，当我们进入特里玻里斯时，我们是一大群蒙选的人。到达时，先前被派去的弟兄们在城门前迎接我们；他们带领我们，把我们带到他预先准备好的各个住处。于是城中起了骚动，一大群渴望见伯多禄的人聚集起来。\par

% structure: p0004 source=h2 target=subsection reason=relative-heading-rank; base=h2

\subsection{第二章 在马罗家的接待}

当我们来到马罗的家——那里已为伯多禄预备好——时，他转向人群，告诉他们他将在后天对他们讲话。于是先前被派去的弟兄们为所有跟我们一起来的人分配了住处。然后，伯多禄进入马罗的家，被请求进食时，他回答说，除非他先确定所有陪伴他的人都已安顿好住处，否则决不进食。然后他从先前被派去的弟兄们那里得知，市民们不仅热情接待了他们，而且由于对伯多禄的爱，对他们非常友善；以至于有几个人因为家里没有客人而感到失望；因为所有人都做了准备，即使来了更多的人，宿主们也仍会缺少客人，而不是客人缺少宿主。\par

% structure: p0006 source=h2 target=subsection reason=relative-heading-rank; base=h2

\subsection{第三章 西满的逃亡}

于是伯多禄非常高兴，称赞弟兄们，祝福他们，并请他们与他同住。然后，他在海里洗了澡，吃了食物，晚上就睡了；照常，在鸡叫时起身，而晚上的灯仍在亮着，他发现我们都醒着。我们总共十六人，即伯多禄和我克莱孟，尼塞塔和阿奎拉，以及先前的那十二人。于是，伯多禄照常向我们致意后，说：\Quote{既然今天我们没有被别人占用，就让我们自己占用自己吧。我将告诉你们你们离开后凯撒勒雅发生的事，而你们要告诉我们西满在这里的所作所为。} 当这些话题正在谈论时，天亮时，一些家人进来告诉伯多禄，西满一听说伯多禄到达，就连夜离开了，前往叙利亚的路上。他们还报告说，群众认为他所说的要间隔一天的时间对他们的爱来说是太长了，他们不耐烦地站在门前，互相谈论他们渴望听到的事，他们希望无论如何能在约定的时间之前见到他；而且随着天色渐亮，人群正在增加，他们自信地相信，无论他们可能假设什么，他们会听到从他而来的讲论。\Quote{那么，} 他们说，\Quote{请指示我们告诉他们你认为合适的；因为有这么多的人聚集在一起，却因为没有得到任何回答而忧伤离去，这是不合理的。因为他们不会认为是他们没有等待指定的日子，反而会认为你轻视了他们。}\par

% structure: p0008 source=h2 target=subsection reason=relative-heading-rank; base=h2

\subsection{第四章 庄稼多}

于是伯多禄满心赞叹，说：弟兄们，你们看，主预言性地所说的每一句话是如何应验的。因为我记得他说过：\BibleQuote{庄稼的确多，但工人少；所以你们应当求庄稼的主，派工人来收割他的庄稼。} \BibleRef{玛 9:37} 看，因此，那在奥秘中预言的事应验了。然而他还说过：\BibleQuote{有许多人要从东方和西方，从北方和南方来，在亚巴郎、依撒格和雅各伯的怀中坐席。} 这也正如你们所看见的，同样应验了。因此，我恳求你们，我的同仆和助手，要勤恳学习讲道的次序和赦免的方式，使你们能拯救人的灵魂，他们因天主的隐秘能力在受教之前就已认出他们应当爱谁。因为你们看见，这些人像忠仆一样，渴望那位他们期待向他们宣告他们主人来临的人，以便他们知道后能完成他的旨意。因此，听天主的圣言并探究他旨意的渴望，是来自天主的；这是赐给外邦人的天主恩赐的开始，使他们藉此能接受真理的教导。\par

% structure: p0010 source=h2 target=subsection reason=relative-heading-rank; base=h2

\subsection{第五章 梅瑟与基督}

因为从起初就是这样赐给了希伯来民族，叫他们爱梅瑟，信他的话；因此也记载说：\BibleQuote{百姓信了天主，也信了祂的仆人梅瑟。}\BibleRef{出 14:31}所以，那本来是来自天主对希伯来民族的特别恩赐，我们现在看到也赐给了那些从外邦人中蒙召来相信的人。但行为的方法则放在每个人的能力和意愿中，这是他们自己的事；而对真理导师怀有爱慕之情，则是天上圣父的恩赐。但得救在于：你们遵行你们藉天主的恩赐而对其怀有爱情和爱慕的那位的旨意，免得那段话临到你们身上，就是祂所说的：\BibleQuote{你们为什么称呼我主啊主啊，却不遵行我的话呢？}\BibleRef{路 6:46}所以，天主赐给希伯来人的特别恩赐是叫他们相信梅瑟；而赐给外邦人的特别恩赐是叫他们爱耶稣。因为主也暗示了这一点，当祂说：\BibleQuote{父啊，天地的主宰，我称谢祢，因为祢将这些事瞒住了智慧及明达的人，而启示给了小孩子。}这明确地宣告：受法律训导的希伯来民族不认识祂；但外邦民族却承认了耶稣，并敬拜祂；因此他们也将得救，不仅承认祂，而且遵行祂的旨意。但那个外邦人，他由天主得了恩赐相信梅瑟，也该由自己的意愿去爱耶稣。同样，希伯来人，他由天主得了恩赐相信梅瑟，也该由自己的意愿去信耶稣；这样，他们每个人，在自己身上既有天主的恩赐，又有自己的努力，就可以藉两者成为完全的。因为关于这样的人，我们的主曾论及一个富人时说：\BibleQuote{他从自己的宝库里取出新和旧的东西来。}\BibleRef{玛 13:52}\par

% structure: p0012 source=h2 target=subsection reason=relative-heading-rank; base=h2

\subsection{第六章 一个会众}

\Quote{关于这些事说得够多了，因为时间紧迫，而百姓的宗教热诚邀请我们向他们讲话。}他说完这话后，便询问哪里有一个适合讨论的地方。\Person{马罗}说：\Quote{我有一个很宽敞的大厅，能容纳五百多人，而且屋内还有一个花园；或者您喜欢在某个公共场所，大家都会更喜欢，因为没有人不渴望至少看到您的面容。}于是伯多禄说：\Quote{带我去看看那大厅或那花园。}他看了大厅后，又进去看花园；忽然全体群众，好像有人召唤他们似的，冲进屋子，又从那里涌进花园，伯多禄已经站在那里，正在选择一个适合讨论的地方。\par

% structure: p0014 source=h2 target=subsection reason=relative-heading-rank; base=h2

\subsection{第七章 病人痊愈}

但当他看见群众如同大河的水一般涌过狭窄的通道时，他便登上一根恰好立在花园墙边的柱子，先以宗教的方式向众人致意。但有些在场的人，长久受魔鬼折磨，扑倒在地，而魔鬼恳求允许它们留在这占据的身体中哪怕一天。但伯多禄斥责它们，命令它们离开；它们立即出去了。之后，另一些患有长期疾病的人请求伯多禄让他们得到痊愈；他允诺，等他的教理讲演结束后，就为他们祈求主。但当他允诺时，他们立刻从疾病中解脱了；他便吩咐他们与那些从魔鬼手中被释放的人一同坐在一旁，如同工作后的休息。与此同时，这事正在发生时，一大群人聚集了，不仅出于听伯多禄讲话的渴望，也因所成就的治愈的传闻。但伯多禄向众人挥手示意安静，使人群平静下来，开始如下向他们讲话：——\par

% structure: p0016 source=h2 target=subsection reason=relative-heading-rank; base=h2

\subsection{第八章 天主眷顾的辩护}

在我看来，在关于真崇拜天主的讲论开始时，首先有必要教导那些尚未获得任何这方面知识的人：在整个天主眷顾中，必须认为这眷顾是无可指责的，世界由之被统治和管理。此外，当前任务的理由以及因天主的能力而获得治愈的人所提供的场合，提示了这开端的话题，即：要表明，有很好的理由使许多人附有魔鬼，以便天主的公义得以显示。因为无知将被发现是几乎所有邪恶之母。但现在让我们进入理由。\par

% structure: p0018 source=h2 target=subsection reason=relative-heading-rank; base=h2

\subsection{第九章 无罪状态是享乐状态}

当天主按自己的肖像和模样造了人之后，祂将祂神性的一种气息和馨香植入祂的工程中，好使人成为祂独生子的分享者，藉着祂也成为天主的朋友和义子。因此，祂自己作为真先知，知道什么行为中悦圣父，便教导他们如何获得这一特权。那时，在人间只有一种对天主的敬拜——纯洁的心和无玷的精神。为此缘故，受造物都与人类遵守一项不可侵犯的盟约。由于对造物主的敬畏，没有任何疾病、身体的失调或食物的腐败能管辖他们；因此，千年的寿命也不会陷入老年的脆弱。\par

% structure: p0020 source=h2 target=subsection reason=relative-heading-rank; base=h2

\subsection{第十章 罪恶是痛苦的根源}

但人过着无忧无虑的生活时，开始认为美好事物的持续并非出于天主的恩赐，而是出于事物的偶然，并将无忧无虑地享受天主悦纳的乐趣视为本性应得的，而非天主仁慈的恩赐——人被这些事引向相反且不敬的思想，最终在懒惰的煽动下，认为众神的生活是他们的本性所有，无需任何劳苦或功德。由此他们每况愈下，甚至相信世界不受天主眷顾治理，美德也无容身之处，因为他们知道自己不劳而获地拥有安逸和享乐，且未经任何工作就被当作天主的朋友。\par

% structure: p0022 source=h2 target=subsection reason=relative-heading-rank; base=h2

\subsection{第十一章：苦痛有益}

因此，藉着天主最公义的审判，劳苦和磨难被指定为对那些沉溺于此类虚妄思想之人的医治。当劳苦和磨难临到他们时，他们被逐出安逸和愉悦之地。大地也开始不劳不作就不给他们任何出产；于是，人的思想在他们内转变，他们被警告要寻求造物主的帮助，并以祈祷和誓愿祈求天主的保护。如此，他们因顺遂而忽略的天主敬拜，反因逆境得以恢复；舒适所败坏的对天主的思想，被困苦纠正了。因此，天主眷顾看到这对人更有益，便移除了仁慈和丰裕之道，视之为有害，而引入了烦恼和磨难之道。\par

% structure: p0024 source=h2 target=subsection reason=relative-heading-rank; base=h2

\subsection{第十二章：哈诺客的升天}

但为表明这些事是因忘恩负义者而发生的，祂将最初人类中的一位提升到不朽，因为祂看见他没有忘记祂的恩宠，并且他盼望呼求天主的名；而其余的人，如此忘恩负义，甚至劳苦和磨难也不能使他们改正和纠正，便被判处可怕的死亡。但在他们中间，祂也找到了一位，他和他的全家是义人，\BibleRef{创 6:9}，祂吩咐他建造一只方舟，使他和奉命与他同去的人得以逃脱，当万物都要被洪水毁灭时：好使恶人被洪水淹没而除灭，世界得以净化；而那位为种族延续而被保存的人，经水净化后，能重新修复世界。\par

% structure: p0026 source=h2 target=subsection reason=relative-heading-rank; base=h2

\subsection{第十三章：偶像崇拜的起源}

但这一切之后，人又转向不敬；为此，天主赐下法律来教导他们生活方式。然而随着时间的推移，天主敬拜和正义被不信者和恶人败坏，我们将稍后更充分地展示。此外，引入了乖谬和迷失的宗教，大多数人都沉迷其中，借口节日和庆典，设立宴饮和筵席，追随笛子、长笛、竖琴和各种乐器，纵容于各种酗酒和奢华。由此产生了各种错误；由此他们发明了丛林和祭台、带子和牺牲，酗酒之后便如癫狂般激动。由此，魔鬼得以进入这一类人的心灵，使他们看似癫狂跳舞，如酒神狂女般狂乱；因而发明了咬牙切齿和腹中咆哮；因而人面容可怖、面目狰狞，以致被醉酒和魔鬼驱使的人，被受骗和迷失者相信是充满了神性。\par

% structure: p0028 source=h2 target=subsection reason=relative-heading-rank; base=h2

\subsection{第十四章：天主既善且义}

因此，既然世上引入了如此多虚假和迷失的宗教，我们被派遣如同良善的商人，向你们带来从祖先传承并保存的真天主敬拜；我们将这些话语如同种子撒在你们中间，并交由你们选择认为正确的事。因为如果你们接受我们带来的这些，你们不仅能使自己逃脱魔鬼的侵袭，还能将它们从他人身上驱走；同时你们将获得永恒美好的赏报。但那些拒绝接受我们所讲论的人，在今世将受制于各种魔鬼和疾病的混乱，他们的灵魂离世后将永远受折磨。因为天主不仅是善的，也是公义的；因为如果祂永远是善的，而从未公义地按各人的行为回报，那么善便成了不义。因为若祂对不敬者和敬虔者同等对待，那便是不义。\par

% structure: p0030 source=h2 target=subsection reason=relative-heading-rank; base=h2

\subsection{第十五章：魔鬼如何获得控制人的能力}

因此，正如我们刚才所说，魔鬼一旦能藉着所提供的机会，通过卑劣和邪恶的行为潜入人的身体，若它们因人的疏忽而长久留在其中——因为人不寻求对自己灵魂有益的事——它们必然迫使这些人日后满足住在其内的魔鬼的欲望。而最糟的是，在世界终结时，当那魔鬼被投入永火时，那顺从它的灵魂也必定与它一同在永火中受折磨，连同它所玷污的身体。\par

% structure: p0032 source=h2 target=subsection reason=relative-heading-rank; base=h2

\subsection{第十六章：魔鬼为何想占有人的身体}

现在，魔鬼渴望占据人的身体，原因如下。他们是意志转向邪恶的灵。因此，通过过度的饮食和情欲，他们怂恿人犯罪，但只限于那些怀有犯罪意图的人，这些人看似单纯地想要满足自然必要的需求，却给魔鬼机会进入他们，因为过度而失去节制。因为只要保持自然的尺度，守住合法的节制，天主的仁慈不会让他们进入人内。但当心灵陷入不敬，或身体充满过度的饮食时，就像被那些忽略自身的人意志和意图所邀请，他们获得权柄对抗那些违反天主法律的人。\par

% structure: p0034 source=h2 target=subsection reason=relative-heading-rank; base=h2

\subsection{第十七章　福音赐予制服魔鬼的能力。}

你们看，认识天主和遵守神圣信仰是何等重要！它不仅保护信者免受魔鬼的攻击，还赐予他们管辖那些统治他人者的能力。因此，你们这些外邦人必须投奔天主，远离一切不洁，使魔鬼被驱逐，天主住在你们内。同时，通过祈祷，将你们交托给天主，呼求祂的帮助对抗魔鬼的无耻；因为\Quote{凡你们祈求，相信的，必得着。}\BibleRef{玛 21:22} 但即使是魔鬼本身，当它们看见一个人的信德增长时，也会相应地离开他，仅停留在此人仍存有某种不信的部分；但对于全心相信的人，它们会毫不延迟地离去。因为当灵魂归向天主的信德时，它获得天上之水的能力，以此像灭火星一般熄灭魔鬼。\par

% structure: p0036 source=h2 target=subsection reason=relative-heading-rank; base=h2

\subsection{第十八章　这能力与信德成正比。}

因此，信德有一个尺度，若它完满，便彻底将魔鬼从灵魂中驱逐；但若有任何缺陷，魔鬼的部分仍留存在不信的部分；灵魂最难理解的是，何时或如何，完全或部分地，魔鬼已被驱逐。因为若它留在任何角落，当它得到机会时，就向人心暗示思想；他们不知这些思想来自何处，便相信魔鬼的暗示，仿佛是自己灵魂的感知。因此，魔鬼暗示一些人因身体必要而追求快乐；以胆汁过多为由，为另一些人的暴躁开脱；以忧郁强烈为由，粉饰一些人的疯狂；甚至以痰多为由，减轻一些人的愚昧。但即使如此，这些事没有一件对身体有害，除非由于饮食过量；因为当这些被过量摄取，其过多以致自然热度不足以消化时，凝结成某种毒物，它像公共下水道一样流过脏腑和所有血管，使身体活动不健康且卑劣。因此，凡事都要达到节制，以免给魔鬼留地步，也免得灵魂被它们占据，与它们一同被投入永恒的火中受折磨。\par

% structure: p0038 source=h2 target=subsection reason=relative-heading-rank; base=h2

\subsection{第十九章　魔鬼煽动偶像崇拜。}

魔鬼还有另一种错误，它们暗示给人的感官，使人以为他们所遭受的苦难来自那些所谓的神，以便借此献祭和礼物，似乎要平息它们，从而加强虚假宗教的敬拜，并躲避我们这些关心他们得救之人，使他们免于错误；但正如我所说，他们这样做，是不知道这些事是魔鬼暗示给他们的，因为他们害怕得救。因此，既然人被赋予了自由意志，每个人都有能力决定是听从我们归于生命，还是听从魔鬼归于灭亡。此外，魔鬼有时以各种形状显形，发出威胁，有时承诺解除痛苦，以便向受骗者灌输它们是神的观念，并使人不知道它们是魔鬼。但对我们这些认识创造奥秘，并知道为何允许魔鬼在今世做那些事的人来说，它们并不隐藏；它们被允许变形为任何形状，暗示邪念，并借着献祭给它们的饮食，潜入领受者的心灵或身体，编造虚妄的梦境以促进对某个偶像的敬拜。\par

% structure: p0040 source=h2 target=subsection reason=relative-heading-rank; base=h2

\subsection{第二十章　偶像崇拜的愚昧。}

然而，谁能如此愚昧，竟被说服去崇拜一个偶像，无论它是金制的还是其他金属？谁不明白，金属只是工匠所喜欢的形状？那么，如果工匠不喜欢，它根本不存在，如何能认为神性在其中？或者，它们如何能指望将来之事向它们启示，而它们对眼前之事毫无感知？因为即使它们能预言一些事，也不应立刻被视为神；占卜是一回事，神性是另一回事。因为\OriginalTerm{皮同}{Pythons}似乎也能占卜，但它们不是神；总之，它们被基督徒从人身上驱逐。那被一个人赶走的东西，怎能是神？但也许你会说，它们行医治，指示人如何痊愈？按此原则，医生也应被当作神来敬拜，因为他们治愈许多人；而且越有技巧，治愈越多。\par

% structure: p0042 source=h2 target=subsection reason=relative-heading-rank; base=h2

\subsection{第二十一章　异教神谕。}

因此很明显，由于它们是魔鬼，它们知道一些事情更快更完美\Emph{超过人}；因为在学习中它们不被身体的沉重所阻碍。因此，作为魔鬼，它们毫不延迟、毫不费力地知道医生经过长时间和大量劳动才获得的东西。因此，它们知道一些比人更多的事并不奇怪；但要注意，它们知道这些并不用于灵魂的救恩，而是用于欺骗灵魂，藉此引导人们崇拜虚假的宗教。但是天主，为了不隐藏如此巨大欺骗的谬误，也为了不让祂自己在允许它们通过占卜、医治和梦境如此大肆欺骗人类上显得是谬误的原因，由祂的仁慈为人类提供了补救，并使那些愿意知道的人可以清楚分辨真假。因此，这就是区别：由真天主所说的话，无论是通过先知还是各种异象，总是真实的；但由魔鬼所预言的不总是真实的。因此，一个明显的标志是：那些在任何时候有虚假的东西都不是由真天主说的；因为在真理中绝无虚假。但在那些说虚假话的人中，可能偶尔会掺杂一点点真理，就像给虚假添加调味料一样。\par

% structure: p0044 source=h2 target=subsection reason=relative-heading-rank; base=h2

\subsection{第二十二章 为何它们有时应验}

若有人说：\Quote{这有什么用？它们甚至有时被允许说真话，从而在人间引入如此多的谬误？}让他接受这个回答：如果它们从来不被允许说任何真话，那么它们根本不会预言任何事；既然它们不预言，人们就不会知道它们是魔鬼。但如果魔鬼不被知道存在于这个世界，我们斗争和争战的原因就会向我们隐藏，我们会在公开中遭受暗中进行的事，也就是说，如果赐给它们的能力只是攻击我们，而非说话。但现在，既然它们有时说真话，有时说假话，我们应当承认，如同我已经说过的，它们的回应来自魔鬼，而非来自天主，在祂那里绝无虚假。\par

% structure: p0046 source=h2 target=subsection reason=relative-heading-rank; base=h2

\subsection{第二十三章 恶不在实体中}

若有人更进一步好奇地追问：\Quote{那么天主造这些恶事有什么用？它们具有如此大的倾向来颠覆人的思想。}对提出这样问题的人，我们回答：我们必须首先探究实体中是否有恶。虽然对他这样说就足够了：受造物审判造物主是不合适的，但判断他人的工作属于具有同等技艺或同等能力的人。然而，直接切入要点，我们绝对地说，实体中没有恶。但如果这样，那么实体的造物主就被徒然地指责了。\par

% structure: p0048 source=h2 target=subsection reason=relative-heading-rank; base=h2

\subsection{第二十四章 为何天主容许恶}

但你将会反驳我说：\Quote{即使这是因自由意志而导致的，难道造物主不知道祂所造的那些人会堕入恶中吗？因此，祂本不该创造那些祂预见到会偏离义路的人。}现在我们告诉那些问这样问题的人，我们所作这类表述的目的是要表明，为什么那些尚未存在之人的邪恶没有胜过造物主的良善。因为如果祂想要充满祂创造的数量和度量，祂害怕将要存在之人的邪恶，并且像一位找不到其他补救和医治方法的人，除了停止祂创造的目的，免得将那些将要存在之人的邪恶归咎于祂；这除了表明造物主方面不配的受苦和不体面的软弱之外还能表明什么？祂如此害怕那些尚未存在者的行为，以至于祂停止了祂所目的要创造的工作？\par

% structure: p0050 source=h2 target=subsection reason=relative-heading-rank; base=h2

\subsection{第二十五章 恶者被转用于善}

但放下这些，让我们认真考虑这一点：宇宙的造物主天主，预见到祂创造物的未来差异，预见到并为祂的每个受造物提供了不同的等级和不同的职分，按照由自由意志产生的独特动向。这样，虽然所有人在创造方式上具有同一实体，但在等级和职分上仍有差异，按照由意志自由产生的思想独特动向。因此祂预见到在祂的受造物中会有过失，而祂公义的方式要求惩罚应随过失而来，为了改正。因此，应当有惩罚的执行者，然而自由意志应当将他们带入那个次序。此外，那些为天上的赏报从事竞赛的人，也必须拥有要征服的敌人。因此，那些被认为是恶的事物也并非毫无用处，因为被征服者不情愿地为征服他们的人获得永恒的赏报。但就此为止，因为随着时间的推移，甚至更隐秘的事也将被揭示。\par

% structure: p0052 source=h2 target=subsection reason=relative-heading-rank; base=h2

\subsection{第二十六章 邪恶天使是诱惑者}

现在，既然你们还不明白多么大的无知黑暗包围着你们，同时我愿向你们解释偶像崇拜在世上从哪里开始。所谓偶像，我指的是你们所崇拜的那些无生命的形象，无论是木制、陶制、石制、铜制或其他任何金属制成。这些的开始是这样的：某些天使，离开了它们本有秩序的轨道，开始偏爱人的恶习，并在某种程度上不配地助长他们的情欲，以便通过这些手段它们更能纵容自己的快乐。然后，为了不显得是自愿倾向不配的服事，教导人：魔鬼能够藉某些技艺——即魔法呼求——被制伏而服从人。于是，如同从邪恶的熔炉和作坊，它们用不虔敬的烟充满了整个世界，虔敬之光被撤去。\par

% structure: p0054 source=h2 target=subsection reason=relative-heading-rank; base=h2

\subsection{第二十七章 含是第一个术士}

因为上述以及一些其他原因，洪水降临世界，正如我们已经说过，并将再次说明；地上所有的人都遭毁灭，只有诺厄的家族幸免，他带着三个儿子和他们的儿媳存活下来。其中有一子名叫\Person{含}，不幸发现了巫术行为，并将它传授给他一个名叫\Person{梅瑟兰}的儿子，埃及人、巴比伦人和波斯人的种族就是从\Person{梅瑟兰}传下来的。当时存在的各民族称他为\Person{琐罗亚斯德}，推崇他为巫术的首创者；在\Person{琐罗亚斯德}名下，也存在许多关于此主题的书籍。因此，他极其频繁地注视星辰，希望被他们尊为神明，便开始从星辰中引出某些火花，向人显示，以使粗鲁无知的人因奇迹而惊愕；为了提高这种声誉，他反复尝试，直至被他自己以过于冒昧的态度接近的魔鬼焚烧并吞噬。\par

% structure: p0056 source=h2 target=subsection reason=relative-heading-rank; base=h2

\subsection{第二十八章 巴贝耳塔}

但当时愚钝的人，本应因看见他死于惩罚而放弃对他的看法，反而更加赞扬他。他们为他建造了一座陵墓以示敬意，甚至崇拜他为天主的朋友，认为他被闪电战车带到了天上，并敬拜他如同活星。因此，在他死后，那些在一代之后学会希腊语的人称他为\Person{琐罗亚斯德}——即\Emph{活星}。总之，以此为例，至今仍有许多人崇拜被闪电击中的人，为他们建陵致祭，敬奉为天主之友。但这个人出生于第十四代，死于第十五代，当时建造了巴贝耳塔，人类的语言被分为许多种。\par

% structure: p0058 source=h2 target=subsection reason=relative-heading-rank; base=h2

\subsection{第二十九章 波斯人的拜火敬礼}

其中首先被提到的是某位名叫\Person{尼默洛}的君王，巫术如同闪电般传给了他，希腊人也称他为\Person{尼努斯}，尼尼微城由此得名。因此，各种纷乱和偏离正轨的迷信都源于巫术。因为很难使人从爱天主转向崇拜又聋又哑无生命的偶像，术士们便使用更高超的手段，试图通过星辰的征兆和仿佛从天降下的运动，以及天主的旨意，将人引向偏离正轨的敬拜。最初受骗的人收集了\Person{琐罗亚斯德}的骨灰——正如我们所说，他被因他过于烦扰而发怒的魔鬼烧死——将这些骨灰带到波斯人那里，让他们以不断的看守保存，当作从天降下的神火，并敬拜为天上的神。\par

% structure: p0060 source=h2 target=subsection reason=relative-heading-rank; base=h2

\subsection{第三十章 英雄崇拜}

同样，其他地方的人为他们所钦佩的人——或因其技艺，或因其德行，或至少对他们怀有极大的爱慕——建造庙宇、设立雕像、创立奥迹、仪式和祭献，并乐于藉着属于神的一切来将他们的名声传于后世；尤其因为，正如我们所说，他们似乎得到一些巫术幻象的支持，以致通过呼召魔鬼，似乎有些事有赖于他们而成就和推动，以欺骗世人。他们还增添了某些庆典和酗酒的宴会，人在其中可以任意放纵；魔鬼乘着饱足的坐骑进入他们体内，与他们五脏六腑混合，盘踞在那里，将人的思想和行动束缚在自己的意志之下。这样的谬误从起初被引入，又加上情欲和酗酒——有血肉之人特别以此取乐——使得天主的宗教，即在于节制和清醒的宗教，在人中逐渐稀少，几乎被废弃。\par

% structure: p0062 source=h2 target=subsection reason=relative-heading-rank; base=h2

\subsection{第三十一章 偶像崇拜导致一切不道德}

因为起初，人敬拜公义且无所不见的天主，既不敢犯罪，也不敢损害邻人，确信天主看见每个人的行动和举止；但后来宗教敬拜转向无生命的偶像，人确知它们不能听、不能看、不能动，便开始放纵地犯罪，进行各类邪恶，因为他们不怕从所敬拜为神的那里遭受任何惩罚。于是战争狂乱爆发；抢劫、掠夺、俘虏、自由沦为奴隶随之而来；各人尽其所能满足自己的情欲和贪得无厌，虽然没有任何力量能满足贪婪。因为正如火，燃料越多，就越发燃烧猛烈；同样，贪欲的疯狂因所获之物而愈加增长猛烈。\par

% structure: p0064 source=h2 target=subsection reason=relative-heading-rank; base=h2

\subsection{第三十二章 邀请}

所以，你们现在要开始以更正确的理解抵抗自己那些不正当的欲望；若你们能以任何方式修复并恢复那起初由天主赋予人的纯洁宗教和无辜生命，那么不朽福佑的希望也能重新赐予你们。要感谢万物仁慈的圣父，藉着祂所立定的和平之王、无可言喻之荣耀的宝库，好使你们现今的罪过也能被泉水、河水、甚至海水的洗涤除去；同时，那三重的真福之名被呼求于你们身上，借此不仅邪恶的鬼神（若存于你们内）被驱逐，而且，当你们弃绝罪过、以全备的信德和纯全的心灵相信天主时，你们也能从他人身上驱逐恶神和魔鬼，并能解救他人脱离痛苦和疾病。因为魔鬼自身认识并承认那些将自己交付给天主的人，有时仅仅因他们的临在就被驱逐，正如你们刚才所见，当我们只向你们致辞问候时，魔鬼因尊重我们的宗教立即开始呼喊，甚至不能忍受我们的临在片刻。\par

% structure: p0066 source=h2 target=subsection reason=relative-heading-rank; base=h2

\subsection{第三十三章　最软弱的基督徒比最强大的魔鬼更有能力}

那么，难道是我们另有一种更优越的本性，所以魔鬼惧怕我们吗？不，我们与你们同属一个本性，但我们在宗教上不同。然而，若你们也愿意像我们一样，我们并不嫉妒，反而劝勉你们，并愿你们确信：当同样的信德、宗教和无辜生命存在于你们内，如同在我们内一样时，你们将拥有同等和同样的能力与德能对抗魔鬼，这是因天主赏报你们的信德。因为正如那有士兵在他属下的人，虽然自己可能较弱，而士兵比他强壮，却\BibleQuote{对这说：去，他就去；对另一说：来，他就来；对另一说：做这事，他就做这事}\BibleRef{玛 8:9}；他能这样做，并非靠自己的能力，而是靠凯撒的威严。同样，每位信者命令魔鬼——虽然它们看起来比人强大得多——并非靠自己的能力，而是靠那使它们服从的天主的能力。因为我们刚才提到的，凯撒为所有士兵、在每个军营和他的整个王国所敬畏，虽然他仅是一人，或许体质软弱，但这唯有藉天主的能力才能实现，是祂使众人恐惧，以致服从一人。\par

% structure: p0068 source=h2 target=subsection reason=relative-heading-rank; base=h2

\subsection{第三十四章　基督受试探}

我们愿你们确切知道，魔鬼对人毫无能力，除非人自愿屈服于它的欲望。因此，那邪恶的首领曾接近那位我们所说由天主立定的和平之王，试探祂，并开始向祂应许世上一切的荣耀；因为他知道，当他以此诱惑他人时，他们曾朝拜他。因此，这邪恶者，不自量力——这实是邪恶的特质——竟以为他会被那他知道必将毁灭他的主所朝拜。所以，我们的主为确认唯一天主的敬拜，回答他说：\Quote{经上记载：\BibleQuote{你要朝拜主，你的天主，惟独事奉祂。}}\BibleRef{玛 4:10}。他被这回答惊吓，惟恐独一真天主的真宗教得以恢复，遂急忙派遣假先知、假宗徒和假教师进入世界，他们虽然以基督之名说话，却成就魔鬼的意愿。\par

% structure: p0070 source=h2 target=subsection reason=relative-heading-rank; base=h2

\subsection{第三十五章　假宗徒}

所以，你们要极其谨慎，不可相信任何教师，除非他从耶路撒冷带来雅各——主的弟兄——的证明，或任何继承他的人。因为除非某人上到那里，并在那里被认可为合适且忠信的教师以宣讲基督的圣言——除非，我说，他从那里带来证明——否则决不可接受。但如今除了我们，你们不必期待先知或宗徒。因为只有一位真先知，我们十二宗徒传讲祂的言语；因为祂是天主悦纳之年，以我们十二宗徒作为祂的十二个月。至于世界为何被造，其间发生了何种变化，以及我们的主为恢复世界而来，为何拣选并派遣了我们十二宗徒，这些将在别时更详细地解释。同时，祂命令我们出去宣讲，并邀请你们参加天上君王的婚宴，这婚宴是圣父为祂圣子的婚礼所预备的，并要我们赐予你们婚宴礼服，即圣洗的恩宠。凡获得此礼服者，如同无瑕的长袍穿着进入君王的婚宴，必须谨慎，免得礼服上任何部分被罪玷污，因而被拒绝为不配和可弃绝的。\par

% structure: p0072 source=h2 target=subsection reason=relative-heading-rank; base=h2

\subsection{第三十六章　无瑕的礼服}

\Quote{但玷污这礼服的方式如下：若有人背离万有的天主圣父和创造者，接受除基督以外别的教师，而基督是唯一忠信真实的先知，并且祂派遣了我们十二宗徒宣讲圣言；若有人对那超越万有的天主性之本体思想不相称——这些事甚至致命地污损了圣洗的礼服。但那些在行为上污损它的事是：凶杀、奸淫、仇恨、贪婪、邪恶的野心。而那些同时污损灵魂和身体的事是：参与魔鬼的筵席，即尝祭物、血、或勒死的动物尸体，以及任何献给魔鬼之物。因此，这应是你们三步中的第一步；第一步产生三十条诫命，第二步六十条，第三步一百条，我们将在别时更详尽地向你们解释。}\par

% structure: p0074 source=h2 target=subsection reason=relative-heading-rank; base=h2

\subsection{第三十七章　会众散会}

他这样说了，并吩咐他们第二天准时到同一个地方来，就遣散了群众；当他们不愿离去时，\Person{伯多禄}对他们说：\Quote{请看在昨天旅途疲惫的份上，帮我这个忙；现在请离去，明天准时集合。}于是他们高兴地离去。但\Person{伯多禄}命令我稍微退开以便祈祷，之后吩咐在花园中有荫凉的部分铺好躺椅；每人按照惯例，认出自己等级的位置，我们便进食。然后，白天还有一段时间，他与我们谈论主的奇迹；到了傍晚，他进入寝室睡觉。\par

\SourceRecord{Translated by Thomas Smith.；Ante-Nicene Fathers；Vol. 8.；Edited by Alexander Roberts, James Donaldson, and A. Cleveland Coxe.；Buffalo, NY: Christian Literature Publishing Co.,；1886.；<http://www.newadvent.org/fathers/080404.htm>.}

% structure: p0001 source=h1 target=section reason=page-title

\section{\WorkTitle{识别录（第五卷）}}

% structure: p0002 source=h2 target=subsection reason=relative-heading-rank; base=h2

\subsection{第一章 伯多禄的问候}

但次日，伯多禄比平时稍早起身，发现我们仍在沉睡。他看见后，便吩咐大家为他保持安静，仿佛他自己还想多睡一会儿，以免打扰我们的休息。待我们睡足醒来，见他已祈祷完毕，正在卧室等候我们。天已破晓，他简短地向我们致意，按照他的习惯问候了我们，随后前往惯常施教之处。他见已有许多人聚集在那里，便依照最初的宗教祷词向他们呼求平安，然后开始讲话如下：\par

% structure: p0004 source=h2 target=subsection reason=relative-heading-rank; base=h2

\subsection{第二章 罪的后果}

万有的创造者天主起初按自己的肖像造了人，赐给他管辖大地、海洋和天空的权力；正如那位真实的先知所告诉我们的，也如事物本身的理性所教导我们的：因为只有人是理性的，而理性应当统治非理性之物。因此，起初当他尚为义人时，他超越一切紊乱和软弱；但一旦犯了罪（正如我们昨日所教导的），成为罪的奴仆，便同时变得软弱可欺。此事被记载下来，是为叫人知道：正如人因不虔敬而变得易于受苦，同样地，因虔诚而能免于受苦；不仅免于受苦，且能凭藉对天主的微小信德治好他人的苦难。因为那真实的先知曾这样应许我们说：\BibleQuote{我实在告诉你们：如果你们有信德像芥子那样大，就是对这座山说：从这里移到那里去，它也会移去。}\BibleRef{玛 17:20} 对这话，你们自己也已有了明证；因为你们昨日看到，当我们出现时，魔鬼连同它们加给人们的那些苦难都逃跑了。\par

% structure: p0006 source=h2 target=subsection reason=relative-heading-rank; base=h2

\subsection{第三章 信与不信}

既然有人受苦，又有人治好受苦者，就必须知道苦难与医治的原因；而这无非是受苦者的不信和医治者的信德。因为不信，不相信有天主的审判，就给予犯罪以自由，而罪使人承受苦难；但信德，相信有天主的审判，就约束人不去犯罪；不犯罪的人不仅不受魔鬼和苦难的困扰，还能驱除别人的魔鬼和苦难。\par

% structure: p0008 source=h2 target=subsection reason=relative-heading-rank; base=h2

\subsection{第四章 无知的根源}

由此可以断定，一切灾祸都源于无知；而无知本身，万恶之母，源于粗心和懈怠，因疏忽而在人的感官中滋长、增强并扎根。若有人教导要将它驱逐，它却像从古老世袭的住所被赶走一样，愤懑地难以剥离。因此，我们必须稍加努力，以探究无知的种种自以为是，藉着知识将它们斩断，特别是在那些已被某些错误观念占据的人身上，无知因着某种知识的假象而更加根深蒂固。因为再没有比一个人以为自己知道他所不知道的、以假为真更糟的了。这好比一个醉汉自以为清醒，而实际上在各方面都像一个\Emph{醉汉}，且自以为清醒，且希望别人也称他为清醒。同样，那些不认识真理、却持有某种知识外观的人，行许多恶事如同行善，奔向毁灭如同奔向得救。\par

% structure: p0010 source=h2 target=subsection reason=relative-heading-rank; base=h2

\subsection{第五章 知识的益处}

因此，我们首先必须切慕认识真理，以便藉此点亮明灯，驱散错误的黑暗。因为如前所述，无知是一大祸患；但它没有实质，对于诚心追求者而言，很容易被驱散。因为无知无非是不知何事对我们有益；一旦知道了，无知便消灭了。所以，应当热切追求真理的知识；而除了那位真实的先知，没有人能赐予这知识。因为这是为那些愿意进入者所设的生命之门，也是为那些走向得救之城者所设的善行之路。\par

% structure: p0012 source=h2 target=subsection reason=relative-heading-rank; base=h2

\subsection{第六章 自由意志}

一个人真正听到那位真实先知的话语后，是愿意还是不愿意接受它，并担起祂的轭（即生活的诫命），这取决于他自己，因为我们的意志是自由的。若情况是听者无法有别于所听的行动，那便有一种自然之力，使他无法转向另一种意见。又若根本没有听众能够接受它，这也是一种自然之力，迫使他做某事，而不留余地给他事。然而，既然心灵可以自由地将其判断转向任何一方，并选择它所认可的道路，就清楚地表明人拥有选择的自由。\par

% structure: p0014 source=h2 target=subsection reason=relative-heading-rank; base=h2

\subsection{第七章 知识的责任}

因此，人在听闻什么对他有益之前，他必定是无知的；因着无知，他渴望并想要做对他无益的事，因此他不会为此受审判。然而，一旦他听见了自己错误的原因，并接受了真理的方法，若他仍停留在早已被占据的错误中，他便理应被传唤审判，承受惩罚，因为他将赐予他为善度生活的那段生命，虚掷在错误的戏弄中。但那听见这些事而甘心接受，并感激善道已被带给他的人，他会更热切地探求，不停地学习，直到查明是否确实有另一个世界，在那里为善人预备了赏报。当他对此确信时，他感谢天主，因为祂向他显示了真理之光；他此后以一切善工引导自己的行为，他确信在将来的世界有为这些善工预备的赏报；同时他不断惊讶并诧异于他人的错误，没有人看见放在他们眼前的真理。然而他自己，因所找到的智慧财富而欢欣，贪得无厌地渴望享受它们，并因实践善工而喜悦；他竭力以洁净的心和纯洁的良心，奔向将来的世界，那时他将能看见万有的君王——天主。\par

% structure: p0016 source=h2 target=subsection reason=relative-heading-rank; base=h2

\subsection{第八章 应克制肉身的欲望}

但我们缺乏并丧失这一切的唯一原因就是无知。因为当人不知道知识中有多大的善时，他们就不容许消除无知之恶；因为他们不知道将一件事物换成另一件会带来多么大的不同。因此，我劝告每位学习者，乐意倾听天主的话语，并以爱真理的心听我们所说的话，好使他们的心接受了最好的种子，能藉善行结出喜乐的果实。因为如果当我在教导关于救恩的事时，有人拒绝接受，并以被恶见占据的心尽力反对，那么他灭亡的原因不在我们，而在他自己。因为他的责任是以公正的判断审视我们所说的话，并明白我们是在说真理之言，好使他知晓实情，以善行引导生活，成为天国有份者，使肉身的欲望服从自己，成为它们的主人，这样最终他自己也能成为万有统治者的中悦产业。\par

% structure: p0018 source=h2 target=subsection reason=relative-heading-rank; base=h2

\subsection{第九章 两个王国}

因为那坚持作恶、作恶的奴仆的人，只要他坚持作恶，就不能成为善的一部分，因为从起初，正如我们所说，天主设立了两个王国，并赐给每人能力，成为他将自己顺服而服从的那个王国的一部分。既然天主已定下，没有一个人能同时事奉两个王国，所以你们要竭尽全力，投靠善君的盟约和法律。因此，当真先知与我们同在时，看见有些富人疏忽了敬拜天主，就如此阐明了此事的真理：祂说：\BibleQuote{没有人能事奉两个主人；你们不能事奉天主而又事奉玛门。}\BibleRef{玛 6:24} 祂用本国语言称呼财富为 \OriginalTerm{玛门}{mammon}。\par

% structure: p0020 source=h2 target=subsection reason=relative-heading-rank; base=h2

\subsection{第十章 耶稣是真先知}

祂就是那位真先知，正如你们所听见的，祂在犹太地显现给我们；祂站在公共场所，仅凭一句话，就使瞎子看见、聋子听见、驱逐魔鬼、使病人康复、死者复生；既然没有一件事对祂是不可能的，祂甚至洞察人心，这是唯独天主才能做到的。祂宣讲了天主之国；我们相信祂是我们在祂所说的一切事上的真先知，我们的信德不仅因祂的话语，也因祂的作为而得到确证；还因为法律的话，在许多世代之前就已预言了祂的来临，都在祂身上应验了；梅瑟以及他之前的先祖雅各伯的行为的预象，在各方面都预表了祂。同样明显的是，祂来临的时间，即祂实际到来的时刻，也被他们所预言；而最重要的是，圣经中包含：祂要为外邦人所期待。这一切都在祂身上同样应验了。\par

% structure: p0022 source=h2 target=subsection reason=relative-heading-rank; base=h2

\subsection{第十一章 外邦人的期待}

但一位犹太人的先知预言祂要为外邦人所期待，这极大地确认了我们在祂身上的真理信德。因为如果他说祂要为犹太人所期待，他就不像预言了什么非凡之事——祂的来临曾应许为了世界的救恩，祂应是同一支派的人、即祂本国的民族所盼望的对象：因为这件事的发生，似乎更像是自然推论，而非需要先知预言之宏伟。然而如今，先知们却说，凡关乎世界救恩所提出的希望、基督将要建立的新王国、以及凡关于祂所宣告的一切，都要转移到外邦人身上；这先知职的宏伟，不是按事物的次序，而是藉着预言令人难以置信的应验而得到确认。因为犹太人从一开始就凭着一个最可靠的传承知道，有一个人将在某个时候来到，万物要藉着祂得以恢复；他们天天默想并期待祂的来临，但当时他们看见祂在他们中间，行出了关于祂所记载的征兆和奇迹，却因嫉妒而瞎了眼，不能认出临在的祂，尽管在祂缺席时他们曾因盼望祂而喜乐；然而我们少数蒙祂拣选的人却明白了这事。\par

% structure: p0024 source=h2 target=subsection reason=relative-heading-rank; base=h2

\subsection{第十二章 外邦人的蒙召}

但这出于天主的眷顾，为的是使这位善者的知识传交给外邦人，那些从未听说过祂、也未从先知学习的人，应承认祂，而那些在日常默想中承认祂的人，反而不认识祂。因为请看，藉着你们这些现在在场、渴望听信祂信德的教导，并要知道祂的来临是什么、如何、怎样的方式，先知预言就应验了。因为这是先知所预言的：祂要被你们这些从未听说过祂的人寻求。\BibleRef{依 65:1} 因此，既然预言的话甚至在你身上应验了，你们正确地唯独信靠祂，你们正确地等候祂，你们正确地询问关于祂，使你们不仅等候祂，而且因信获得祂王国的继承；按照祂自己所说的，每一个人都成为他所服从之人的奴仆。\BibleRef{若 8:34}\par

% structure: p0026 source=h2 target=subsection reason=relative-heading-rank; base=h2

\subsection{第十三章 外邦人的邀请。}

因此，醒来吧，接受我们的主和天主，就是那掌管天地的主，并使自己与祂的肖像和模样相符，如同真先知亲自教导说：\Quote{你们应当慈悲，如同你们的天父是慈悲的，祂使太阳升起，光照善人，也光照恶人，降雨给义人，也给不义的人。}因此，效法祂，敬畏祂，正如给人类的诫命所说：\Quote{你当朝拜主你的天主，唯独事奉祂。}因为唯独事奉这位主对你们有益，藉着祂认识独一的天主，你们可以从那些你们徒然惧怕的众多神祇中得释放。因为那不惧怕创造万有的天主，却惧怕那些他自己亲手所造的东西，他所做的岂不是使自己屈服于虚妄无意义的恐惧，使自己比那些在他心中产生恐惧的东西更为卑贱和可鄙吗？但相反，藉着那邀请你们的天主的良善，回归你们原来的高贵，并以善行表明你们承载着造物主的肖像，好因默观祂的模样，你们可被认为是祂的儿女。\par

% structure: p0028 source=h2 target=subsection reason=relative-heading-rank; base=h2

\subsection{第十四章 偶像无益。}

因此，开始从你们心中驱逐偶像的虚妄观念，以及你们无用而空洞的恐惧，好使你们同时也能脱离不义的奴役状态。因为那些甚至不能成为你们有益仆人的，竟成了你们的主人。因为无生命的偶像如何能显得适合服侍你们？它们既不能听，也不能看，更不能感知任何事物。是的，甚至它们所制成的材料，无论是金、银，甚至是铜或木，虽然可能对你们有实际用途，但你们却将它们做成神像，使其完全失效无用。因此，我们向你们宣告对天主的真正敬拜，同时警告并劝勉敬拜者，要藉着善行效法他们所敬拜的那一位，并赶快回归祂的肖像和模样，如同我们先前所说的。\par

% structure: p0030 source=h2 target=subsection reason=relative-heading-rank; base=h2

\subsection{第十五章 拜偶像的愚妄。}

但我希望那些拜偶像的人告诉我，他们是否愿意变得像他们所拜的那些东西？你们中有人愿意像它们那样看吗？或像它们那样听？或拥有它们那样的理解力？愿这远离我的任何听众！因为这对一个人来说更应被视为诅咒和羞辱，他自身带着天主的肖像，尽管失去了模样。那么，它们应被算作怎样的神祇，模仿它们会受到敬拜者的憎恶，拥有它们的模样会成为羞辱？那么怎样？熔化你们无用的偶像，做成有用的器皿。熔化那无用无活力的金属，做成适合人使用的工具。但有人说：人类的法律不允许我们这样做。他说得好；因为禁止的是人类的法律，而不是它们自身的能力。那么，那些靠人类法律而非自身能力保护的神祇，是怎样的神？同样，它们靠看门狗和门闩的保护免遭盗贼，至少如果它们是银、金甚至是铜做的；而那些石头和陶制的则因自身的无价值而受保护，因为没有人会偷一个石头或陶制的神。因此，那些因更贵重的金属而暴露于更大危险的神似乎更为可怜。既然它们可以被偷，必须被人看守，可以被熔化、称重，用锤子锻造，有理智的人难道应该视它们为神吗？\par

% structure: p0032 source=h2 target=subsection reason=relative-heading-rank; base=h2

\subsection{第十六章 唯独天主是合宜的敬拜对象。}

唉！人的理智陷入了何等可怜的光景！因为如果惧怕死人被视为最大的愚妄，那么对那些惧怕比死人更糟的东西的人，我们又该如何评判呢？因为那些偶像甚至不能算在死人之列，因为它们从未活过。甚至死人的坟墓也比它们更可取，因为死人虽然死了，但曾有过生命；而你们所敬拜的那些，连一切东西中最低贱的生命，如青蛙和猫头鹰的生命，也未曾拥有。但何必再多说呢？因为对敬拜偶像的人说以下话就够了：\Quote{你难道没有看见你所敬拜的那位看不见吗？听，你所敬拜的那位听不见吗？要明白，那位毫无理解力吗？}因为它是人手的作品，必然没有理解力。因此你们敬拜一个没有感觉的神，而每一个有感觉的人都相信，连那些由天主所造且有感觉的东西，如日、月、星辰及天上和地上的一切，也不应受敬拜。因为他们认为合理的是：不应敬拜那些为世界服务的东西，而应敬拜这些东西和整个世界的创造者。因为连这些东西也喜欢祂受敬拜和崇拜，而不悦创造者的尊荣被给予受造物。因为唯独敬拜天主才是\Emph{它们}所能接受的，祂是唯一非受造的，而万物都是祂的受造物。因为正如唯独非受造的那位才有资格是神，一切受造物都不是真正的神。\par

% structure: p0034 source=h2 target=subsection reason=relative-heading-rank; base=h2

\subsection{第十七章 古蛇的暗示。}

所以，你们最应当明白的是那古蛇的欺骗及其狡猾的暗示。他仿佛以明智的方式欺骗你们，像是藉着某种理性爬进你们的感官；他从头部开始，滑入你们的内髓，把欺骗你们视为巨大的收获。因此，他向你们的思想灌输各种神祇的意见，仅是为了使你们远离唯一天主的信仰，因为他知道你们的罪恶是他的安慰。因为他因自己的邪恶从起初就被判罚吃尘，因为他使那从尘土中取来的人又归于尘土，直到你们的灵魂经过火的考验而得以恢复时为止——这我们将在另一时间更充分地教导你们。从他那里，生出了所有错误和怀疑，你们正是由此被从对唯一天主的信德和信仰中驱离。\par

% structure: p0036 source=h2 target=subsection reason=relative-heading-rank; base=h2

\subsection{第十八章：他的第一次暗示。}

首先，他向人的思想暗示不要去听真理的话语，藉此他们本可以驱散对邪恶之事的无知。他这样做，仿佛展示另一种知识，摆出那许多人持有的意见，即认为若无知就不会被定罪，凡未听见的也不必被追究；由此他劝说人们转离听道。但我告诉你们，与此相反，无知本身就是一种最致命的毒药，足以无需外援而毁灭灵魂。因此，没有哪个无知者能因无知而逃脱，反而必定灭亡。因为罪恶的力量自然地摧毁罪人。既然审判将根据理性，无知的缘由和起源，如同每一样罪恶一样，将被查究。因为那不愿知道如何获得生命，且宁愿无知以免被定罪的人，单从这一事实，就被视为明知故犯。因为他知道自己不愿听的是什么；那藉蛇的诡计而得的狡诈对他毫无开脱的用处，因为他将面对的是那洞察内心的一位。但为使你们知道无知本身就带来毁灭，\Emph{我向你们保证}：当灵魂离开身体时，如果它离开时仍不知道创造它的那一位，以及它在此世从祂那里获得了一切所需之物，它就会像忘恩负义、不忠不信者一样，被逐出祂光明的王国。\par

% structure: p0038 source=h2 target=subsection reason=relative-heading-rank; base=h2

\subsection{第十九章：他的第二次暗示。}

再者，那邪恶的蛇向人暗示另一种意见，你们许多人惯常提出：正如我们所说，有一位天主，是万有的主；但这些人又说，这些也是神。因为正如有一位凯撒，他下面有许多官长——例如总督、执政官、护民官和其他官员——同样，我们认为，有一位较一切更大的天主，但此外，这些神祇被安置在此世，像我们所说的那些官员，确实服从那位更大的天主，却管理我们和此世的事物。针对此，我将向你们表明，就在你们为欺骗而提出的那些事上，你们如何被真理的理由所驳斥。你们说天主占据凯撒的位置，那些被称为神的代表祂的法官和官长。那么，就按你们所引用的凯撒的例证来认定：要知道，正如凯撒的法官或行政官，如总督、代理总督、将军或护民官，若擅取凯撒的名号，则取用者和授予者都当同被处死，同样，在此事上你们也当注意：若有人将天主的名称给任何不同于祂者，而那人接受了，则他们将同受一种更可怕的灭亡，比凯撒的仆人更甚。因为冒犯凯撒者受暂时的毁灭；但冒犯那位唯一真天主者，将受永罚，而且理所当然，因为以不正当的条件损伤了那独一无二的名称。\par

% structure: p0040 source=h2 target=subsection reason=relative-heading-rank; base=h2

\subsection{第二十章：埃及人的偶像崇拜。}

\Quote{虽然\InnerQuote{神}这个词并非天主的名称，但世人暂且用这个词作为祂的名称。因此，如我所说，当这词被侮慢地使用时，这羞辱便指向对真名的伤害。总之，古埃及人自以为发现了天体运行的理论和星辰的本性，但因魔鬼蒙蔽了他们的感官，他们却使这不可通传的名称遭受各种侮辱。因为有些人教导应崇拜他们的牛，即称为\OriginalTerm{阿匹斯}{Apis}的；另一些人教导应崇拜公山羊、猫、\OriginalTerm{朱鹭}{ibis}、鱼、蛇、洋葱、下水道、\OriginalTerm{腹鸣}{crepitus ventris}，以及无数其他我甚至羞于提及的东西。}\par

% structure: p0042 source=h2 target=subsection reason=relative-heading-rank; base=h2

\subsection{第二十一章：埃及人的偶像崇拜比其他人更合理。}

伯多禄这样说时，我们所有听他说话的人都笑了。于是伯多禄说：你们嘲笑别人的荒谬，因为长期的习惯使你们看不见自己的荒谬。其实，你们嘲笑埃及人的愚昧并非没有理由，他们崇拜哑巴动物，而他们自己却有理性。但我要告诉你们，他们也在嘲笑你们；因为他们说：我们崇拜活的动物，虽然是必死的；但你们却崇拜和敬奉从未活过的东西。他们还加上说，这些是某些力量的象征和比喻，人类的种族靠这些力量的帮助得以管理。他们因羞愧而躲进这种托词，捏造这些和类似的借口，试图遮掩他们的错误。但现在不是回答埃及人的时候，也不是放下对在场者之关怀去医治缺席者之疾病的时候。因为你们并没有为这种疾病而悲伤，反而嘲笑他人的疾病，这清楚表明你们自己并未染上此病。\par

% structure: p0044 source=h2 target=subsection reason=relative-heading-rank; base=h2

\subsection{第二十二章：第二次暗示的继续。}

但是，让我们回到你们身上，你们的主张是：应将天主视为\Person{凯撒}，而将众神视为\Person{凯撒}的臣仆和代理人。请留心听我讲，我立刻就要向你们指出那蛇的隐藏之处，它们藏在这论证的曲折回环之中。所有人都应将以下这点视为确定无疑的：没有任何受造物能与天主同等，因为祂不是由谁造的，而是祂自己创造了万有；事实上，也找不到任何如此不合理的人，竟以为受造物能与造物主相提并论。因此，如果人的心智——不单凭理性，甚至凭某种自然本能——正确持有这种观点，即那被称为天主的，是没有任何事物可以与祂相比或相等的，而是超越一切、凌驾一切，那么，怎么能认为那被相信是至高无上的名号，可以正当地给予你们认为是为人类生活的服务与安逸而受差遣的那些存在呢？但我们还要再加一点。这个世界无疑是被造的，并且是可朽坏的，我们稍后会更充分地展示这一点；同时，我们承认它既是受造的，又是可朽坏的。因此，如果世界不能被称为天主，而且理应如此，因为它是可朽坏的，那么世界的各部分又如何能承当天主的名号呢？因为既然整个世界都不能是天主，它的各部分就更不能了。所以，如果我们回到\Person{凯撒}的例子，你们会看到你们的错误有多大。任何人，即使与\Person{凯撒}有同样的本性，也不得与\Person{凯撒}相比；那么，你们认为有人应当与那在这点上超越一切——即祂不是由谁造的，而是祂自己创造了万有——的天主相比吗？然而，你们确实不敢将\Person{凯撒}的名号给予任何其他人，因为他会立即惩罚冒犯他的人；而你们却敢将天主的名号给予其他存在，因为祂延迟对冒犯者的惩罚，是为了给他们悔改的机会。\par

% structure: p0046 source=h2 target=subsection reason=relative-heading-rank; base=h2

\subsection{第二十三章　第三点建议}

那蛇也常借他人的口这样说：\Quote{我们敬拜可见的像，以敬礼无形可见的天主。}这无疑是错误的。如果你们真的愿意敬拜天主的肖像，你们就该对人行善，如此在他身上敬拜天主的真实肖像。因为天主的肖像在每一个人内，但祂的模样却不在所有人内，而只在灵魂良善、心思纯洁的人内。因此，如果你们真正愿意尊崇天主的肖像，我们向你们宣告真理：你们应当对那按天主的肖像所造的人行善，并给予尊崇与敬意；应当给饥饿的人食物，给口渴的人水喝，给赤身露体的人衣服穿，给旅客提供住宿，给囚犯所需用品；这些才是真正献给天主的。这些事足以尊崇天主的肖像，以至于不行这些事的人被视为侮辱那神圣的肖像。那么，那种从一个石像或木像跑到另一个，将空虚无生命的像当作神明敬拜，却轻视那些确实拥有天主肖像的人，这算是什么对天主的敬礼呢？相反，你们要确信，凡杀人、行奸淫，或任何给人带来痛苦或伤害的事，在这一切中，天主的肖像都被侵犯了。因为伤害人是对天主的极大不敬。因此，每当你们对别人做你们不愿别人对你们做的事时，你们就用不应有的痛苦玷污了天主的肖像。所以，你们要明白，这正是潜伏在你们内的那蛇的暗示，它说服你们，使你们以为敬拜无知觉的事物就可以显得虔诚，而伤害有知觉、有理性的人却不显得不敬。\par

% structure: p0048 source=h2 target=subsection reason=relative-heading-rank; base=h2

\subsection{第二十四章　第四点建议}

但这些事，那蛇又用另一张嘴回答我们，说：\Quote{如果天主不希望这些事物存在，那么它们就不该存在。}我不是在告诉你们，为什么在这个世界上，为了考验每个人的心思，许多相反的事物被允许存在。但现在适合说的是：如果按你们所说，一切本该受敬拜的事物都不该存在，那这个世界就几乎什么也没有了。因为你们还有什么东西没有敬拜过呢？太阳、月亮、星辰、水、大地、山岳、树木、石头、人；这些之中没有一样是你们未曾敬拜的。因此，按照你们的说法，这些事物都不该由天主造出来，好叫你们没有任何可敬拜的东西！是的，祂甚至不该造人本身来作敬拜者！但这正是潜伏在你们内的那蛇所渴望的：它不放过你们中的任何人；它不愿你们任何一人逃脱毁灭。但事实并非如此。因为我告诉你们，有罪的并不是那受敬拜的东西，而是那敬拜的人。因为天主的审判是公义的；祂以某种方式审判那受害者，又以另一种方式审判那作恶者。\par

% structure: p0050 source=h2 target=subsection reason=relative-heading-rank; base=h2

\subsection{第二十五章　第五点建议}

但你们说：那么，那些敬拜不当敬拜之物的，应该立即被天主消灭，以免他人效尤。难道你们比天主更有智慧，竟敢向祂进谏？\BibleRef{罗 11:34}祂知道该做什么。因为对于所有处于无知中的人，祂都忍耐，因为祂是仁慈而宽厚的；祂预见许多不敬虔的人会变得虔敬，甚至有些敬拜不洁雕像和被玷污偶像的人，也会归向天主，离弃罪恶，行善立功，获得救恩。但有人会说：我们根本不该动念去做这些事。你们不明白什么是自由意志，而且忘记了：那出于自己意愿行善的才是善人；但若因不得已而持守善，就不能称为善，因为那不是出于他自己。因此，既然每个人都有选择善恶的自由，他或获得赏报，或招致毁灭。不，有人说，天主把我们所想的都放在我们心里。诸位，你们这是什么意思？你们在亵渎。因为如果祂把我们的所有思想都放进我们心里，那么奸淫、贪婪、亵渎和各种放荡的思想也都是祂所暗示的了。我恳求你们，停止这些亵渎，明白什么是合乎天主尊荣的敬意。不要像你们有些人惯常说的那样，说天主不需要人的尊荣。的确，祂什么都不需要；但你们应当知道，你们献给天主的尊荣，对你们自己有益。因为还有什么比人不对他的造物主感恩更可恶的呢？\par

% structure: p0052 source=h2 target=subsection reason=relative-heading-rank; base=h2

\subsection{第二十六章 第六个建议}

但有人说：我们做得更好，因为我们既感谢祂，也感谢与祂同在的一切。在这一点上你们不明白这正是你们救恩的毁灭。就好像一个病人同时请来医生和下毒者为他治病；下毒者确实能伤害他，却不能治愈他；而真正的医生会拒绝将他的药物与毒药混合，以免这人的毁灭被归于良善者，或他的康复被归于有害者。但你们说：难道天主就因此生气或嫉妒，如果祂施恩于我们，而我们将感谢献给他人？即使祂不生气，祂至少也不愿成为谬误的根源，让祂的功劳归于虚妄的偶像。还有什么比从天主那里获得恩惠，却向木头石头感谢更不敬、更忘恩负义的呢？所以，起来，明白你们的救恩吧。因为天主不需要任何人，也不需求什么，也不受任何伤害；但我们或得帮助或受伤害，在于我们感恩或忘恩。因为我们的赞美给天主带来什么增益，或我们的亵渎使祂损失什么？只须记住这一点：天主使那向祂感恩的灵魂接近并成为祂的朋友。而邪恶的魔鬼拥有忘恩负义的灵魂。\par

% structure: p0054 source=h2 target=subsection reason=relative-heading-rank; base=h2

\subsection{第二十七章 受造物报复罪人}

但这一点我也愿你们知道：对于这样的灵魂，天主不直接报复，而是祂的整个受造界起来惩罚不敬虔的人；虽然在今世，天主的良善同样将世界之光和土地的出产赐给虔敬和不虔敬的人，但太阳并非毫无悲伤地向不虔敬的人提供光芒，其他元素也并非毫无悲伤地为他们服务。总之，有时甚至违逆造物主的良善，元素因恶人的罪行而疲乏；因此，要么土地的果实枯萎，要么空气的构成变坏，要么太阳的热度过度增加，要么雨水或寒冷过多。于是瘟疫、饥荒和各种形式的死亡相继而来，因为受造物急于向恶人复仇；然而天主的良善约束它，抑制它对恶人的怒愤，迫使它服从祂的仁慈，而不是被人的罪恶和罪行激怒。因为天主的耐心等待人的悔改，只要他们还在这个身体内。\par

% structure: p0056 source=h2 target=subsection reason=relative-heading-rank; base=h2

\subsection{第二十八章 惩罚的永恒性}

如果有人执迷不悟直到生命的尽头，那么，当不朽的灵魂离开身体时，它就要为它的执迷不悟受罚。因为甚至不敬虔者的灵魂也是不朽的，虽然他们自己也许希望灵魂与身体一同结束。但并非如此；因为他们承受着永不熄灭的火焰的折磨，而且他们的毁灭没有死亡的性质。但也许你会对我说：你在恐吓我们，伯多禄。我们如何对你们讲述真实的事呢？难道我们能保持沉默来向你们宣告真理吗？我们不能不按照实际情况叙述。如果我们保持沉默，我们就成了导致你们毁灭性无知的原因，就会满足那潜伏在你们里面、堵塞你们感官的蛇，它狡猾地向你们暗示这些事，好使你们永远成为天主的敌人。但我们被派遣正是为此：向你们揭露它的伪装，融化你们的敌意，使你们与天主和好，让你们归向祂，并通过善工得祂欢心。因为人与天主为敌，对祂怀着不理智、不虔敬的心态和邪恶的倾向，尤其当他自以为知道什么，实际却一无所知时。但当你们放下这些，开始以天主所喜悦和厌恶的为喜悦和厌恶，并愿意天主所愿意的，那时你们才可真正被称为祂的朋友。\par

% structure: p0058 source=h2 target=subsection reason=relative-heading-rank; base=h2

\subsection{第二十九章 天主对人事的眷顾}

\Quote{但你们中或许有人说：天主不关心人间之事；我们连认识祂都无法达到，又怎能达到与祂的友谊呢？天主的治理见证祂确实关心人类的事务：因为太阳每日服侍于它，雨水浇灌它；泉源、河流、风以及所有元素都服侍它；这些事物越为人所知，就越彰显天主对人的眷顾。因为除非藉至高者的能力，更强大的永不会服侍更弱小的；借此显示天主不仅对人有关怀，更有一份深厚的爱，因祂派遣如此高贵的事物为人服务。但人也能达到与天主的友谊，这由那些祂应允其祷告的人的例子所证明：他们愿时，祂使天不下雨；他们祈祷时，祂又开启天空。此外，祂也将许多其他事物赐予那些遵行祂旨意的人，这些恩赐只能赐给祂的朋友。但你们或许会说：我们也崇拜这些事物，对天主有何损害？倘若你们中有人将本应归于他父亲（他从中领受了数不尽的恩惠）的荣誉归给另一个人，并且尊敬一个陌生人和外邦人如父亲，难道你们不认为他对父亲不孝，最应被剥夺继承权吗？}\par

% structure: p0060 source=h2 target=subsection reason=relative-heading-rank; base=h2

\subsection{第三十章 要抛弃祖传的宗教}

另有人说：\Quote{如果我们不崇拜那些从我们祖先传下来的\Emph{偶像}，并背弃祖先留给我们的宗教，那便是邪恶的。}按此原则，若一个人的父亲是强盗或无赖，他也不应改变父亲传给他的生活方式，也不应从父亲的错误中被召回更好之道；并且，如果一个人不与父母一同犯罪，或不与他们一同坚持不敬，便被视为不虔敬。另有人说：\Quote{我们不应打扰天主，不应总以苦难的抱怨或祈求的需要来烦扰祂。}何等愚昧无知之回答！你们以为，若你们为祂的恩惠而感谢祂，便是打扰天主，而你们为祂的恩赐感谢木头石头，却不认为那是打扰祂？并且，当久旱无雨时，我们所有人都转眼望天，向全能天主恳求雨水的恩赐，我们全体与幼童一同向天主倾心祈祷，恳求祂的怜悯？但真正忘恩负义的人，一旦获得祝福，便迅速忘记：因为他们一收完庄稼或葡萄，就立刻将初果献给聋哑的像，并在庙宇或丛林中对天主所赐之物还愿，然后向魔鬼献祭；接受了恩惠，却否认施恩者。\par

% structure: p0062 source=h2 target=subsection reason=relative-heading-rank; base=h2

\subsection{第三十一章 异教及其罪恶}

但有人说：\Quote{这些事物是为喜乐和舒心安神而设立；它们被设计的目的，是使人暂时从忧虑和悲伤中放松。}现在你们看，你们为自己所行之事招致何等指责！若这些事物为的是减轻忧愁、提供享受，那么为何对魔鬼的呼求要在丛林和树林中举行？那疯狂的旋转、肢体的划伤、肢体的割裂是什么意思？何以在其中产生狂怒？何以产生疯狂？何以妇女疯狂地狂跑、披散头发？何处传来尖叫和磨牙？何处传来心和肠的咆哮？以及所有那些不论是假装还是魔鬼的运作所制造的，都用来惊吓愚昧无知之人的事物？这些事是为减轻思想而行，还是为压迫它？你们岂不仍然知觉和领悟，这是潜伏在你们内的蛇的计谋，它藉非理性的错误建议将你们从真理的领悟中引开，为要使你们成为情欲、贪欲和一切可耻之事的奴隶和仆人？\par

% structure: p0064 source=h2 target=subsection reason=relative-heading-rank; base=h2

\subsection{第三十二章 真宗教召唤人归于清醒和谦逊}

但我以清晰的宣讲之声向你们宣告：相反，天主的宗教召唤你们归于清醒和谦逊；命令你们戒除放纵和疯狂，以忍耐和温和阻止愤怒的侵袭；满足于自己的所有和节俭的德行；甚至穷困所迫也不掠夺他人的财物，而在一切事上遵守正义；完全远离偶像祭献：因为藉这些事，你们引出魔鬼到你们身边，并自愿给予它们进入你们的能力；因此你们接纳了那成为疯狂或非法之爱的根源的东西。\par

% structure: p0066 source=h2 target=subsection reason=relative-heading-rank; base=h2

\subsection{第三十三章 不敬的起源}

因此一切不敬的起源由此而来；由此而来谋杀、奸淫、偷盗；当你们沉溺于亵渎的奠祭和香气，并给恶灵以统治和在你们身上取得某种权威的机会，便形成了所有邪恶和恶行的温床。因为当它们侵入你们的感官时，它们除了行属于情欲、不义和残忍的事，并迫使你们服从一切取悦它们的事，还能做什么呢？天主确实凭某种公义的审判，允许你们在它们手下受此苦楚，为要使你们从自己行为和感受的羞耻中，明白屈服于魔鬼而不屈服于天主是何等不配。同样，因与魔鬼的友谊，人被引向羞耻和卑劣的行为；因此，人甚至走向生命的毁灭，或通过情欲之火，或通过愤怒的疯狂，或通过过度的悲伤，以致众所周知，有些人甚至对自己下毒手。并且，正如我们所说，凭天主公义的判决，他们没有被阻止如此行，为使他们都明白自己屈服于谁，也知道自己抛弃了谁。\par

% structure: p0068 source=h2 target=subsection reason=relative-heading-rank; base=h2

\subsection{第三十四章 谁是崇拜天主的人？}

但是有人会说：这些苦难有时也临到那些敬拜天主的人。这不是真的。因为我们说，一个敬拜天主的人，是承行天主的旨意、遵守祂法律诫命的人。因为在天主眼中，他不是犹太人——那在人前被称为犹太人的（也不是外邦人——那被称为外邦人的），而是那相信天主、履行祂的法律、承行祂的旨意的人，即使未受割礼。那才是真正的敬拜天主的人，他不仅自身摆脱苦难，还使他人摆脱苦难；即便这些苦难重如高山，他也藉着对天主的信德将它们移开。是的，因着信德，他真的能移山，连带树木，如果必要的话。但是那些看似敬拜天主的人，若没有被全备的信德和对诫命的服从所坚固，反而是一个罪人，他便因自己的罪，在自身内为那些苦难留出了位置。这些苦难是天主为惩罚犯罪的人而定的，好让他们藉着惩罚，索取他们应得的罪债，并将他们净化后，带到众人的普遍审判面前，只要他们在受罚时其信德不失落。因为现世中不信者的惩罚是一种审判，藉此他们开始与未来的福乐分离；但敬拜天主者的惩罚，虽因他们陷入的罪恶而加于他们，却会向他们索取其所应得的，好使他们避免审判，在现世偿还罪债，至少一半，免去那里所预备的永罚。\par

% structure: p0070 source=h2 target=subsection reason=relative-heading-rank; base=h2

\subsection{第三十五章 将来的审判}

但是，凡是不相信将有天主的审判的人，不会把这些事当作真实的；因此，他因被现世生活的快乐所束缚，被关在永恒的美事之外；所以我们不疏忽向你们宣讲我们认为对你们的救恩必要的事，并指给你们对天主的真敬拜是什么，好使你们相信天主，能藉着善工，与我们一起成为来世的承继人。但若你们尚未相信我们所说的是真的，那么首先，你们不当因我们向你们宣告我们认为好的事而恼怒或敌对我们，也不当因我们不吝惜也赐给你们那我们认为能给我们自己带来救恩的事——正如我所说，我们热切努力，好使你们与我们共同承继那我们认为将临到我们身上的福乐。但你们所宣告的这些事是否确实真实，除非你们服从所吩咐的，否则无从得知；好使你们从事物的结局以及最确实的福乐终点而受教。\par

% structure: p0072 source=h2 target=subsection reason=relative-heading-rank; base=h2

\subsection{第三十六章 讲论的结束}

\Quote{所以，尽管潜伏在你们内的蛇，用千种败坏的手段占据你们的感官，并在你们路上设下千种障碍，藉此使你们远离救恩教诲的听闻，你们更应当抵抗它，藐视它的暗示，更加频繁地聚集来听道并接受我们的教导，因为没有人不经过教导就能学到任何东西。}\par

当他说完话后，他命人把那些被疾病或魔鬼折磨的人带到自己跟前，按手在他们身上祈祷；就这样他遣散了人群，嘱咐他们在他留在此地的日子里前来听道。因此，当人群散去后，\Person{伯多禄}与许多愿意这样做的人一起，在穿过园子的水中洗了身体；然后他命人将卧榻铺在一棵非常荫凉的树下，并指示我们按在\Place{凯撒勒雅}所定的次序斜靠。就这样，在取了食物并按照希伯来人的方式感谢天主之后，由于白天还剩一些时间，他命我们就任何我们喜欢的事质问他。虽然我们与他共有二十人，他回答了每个人任意提问的问题；这些细节我曾记在书中，前些时已寄给你们。到了晚上，我们与他一同进入住所，各归各处，入睡。\par

\SourceRecord{Translated by Thomas Smith.；Ante-Nicene Fathers；Vol. 8.；Edited by Alexander Roberts, James Donaldson, and A. Cleveland Coxe.；Buffalo, NY: Christian Literature Publishing Co.,；1886.；<http://www.newadvent.org/fathers/080405.htm>.}

% structure: p0001 source=h1 target=section reason=page-title

\section{《认亲记》第六卷}

% structure: p0002 source=h2 target=subsection reason=relative-heading-rank; base=h2

\subsection{第一章 第六卷 勤于研究}

但天色初露，黎明驱散退去的黑暗，\Person{伯多禄}进入花园祈祷，从那里返回我们这里时，为比平时稍晚唤醒并来到我们这里而道歉，这样说：\Quote{既然春天延长了白昼，夜晚当然更短；因此，如果有人想在夜晚花些时间研究，就不能在所有季节都保持相同的醒来时间，而应该无论夜晚长短，都花相同的时间睡眠，并且非常小心，不要从他习惯用于研究的时间中削减，从而增加睡眠并减少清醒的时间。这也需要注意，以免如果在食物尚未消化时打断睡眠，未消化的团块会主导心神，并因粗劣蒸气散发使内在感官混乱不安。因此，也要用足够的休息来照顾身体的那一部分，以便充分完成对身体所当尽的那些事之后，身体能在其他方面向心灵提供应有的服务。}\par

% structure: p0004 source=h2 target=subsection reason=relative-heading-rank; base=h2

\subsection{第二章 要在短时间内完成许多事}

他说完这些话，因为已有许多人聚集在花园惯常的地方听他讲话，\Person{伯多禄}便走出去；照常问候人群后，开始讲话如下：的确，正如被农夫忽视的土地必然长出荆棘和蒺藜，同样，你们的感官因长期疏忽，对事物和虚假科学的教条产生了大量有害意见；现在需要用许多关怀来耕耘你们心灵的土地，使真理之言——它是心灵真正勤劳的农夫——用不断的教训来耕耘它。因此，你们应当顺服它，并剪除多余的操劳和忧虑，以免有害的杂草窒息好种子的话语。因为短暂而热切的勤勉可以弥补长时间的疏忽；每个人的生命时间都不确定，所以我们必须赶紧获得救恩，以免突然的死亡捕捉到迟延的人。\par

% structure: p0006 source=h2 target=subsection reason=relative-heading-rank; base=h2

\subsection{第三章 义怒}

因此，我们更要热切努力，趁着还有时间，恶习的积弊可以切除。除非你们对自己无益和卑劣的行为发怒，否则你们不能做到这一点。因为这是一种正义而必要的忿怒，藉此每个人都对自己感到愤慨，并为自己犯错和做错的事责备自己；藉着这种愤慨，我们心中燃起某种火，它如同浇在荒芜的田地上，烧毁和焚尽卑劣快乐的根，使心灵的土地对\Person{天主}的圣言的好种子更加肥沃。我认为你们有足够正当的理由发怒，由此可以点燃那最正义的火，如果你们考虑无知之恶如何把你们引入怎样的错误，它如何使你们坠落并一头栽进罪中，它把你们从什么好东西中拉走，驱入什么邪恶，而且，比一切更重要的是，它如何使你们在来世永罚中受罚。现在真理之光照耀你们，难道对于这一切，内心不会燃起最正义的愤慨之火吗？难道那令\Person{天主}喜悦的忿怒之焰不在你们内升起，烧尽并毁灭一切从根长出的枝条，以防任何邪恶贪欲的嫩芽在你们内萌发吗？\par

% structure: p0008 source=h2 target=subsection reason=relative-heading-rank; base=h2

\subsection{第四章 不是和平，而是刀剑}

因此，那派遣我们者，当祂来到，看见全世界都陷入邪恶时，并没有立即给在错误中的人和平，以免祂使他停留在恶中；而是把真理的知识放在无知之恶的废墟对面，以便如果人们悔改并仰望真理之光，他们可以正确地悲伤自己曾被欺骗并被引入错误的悬崖，并可以点燃有益忿怒之火来对抗欺骗他们的无知。因此，祂说：\BibleQuote{我来是把火投在地上；我多么希望它已经燃起！} \BibleRef{路 12:49} 所以，在这场生命中，有一场战斗要我们打；因为真理和知识的话必然使人从错误和无知中分离出来，如同我们常看到身体中腐烂的死肉被刀从活肢体的连接处割开。真理的知识所产生的效果正是如此。因为为了救恩，例如，接受了真理之言的儿子必须与他不信的父母分离；或者，父亲与儿子分离，女儿与母亲分离。就这样，知识与无知、真理与错误的战斗在相信与不信的亲属和关系人之间兴起。因此，那派遣我们者又说道：\BibleQuote{我来不是为地上送和平，而是送刀剑。} \BibleRef{玛 10:34}\par

% structure: p0010 source=h2 target=subsection reason=relative-heading-rank; base=h2

\subsection{第五章 战斗如何开始}

但若有人说：人脱离双亲，这岂合理？我将告诉你们是何缘故。因为若他们与父母同留于谬误之中，既不能对父母有益，自己也要同他们一起丧亡。因此，那将要得救的人与那不愿得救的人分离，这是合理的，且十分合理。但也要注意：这种分离并非来自那些明悟正道的人；因为他们愿意与亲属同在，向他们行善，并教导他们更美的事。然而，无知有这样一个独特的恶习：它不能容忍那驳斥它的真理之光靠近；所以，分离是由无知者发起的。因为那些领受真理知识的人，由于真理充满良善，便渴望——若有可能——将之分享给所有的人，如同由良善的天主所赐一样；甚至分享给那些仇恨和迫害他们的人：因为他们知道，无知是他们犯罪的根源。因此，简言之，当主自己被那些不认识祂的人带向十字架时，祂曾为杀害祂的人祈求圣父，说：\Quote{父啊，赦免他们的罪，因为他们不知道他们做的是什么！}\BibleRef{路 23:34}宗徒们效法主，当他们自己受苦时，也同样为杀害他们的人祈祷。\BibleRef{宗 7:60}如果我们被教导要甚至为杀害我们的人和迫害我们的人祈祷，那么我们岂不更应该忍受父母和亲属的迫害，并为他们的悔改祈祷吗？\par

% structure: p0012 source=h2 target=subsection reason=relative-heading-rank; base=h2

\subsection{第六章　爱天主应胜过爱父母}

那么，让我们接下来仔细思考：我们爱父母有什么理由？经上说，我们爱他们，因为他们似乎是我们生命的来源。但我们的父母并不是我们生命的来源，而只是途径。因为他们并非赐予生命，而是提供我们进入此生的途径；而生命的惟一来源是天主。因此，若我们愿意爱那生命的来源，就当知道应当爱祂。但有人会说：我们不能认识祂；但我们认识并爱慕他们。就算是这样：你们不能认识天主\Emph{是什么}，但你们很容易认识天主\Emph{不是什么}。因为，有谁不能知道木头、石头、铜或其他类似物质不是天主呢？但如果你们不愿用心思考那些你们容易领会的事，那么，你们在认识天主上受阻，并非因为不可能，而是因为怠惰；因为若你们愿意，即便从这些无用的偶像，你们也可以被引上理解之路。\par

% structure: p0014 source=h2 target=subsection reason=relative-heading-rank; base=h2

\subsection{第七章　大地是为人类受造的}

因为显然这些偶像是用铁制工具造成的；铁经火锻炼；火被水熄灭；水由气息推动；气息源自天主。因为先知梅瑟这样说：\Quote{在起初，天主创造了天地。大地是无形且无序的；黑暗笼罩深渊；天主的圣神运行在水面上。}\BibleRef{创 1:1}这圣神如同造物主的手，按天主的命令将光与黑暗分开；之后，那不可见的天产生了这个可见的天，为使高处成为天使的居所，低处成为人的居所。因此，为了你们的缘故，按天主的命令，水从地面退去，使大地能为你们结出果实；祂也在大地中埋藏了湿气之源，使泉水与河流从其中涌出，供你们使用。为了你们的缘故，大地奉命生出活物，以及一切可以供你们使用和享受的事物。难道不是为你们，风才吹动，使大地受其滋润而生出果实吗？难道不是为你们，雨水降下，季节更替吗？难道不是为你们，太阳升起又落下，月亮盈亏变化吗？为你们，海洋献出它的服务，使万物都服从于你们——你们这些忘恩负义的人！对于这一切，难道不会有一个公义的惩罚报应吗？因为你们比什么都更不认识那赐予这一切恩惠者，而你们本该高于一切地认识并敬拜祂。\par

% structure: p0016 source=h2 target=subsection reason=relative-heading-rank; base=h2

\subsection{第八章　圣洗的必要性}

但现在，我藉由同样的途径引领你们走向理解。因为你们看到万物都源自水。但水是由\OriginalTerm{独生子}{Only-begotten}首先造的；全能的天主是独生子的元首，我们藉着祂以我们上述的顺序到达圣父。但当你们到达圣父时，你们将知道祂的旨意是：你们要藉着那首先受造的水重生。因为那由水重生的人，在圆满善工之后，便成为那使他在不朽中重生者的嗣子。因此，要怀着准备好的心，如同儿子走向父亲一样，使你们的罪得以洗除，并在天主面前证明无知是它们惟一的根源。因为若在得知这些事后，你们仍留在不信中，那么你们灭亡的原因就要归咎于你们自己，而非无知了。难道你们以为，即使你们培养一切虔诚和一切义德，若不领受圣洗，就能对天主怀有希望吗？不，反倒要受更大的惩罚：因为人从善工中确实获得功德，但这仅当这些善工是按天主的命令而行。天主已命令每个敬拜祂的人，藉圣洗接受印记；但若你们拒绝，顺从自己的意愿而非天主的，你们无疑是在反对并敌对祂的旨意。\par

% structure: p0018 source=h2 target=subsection reason=relative-heading-rank; base=h2

\subsection{第九章　圣洗的效用}

但你或许会说：\Quote{水的圣洗对敬拜天主有何贡献？}首先，因为那悦乐天主的事得以成就。其次，因为你由水和天主重生再生时，你从前经由人而有的那脆弱出生便被割断，这样你终能获得救恩；否则是不可能的。因为真先知曾以誓词向我们作证说：\BibleQuote{我实在告诉你们：人除非由水重生，不能进入天国。}因此，赶紧吧！因为这些水中含有某种仁慈的德能，它在起初就运行在水上，并承认那些因三重圣事之名受洗的人，救他们脱离将来的惩罚，将那些因圣洗而祝圣的灵魂作为礼物献给天主。所以，投奔这些水吧！因为唯独它们能扑灭将来之火的暴烈；凡迟延去投奔它们的人，显然不信的偶像仍留在他内，并借此阻止他奔向那赐予救恩的水。因为无论你是义人或是不义的人，圣洗在各方面都是你必需的：对义人，为使成全在他内完成，并使他为天主而生；对不义的人，为使他在无知中所犯的罪获得赦免。因此，众人都应毫不迟延地为天主而重生，因为每个人的生命终点都是不确定的。\par

% structure: p0020 source=h2 target=subsection reason=relative-heading-rank; base=h2

\subsection{第十章  善工的必要}

但你藉水重生后，要以善工在你内显示那生了你的父的相似。现在你认识天主，就当以祂为父尊敬祂；而尊敬祂，就是按祂的旨意生活。祂的旨意是：你们要这样生活，以致不知何为凶杀或奸淫，逃避仇恨和贪恋，除去忿怒、骄傲和夸耀，憎恶嫉妒，并视这一切完全与你们不相称。确实，我们宗教有一种特殊的操守，与其说是加给人，不如说因它的纯洁而为每位崇拜天主者所追求。我说的是洁德，它有许多种，但首先是：人人都要小心，\Quote{不可接近经期的妇人}；因为天主的法律视此为可憎。但即使法律未曾就这些事给予训诫，难道我们愿意像甲虫一样在污秽中打滚吗？因为作为有理性、能体悟天理的人，我们应当比动物更有所持守，其首要之务是守护良心免受内心的一切玷污。\par

% structure: p0022 source=h2 target=subsection reason=relative-heading-rank; base=h2

\subsection{第十一章  内在与外在的洁净}

再者，用水洗身也是好的，且有助于纯洁。我称之为好，并不是因为它像心灵净化那首要之善，而是因为身体的洗涤是那善的结果。因为我们的师傅也曾责备某些法利塞人和经师，他们似乎比别人更好，且与民众分离，称他们为伪善者，因为他们只洁净那些被人看见的事物，却让他们的心——唯独天主看见——污秽不堪。因此，祂对他们中的一些人——并非所有人——说：\BibleQuote{祸哉，你们经师和法利塞人，伪善者！因为你们洁净杯盘的外面，里面却充满污秽。瞎眼的法利塞人！先清洁里面，外面也就清洁了。}\BibleRef{玛 23:25}因为，实在说来，如果心思因知识之光而得以净化，一旦它洁净明朗，就必然关心人外在的事，即他的肉身，使之也得以洁净。但若人对外在的事——肉身的洁净——不加理会，那就说明他对心灵的纯洁和内心的清洁毫不关心。因此，这就发生：内心洁净的人无疑也外在洁净，但外在洁净的人却不一定内心也洁净——即当他做这些事是为了取悦人时。\par

% structure: p0024 source=h2 target=subsection reason=relative-heading-rank; base=h2

\subsection{第十二章  洁德的重要性}

此外，还应遵守这种洁德：性交不可轻率地为纯快乐而进行，而应为生育子女。既然这操守甚至在某些低级动物中也能见到，那么有理性、崇拜天主的人若不行此，岂不愧对？然而，还有进一步的理由，为何那些持守真天主崇拜的人应在我们所谈及的那些形式及其他类似形式中遵守洁德：因为甚至在那些仍被魔鬼的错误所束缚的人中，这洁德也在某种程度上被遵守。那么，既然你们现已悔改，难道你们不愿遵守你们在错误中曾遵守的规矩吗？\par

% structure: p0026 source=h2 target=subsection reason=relative-heading-rank; base=h2

\subsection{第十三章  基督徒道德的优越性}

但你们中或许有人会说：那么我们必须遵守我们敬拜偶像时所做的一切事吗？不是全部。但凡是做得好的事，你们如今也该遵守；因为，如果任何事由在错误中的人做得正确，那么这事必定来自真理；反之，如果在真宗教中有什么做得不正确，那无疑是借自错误。因为善就是善，即使由在错误中的人做出；恶就是恶，即使由追随真理的人做出。或者我们竟会如此愚昧，以致看见一个敬拜偶像的人清醒，我们就拒绝清醒，以免显得与敬拜偶像的人做同样的事？并非如此。但让我们以此作为我们的操练：如果那些犯错的人不杀人，我们甚至不该动怒；如果他们不犯奸淫，我们甚至不该贪恋别人的妻子；如果他们爱邻人，我们甚至该爱仇敌；如果他们借钱给有偿还能力的人，我们该给予那些我们不指望收回的人。在一切事上，我们这些盼望承受永恒世界遗产的人，应当超越那些只知道现世的人；因为我们知道，如果他们的工作，在审判之日与我们的工作相比，被发现相似且相等，那么我们就会羞愧，因为我们被发现在工作上与那些因无知而被定罪、对来世没有希望的人相等。\par

% structure: p0028 source=h2 target=subsection reason=relative-heading-rank; base=h2

\subsection{第十四章 知识增加责任}

\Quote{羞愧确实是我们应得的份，如果我们并不比那些在知识上不如我们的人做得更多。但如果与他们在工作上相等尚且让我们羞愧，那么，如果将来的审查发现我们比他们更差更糟，我们将会怎样呢？因此，请听我们的真先知是如何教导我们这些事的；因为关于那些忽视聆听智慧话语的人，祂这样说：\InnerQuote{南方女王要在审判时与这一代人一同起来，并定他们的罪，因为她从地极而来，要听撒罗满的智慧；看，比撒罗满更大的就在这里，他们却不听祂。} 至于那些拒绝悔改自己恶行的人，祂这样说：\InnerQuote{尼尼微人要在这审判时与这一代人一同起来，并定他们的罪，因为他们因约纳的宣讲而悔改了；看，比约纳更大的就在这里。} 你看，祂如何通过引用那些来自外邦人无知之人的例子，来谴责那些受法律教导的人，并表明前者甚至不比那些似乎生活在错误中的人强。从这一切事上，祂所提出的声明就被证明了，即贞洁——它甚至在某种程度上被生活在错误中的人所遵守——应当被我们这些追随真理的人以更纯洁、更严格的方式，在其所有形式中遵守，正如我们上面所展示的；而且更因为在我们这里，永恒的奖赏被分配给对它的遵守。}\par

% structure: p0030 source=h2 target=subsection reason=relative-heading-rank; base=h2

\subsection{第十五章 在特里波利斯任命主教、司铎、执事及寡妇}

他说完这些话以及其他类似的话后，就遣散了人群；然后照常与朋友们共进晚餐后，便就寝了。就这样，他连续三个月整教导天主圣言，使许多人归信。最后，他命令我禁食；禁食后，他在靠近海边的泉源中为我施行了永流之水的圣洗。当我因天主赋予我的重生恩宠而喜乐地与弟兄们共度假日时，\Person{伯多禄}命令那些已被指定先行的人前往\Place{安提约基雅}，并在那里再等候三个月。他们出发后，他亲自带领那些已完全接受主信仰的人到我所提到的靠近海边的泉源，为他们施洗；并为他们举行圣体圣事后，任命了曾在家中款待他、如今在一切事上已成全的\Person{玛洛}为他们的主教，同时与他一起任命了十二位司铎和执事。他还设立了寡妇团体，并安排了教会的所有服务；他嘱咐他们都服从他们的主教\Person{玛洛}，在一切他所命令的事上。就这样，当一切事务妥善安排，三个月期满后，我们告别了在\Place{特里波利斯}的人，启程前往\Place{安提约基雅}。\par

\SourceRecord{Translated by Thomas Smith.；Ante-Nicene Fathers；Vol. 8.；Edited by Alexander Roberts, James Donaldson, and A. Cleveland Coxe.；Buffalo, NY: Christian Literature Publishing Co.,；1886.；<http://www.newadvent.org/fathers/080406.htm>.}

% structure: p0001 source=h1 target=section reason=page-title

\section{\WorkTitle{认罪录（第七卷）}}

% structure: p0002 source=h2 target=subsection reason=relative-heading-rank; base=h2

\subsection{第一章　从的黎波里出发}

终于我们离开了腓尼基城的黎波里，在离的黎波里不远的奥尔托西亚作了第一次停留；次日我们也在那里住下，因为几乎所有信主的人都舍不得与伯多禄分离，一直跟着他走到这里。然后我们来到安塔拉杜斯。但由于我们一行人众多，伯多禄对尼塞塔和阿奎拉说：\Quote{既然有这么多弟兄与我们同行，每进一城都会招致不少嫉妒，我想我们必须采取办法，既要避免因排场而激起恶者的妒恨，又不至于做出阻止他们跟随我们这样令人不快的事。因此，我希望你们，尼塞塔和阿奎拉，带着他们先走，把群众分成两队，这样我们进城时可以分散开，而不是聚在一起进入每一座外邦人的城邑。}\par

% structure: p0004 source=h2 target=subsection reason=relative-heading-rank; base=h2

\subsection{第二章　门徒分成两队}

\Quote{但我知道，你们认为至少两天见不到我是一件伤心的事。请相信，无论你们多么爱我，我对你们的感情比你们大十倍。可是，如果我们因着彼此相爱而不去做合理和光荣的事，这样的爱就显得不合理了。因此，在不减少我们彼此相爱的前提下，让我们关注那些似乎有用且必要的事；尤其是因为每一天都可能没有你们在场听我讲论。因为我打算逐一经过各省最著名的城邑，你们也知道，每城住三个月以施教。所以现在你们先我到最近的城邑劳狄刻雅去，我将在两三天后跟上，按我的计划。但你们要在离城门最近的那家客栈等我；在那里再住几天后，你们再先我到更远的城邑去。我希望你们在每座城都这样做，为的是尽量避嫌，同时也让与我们同行的弟兄们，因你们的预先安排在各城找到住处，不至于显得像流浪者。}\par

% structure: p0006 source=h2 target=subsection reason=relative-heading-rank; base=h2

\subsection{第三章　行程安排}

伯多禄这样说，他们当然服从了，说道：\Quote{做这事并不令我们十分伤心，因为这是你——被基督的远见所拣选，在一切事上行善并善于劝告的人——所吩咐的；而且，虽然一天或两天见不到我们的主伯多禄是一大损失，但并非不能忍受。我们还想到那十二位先行的弟兄，他们每三个月在你停留每座城期间，有一个月不能听你讲论、见你面。因此我们不会延迟执行你的命令，因为你的一切吩咐都是合宜的。}说着，他们便前行了，领受了指示：要对同行的弟兄们在城外说话，请他们不要成群结队、喧闹地进城，而是分散、分开进入。\par

% structure: p0008 source=h2 target=subsection reason=relative-heading-rank; base=h2

\subsection{第四章　克莱孟因与伯多禄同留而喜乐}

他们走后，我克莱孟非常欢喜，因为他把我留在他身边。我对他说：\Quote{我感谢天主，你没有把我同其他人一起打发走，否则我会因忧伤而死。}伯多禄说：\Quote{那么，如果有必要派你到某处去施教，你会怎样？如果为了善工而与我分离，你会死吗？难道你不能克制自己，耐心承受必然之事吗？难道你不知道，朋友们即使身体分离，也在记忆中常在一起；正如另一些人，身体相近，思想却分离吗？}\par

% structure: p0010 source=h2 target=subsection reason=relative-heading-rank; base=h2

\subsection{第五章　克莱孟对伯多禄的深情}

我回答说：\Quote{主，请不要以为我这样受苦是没有理由的；我对你的这种爱是有某种原因和理由的。因为我把你当作我所有情感的唯一对象，以代替父亲、母亲和兄弟；但最重要的是，你是我的救恩和认识真理的唯一原因。而且我也不认为这是小事：我的青春年华容易陷入情欲的陷阱；我害怕没有你在身边，因为你的同在就能使一切不合理的软弱感到羞愧。虽然我信赖天主的仁慈，甚至我的思想，由于从你的教诲中所领受的，已经不能再接受任何其他事物。此外，我记得你在凯撒勒雅说过：\InnerQuote{若有人愿意跟我同行，而不违反孝道，就让他跟我同行。}你的意思是说，他不可使那按天主的安排他所应依附的人伤心；例如，不可离弃忠信的妻子，或父母等。如今我完全摆脱了这些，因此我适合跟随你；我希望你准我以仆人的身份服侍你。}\par

% structure: p0012 source=h2 target=subsection reason=relative-heading-rank; base=h2

\subsection{第六章　伯多禄简朴的生活}

于是\Person{伯多禄}笑着说：\Quote{\Person{克肋孟}，难道你不认为，极度的需要必然使你成为我的仆人吗？因为还有谁能铺开我的被单，整理我漂亮的被褥？谁能随时保管我的戒指，准备我的长袍——这些我必须不断更换？谁又能监督我的厨师，用最深奥而多样的技艺准备各种精选的肉食？所有这些东西，花费巨大，都是为了那些娇生惯养的人，甚至可以说，为了他们的食欲，如同为了某种巨大的野兽。但也许，虽然你与我同住，你并不了解我的生活方式。我只靠面包和橄榄为生，甚至很少用蔬菜；而我的衣着就是你看到的，一件长衣和一件披肩：有了这些，我别无他求。这些对我已足够，因为我的心思不在现世的事物上，而在永恒的事物上，因此没有现世可见的事物能使我喜悦。正因如此，我确实接纳并钦佩你对我的一片好心；我更加称赞你，因为你虽然习惯了如此丰裕，却能如此迅速地放弃它，并适应我们这种只使用必需品的生活。因为我和我的兄弟\Person{安德肋}从小不仅成为孤儿，而且极其贫穷，由于需要而习惯了劳动，因此我们现在也能轻易承受旅途的辛劳。但如果你同意并允许的话，我这个劳动者反而更容易尽仆人的职责服侍你。}\par

% structure: p0014 source=h2 target=subsection reason=relative-heading-rank; base=h2

\subsection{第七章  \Person{伯多禄}的谦卑}

但我听到这话时战栗不已，眼泪立刻夺眶而出，因为如此伟大的一位人物——他比整个世界更宝贵——竟向我提出这样的建议。他见我在哭泣，便询问原因。我回答他说：\Quote{我究竟如何得罪了你，使你要用这样的提议来折磨我？}伯多禄说：\Quote{如果我说应该服侍你是一件坏事，那么你首先就错了，因为你对我说了同样的话。}我说：\Quote{情况不同：我这样做是合适的；而你，作为至高天主的传报者，来拯救人的灵魂，对我说这样的话，实在令人悲伤。}伯多禄说：\Quote{我会同意你的看法，若非我们的主——祂为拯救整个世界而来，高于一切受造物——曾屈尊作了仆人，为要劝服我们，不要羞于向我们的弟兄履行仆人的职责。}我说：\Quote{我若以为能说服你，那便是愚昧；但我仍感谢天主的眷顾，因为我堪当以你代替我的父母。}\par

% structure: p0016 source=h2 target=subsection reason=relative-heading-rank; base=h2

\subsection{第八章  \Person{克肋孟}的家世}

\Person{伯多禄}说：\Quote{那么，你家里就没有其他人存活了吗？}我回答说：确实有许多显赫的人物，出自\Person{凯撒}的家族；因为凯撒本人曾将一位女子嫁给我父亲，这女子是他的亲戚，与他一同受过教育，出身于相当显赫的家族。我的父亲与她生了两个双胞胎儿子，在我之前出生，他们彼此不太相像——这是我父亲告诉我的，因为我从未见过他们。事实上，我对母亲也没有清晰的记忆；但我珍藏着她面容的印象，仿佛在梦中见过。我母亲名叫\Person{玛提狄亚}，我父亲名叫\Person{福斯蒂尼亚诺斯}：我的两个兄弟名叫\Person{福斯蒂诺斯}和\Person{福斯图斯}。当时我只有五岁，母亲见到一个异象——这是我从父亲那里听来的——警告她除非她迅速带着双胞胎儿子离开这座城市，并在外居住十年，否则她和孩子们将遭遇悲惨的命运。\par

% structure: p0018 source=h2 target=subsection reason=relative-heading-rank; base=h2

\subsection{第九章  母亲和兄弟的失踪}

于是我的父亲——他非常疼爱他的儿子——把他们连同他们的母亲送上一艘船，打发他们去\Place{雅典}受教育，并配备了男女仆人以及充足的钱财；只留下我作为他的安慰，并且为此庆幸，因为异象没有命令我也与母亲同去。一年后，我父亲派人带钱去雅典给他们，并想知道他们的近况；但派去的人再也没有回来。第三年，我悲伤的父亲又派了其他人带钱去，他们在第四年返回，说既没有见到我的母亲，也没有见到我的兄弟，他们从未到达雅典，也没有发现任何与他们同行之人的踪迹。\par

% structure: p0020 source=h2 target=subsection reason=relative-heading-rank; base=h2

\subsection{第十章  父亲的失踪}

\Quote{我父亲听了这话，因过度的悲伤而不知所措，不知该往何处去，也不知该往何处寻找，便带着我下到港口，开始询问水手们，是否有谁在四年前曾在某处看见或听说过一具母亲和两个幼童的遗体被冲上岸。一个人说一件事，另一个人说另一件事，但在这茫茫大海中，我们查问不出任何确切的消息。然而，我父亲由于对妻子和儿女的深厚感情，仍被虚幻的希望所喂养，直到他决定把我托付给监护人，留在\Place{罗马}——当时我已十二岁——而他自己去寻找他们。于是他哭着下到港口，上了一艘船，便出发了；从那时起直到现在，我再也没有收到过他的任何信件，也不知道他是死是活。但我更倾向于认为他也已丧生，或是由于心碎，或是由于海难；因为自那时起已过了二十年，我再也没有得到过他的任何音讯。}\par

% structure: p0022 source=h2 target=subsection reason=relative-heading-rank; base=h2

\subsection{第十一章  受苦对外教人和基督徒的不同影响}

\Person{伯多禄}听了这话，流下了同情的眼泪，对在座的朋友们说：\Quote{如果一位敬拜天主的人遭受了这个人的父亲所遭受的一切，人们会立刻将他宗教信仰归为他灾祸的原因；但当这些事发生在可怜的外邦人身上时，他们便把自己的不幸归咎于命运。我称他们为可怜的人，是因为他们在这里既被错误所困扰，又被剥夺了未来的希望；然而，当天主的敬拜者遭受这些事时，他们的耐心忍受有助于洗刷他们的罪恶。}\par

% structure: p0024 source=h2 target=subsection reason=relative-heading-rank; base=h2

\subsection{第十二章  前往\Place{阿剌多斯}的远行}

此后，在场的一人开始请求\Person{伯多禄}，在翌日清晨我们应前往一个邻近的岛屿，名为\Place{阿拉杜斯}，距离不超过六弗隆，为观看岛上一件奇妙的作品，即巨大的葡萄木柱子。\Person{伯多禄}同意了，因为他十分随和；但他吩咐我们，当我们离开船时，不可一拥而上地去看：他说：\Quote{因为我不愿你们被群众注意到。}因此，次日，我们在一小时内乘船到达该岛，立刻赶往那奇妙柱子的所在地。它们被安置在一座神庙中，里面有\Person{菲迪亚斯}极其宏伟的作品，我们每个人都认真凝视。\par

% structure: p0026 source=h2 target=subsection reason=relative-heading-rank; base=h2

\subsection{第十三章　讨饭的妇人。}

但\Person{伯多禄}只欣赏了柱子，丝毫未被绘画的优美所迷醉，他走出去，在门前看见一个穷妇人正向进去的人讨饭；他认真地看着她说：\Quote{告诉我，妇人，你身体缺少哪一部分，以致你屈尊乞讨，而不用天主给你的双手劳动谋生呢？}她叹息着说：\Quote{但愿我有能活动的手；但现在只保留了手的样子，因为它们已经毫无生气，被我咬得软弱无力，毫无感觉。}\Person{伯多禄}说：\Quote{什么原因让你对自己造成如此大的伤害？}\Quote{缺乏勇气，}她说，\Quote{别无他因；如果我有一点勇敢，我早就从悬崖跳下，或投身深海，结束我的痛苦了。}\par

% structure: p0028 source=h2 target=subsection reason=relative-heading-rank; base=h2

\subsection{第十四章　妇人的悲伤。}

\Person{伯多禄}说：\Quote{妇人，你以为那些自杀的人会免除痛苦，而不是那些对自己施加暴力的人的灵魂要受更大的惩罚吗？}她说：\Quote{但愿我确知灵魂在阴间活着，因为我乐意接受自杀的惩罚，只要能见到我的亲爱的孩子们，哪怕只有一小时。}\Person{伯多禄}说：\Quote{究竟是什么事使你如此悲伤？我很想知道。因为如果你告诉我原因，我或许能向你证明，妇人，灵魂确实在阴间活着；而且我能给你一些补救，代替悬崖或深海，使你能毫无痛苦地结束生命。}\par

% structure: p0030 source=h2 target=subsection reason=relative-heading-rank; base=h2

\subsection{第十五章　妇人的故事。}

那妇人听到这受欢迎的应许，便开始说：\Quote{我的出身或我的国家，既不易相信，也无需讲述。只须解释我悲伤的原因，即我为何咬自己的手以致无力。我生自高贵父母，嫁给一个相当有势力的丈夫，生了一对双胞胎儿子，之后又生了一个。但我丈夫的兄弟对我怀有强烈的非法情爱；因我将贞洁视为高于一切，既不愿同意如此大恶，也不愿向我丈夫揭露他兄弟的卑鄙，我考虑如何能不沾染污秽，又不使兄弟反目，从而将整个高贵家族的名声拖入耻辱。因此，我下定决心带着我的双胞胎离开我的国家，直到那种因我的容貌而助长和激发的乱伦之爱消退；我认为我们的另一个儿子应留下，以某种程度安慰他的父亲。}\par

% structure: p0032 source=h2 target=subsection reason=relative-heading-rank; base=h2

\subsection{第十六章　妇人的故事（续）}

\Quote{为了实施这个计划，我假装作了一个梦，在梦中有一位神祇显现站在我身边，告诉我应立即带着双胞胎离开那城，直到他命令我返回时才回来；如果我不这样做，我将与所有孩子一同灭亡。事情就这样做了。因为我一告诉丈夫这个梦，他就惊恐万状；他派我的双胞胎儿子以及男女奴仆与我同行，给了我许多钱，命我航行到\Place{雅典}，在那里教育我的儿子，并待在那里直到那命令我离开的神允许我返回。当我与儿子们航行时，夜间因狂风大作而遭遇海难，我这个可怜人被吹到这里；当所有人都丧生时，一个巨浪抓住我，将我抛到岩石上。当我坐在那里，仅存一线希望，或许能找到我的儿子们时，我没有投身深海，尽管那时我的灵魂因悲伤而纷乱沉醉，既有勇气也有力量这样做。}\par

% structure: p0034 source=h2 target=subsection reason=relative-heading-rank; base=h2

\subsection{第十七章　妇人的故事（续）}

\Quote{但天亮时，我哭喊着环顾四周，看我那不幸的儿子们的尸体是否被冲上岸边，一些看见我的人动了怜悯之心，先在海面上搜索，然后沿着海岸，看能否找到我的任何一个孩子。但哪里都找不到他们，当地的妇人们怜悯我，开始安慰我，每个人都讲述自己的悲伤，希望我能从她们与我相似的灾难中获得安慰。但这反而使我更加悲伤；因为我的性格不是那种能把别人的不幸当作安慰的人。当许多人想热情接待我时，一位住在这里的穷妇人硬要我进入她的小屋，说她的丈夫是个水手，年轻时就死于海中，虽然后来许多人向她求婚，但她因为爱丈夫而宁愿守寡。她说：\InnerQuote{因此，我们将分享用我们双手劳动所能获得的一切。}}\par

% structure: p0036 source=h2 target=subsection reason=relative-heading-rank; base=h2

\subsection{第十八章　妇人的故事（续）}

\Quote{而且，为了不用一个冗长无益的故事耽搁您，我因为她对丈夫所持的忠实情感而乐意与她同住。但不久之后，我的手（我真是个不幸的女人！），长期被啃咬而撕裂，变得无力；收留我的那个女人患了瘫痪，现在躺在床上；那些从前怜悯我的女人的情感也冷淡了。我们两人都无助。如您所见，我坐着乞讨；当我得到一点东西时，一餐饭两个人分。看，现在您已听够了我的事；为什么您不履行您的诺言，给我一种补救方法，使我们俩都能无痛苦地结束这悲惨的生命？}\par

% structure: p0038 source=h2 target=subsection reason=relative-heading-rank; base=h2

\subsection{第十九章。\Person{伯多禄}对故事的反思。}

当她说话时，\Person{伯多禄}因思绪纷乱而如同被雷击一般站着；我\Person{克莱孟}走上前说：\Quote{我到处找你，现在我们该怎么办？}但他命令我先行到船那里等他；由于不可违抗，我照做。但\Person{伯多禄}（后来他将一切告诉了我）起了某种疑心，向那妇人询问她的家庭、她的家乡和她儿子的名字；\Quote{你立刻，}他说，\Quote{如果你告诉我这些事，我就会给你补救。}但她像是受着逼迫，因为不愿承认这些事，却又渴望补救，便编造谎言，说她是\Place{厄弗所}人，她的丈夫是\Place{西西里}人，并给她儿子们起假名字。于是\Person{伯多禄}以为她回答了实情，便说：\Quote{哎呀！妇人，我想今天会有极大的喜乐临于我们；因为我怀疑你是我最近得知某些类似事情的那个妇人。}但她恳求他说：\Quote{我求你把那些事告诉我，好让我知道在女人中是否有比我更不幸的。}\par

% structure: p0040 source=h2 target=subsection reason=relative-heading-rank; base=h2

\subsection{第二十章。\Person{伯多禄}对妇人的陈述。}

于是，\Person{伯多禄}——他不能欺骗，且被怜悯所动——开始说：\Quote{在我的追随者中（他们因宗教和教派而跟随我）有一个年轻\Place{罗马}公民，他告诉我，他有一位父亲和两个孪生兄弟，但他们一个也不在了。\InnerQuote{我的母亲，}他说，\InnerQuote{正如我从父亲那里得知，她看见一个异象，她必须带着她的孪生儿子暂时离开\Place{罗马}城，否则他们会死于可怕的死亡；当她离开后，就再也没有人见过她。}后来，他的父亲动身去寻找他的妻子和儿子，也失踪了。}\par

% structure: p0042 source=h2 target=subsection reason=relative-heading-rank; base=h2

\subsection{第二十一章。一个发现。}

当\Person{伯多禄}这样说时，那妇人惊愕得昏了过去。于是\Person{伯多禄}开始扶起她，安慰她，问她发生了什么事，或她受了什么苦。但她终于艰难地恢复了呼吸，振作起精神，迎接她所希望的巨大喜乐，同时擦着脸说：\Quote{你所说的那个年轻人在这里吗？}\Person{伯多禄}明白事情后，说：\Quote{你先告诉我，否则你不能见他。}于是她说：\Quote{我是那年轻人的母亲。}\Person{伯多禄}说：\Quote{他叫什么名字？}她回答：\Quote{\Person{克莱孟}。}于是\Person{伯多禄}说：\Quote{正是他；刚才和我说话的就是他，我已经吩咐他先到船那里等我。}于是她俯伏在\Person{伯多禄}脚前，开始恳求他赶快到船那里去。\Person{伯多禄}说：\Quote{行，只要你答应我按我说的去做。}她说：\Quote{我什么都做；只是让我见我的独生子，因为我想在他身上也会看到我的孪生子。}\Person{伯多禄}说：\Quote{你见到他后，暂且掩饰一下，直到我们离开这座岛。}她回答说：\Quote{我会这样做的。}\par

% structure: p0044 source=h2 target=subsection reason=relative-heading-rank; base=h2

\subsection{第二十二章。幸福的相遇。}

于是\Person{伯多禄}拉着她的手，领她到船那里。当我看见他伸手给那妇人时，我开始笑起来；然而，为了表示尊敬，我试图用我的手替换他的手，去扶那妇人。但一碰到她的手，她就发出一声大叫，冲进我的怀抱，开始以母亲的吻吞噬我。但我，不知情，把她当作疯女人推开；同时，尽管怀着敬畏，我对\Person{伯多禄}有些生气。\par

% structure: p0046 source=h2 target=subsection reason=relative-heading-rank; base=h2

\subsection{第二十三章。一个奇迹。}

但\Person{伯多禄}说：\Quote{住手：你这是什么意思，哦，\Person{克莱孟}，我的儿子？不要推开你的母亲。}我听到这话，立刻泪流满面，扑倒在我母亲身上（她已倒在地上），开始吻她。因为我一听，就渐渐记起了她的面容；我注视得越久，就越觉得熟悉。与此同时，一大群人聚集起来，听说那个常坐着乞讨的妇人被她的儿子认出来了（他是个好人）。当我们想迅速离开这座岛时，我母亲对我说：\Quote{我亲爱的儿子，我应该向收留我的妇人告别；因为她贫苦、瘫痪、卧床不起。}\Person{伯多禄}和所有在场的人听到这话，都赞叹那妇人的善良和明智；于是\Person{伯多禄}立刻派人去把那妇人连她的床一起抬来。当她被抬来放在人群中间时，\Person{伯多禄}当着众人说：\Quote{如果我是真理的宣讲者，为了坚振所有在场者的信德，使他们知道并相信有一位\OriginalTerm{天主}{God}，祂创造了天地，因祂的圣子\Person{耶稣基督}之名，让这妇人起来吧。}他一说完，她就完全好了，起来俯伏在\Person{伯多禄}脚前；她以亲吻问候她的朋友和熟人，问她这一切是什么意思。但\Person{伯多禄}简短地向她叙述了\WorkTitle{互认}的整个经过，使周围的群众感到惊讶。\par

% structure: p0048 source=h2 target=subsection reason=relative-heading-rank; base=h2

\subsection{第二十四章。离开\Place{阿辣多斯}。}

于是伯多禄尽可能利用时间，向群众宣讲对天主的信仰和宗教的规条，然后补充说，若有人想知道更多关于这些事的确切道理，他该到安提约基雅来，\Quote{那里}，他说，\Quote{我们决意停留三个月，并要完备地讲授有关救恩的事。因为假如}，他说，\Quote{人们为了经商或从军而离开家乡和父母，并不害怕远途航行，那么为了永生而离家三个月，又怎会觉得沉重或困难呢？}他说完这些及类似的话后，我向那位曾款待我母亲、并藉伯多禄恢复健康的妇女呈上一千德拉克马，并在众人面前将她托付给该城的一位首领，此人是个好人，他承诺乐意办理我们所请求的事。我也向其他一些人、以及那些据说曾在我母亲患难时安慰过她的妇女们分发了一些钱，并向她们表达了感谢。此后，我们与母亲一同起航，前往\Place{安塔拉杜斯}。\par

% structure: p0050 source=h2 target=subsection reason=relative-heading-rank; base=h2

\subsection{第二十五章 旅途}

及至我们来到住所，我母亲开始问我父亲的下落，我告诉她父亲前去寻找她，再也没有回来。她听了只是叹息，因为因我而有的巨大喜乐减轻了她其他的悲伤。次日，她与我们同行，坐在伯多禄的妻子身旁；我们到了\Place{巴兰内亚}，在那里停留三天，然后前往\Place{帕托斯}，又到\Place{加巴拉}，最后抵达\Place{劳迪西亚}。尼塞塔和阿奎拉在城门前迎接我们，亲吻我们后，领我们到住所。伯多禄见这是一座宏大而辉煌的城市，便说值得在此停留十天，甚至更久。于是尼塞塔和阿奎拉问我这位陌生的妇人是谁，我回答：\Quote{这是我的母亲，是天主藉我主伯多禄赐还给我的。}\par

% structure: p0052 source=h2 target=subsection reason=relative-heading-rank; base=h2

\subsection{第二十六章 综述}

我这样说了之后，伯多禄便开始按顺序向他们讲述整个事情，他说：我们到了\Place{阿拉杜斯}，我曾吩咐你们先走，你们离开后的同一天，克莱孟在谈话中向我讲述了他的出身和家世，以及他如何失去父母，还有两个比他年长的孪生兄弟，据他父亲说，他母亲曾见过一个异象，命令她带着孪生兄弟离开罗马城，否则她和他们会突然灭亡。她把梦告诉了父亲，父亲出于对儿子的慈爱，害怕他们遭遇不测，便将妻子和儿子送上船，备齐一切必需品，送往雅典接受教育。后来他屡次派人打听他们，却毫无踪影。最后父亲亲自去寻找，至今下落不明，\Emph{毫无音讯}。克莱孟向我叙述这些事后，有人来请我们到附近的\Place{阿拉杜斯}岛，去看一些非常巨大的葡萄木柱子。我同意了；到达后，其他人都进入神庙内部，而我——不知为何——无心再往前走。\par

% structure: p0054 source=h2 target=subsection reason=relative-heading-rank; base=h2

\subsection{第二十七章 综述续}

\Quote{但我正在外面等候他们时，开始注意到这位妇人，并奇怪她身体哪部分残疾，以致不靠双手劳动谋生，却忍受乞讨的羞辱。我便问她其中的原因。她承认自己出身高贵，嫁给了一位同样高贵的丈夫，\InnerQuote{他的兄长}，她说，\InnerQuote{因对我怀有不法的爱，想玷污兄弟的床榻。我厌恶此事，但又不敢将如此大恶告诉丈夫，唯恐引起兄弟间的争斗，使家族蒙羞，便决定带我的两个孪生子离开家乡，留下小儿子安慰父亲。为使这事显得正当，我认为最好假托一个梦，告诉丈夫说，我在异象中看见一位神明，命我立即带着两个孪生子离开该城，直到他指示我返回为止。}她告诉我，她丈夫听后相信了她，便将她和孪生儿子送往雅典接受教育。但途中遭遇可怕的风暴，被冲到那座岛上，船只破碎，她被一个浪头托到岩石上，唯独为此而拖延自杀，\InnerQuote{直到}，她说，\InnerQuote{我能至少拥抱我可怜儿子的尸体，并埋葬他们。但天明时，许多人聚集起来，可怜我，给我披了一件衣服。但我这不幸的人，流着许多眼泪恳求他们，看是否能找到我可怜儿子的尸体。我甚至用牙齿撕裂全身，哀嚎哭喊，不住地喊道：我这苦命的女人，我的法乌斯都在哪里？我的法乌斯提努斯在哪里？}}\par

% structure: p0056 source=h2 target=subsection reason=relative-heading-rank; base=h2

\subsection{第二十八章 更多的相认}

当伯多禄这样说时，尼塞塔和阿奎拉突然跳起来，惊异不已，开始极其激动地说：\Quote{主啊，万有的主宰和天主，这些事是真的吗？还是我们在做梦？}然后伯多禄说：\Quote{除非我们疯了，这些事是真的。}他们稍停片刻，擦了擦脸，说：\Quote{我们是福斯蒂努斯和福斯图斯。从一开始你讲述这事时，我们立刻怀疑你所说的事或许与我们有关；但再一想，人生中许多类似的事常常发生，我们便保持了沉默，尽管我们的心被某种希望所触动。因此我们等待你的故事结束，如果完全清楚是关乎我们，那时我们就承认。}他们这样说后，哭着走到我们母亲那里。见她睡着了，想要拥抱她，伯多禄阻止他们说：\Quote{先让我预备你母亲的心，免得突然的极大喜乐使她失去理智，神志混乱，尤其她现在因睡眠而昏沉。}\par

% structure: p0058 source=h2 target=subsection reason=relative-heading-rank; base=h2

\subsection{第二十九章．\Quote{没有什么是不洁的与不洁净的}}

因此，当我们母亲睡醒后，伯多禄开始对她说：\Quote{妇人，我希望你知道我们宗教的一项规矩。我们敬拜创造世界的唯一天主，并遵守祂的法律，其中祂首先命令我们敬拜祂，尊敬祂的圣名，孝敬父母，保持贞洁和正直。但我们也遵守这一点：不与外邦人共桌，除非他们信了，领受了真理，受了洗，并藉着对圣名的一种三重呼求而被祝圣；然后我们才与他们共食。否则，即使是父亲或母亲、妻子、儿子或兄弟，我们也不能与他们共桌。因此，既然我们为宗教的特殊原因这样做，请你不觉得难以接受：你的儿子不能与你共食，除非你对他所信的道理有同样的判断。}\par

% structure: p0060 source=h2 target=subsection reason=relative-heading-rank; base=h2

\subsection{第三十章．\Quote{谁能阻止水呢？}}

她听了这话，说：\Quote{今天有什么事阻止我受洗呢？因为甚至在见到你之前，我就完全远离了他们称之为神明的那些东西，因为他们不能为我做任何事，尽管我常常几乎天天向他们献祭。至于贞洁，我该说什么呢？既不是从前享乐欺骗过我，也不是后来的贫穷迫使我去犯罪。但我想你很清楚我多么热爱贞洁，因为我假托那个梦，逃避了不洁之爱的陷阱，并得以带着我的双胞胎外出，且留下我这儿子克莱孟作为他父亲的安慰。因为如果两个孩子对我来说都不够，那么如果他们父亲一个都没有，岂不更令他伤心？他因为对儿子的深情而可怜，连梦的权威都几乎不能说服他把福斯蒂努斯和福斯图斯——这克莱孟的兄弟——交给我，而他自己只满足于克莱孟一人。}\par

% structure: p0062 source=h2 target=subsection reason=relative-heading-rank; base=h2

\subsection{第三十一章．过度喜乐}

她还在说话时，我的兄弟们再也忍不住，流着泪冲进母亲的怀抱，亲吻她。但她说：\Quote{这是怎么回事？}伯多禄说：\Quote{妇人，不要惊慌；要坚强。这些是你的儿子福斯蒂努斯和福斯图斯，你原以为他们已葬身深海；但他们如何活着，如何在那可怕之夜逃脱，以及其中一个如何叫尼塞塔，另一个叫阿奎拉，他们自己会向你解释，我们也将与你一同听。}伯多禄这样说后，我们的母亲因过度喜乐而昏厥；过了一会儿，她苏醒过来，说：\Quote{亲爱的儿子们，我恳求你们，告诉我自从那悲惨残酷的夜晚以来，你们遭遇了什么。}\par

% structure: p0064 source=h2 target=subsection reason=relative-heading-rank; base=h2

\subsection{第三十二章．\Quote{他领他们到达渴望的港湾}}

于是尼塞塔开始说：\Quote{母亲，在那船破碎的夜晚，我们被抛在海中，靠着一块破船的碎片支撑，一些以海上抢劫为生的人发现了我们，把我们放在他们的船上，通过奋力划桨战胜波涛，经过各种曲折，把我们带到凯撒勒雅·斯特拉托尼斯。在那里，他们使我们挨饿，殴打我们，恐吓我们，使我们不敢说出真相；他们改了我们的名字，把我们卖给一个非常尊贵的寡妇，名叫犹斯塔。她买了我们，待我们如同儿子，精心教育我们学习希腊文学和自由艺术。长大后，我们也从事哲学研究，以便能驳倒外邦人，用哲学辩论支持神圣宗教的教义。}\par

% structure: p0066 source=h2 target=subsection reason=relative-heading-rank; base=h2

\subsection{第三十三章．另一场海难被阻止}

\Quote{但出于友谊和少年同伴关系，我们依附于一个与我们一同受教育的术士西满，几乎被他欺骗。因为在我们宗教中提及一位先知，所有遵守那宗教的人都盼望祂的来临，通过祂，永生和幸福的生命应许给信祂的人。我们曾以为这西满就是祂。但母亲，这些事将在更合适的时候向你解释。与此同时，当我们几乎被西满欺骗时，我主伯多禄的一位同工，名叫匝凯，警告我们不要被那术士迷惑，并在伯多禄抵达时把我们引荐给他，好由他教导我们正确和完全的事。我们希望这事也发生在你身上，正如天主恩赐给我们一样，使我们能与你共桌共食。母亲，因此，你相信我们淹死在海里，而我们却被海盗偷走了。}\par

% structure: p0068 source=h2 target=subsection reason=relative-heading-rank; base=h2

\subsection{第三十四章．圣洗之前必须守斋}

我们的母亲跪在伯多禄脚前，恳求并哀求他，让自己和她的女主人即刻受洗。她说：\Quote{那样，我连一天也不能失去与我儿子们的相伴和共处。}同样，我们她的儿子们也与她一同恳求伯多禄。但他说：\Quote{怎么？你们以为只有我是无怜悯的，不愿你们享受与母亲同桌的欢聚吗？但她必须先至少禁食一天，然后才受洗；这是因为我从她口中听到一个宣告，这宣告向我显明了她的信德，并证实了她的信仰；否则，她本应受教导和教训许多日子，才可受洗。}\par

% structure: p0070 source=h2 target=subsection reason=relative-heading-rank; base=h2

\subsection{第三十五章 渴望他人得救}

于是我说：\Quote{我请求您，我主伯多禄，告诉我们那是什么宣言，您说那为您提供了她信德的证据？}接着伯多禄说：\Quote{她要求她所爱的女主人与她一同受洗，以报答其恩惠。除非她相信洗礼中有某种伟大的恩赐，否则她不会为她所爱之人请求这份恩宠。因此，我也责备许多人，他们自己受洗并相信，却没有为他们所爱的人——如妻子、儿女或朋友——做出与信仰相称的事，他们没有劝勉那些人追求自己所获得的，\Emph{一如他们若真相信永生由此赐予就会做的}。简言之，若他们看见这些人患病或面临任何身体危险，他们便忧伤哀悼，因为确信这威胁着他们的毁灭。所以，若他们确知不敬拜天主的人要受永火惩罚，他们何时才会停止警告和劝勉呢？或倘若他们拒绝，他们又怎能不忧伤哀悼，确知永罚在等待他们呢？如今，我们立刻派人去请那妇人，看她是否喜爱我们宗教的信仰；我们按所见行事。但既然你的母亲对洗礼判断得如此忠诚，让她在受洗前只禁食一天。}\par

% structure: p0072 source=h2 target=subsection reason=relative-heading-rank; base=h2

\subsection{第三十六章 儿子们的恳求}

但她在我主伯多禄的妻子面前发誓说，自从认出儿子以来，因过度喜乐，她无法进食任何食物，除了昨天喝了一杯水。伯多禄的妻子也作证说，确实如此。于是阿基拉说：\Quote{那么，有什么妨碍她受洗呢？}接着伯多禄微笑着说：\Quote{但这并不是洗礼的禁食，因为她不是为了受洗而禁食。}于是尼塞塔说：\Quote{但或许天主愿意：在我们认母之后，我母亲甚至连一天也不应与我们分开共桌，为此预定了这禁食。因为正如她在无知中保持了贞洁，以助她获得洗礼的恩宠；同样，她在知道禁食的理由之前禁食，以助她获得洗礼，并使我们从相识之初，她就能立刻享受与我们共桌的交往。}\par

% structure: p0074 source=h2 target=subsection reason=relative-heading-rank; base=h2

\subsection{第三十七章 伯多禄不为所动}

于是伯多禄说：\Quote{愿那邪恶者不得利用母爱胜过我们；但愿你们和与你们同处的我，今天与她一同禁食，明天她将受洗；因为为了任何人或友谊而松弛弱化真理的规条是不对的。那么，让我们不要退缩与她一同受苦，因为违反任何诫命都是罪。但我们要教导我们外在的肉体感官，使它们服从我们内在的感官；不要强迫我们那品味天主之事的内心感官，去跟随那品味肉体之事的外在感官。为此，主也命令说：\BibleQuote{凡注视妇女，有意贪恋她的，他已在心里奸淫了她。}并且又加上：\BibleQuote{若你的右眼使你跌倒，剜出它来，从你身上扔掉；因为丧失你一个肢体，比你全身投入地狱之火，为你有益。} \BibleRef{玛 5:28}祂不是说\Emph{已经使你跌倒}，叫你犯罪之后才除去罪的缘由；而是说\Emph{若它使你跌倒}，即在你犯罪之前，砍断那激惹、刺激你的罪因。但兄弟们，你们谁也不可认为主赞同切除肢体。祂的意思是，应当切除意向，而非肢体，以及那些诱人犯罪的缘由，使我们的思想，乘着视觉之车，在肉体感官的支持下，趋向于天主的爱；而不让肉体的眼睛像放荡的马匹一样松开缰绳，急于在诫命之道外奔跑，而应使肉体的视觉服从心灵的判断，不许我们的双眼——天主使之成为祂工程的观看者和见证者——成为邪恶欲望的淫媒。因此，让肉体的感官和内心的思想都服从天主的律法，让它们事奉祂的旨意，承认自己是祂的工程。}\par

% structure: p0076 source=h2 target=subsection reason=relative-heading-rank; base=h2

\subsection{第三十八章 贞洁的赏报}

因此，正如奥迹所要求的秩序与理智，次日她在海中领受了圣洗，回到住所后，按序领受了宗教的一切奥迹。我们她的儿子，\Person{尼塞塔}、\Person{阿奎拉}和我\Person{克肋孟}，都在场。之后，我们与她一同用餐，与她一起光荣天主，衷心感谢\Person{伯多禄}的热忱和教导，他通过我们母亲的例子向我们显示，贞洁的善不会在天主面前失落；\Quote{相反，}他说，\Quote{不贞洁不会逃脱惩罚，虽然它可能不会立即受罚，而是缓慢地。但如此令天主喜悦的，}他说，\Quote{是贞洁，以至于它甚至在现世生活中赐予那些仍在谬误中的人一些恩宠；因为未来的福乐只为那些藉圣洗的恩宠保持贞洁与正义的人而存留。简言之，发生在你们母亲身上之事就是此例，因为所有这些福祉都已恢复给她，作为她贞洁的赏报；为守护和保持此贞洁，仅靠节欲并不足够；但当任何人察觉到陷阱和欺骗正在预备时，他必须立刻逃离，如同躲开烈火或疯狗的攻击，而不应相信他可以通过哲学思辨或迎合它们来轻易挫败此类陷阱；但是，如我所说，他必须逃离并远离，正如你们的母亲也通过她对贞洁真实而完整的爱所行的。因此，她得以保全给你们，而你们也得以保全给她；此外，她还被赋予了永生的知识。}他说了这些以及许多类似的话后，傍晚来临，我们就去睡了。\par

\SourceRecord{Translated by Thomas Smith.；Ante-Nicene Fathers；Vol. 8.；Edited by Alexander Roberts, James Donaldson, and A. Cleveland Coxe.；Buffalo, NY: Christian Literature Publishing Co.,；1886.；<http://www.newadvent.org/fathers/080407.htm>.}

% structure: p0001 source=h1 target=section reason=page-title

\section{\WorkTitle{认亲录（卷八）}}

% structure: p0002 source=h2 target=subsection reason=relative-heading-rank; base=h2

\subsection{第一章 老工人}

如今第二天早晨，伯多禄带着我的兄弟和我同去，我们下到港口在海中沐浴，此后我们退到一处隐秘的地方祈祷。但有一个贫穷的老人，一个工人——从他的衣着可以看出来——开始热切地观察我们，而我们没有看见他，他想要看看我们在秘密中做什么。当他看见我们祈祷时，他等我们出来，然后向我们致意，说：\Quote{如果你们不介意，不把我当作一个好奇而好管闲事的人，我希望和你们交谈；因为我怜悯你们，不愿你们在真理的外表下犯错，惧怕不存在的事物；或者如果你们认为其中有什么真理，就请告诉我。因此，如果你们耐心地接受，我可以用几句话教导你们什么是正确的；但如果你们不喜欢，我就离开去做我的事。} 伯多禄回答他说：\Quote{请说出你认为是好的话，我们将乐意倾听，无论是真是假；因为你应当受到欢迎，就像一个为子女担忧的父亲，愿意把你看为好的东西赐给我们。}\par

% structure: p0004 source=h2 target=subsection reason=relative-heading-rank; base=h2

\subsection{第二章 命数}

然后老人接着说：\Quote{我看见你们在海中沐浴，然后退到一处隐秘的地方；因此，在你们没有察觉的情况下，我观察你们在做什么，看见你们祈祷。所以，我怜悯你们的错误，等你们出来，好对你们说话，教导你们不要在这种仪式上犯错；因为世上既没有天主，也没有任何敬拜，也没有天主眷顾，一切都是由偶然的机遇和\Emph{\OriginalTerm{命数}{genesis}}所做和所支配的，正如我通过学习最清楚地向自己发现的。所以，不要犯错：因为无论你们祈祷与否，无论你们的\Emph{\OriginalTerm{命数}{genesis}}包含什么，那必会临到你们。} 当时我克肋孟不知为何心中感动，回想起他许多似乎熟悉的事；因为有人说得对，凡是出于一个人的，即使久别，亲情的火花也永不熄灭。因此我开始问他是谁，从何处来，家世如何。但他不愿回答这些问题，说：\Quote{这与我告诉你们的有什么关系？但首先，如果你们愿意，让我们讨论我们提出的那些问题；之后，如果情况需要，我们可以像朋友对朋友那样，彼此透露姓名、家族、国籍以及其他与此相关的事。} 然而我们所有人都钦佩这个人的口才、品性的庄重和言谈的平静。\par

% structure: p0006 source=h2 target=subsection reason=relative-heading-rank; base=h2

\subsection{第三章 友好会谈}

但伯多禄在漫步交谈时，正寻找适合会谈的地方。当他看到港口附近一处安静的角落时，就让我们坐下；于是他自己先开始。他并没有轻视那位老人，也没有因他衣着破旧而看不起他。因此，他说：\Quote{既然在我看来你是个有学问且有同情心的人，因为你来到我们这里，希望把你看为好的东西告诉我们，我们也愿意向你们阐明我们认为良善和正确的事；如果你不认为它们是真实的，你会善意地接受我们对你的善意，如同我们接受你的善意一样。} 伯多禄正这样说着的时候，一大群人聚集起来。然后老人说：\Quote{也许人群的在场使你不安。} 伯多禄回答说：\Quote{一点也不，仅此一点：恐怕在我们讨论中真理显明时，你在众人面前羞于承认并同意你听明白的真实之事。} 对此老人回答说：\Quote{我年老还不至于如此愚蠢，以致明白了什么是真实的，却为了讨好群众而否认它。}\par

% structure: p0008 source=h2 target=subsection reason=relative-heading-rank; base=h2

\subsection{第四章 问题陈述}

于是伯多禄开始说：\Quote{那些宣讲真理之言、启迪人灵魂的人，在我看来就像太阳的光芒，一旦出现并向世界显现，就不再能被隐藏或遮掩，它们与其说是被人们看见，不如说是赐予所有人视力。因此，有一位对真理的宣告者说得很好：\BibleQuote{你们是世界的光，建在山上的城是不能隐藏的；人点灯，并不是放在斗底下，而是放在灯台上，照耀屋中所有的人。}\BibleRef{玛 5:14}} 然后老人说：\Quote{他说得好，无论他是谁。但你们中一人应说明，根据他的意见，应当遵循什么，以便我们将言论导向明确的目标。因为为了找到真理，仅仅推翻对方所说的话是不够的，而且自己也应该提出对方可能反对的东西。因此，为了使双方处于平等地位，我认为合理的是我们每人先陈述自己所持的意见。如果你们愿意，我先开始。那么，我说，世界不是按照天主的眷顾所治理的，因为我们看到世上许多事情做得不义和无序；但我说，是\OriginalTerm{命数}{genesis}在做并支配一切。}\par

% structure: p0010 source=h2 target=subsection reason=relative-heading-rank; base=h2

\subsection{第五章 允许自由讨论}

当\Person{伯多禄}正要回答时，\Person{尼塞塔}抢先说道：\Quote{我的主\Person{伯多禄}，请允许我回答这个问题；不要认为我这样一个年轻人与一位老人争论是失礼，而是让我如同儿子与父亲交谈一样。} 老人说：\Quote{我的儿子，我不仅希望你陈述你的见解，而且如果你的同伴中任何人，甚至旁观者中任何人，认为自己知道什么，也让他毫不犹豫地陈述出来：我们将乐意聆听；因为通过许多人的贡献，未知之事更容易被发现。} 于是\Person{尼塞塔}回答说：\Quote{不要以为我鲁莽行事，我的父亲，因为我打断了我主\Person{伯多禄}的发言；相反，我这样做是为了尊敬他。因为他是一位充满一切知识的天主之人，甚至不缺乏希腊学问，因为他充满了天主之神，对祂而言没有什么是未知的。但由于适合他谈论天上的事，我将回答那些关于希腊人空谈的事。但在我们以希腊方式辩论之后，当我们到达似乎没有结果的时刻，那时他本人，充满天主的智慧，将公开而清楚地向我们揭示所有事情的真理，这样不仅我们，而且所有在周围聆听的人，都将学到真理之路。因此现在让他作为裁判；当我们两人中有一人屈服时，让他接手此事，给出无可置疑的判决。}\par

% structure: p0012 source=h2 target=subsection reason=relative-heading-rank; base=h2

\subsection{第六章 问题的另一面陈述。}

当\Person{尼塞塔}这样说了以后，那些聚集的人彼此交谈说：\Quote{这就是我们听说的那位\Person{伯多禄}吗？他是在\Place{犹太}显现并行了许多征兆和奇迹的那位最受称许的门徒？} 他们怀着极大的敬畏和尊敬注视着他，将主的忠仆的荣誉归于主。\Person{伯多禄}看到这情形，对他们说：\Quote{让我们全神贯注地聆听，对每人所说的话保持公正的判断；在他们的辩论之后，我们也要添加似乎必要的内容。} \Person{伯多禄}说了这话，群众就欢喜起来。于是\Person{尼塞塔}开始如下发言：\Quote{我的父亲，你提出了世界不受天主眷顾治理，而是所有事物——无论是与每个人的性情相关，还是与行为相关——都受\OriginalTerm{宿命}{genesis}支配。这一点我可以立即回答；但因为遵守秩序是正当的，我们也陈述我们的立场，正如你自己所请求的那样。我说世界受天主眷顾治理，至少在其需要祂治理的事务上。因为唯有祂——那创造了世界的公义的天主——亲手掌握万有，祂必将在某一时刻按照每个人的行为给予报应。现在，你有了我们的立场；请继续，随你喜好，要么推翻我的立场，要么建立你自己的立场，这样我就能应对你的陈述。或者如果你希望我先发言，我将毫不犹豫。}\par

% structure: p0014 source=h2 target=subsection reason=relative-heading-rank; base=h2

\subsection{第七章 道路扫清。}

老人回答说：\Quote{我的儿子，你愿意先发言，还是更愿意我先说，这没有区别，尤其是在友好讨论的人之间。然而，你先说，我将乐意聆听；我希望你甚至能够在我所说的那些事情之后，提出与它们相反的观点，并通过两者的比较来展示真理。} \Person{尼塞塔}回答：\Quote{如果您愿意，我甚至可以陈述您一方的论点，然后进行答辩。} 老人说：\Quote{先向我展示你如何能知道我尚未说出的事情，这样我就能相信你能把握我一方的论点。} \Person{尼塞塔}说：\Quote{您的学派，甚至从您所提出的命题中，就对那些精通此类学说的人显明出来；其后果也是确定的。因为我不无知于哲学家们的学说，所以我知道从您所提出的事物中会产生什么；尤其因为我更常光顾\Person{伊壁鸠鲁}学派，胜过其他哲学家的学派。但我的兄弟\Person{阿奎拉}更关注\OriginalTerm{皮浪主义者}{Pyrrhonists}，我们的另一位兄弟关注\OriginalTerm{柏拉图主义者}{Platonists}和\OriginalTerm{亚里士多德主义者}{Aristotelians}；因此，您面对的是有学识的听众。} 老人说：\Quote{你们已经很好地、合乎逻辑地告诉我们，你们如何察觉到从所陈述的命题中得出的结论。但我所主张的比\Person{伊壁鸠鲁}的教条更多；因为我引入了\OriginalTerm{宿命}{genesis}，并断言它是人类所有行为的起因。}\par

% structure: p0016 source=h2 target=subsection reason=relative-heading-rank; base=h2

\subsection{第八章 本能。}

我\Person{克莱孟}对他说：\Quote{听，我的父亲：如果我的兄弟\Person{尼塞塔}使您承认世界并非没有天主眷顾而治理，那么我就能回答您那部分关于\OriginalTerm{宿命}{genesis}的剩余问题；因为我非常熟悉这一学说。} 我这样说了以后，我的兄弟\Person{阿奎拉}说：\Quote{我们称他为父亲有什么用处？因为我们被告诫不要称地上任何人为父亲。}\BibleRef{玛 23:9} 然后，他转向老人说：\Quote{不要见怪，我的父亲，因为我对我的兄弟称您为父亲提出了批评；因为我们有一条诫命，不要称任何人为父亲。} 当\Person{阿奎拉}说这话时，所有旁观的集会者、老人和\Person{伯多禄}都笑了。\Person{阿奎拉}问他们为何都笑时，我对他说：\Quote{因为你自己正做了你在别人身上所批评的事；因为你称老人为父亲。} 但他否认说：\Quote{我没有意识到我称他为父亲。} 同时，\Person{伯多禄}被某种怀疑所触动，正如他后来告诉我们的；他看向\Person{尼塞塔}，说：\Quote{继续你提出的内容。}\par

% structure: p0018 source=h2 target=subsection reason=relative-heading-rank; base=h2

\subsection{第九章 单纯与复合。}

于是\Person{尼塞塔}开始如下说：\Quote{一切存在的事物，要么是单纯的，要么是复合的。单纯的事物没有数目、可分性、颜色、差异、粗糙、光滑、重量、轻便、性质、数量，因此也没有终结。但复合的事物则由两个、三个、甚至四个\Emph{元素}，或者至少由几个元素复合而成；而复合的事物也必然能被分割。}老人听了这话，说：\Quote{你说得极好，极有学识，我的孩子。}然后\Person{尼塞塔}继续说：\Quote{因此，那单纯的事物，没有任何使存在物可被分解的东西，无疑是无法理解、无限的，不知起始也不知终结，所以它是独一、唯一的，且无需作者而存在。但复合的事物则受制于数目、多样性和可分性——它必然由某个作者复合而成，是一种多样性集合为一种类。因此，无限者，就善而言是父，就能力而言是造物主。在无限者中，创造的能力不能停止，善也不能止息；而是被善推动以改变现有之物，被能力推动以安排并坚固它们。因此，正如我们所说，有些事物被改变，由两个或三个元素组成，有些由四个，另一些由更多元素组成。但由于我们目前探讨的是世界的方法和实质，而公认世界由四个元素复合而成，以上所提的十种区别都属于它，让我们从这些较低的阶梯开始，并上升到更高的。因为从我们眼见手触的事物，有一条路通向理智和不可见的事物；正如算术教导中所包含的：当探求神性事物时，我们从较低数字上升到较高数字；但当阐释现存和可见事物的方法时，顺序则从较高数字导向较低数字。难道不是这样吗？}\par

% structure: p0020 source=h2 target=subsection reason=relative-heading-rank; base=h2

\subsection{第十章 创造暗示眷顾}

于是老人说：\Quote{你论证得非常好。}然后\Person{尼塞塔}说：现在，我们必须探求世界的方法；其中首要的探求分为两部分。问题在于：世界是被造的的还是未被造的？如果它未被造，那么它本身就是那\OriginalTerm{非受生者}{Unbegotten}，万物从它而来。但如果它是被造的，那么这个问题又分为两部分：它是自己造的还是由另一者所造？如果它真是自己造的，那么无疑眷顾就被排除了。如果不承认眷顾，那么心智就徒然被激发向德，正义徒然被持守，因为没有谁按功劳回报义人。甚至灵魂本身也似乎不是不朽的，如果没有眷顾的安排来接引它离开身体。\par

% structure: p0022 source=h2 target=subsection reason=relative-heading-rank; base=h2

\subsection{第十一章 普遍或特殊的眷顾}

如果教导说有眷顾存在，且世界由它造成，那么另一些问题就会迎向我们，必须加以讨论。因为人们会问：眷顾以何种方式运作？是普遍地对待整体，还是特殊地对待部分？或是普遍地也对待部分？或是既普遍地对待整体，又特殊地对待部分？我们所说的普遍眷顾，意思是：仿佛\Person{天主}在起初创造世界时，就给事物指定了一种秩序和进程，然后就不再关心所发生的事了。但特殊地对待部分的眷顾是这样：祂眷顾某些人或地方，而非其他。但普遍地对待一切，同时又特殊地对待部分，则是这样：如果\Person{天主}起初造了万物，并且对待每一个个体直到终末，按各人的行为施行赏报。\par

% structure: p0024 source=h2 target=subsection reason=relative-heading-rank; base=h2

\subsection{第十二章 祈祷与命运相悖}

因此，第一种说法宣称\Person{天主}在起初造了万物，并给事物施加了进程和秩序，此后就不再过问，这种说法认为一切都是按\OriginalTerm{命运}{genesis}发生的。对此，我们将首先予以回应；尤其是针对那些敬拜神灵并维护\OriginalTerm{命运}{genesis}的人。无疑，这些人向神灵献祭并祈祷时，希望获得某种与命运相反的东西，从而取消了命运。但当他们嘲笑那些激励德行、劝勉节制的人，并说没有人能做或遭受任何不是命运所注定的事时，他们确实根除了对神性的所有敬拜。因为如果你不能从那些神灵那里获得任何命运所定之方法不允许的东西，你又何必敬拜他们？暂且以此反驳这些人。但我说，世界由\Person{天主}所造，并且终将被祂毁灭，为要使那永恒的世界显现出来；这永恒的世界正是为此目的而造的：为了永远存在，并接纳那些在\Person{天主}审判中配得上它的人。但还有另一个不可见的世界，它包含这个可见的世界于自身之内——在我们完成关于可见世界的讨论之后，我们也将论及它。\par

% structure: p0026 source=h2 target=subsection reason=relative-heading-rank; base=h2

\subsection{第十三章 造物主的存在是必然的}

然而，在此期间，许多哲学家中的智者都见证了这个有形可见的世界是被造而成的。但为了避免我们似乎将断言用作证人，仿佛我们需要它们似的，请允许我们探讨它的原理。这个有形可见的世界是物质的，从它可见这一事实便足够明显。然而，每一个物体都接受两种\OriginalTerm{差异}{Differentiæ}中的\Emph{一种}：因为它是紧密固体的，或是分成而离散的。如果构成世界的物体是紧密固体的，且该物体根据其差异被分割和划分为不同的种类和部分，那么必须理解有一位存在者将那紧密固体的物体分开，并把它引向许多部分和不同形式；或者，如果这个世界的全部质量是由众多分散的物体部分复合和凝聚而成，那么仍然必须理解有一位存在者将分散的部分收集为一，并赋予这些事物以不同的种类。\par

% structure: p0028 source=h2 target=subsection reason=relative-heading-rank; base=h2

\subsection{第十四章 受造的方式。}

而且，我确实知道，有几位哲学家更持此见解：造物主天主从一体中作出了区分和分别，他们将此一体称为物质，它由四种元素构成，藉天主的眷顾之某种调和而混合为一。因为我认为某些人所说的——世界的身体是单纯的，即没有任何结合——是虚妄的；因为很明显，单纯的东西既不能是身体，也不能被混合、繁衍或分解；而所有这些，我们看见都发生在世界的身体上。因为若它是单纯的，且内中没有可被分解和分开的东西，它又怎么可能被分解呢？但如果身体似乎是由两种、三种甚至四种元素构成——凡是稍有知觉的人，谁不察觉必定有一位将众元素收集为一，且保持调和的尺度，从不同部分做成一个固体身体的\Emph{那位}？因此，我们称\Person{这位}为天主，世界的造物主，并承认祂是宇宙的创造者。\par

% structure: p0030 source=h2 target=subsection reason=relative-heading-rank; base=h2

\subsection{第十五章 受造的理论。}

因为希腊哲学家们在探求世界起源时，已有不同途径。简言之，毕达哥拉斯说数目是其源始元素；卡利斯特拉托斯说是性质；阿尔克迈翁说相反；阿那克西曼德说无限；阿那克萨戈拉说部分均等；伊壁鸠鲁说原子；狄奥多罗斯说\OriginalTerm{无部分者}{\GreekText{ἀμερῆ}}，即无部分之物；阿斯克勒庇俄斯说\OriginalTerm{肿块}{\GreekText{ὄγκοι}}，我们可称之为肿瘤或肿胀；几何学家说端点；德谟克利特说理念；泰勒斯说水；赫拉克利特说火；第欧根尼说气；巴门尼德说土；芝诺、恩培多克勒、柏拉图说火、水、气、土。亚里士多德还引入了第五元素，他称之为\OriginalTerm{不可名状者}{\GreekText{ἀκατονόμαστον}}，即不可命名者；无疑是指那将四种元素联合为一、创造世界的那一位。因此，无论构成世界的元素是两种、三种、四种、更多还是无数，在每种假设中都显明有一位天主，祂将众多收集为一，并再次将已收集的引向不同种类；由此证明，世界的机器若没有一位制造者和安排者，便不能存在。\par

% structure: p0032 source=h2 target=subsection reason=relative-heading-rank; base=h2

\subsection{第十六章 世界由创造者从无中造成。}

但由另一事实也可看出它们从无开始：在元素的结合中，若一个缺乏或过剩，其他便会松弛而衰落。例如，若任何物体中缺乏水分，干燥便不能维持；因为干燥由水分滋养，正如寒冷由热滋养一样；在其中，如我们所说，若一个有缺陷，整体便会分解。在此，它们给出其起源的迹象：即它们是从无中造成的。既然物质本身被证明是被造的，那么构成世界的各部分和种类，怎能被认为非受造呢？但关于物质及其性质，现在不是谈论的时候：只需教导这一点就足够了：天主是万物的创造者，因为，若构成世界的物体是固体且联合的，没有创造者它就不能被分离和区分；或者，若它是由分散且分离的部分收集为一，没有制造者它就不能被收集和混合。因此，若天主如此清楚地被显明为世界的创造者，伊壁鸠鲁还有什么余地引入原子，并断言不仅可感觉的身体，甚至理智和理性的心灵，都是由无感觉的微粒构成的呢？\par

% structure: p0034 source=h2 target=subsection reason=relative-heading-rank; base=h2

\subsection{第十七章 原子学说难以成立。}

但你会说，按伊壁鸠鲁的意见，一连串的原子无休止地到来，彼此混合，经过无限且无尽的时间而聚集，形成了固体物体。我并不将此意见视为纯粹的虚构，而且还是一个糟糕的虚构；但让我们审视它，无论它是什么性质，看看所说的能否成立。因为他们说那些微粒，他们称为原子的，具有不同的性质：有些是潮湿的，因而沉重，向下倾向；另一些是干燥和土性的，因而依然沉重；但另一些是火性的，因而总是向上推；另一些是冷而惰性的，总是停留在中间。既然有些如火性总是向上倾向，另一些如潮湿和干燥总是向下，还有一些保持不等的中间路线，它们如何能相遇并形成一个身体呢？如果有人从高处扔下小块稻草和同样大小的铅块，轻的稻草能跟上铅块吗？虽然同样大小，但重的到达底部要快得多。同样，原子尽管大小相等，但重量不等，较轻的永远赶不上较重的；但如果它们不能同步，当然也不能混合或形成一个身体。\par

% structure: p0036 source=h2 target=subsection reason=relative-heading-rank; base=h2

\subsection{第十八章 原子的汇聚不能造成世界。}

然后，其次，若它们无休止地浮动，不断到来，并被添加到已经满盈的事物上，当新重量不断堆积在如此巨大的重量上时，宇宙如何能站立得住？我也问这个问题：我们所见的天穹，若是由原子逐渐聚合而成，那么在建造过程中，若没有支柱捆绑那敞开的结构顶部，它如何能不坍塌？因为正如建筑圆形穹顶的人，若不捆绑中央顶部的固定物，整个建筑便立即倒塌；同样，我们所见的以如此优美形式聚合起来的世界圆形，若不是由造物主的能力透过一次性的神圣能力表达而造成，而是由原子逐渐聚合、建造，不是按理性的要求，而是按偶然的结果发生，那么它如何在被聚合和固定之前没有倒塌和粉碎呢？此外，我问：如此巨大质量的基础铺设在哪里？还有，你所谓的基础，它又立于何物？而那个基础，又是什么支撑它？我这样不断追问，直到答案归于虚无和空无！\par

% structure: p0038 source=h2 target=subsection reason=relative-heading-rank; base=h2

\subsection{第十九章 原子论的更多困难}

但若有人说，火性的原子结合形成了物体，且因为火的性质不是向下，而是向上，所以火的本性总是向上推，支撑着置于其上的世界质量；对此我们回答：火性的原子总是趋向最高处，如何能下降到较低处，并居于一切最低处，以形成万物的基础？而更重的性质，即土或水的性质，总是先于较轻的，如我们已说过的；因此，他们也断言，天作为更高的结构，是由较轻且总是飞升向上的火性原子组成的。因此世界不能以火或任何其它东西为基础；较重的原子与较轻的原子之间也不可能有任何结合或凝聚，即那些总是向下沉的与那些总是向上飞的不可能结合。由此已充分证明，世界的物体是由原子结合而巩固的；并且无感觉的物体，即使能以某种方式相遇并结合，也不能赋予物体形式和尺寸，形成肢体，或产生性质，或表达数量；所有这些，凭借其精确性，都证明了造物主的手，并显示了理性的运作，这理性我称之为圣言和天主。\par

% structure: p0040 source=h2 target=subsection reason=relative-heading-rank; base=h2

\subsection{第二十章 柏拉图的见证}

但有人会说，这些事是由自然所成。现在，这里争论的是名称。因为虽然明显那是理智和理性的作品，但你所谓自然，我称之为天主造物主。很明显，无论是按如此必要的区分安排的各类物体，还是心智的各种能力，都不能或不可能由无理性的、无感觉的工作造成。但若你视哲学家为合适的证人，柏拉图在\WorkTitle{蒂迈欧篇}中为此作证，他在讨论世界的创造时，问道：世界是否一直存在，还是有开始，并断定它是被造的。他说：\Quote{因为它是可见、可触、有形的；但显然，所有这类事物都是被造的；而被造之物无疑有作者，由他造成。然而，这位万有的造物主和父，很难发现；即使发现了，也不可能向群众宣示他。}这就是柏拉图的声明；但即使他和其它希腊哲学家选择对世界的创造保持沉默，难道对有理智的人不都显而易见吗？因为哪个人，有一丝理性，当他看见一所房子具备一切有用之需，其屋顶做成球体形状，饰以各种绚丽色彩和不同图案，并用巨大而光辉的灯具装饰；我说，谁看见这样的建筑，能不立刻断言，它是由最智慧、最有力的工匠建造的？同样，谁能如此愚蠢，当他仰望天穹，看见太阳和月亮的光辉，看见星辰的运行和美丽，以及它们被固定法则和周期所指定的路径，竟不大声呼喊这些作品与其说是由智慧而有理性的工匠造成，不如说是由智慧和理性本身造成？\par

% structure: p0042 source=h2 target=subsection reason=relative-heading-rank; base=h2

\subsection{第二十一章 机械论}

但若你更愿听取其它希腊哲学家的意见——而且你熟悉机械科学——你当然知道他们关于天的论述。因为他们设想一个球体，向每一边都同等圆润，平等地朝向所有点，与地心在所有方向上距离相等，并因其自身的对称而如此稳定，以致其完美的均衡不允许它向任何一边倾斜；因此球体虽无支柱支撑，却能维持。若世界结构真有此形态，其中明显有神圣的作品。但若如其他人所想，球体置于水上，由水支撑，或浮于水中，即便这样，其中也显示出伟大设计者的作品。\par

% structure: p0044 source=h2 target=subsection reason=relative-heading-rank; base=h2

\subsection{第二十二章 星辰的运动}

但为避免对并非众人皆知之事所作的断言可能显得可疑，让我们转向那些无人不知的事。是谁以如此伟大的理性安排了星辰的运行，规定了它们的升起与落下，并指定每一颗星在确定而有规律的时间完成天体的运行？是谁赋予一些星总是靠近落处，另一些星则回归升起？是谁为太阳的行程设定了度量，使其通过不同的运动来划分时辰、日子、月份和季节的变化？——使其通过行程的精确度量来区分冬、春、夏、秋，并始终借着年岁相同的变化，有条不紊、循环往复？我要说，谁会不宣称如此秩序的主宰者正是天主的智慧呢？我们所说这些，乃是根据希腊人向我们提供的有关天体科学的知识。\par

% structure: p0046 source=h2 target=subsection reason=relative-heading-rank; base=h2

\subsection{第二十三章 眷顾尘世之事}

至于我们在陆地上或海洋中看到的事物，又当如何？它们难道不明确教导我们：其中不仅有天主的化工，也有天主的眷顾吗？因为在某些地方有高耸的山脉，\Emph{其目的在于}通过天主的安排，空气被压缩和限制在里面，被迫挤压成风，使果实得以萌芽，并在昴宿星团因太阳的烈焰而炽热时，调节夏天的暑热。但你仍然会说：为何要有太阳的烈焰，以致需要调节呢？那么，为人类使用所必需的果实又如何成熟呢？但也要注意这一点：在正午的轴线上，即炎热最甚的地方，并没有大量的云层聚集，也没有丰沛的降雨，以免给居民带来疾病；因为水汽云层若受骤热作用，会使空气变得不洁、有害。而大地承接暖雨，也不会给庄稼带来滋养，而是毁灭。谁能怀疑这其中有着天主眷顾的运作？简言之，埃及靠近埃塞俄比亚的炎热地区，为避免因降雨而使其空气无法救治地恶化，其平原不接纳来自云层的雨水，而是接受来自尼罗河泛滥的、如同地上的降雨。\par

% structure: p0048 source=h2 target=subsection reason=relative-heading-rank; base=h2

\subsection{第二十四章 河流与海洋}

对于泉源和河流——它们以永不停息的流动汇入大海——我们又能说什么？由于天主的眷顾，它们丰沛的供应并不枯竭，而海洋虽然收纳如此大量的水，却不见增长；无论是注入的元素还是被注入的元素，都保持同样的比例。但你会对我说：咸水自然会消耗注入其中的淡水。然而，这正显明天主眷顾的化工：天主使那汇聚它为自己所预备、供人类使用的所有水流的元素成为咸的。如此，经过如此漫长的岁月，海洋的河床没有被填满，也没有产生毁灭大地和人类的洪水。没有人会愚蠢到认为如此伟大的理性与如此伟大的眷顾是由无理性的自然所安排的。\par

% structure: p0050 source=h2 target=subsection reason=relative-heading-rank; base=h2

\subsection{第二十五章 植物与动物}

关于植物，关于动物，我又能说什么？难道不是眷顾规定了植物在因衰老而腐朽时，通过自身产生的侧芽或种子得以再生，而动物则通过繁殖延续？并且借由眷顾的一种奇妙安排，在动物出生之前，乳汁已在母畜的乳房中预备好；它们一出生，无需任何引导，便会寻找为它们预备的滋养储备。不仅产生雄性，也产生雌性，以便通过两者使种族得以延续。但为了避免这一点，如有些人所认为的，是由某种自然的秩序而非造物主的安排所成就，天主便作为一种眷顾的证明和标记，以例外的方式规定了少数动物在地上保存其种类：例如，乌鸦通过口受孕，鼬鼠通过耳生产；一些鸟类，如母鸡，有时产下由风或尘埃感应的蛋；其他动物则将雄性转变为雌性，并每年改变性别，如野兔和鬣狗——它们被称为怪物；有些从泥土中生出，从土中获得身体，如鼹鼠；另一些从灰烬中生出，如毒蛇；另一些从腐烂的肉中生出，如马肉中的黄蜂、牛肉中的蜜蜂；另一些从牛粪中生出，如甲虫；另一些从植物中生出，如罗勒上的蝎子；反之，植物也从动物生出，如从雄鹿或母山羊的角上生出欧芹和芦笋。\par

% structure: p0052 source=h2 target=subsection reason=relative-heading-rank; base=h2

\subsection{第二十六章 种子的萌芽}

还有什么事例，天主的眷顾以种种方式规定了受造物的产生，而那被视为自然所赋予的秩序被超越，由此表明并非无理性的事物过程，而是由祂自己的理性所安排的呢？在这当中，岂不也显示出眷顾的完满工程吗？当种子被播种，通过土地和水为人的滋养而预备时，这些种子被交付给大地，土壤便如同从它的乳房中，将因天主的旨意而吸收的水分泌入种子。因为水中含有天主从起初就赋予的某种神能，借着这神能的运作，那将要形成的身体结构在种子自身内开始形成，并通过芽和穗发育；因为种子因水分膨胀，那被安放在水中的神能，作为一种非物质的实体，通过某些狭窄的脉络通道运行，激发种子生长，并形成生长植物的种类。因此，借着那含有并内蕴生命之灵的湿润元素，不仅使种子复活，而且还再现了与所播种子完全相像的外观和形态。如今，谁若还稍有理智能认为这过程取决于无理性的自然，而非天主的智慧呢？最后，这些事也与人类的出生相似；因为大地似乎取代了子宫的位置，种子投入其中，借着水和神能的力量被形成和滋养，正如我们上面所说的。\par

% structure: p0054 source=h2 target=subsection reason=relative-heading-rank; base=h2

\subsection{第二十七章　水的力量}

但在此，天主的眷顾也令人钦佩：它允许我们看见并知道所造之物，却将造物的方式和途径置于隐秘与隐蔽之中，使不配知晓的人无法得知，而只向配得忠信的人启示，当他们堪当领受时。然而，借着事实和实例证明，没有任何东西是从土地的实体传递给种子的，一切都取决于水的元素及其内蕴的神能——假设，例如，将一百塔连特的土放入一个非常大的槽中，在其中播撒几种不同的种子，如草药或灌木的种子，并供给足够的水来浇灌，如此管理数年，然后将收获的种子储存起来，例如谷粒或大麦等，逐年分别保存，直至每种种子达到一百塔连特的重量，然后也将茎秆连根拔起并称重；在所有这些都从槽中取出后，再将土称量，它仍会完整无缺地归还一百塔连特的重量。那么，我们该说所有这些重量、所有不同种子和茎秆的数量是从哪里来的呢？这不是显然来自水吗？因为大地保留了自己的全部，而所倒进去的水却无处可寻，这是由于天主规定的强大德能，它借着同一种水，既预备了如此多种种子和灌木的质料，又形成了它们的种类，并在倍增收成时保存了种类。\par

% structure: p0056 source=h2 target=subsection reason=relative-heading-rank; base=h2

\subsection{第二十八章　人的身体}

从这一切，我认为已足够充分明白地证明：万物都是藉着设计的理智而产生，宇宙也是由它构成，而非藉着自然的无理性运作。但请允许我们现在来讨论我们自身的实体，即人的实体，他是大世界中的小世界，即\OriginalTerm{小宇宙}{microcosm}；让我们思考它是如何以何等理智被组成的。尤其由此你将理解造物主的智慧。因为尽管人由不同的实体构成，一为有死的，另一为不朽的，然而由于造物主的巧妙设计，它们的差异并不妨碍其结合，尽管实体彼此不同且相异。因为一个取自土地并由造物主形成，另一个则从不朽的实体赐予；然而，它的不朽尊荣并未因这结合而受损。它并非如某些人所想，由理性、贪欲和情欲组成，而是这些情感似乎存在于其中，使它能在这些方向的每一方向被推动。因为身体由骨骼和肉组成，其开端来自男人的种子，这种子因温暖从骨髓中提取，并被送入子宫作为土壤，附于其上，逐渐从血液的泉源得到滋润，从而转化为肉和骨，并形成与注入种子者相似的模样。\par

% structure: p0058 source=h2 target=subsection reason=relative-heading-rank; base=h2

\subsection{第二十九章　身体的匀称}

且注意在此设计者的工程：祂如何将骨骼像柱子一样插入，使肉得以在其上支撑和承载。然后，两侧——即左右——又如何保持相等的尺度，使脚与脚相应，手与手相应，甚至指与指相应，每部分都与另一部分完全相等；眼与眼，耳与耳，不仅彼此适合匹配，而且还为必要的用途而造。例如，手被造得适于工作；脚适于行走；眼睛，有眉毛如哨兵般保护，以服务于视觉；耳朵被造得适于听觉，像铙钹一样振动落在其上的声音，并将其向内传送，甚至传递到内心的理解；而舌头在说话时叩击牙齿，如同琴弓。牙齿也被形成，有些用于切割和分解食物，并将其递给内齿；而内齿则像磨一样将其碾磨，使其更易于消化并被送入胃中；因此它们也被称为磨齿。\par

% structure: p0060 source=h2 target=subsection reason=relative-heading-rank; base=h2

\subsection{第三十章　呼吸与血液}

鼻孔也被造来收集、吸入和呼出空气，以便通过呼吸的更新，使心中固有的天然热力凭借肺部得以按需加温或冷却；而肺被造来存于胸腔，借其柔软性安抚和滋养心脏的活力——生命似乎寄居于此；我说的是生命，不是灵魂。至于血液的本质，它如同河流从源头涌出，先沿一根管道流动，然后分散到无数静脉中，如同渠道，以生命之流灌溉整个人体的疆域，此过程由位于右侧的肝脏提供支持，以消化食物并将其转化为血液；而在左侧则安放了脾，它吸引并在某种程度上净化血液中的杂质。\par

% structure: p0062 source=h2 target=subsection reason=relative-heading-rank; base=h2

\subsection{第三十一章 肠腑}

此外，肠脏的构造中又运用了何等理性呢？它们被排列成长长的环形盘绕，以便逐步排出食物残渣，既不致使腹腔突然空虚，也不致被随后摄入的食物所阻碍。而且它们被造成膜状，使得外部部分能够逐渐吸收水分；若水分突然涌出，会使内部空虚；而若被厚皮阻碍，则会使外部干燥，并以令人痛苦的口渴扰乱整个人体结构。\par

% structure: p0064 source=h2 target=subsection reason=relative-heading-rank; base=h2

\subsection{第三十二章 生育}

再者，女性的体型和子宫的腔室最适合接受、滋养并赋予胚种生命，谁不相信这是凭理性和预见造成的呢？因为女性身体只有在此部位才与男性不同，胎儿在此被安置、保存和滋养。同样，男性与女性的区别也仅在于身体中具有注入精子并繁衍人类能力的那一部分。由此，从肢体必要的差异中可见天主眷顾的伟大证明；尤其在于，当形态相似时，却发现有用途的差异和功能的不同。因为男女同样有乳房，但只有女性的乳房充满乳汁，以便她们一旦生产，婴儿就能找到适合他的营养。如果我们看到人体器官的安排如此有序：所有其余部分形态相似，只有那些用途需要差异的部分才有所不同；我们在男人身上既看不到多余，也看不到欠缺，在女人身上也未见不足或过剩——那么，有谁不会从这一切中认出理性的运作和造物主的智慧呢？\par

% structure: p0066 source=h2 target=subsection reason=relative-heading-rank; base=h2

\subsection{第三十三章 受造物中的对应}

与此一致的还有其它动物合理的差异，每一种都适应其自身的用途和服务。树木的多样性、草药的差异（无论是在形态还是汁液方面）也同样证实了这一点。季节的变迁也证明了这一点：分为四个时期，年轮以确定的时辰、日子、月份闭合，且不偏离既定的计算一个时辰。简言之，世界本身的年代也由确定不变的计算和明确年数来计算。\par

% structure: p0068 source=h2 target=subsection reason=relative-heading-rank; base=h2

\subsection{第三十四章 创造世界的时间}

\Quote{但你会说：世界是何时创造的？为什么这么晚？即使它被造得更早，你也会提出这个异议。因为你可能会说：为什么不更早呢？这样，通过回溯无限的年代，你仍会问：为什么不更早呢？然而我们现在并非讨论为什么它没有被造得更早，而是它究竟是否被造。因为如果它显然是被造的，那么它必然是一位强大而至高的工匠的作品；如果这一点是清楚的，那么就必须由智慧的工匠按其喜悦选择并决定创造的时间。除非你认为，那构建了世界宏伟结构、赋予了各个对象以形式和种类、不仅赋予它们美丽的外观，而且为它们未来的用途赋予最适宜且必不可少之习性的全部智慧，唯独忘记了为如此宏伟的创造工程选择一个合适的时间。祂无疑有其确定的理由和明显的原因，说明祂在何时、为何以及如何创造了世界；但这些理由不适合向那些不情愿审视和理解摆在眼前、见证祂眷顾的事物的人揭示。因为那些隐藏于智慧感官之内、如同藏在王室宝库中的秘密，只向那些从祂那里学习的人开启——这些秘密在祂那里被封存和收藏。因此，是天主创造了万物，祂本身并非受造者。而那些用自然代替天主、宣称万物由自然造成的人，并不觉察他们所用名称的错误。因为如果他们认为自然是不具理性的，那么设想理性的受造物可以从不具理性的创造者产生，便是极其愚蠢的。但如果它是理性——即\OriginalTerm{圣言}{Logos}——万物显然由它造成，那么他们在谈论创造者的理性时，便是无故更改名称了。我的父亲，若你对这些事有什么要说的，就请说吧。}\par

% structure: p0070 source=h2 target=subsection reason=relative-heading-rank; base=h2

\subsection{第三十五章 好客之争}

尼塞达这样说时，老人回答说：\Quote{我儿，你的论证确实明智而有力；以至我认为关于天主眷顾的题目已无法处理得更好了。但天色已晚，我希望明天就你所论证的答复一些话；如果你能在这上面令我满意，我将承认我亏欠你的恩惠。}老人这样说时，伯多禄站起来。那时，在场的一位劳狄刻雅人的首领，请求伯多禄和我们，让他给老人换上别的衣服，而不是他所穿的破烂衣服。伯多禄和我们拥抱了他，称赞他可敬而美好的意向，说：\InnerQuote{我们并非如此愚昧不敬，竟不将身体所需之物赐给那托付了如此宝贵话语的人；我们盼望他乐意接受，如同父亲从儿子们手中接受，并且我们相信他会与我们共享我们的房屋和生活。}我们这样说时，那位城中的首领极力以最大的热情和许多甜言蜜语试图将老人从我们这里带走，而我们更加热切地设法留住他，全体民众喊道，应该照老人自己所喜悦的去做；得到静默后，老人起誓说：\InnerQuote{今天我不与任何人停留，也不从任何人取什么东西，免得一人的选择成为另一人的忧伤；以后这些事可以办，如果这样看来合适。}\par

% structure: p0072 source=h2 target=subsection reason=relative-heading-rank; base=h2

\subsection{第三十六章  安排次日事宜}

老人这样说时，伯多禄对城中的首领说：\Quote{既然你在我等面前显出了善意，你不应忧愁离去；但我们要以恩惠回报恩惠。将你的房屋指给我们，并预备好，使明日的讨论可在那里举行，并容许任何愿意在场聆听的人进入。}城中的首领听了这话，非常欢喜；全体民众也欣然听见。人群散去后，他指出了自己的房屋；老人也准备离开。但我吩咐一个仆从暗中跟随老人，查明他在何处停留。我们回到寓所后，将我们与老人周旋的一切告诉了弟兄们；然后照常用过晚餐，便就寝了。\par

% structure: p0074 source=h2 target=subsection reason=relative-heading-rank; base=h2

\subsection{第三十七章  \Quote{你从我听到的那纯正话语的规模}}

次日，伯多禄早早起来呼唤我们，我们一同前往前一天我们去过的隐密处祈祷。祈祷后，我们从那里前往指定的地方时，他沿路劝诫我们说：\Quote{最亲爱的同仆们，请听我说：你们每人按自己的能力，为那些前来归向我们宗教之信仰的人提供益处，这是好的；因此，不要畏缩于教导无知者，并按照天主眷顾所赐予你们的智慧施教，只是要使你们言论的雄辩与你们从我听到、受托付的那些事相结合。但不要说出任何自己的、未经托付的事，尽管你们自己可能认为真实；而要持守那些正如我所说的、我自己从真先知领受并交付给你们的事，尽管它们可能显得不那么权威。因为常常发生这样的事：人因相信自己凭自己的思想发现了更真实、更有力的真理形式，而偏离了真理。}\par

% structure: p0076 source=h2 target=subsection reason=relative-heading-rank; base=h2

\subsection{第三十八章  首领的房屋}

我们欣然同意伯多禄的这些劝诫，向他说，我们将只做使他喜悦的事。于是他说道：\Quote{为使你们能安全地操练，你们各人当在我面前依次主持讨论，每人阐明自己的问题。那么，既然尼塞达昨天已充分论述，就让阿基拉今天主持讨论；阿基拉之后，克肋孟；然后我，若情况需要，再加添一些。}我们这样交谈时，到了那房屋；屋主迎接我们，领我们到一个布置成剧场式样、建造美丽的房间。我们在那里看到大群人在等候我们，他们是夜间来的，其中也有昨天与我们辩论的老人。因此我们进入，伯多禄在我们中间，环顾四周看能否看见老人；伯多禄看见他隐藏在人群中间，便叫他到跟前，说：\Quote{既然你拥有比大多数人更光明的灵魂，为何隐藏自己，因谦逊而躲避？不如到这里来，提出你的见解。}\par

% structure: p0078 source=h2 target=subsection reason=relative-heading-rank; base=h2

\subsection{第三十九章  重述昨日的论证}

当\Person{伯多禄}这样说的时候，人群立刻开始为老人让路。老人走上前来，这样开始：虽然我不记得那年轻人昨天发表的谈话中的措辞，但我记得它的主旨和次序；因此，我认为有必要，为了那些昨天未在场的人，唤起已说过的话，并简短地重复一切，以便即使我遗漏了什么，也可由那位发表谈话且此刻在场的人提醒我。这就是昨天讨论的主旨：我们所看见的一切，既然它们由某种比例、技艺、形式和种类构成，就必须相信是由明智的能力所造；但如果是理智和理性形成了它们，那么世界也是由同一理性的眷顾所治理，尽管世界中所行的事在我们看来可能并不完全正确。但由此推论，如果天主和理智是万物的创造者，祂也必须是正义的；但若祂是正义的，祂必然审判。若祂审判，人就必须按其行为受审判；若每人按其行为受审判，终有一天义人与罪人之间会有公义的分离。我想，这就是整个谈话的实质。\par

% structure: p0080 source=h2 target=subsection reason=relative-heading-rank; base=h2

\subsection{第四十章 生成}

\Quote{因此，若能证明理智和理性创造了万物，那么后来的事物也必由理性和眷顾管理。但若无理智和盲目的本性产生万物，审判的理性无疑被推翻；若没有审判者，就没有期望罪的惩罚或善行的赏报。既然整件事依赖于此，系于这一要点，不要见怪，我希望这一点能更充分地讨论和处理。因为在这点上，仿佛第一道门对所有提出的事物关闭了，所以我希望首先为我开启它。现在请听我的教义；若你们中有人愿意，就回答我：因为若我听到真理，我将不耻于学习，并赞同发言正确者。那么，你昨天发表的谈话，主张万物凭借技艺、尺度和理性而存在，并未完全说服我世界的创造者是理智和理性；因为我能指出许多事物，它们凭借合适的尺度、形式和种类而存在，却显然不是由理智和理性所造。此外，我还看到世界上许多事情没有秩序、连续或正义，并且没有生成的进程，什么也不能发生。我将在后续以我自己的情况最清楚地证明这一点。}\par

% structure: p0082 source=h2 target=subsection reason=relative-heading-rank; base=h2

\subsection{第四十一章 彩虹}

当老人这样说了之后，\Person{阿奎拉}回答说：\Quote{正如您自己提出的，任何人若愿意，可以回答您说的话，我的兄弟\Person{尼塞塔}允许我今天主持辩论。}然后老人说：\Quote{继续，我儿，随你便。}\Person{阿奎拉}回答说：\Quote{您承诺过要指出世界上有许多事物具有由同等理性安排的形式和种类，却显然不是由天主作为造物主造成的。那么，现在按您的承诺，指出这些事情吧。}\par

% structure: p0084 source=h2 target=subsection reason=relative-heading-rank; base=h2

\subsection{第四十二章 类型与形式}

\Person{阿奎拉}说：如果某物是从类型和形式表达出来的，就立刻可知它源于理性，且缺乏理智无法制成；因为表达图形和形式的类型本身，也不是缺乏理智制成的。例如，如果将蜡按在刻纹的戒指上，它就会从戒指上取得印记和图形，而戒指无疑是无感觉的；但表达图形的戒指是工匠的手所雕刻的，是理智和理性赋予了戒指类型。同样，彩虹也在空中被表达出来；因为太阳在云彩稀薄化的过程中将光线印在云上，将其圆环的类型附在潮湿的云上，就像附在软蜡上，产生了彩虹的外观；正如我所说，这是通过太阳的光芒在云上的反射，并从中再现其圆环的光芒而实现的。但这并不总是发生，只有当湿润云彩稀薄化提供了机会时才发生。因此，当云彩再次凝聚结合时，彩虹的形式就会消解消失。最后，没有太阳和云彩就看不到彩虹，正如除非有类型和蜡或其他材料，否则不会产生影像。如果天主创造者在起初就创造了类型，如今形式和种类可从其表达出来，但这并不奇怪。这类似于此：在起初天主创造了无感觉的元素，祂可以用它们来形成和发展所有其他事物。即使那些塑造雕像的人，也先用黏土或蜡制作模具，然后雕像的图形从其中产生。之后，雕像还会产生一个影子，这个影子总是带着雕像的形式和相似。那么我们要说什么？无感觉的雕像形成了一个与雕像本身同样精心完成的影子？还是应当毫不犹豫地将影子的完成归功于塑造雕像的那位？\par

% structure: p0086 source=h2 target=subsection reason=relative-heading-rank; base=h2

\subsection{第四十三章 天主所造看似无用和卑贱之物}

\Quote{若这样，这似乎是正确的，关于此事所说的已足够，让我们来探讨其他问题；或者如果你认为还有所缺，我们再重述一遍。}那老人说：\Quote{我希望你再重述一遍，因为还有许多其他事物，我看到是以同样方式被造的：树木的果实也是这样被造，形态优美，圆润奇妙；叶子的外观以极大的优雅形成，绿色的薄膜以精巧的技艺编织；此外，跳蚤、老鼠、蜥蜴之类，我们能说这些也是天主造的吗？因此，从这些卑贱之物，人们推导出关于更高事物的猜想，认为它们绝不是由理性之技巧形成的。}\Quote{你推断得好，}\Person{阿基拉}说，\Quote{关于叶子的纹理，以及微小动物，从这些，对更高受造物的信仰被撤去；但不要让这些事欺骗你，以为天主仿佛只凭两手工作，不能完成一切被造之物；但要记住，我的兄弟\Person{尼塞塔}昨日如何回答你，并如同儿子对父亲说话一样，真诚地提前揭示了奥迹，解释了那些看似无用的事物为何及如何被造。}\par

% structure: p0088 source=h2 target=subsection reason=relative-heading-rank; base=h2

\subsection{第四十四章 有序与无序}

那老人说：\Quote{我希望从你那里听到，为什么那些无用之物是由那至高理性的意愿所造？}\Quote{如果，}\Person{阿基拉}说，\Quote{你完全明了它们之中有理性和理智的工作，那么你也会毫不迟疑地说出它们为何被造，并宣告它们被造得正确。}对此，那老人回答说：\Quote{我的孩子，我不能说那些看似由技艺形成之物是由理性所造，因为世界上我们看到其他事物被做得不义和无序。}\Quote{如果，}\Person{阿基拉}说，\Quote{那些无序所做之事不允许你说它们是由天主眷顾所做，那么为什么那些有序所做之事不迫使你说它们是由天主所做，而非理性之性不能产生理性之作？因为确实，我们毫不否认，在这世界上，有些事做得有序，有些无序。因此，那些理性地所做之事，相信它们是由眷顾所做；但那些非理性地、无序地所做之事，相信它们是自然发生、偶然发生的。但我奇怪人们没有察觉到，在有感觉之处，事物可以做得有序或无序，但在没有感觉之处，两者都不能做到；因为理性创造秩序，而秩序的过程必然产生某种无序，如果有相反之事发生扰乱了秩序。}然后那老人：\Quote{这正是我希望你向我展示的。}\par

% structure: p0090 source=h2 target=subsection reason=relative-heading-rank; base=h2

\subsection{第四十五章 太阳和月亮的运行}

\Person{阿基拉}说：我将毫不迟疑地这样做。天上显示两个可见的征象——一个是太阳的，另一个是月亮的；还有其他五颗星紧随其后，每颗星描绘自己单独的轨道。因此，天主将它们安置在天上，藉此空气的温度可按季节调节，交替和变化的秩序得以保持。但藉着同样的\Emph{征象}，若有时因人的罪恶，瘟疫和败坏被降在地上，空气被扰乱，瘟疫临到动物，枯萎临到庄稼，一个毁灭性的年份在各方面临到人类；因此，藉着同一种手段，秩序既被保持也被摧毁。因为即使对不信者和无知者，也是明显的：太阳的行程对世界有用且必要，由天主眷顾所指定，总是保持有序；但月亮的行程，与太阳的行程相比，在其盈亏中，对无知者似乎是失序和不安定的。因为太阳以固定的、有规律的时期运行：从它而来的是时辰，从它而来的是它升起时的白日，从它而来的还有它落下时的黑夜；从它计算月份和年份，从它产生季节的变化；当它升到更高区域时，它调和春天；但到达天顶时，它点燃夏天的炎热；再次下降时，它产生秋天的温和；当它回到最低圈时，由于天空的冰封，它将严冬的寒冷留给我们。\par

% structure: p0092 source=h2 target=subsection reason=relative-heading-rank; base=h2

\subsection{第四十六章 太阳和月亮：善恶两者的仆役}

但我们在另一时间将更详细地讨论这些主题。现在，同时，\Emph{我们注意到}：虽然他是调节季节变化的那个好仆役，但当惩罚按照天主的意愿加在人身上时，他燃烧得更猛烈，以更炽热的火焰烧毁世界。同样，月亮的行程，以及那对无知者似乎失序的变化，适应于庄稼、牲畜和所有生物的成长；因为藉着她的盈亏，藉着天主眷顾某种奇妙的安排，一切出生的都被滋养和成长；关于这一点，我们可以更详细地谈论并展开细节，但所提问题的方向召唤我们回去。然而，藉着产出它们的同样工具，一切事物都被滋养和增长；但当因任何正当原因，既定秩序的调节被改变时，败坏和失调产生，以便惩罚通过天主的意愿临到人，正如我们上面所说的。\par

% structure: p0094 source=h2 target=subsection reason=relative-heading-rank; base=h2

\subsection{第四十七章 对义人和恶人的惩罚}

但也许你会说，在共同的惩罚中，虔敬和不虔敬的人遭遇相同的事，这又如何解释？确实如此，我们也承认；但惩罚对虔敬的人有益，使他们在今生受苦，从而更纯洁地进入来世，那里为他们准备了永久的安息；同时，即使不虔敬的人也可能从惩罚中获益，或者将来的审判可对他们作出公正的判决；因为在相同的惩罚中，义人感谢天主，而不义的人则亵渎。因此，既然对事物的看法分为两部分，即有些事是按秩序发生的，有些是违背秩序的，那么从那些按秩序发生的事，就应当相信有天主眷顾；至于那些违背秩序发生的事，我们应向那些通过先知教导得知其缘由的人请教原因：因为那些熟悉先知话语的人知道，在每一世代，枯萎、冰雹、瘟疫等何时发生，为何发生，以及因何罪这些作为惩罚降下；由此悲伤、哀叹和痛苦的原因临到人类；颤抖的疾病也随之而来，这从一开始就是对弑亲罪的惩罚。\par

% structure: p0096 source=h2 target=subsection reason=relative-heading-rank; base=h2

\subsection{第48章　为罪而施的惩罚}

\Quote{因为在世界之初，并没有这些灾祸，它们是从人的不虔敬兴起的；从此，随着不义的不断增加，灾祸的数量也增加了。但天主眷顾之所以对所有人宣判定案，是因为今生并非人人都能按其应得受报。那些从起初、在没有邪恶原因时妥善有序设立的事，不应以因人的罪而降临世界的灾祸来评判。总之，作为起初之事的标识，有些民族被发现与这些灾祸无缘。例如\OriginalTerm{赛里斯人}{Seres}，因为他们贞洁地生活，得以免受这一切；在他们那里，妇女怀孕后或正在洁净时，不得接近她。没有人吃不洁的肉，没有人知道祭祀；所有人都依据正义自我裁判。因此，他们不受我们所说的那些瘟疫惩罚；他们活到极老，无疾而终。但我们这些可怜人，如同与致命毒蛇同住\BibleRef{则 2:6}——我指与恶人同住——必然在今生与他们同受苦难的瘟疫，但我们从将来善事的安慰中抱有希望。}\par

% structure: p0098 source=h2 target=subsection reason=relative-heading-rank; base=h2

\subsection{第49章　天主的诫命被轻视}

\Quote{如果，}老人说，\Quote{即使义人因他人的不义而受苦，天主作为预知此事者，本应命令人不要做那些必要使义人与不义者一同受苦的事；或者，如果他们做了，祂本应对世界施以某种纠正或净化。} \Quote{天主，}阿奎拉说，确实这样命令了，并且通过先知赐下关于人应如何生活的诫命；但连这些诫命他们也轻视了：是的，若有人愿意遵守，他们就用各种伤害折磨他，直到把他从立志遵守中驱离，使他转从无信仰的乌合之众，变得和他们一样。\par

% structure: p0100 source=h2 target=subsection reason=relative-heading-rank; base=h2

\subsection{第50章　洪水}

因此，简言之，起初当全地都被罪恶玷污时，天主降洪水于世界，就是你们所说的\OriginalTerm{德乌卡利翁}{Deucalion}时期的洪水；那时祂在方舟中救了一位义人及其儿子们，并与他一同保存了所有植物和动物的种类。然而，即使从他们而出的人，过了一段时间后，又做了如同他们前人一样的事；因为他们忘记了所遭遇的事，以至于后代甚至不相信洪水曾发生过。为此天主也决定，在今世不再有第二场洪水，否则每一世代都该有一场，按他们因不信而犯的罪来说；但祂宁愿让某些以邪恶为乐的天使统治各个民族——他们被赐予对个人的权力，但只有一个条件，即若有人先因犯罪而自甘服从他们——直到那以善为乐者来临，由祂完成义人的数目，并因全世界虔敬人数的增加，使不虔敬在某种程度上受到抑制，并使所有人都知道，一切善事都是由天主做的。\par

% structure: p0102 source=h2 target=subsection reason=relative-heading-rank; base=h2

\subsection{第51章　罪恶带来的灾祸}

但因自由意志，每个人在不信将来之事时，通过恶行陷入灾祸。世上那些看似违背秩序发生的事，正是由不信而产生的。因此，天主眷顾的安排也值得赞叹：祂起初赐予走在良善生活之路的人享受不朽的美物；但当他们犯罪时，他们因罪生出恶。每一善事都与恶相连，仿佛因罪之盟约而结合，因为大地已被人的血污染，祭台已为魔鬼点燃，他们用祭祀的污秽烟雾污染了空气；于是元素最终先被败坏，又将败坏的过失传给人，如同树根\Emph{将其品质传予}树枝和果实。\par

% structure: p0104 source=h2 target=subsection reason=relative-heading-rank; base=h2

\subsection{第52章　\Quote{没有不带刺的玫瑰}}

因此，请留意在此事上，正如我所说过的，天主眷顾是如何公正地救助了那些变质的事物；既然源于罪恶的种种邪恶与天主的恩赐混杂在一起，祂便将两个首领分派给这两个领域。\Emph{相应地}，祂将善事的安排交给了那以善为乐者，为的是使那些信靠\Emph{祂}的人能归信祂的眷顾；而将那无序且无益的事务交给了那以恶为乐者，这些事当然会使人对祂的眷顾产生疑虑；这样，一位公正的天主便作出了公正的划分。故此，一方面，星辰的有序运行使人确信世界是由一位设计者之手所造；另一方面，大气的扰动、致病的微风、不受约束的闪电之火，却使人对天主眷顾的工程产生怀疑。因为，正如我们所说过的，每一件美善之事都有其相应的对立邪恶之事与之相连：正如冰雹与滋养的降雨相对，霉烂的腐败与温柔的露水相连，风暴的旋风与和风相伴，不结果实的树与结果实的树为伍，有害的草药与有益的草药并存，凶猛的野兽与温顺的动物共处。但这一切都出自天主的安排，因为人的意志选择已经偏离了善的目的，堕入了邪恶。\par

% structure: p0106 source=h2 target=subsection reason=relative-heading-rank; base=h2

\subsection{第五十三章 万物皆有对应的对立面}

\Quote{因此，这一划分存在于世界上所有事物之中；既然有虔诚的人，也就有不虔敬的人；既然有先知，也就有假先知；在外邦人中，有哲学家和假哲学家。此外，阿拉伯民族及其他许多民族，为了服务于他们的不虔敬，也模仿了犹太人的割损。同样，对邪魔的敬拜与对天主的敬拜相对，圣洗与圣洗相对，律法与法律相对，假宗徒与宗徒相对，假教师与教师相对。因此，在哲学家中，有人肯定天主眷顾，有人否定它；有人主张只有一位天主，有人主张多位。简言之，事情已经到了这种地步：邪魔被天主的言语所驱逐，这言语宣告了天主眷顾的存在，而巫术为了确证不信，却找到了以相反方式模仿此事的途径。于是，便发现了用咒语化解蛇毒的方法，以及行使与天主的言语和能力相反的治愈术。巫术也发现了与天主的众使者相对的侍奉，将招魂和邪魔的幻象置于对立面。不再赘述，凡是有助于相信天主眷顾之事，无不有另一方面为不信预备了对应之物。因此，那些不认识这一划分的人，由于世界上那些与自身不协的事物，便认为没有天主眷顾。但您，我的父亲，作为一位智者，请从这一划分中选取那维护秩序、有助于相信天主眷顾的部分，而不要只追随那违背秩序、抵消天主眷顾信仰的部分。}\par

% structure: p0108 source=h2 target=subsection reason=relative-heading-rank; base=h2

\subsection{第五十四章 一个比喻}

对此，老人回答说：\Quote{我的儿子，请指给我一条路，使我能在心中确立这两种秩序中的一种，一种肯定天主眷顾，另一种否定它。} \Quote{对于一个有正确判断力的人来说，}\Person{阿奎拉}说，\Quote{决定是容易的。因为您所说的秩序与无序，恰恰可以被一位设计者所产生，而不能被无感觉的自然所产生。让我们用一个比喻来说明：假设一大块石头从高岩上脱落，直坠而下，摔在地上碎成许多碎片，那么在那一堆碎片中，是否有可能找到一块具有完全形状和样貌的呢？}老人回答说：\Quote{不可能。} \Quote{但是，}\Person{阿奎拉}说，\Quote{如果有一位雕刻家在旁，他可以用他灵巧的手和理性的心智，把从山上切下的石头塑造成他喜欢的任何形状。}老人说：\Quote{确实如此。} \Quote{因此，}\Person{阿奎拉}说，\Quote{当没有理性心智时，石料无法成形；但当有设计的心智时，便可以有形式和畸形。例如，如果一个工匠从山上切下一块他想要赋予形状的石料，他必须先切割出未成形、粗糙的毛坯；然后，通过逐步锤击和雕琢，遵循他艺术的规则，表达出他心中构思的形状。这样，从无形状或畸形，经工匠之手，形状得以实现，两者都来自工匠。同样，世界上所发生的事，也是由一位设计者的天主眷顾所完成，尽管它们看起来可能不太有序。因此，既然这两条路已经向您显明，您也听到了它们的划分，就请逃离不信之路，免得它引您到那以恶为乐的首领那里；而要追随信德之路，以便您能到达那以善人为乐的王那里。}\par

% structure: p0110 source=h2 target=subsection reason=relative-heading-rank; base=h2

\subsection{第五十五章 两个国度}

对此，老人回答说：\Quote{但那位喜欢邪恶的王子是如何被造的？由什么造出？还是没有造出？} \Person{阿奎拉}说：\Quote{这个话题属于另一时间；但为了使你不至于完全得不到回答就离开，我也对这个问题略陈己见。天主在创世之前预知一切，知道将来有些人会趋向善，另一些人则相反，就将那些选择善的归于祂自己的管理和照顾，称他们为祂独特的产业；但将那些转向邪恶的人的管理权交给了那些天使，他们不是因本性，而是因反对而不愿与天主同在，被嫉妒和骄傲的恶习所败坏。因此，祂使那些人成为相称臣民的相称王子；然而，祂将那些臣民交托给那些天使，使他们没有能力随心所欲地对待他们，除非他们越过了从起初就被指定的界限。而指定的界限就是：除非一个人首先执行魔鬼的意志，否则魔鬼对他没有权柄。}\par

% structure: p0112 source=h2 target=subsection reason=relative-heading-rank; base=h2

\subsection{第56章 邪恶的起源}

然后老人说：\Quote{你说得很好，我的孩子。现在只差你告诉我邪恶的实质来自哪里：因为如果它是由天主造的，邪恶的果实表明根有毛病；因为它似乎也是邪恶的本性。但如果这实质与天主同等永恒，那么那同样无始且同等永恒的东西怎能服从于另一位呢？} \Quote{它并非一直存在，}\Person{阿奎拉}说，\Quote{但也不必然推出，如果它是由天主造的，就认为它的造物主如同祂所造之物那样。因为天主确实造了万物的实质；但如果一个合理的理性，即使由天主所造，不认同造物主的法律，超出所规定的节制的界限，这怎能归咎于造物主呢？或者如果有比这更高的理由，我们则不知道；因为我们不能完全知道任何事，尤其对于那些我们无知不会受审判的事。但那些我们将会受审判的事是最容易理解的，几乎一句话就可完成。因为我们行为的所有规范几乎都归结于此：我们不愿承受的事，就不该对他人做。就如你不想被杀，就必须注意不杀别人；你不想自己的婚姻受玷污，就不可玷污别人的床；你不想被偷，就不可偷窃；人的一切行为都包含在这个规则之内。}\par

% structure: p0114 source=h2 target=subsection reason=relative-heading-rank; base=h2

\subsection{第57章 老人未被说服}

然后老人说：\Quote{不要见怪，我的孩子，我将要说的话。你的话虽然有力，却不能使我相信任何事可以脱离\OriginalTerm{生成}{Genesis}而做成。因为我知道一切事都是因\OriginalTerm{生成}{Genesis}的必然性而临到我，所以我不能相信行善或行恶在于我们的能力；如果我们不能掌握自己的行为，就不能相信将来有审判，使恶人受罚，善人得赏。总之，既然我看你精通这类学问，我就从这门\OriginalTerm{技艺}{art}本身略陈几件事。} \Quote{如果，}\Person{阿奎拉}说，\Quote{你想从那门科学添加什么，我的兄弟\Person{克莱孟}会很小心地回答你，因为他更详细地研习了\OriginalTerm{数学科学}{science of mathematics}。因为我可以用其他方式主张我们的行为在乎我们自己；但我不该僭越那些我没有学过的。}\par

% structure: p0116 source=h2 target=subsection reason=relative-heading-rank; base=h2

\subsection{第58章 审判天主}

当\Person{阿奎拉}这样说了之后，\Person{克莱孟}说：\Quote{明天，我的父亲，你可以随意发言，我们乐意听；因为我想，你所面对的并非不了解你所宣称的科学的人，这也会让你高兴。}因此，当老人与我约定次日就\OriginalTerm{生成}{Genesis}的问题——是否万物都受其影响，还是我们里面有不受\OriginalTerm{生成}{Genesis}支配而由心智判断决定的成分——举行讨论时，\Person{伯多禄}站起身来，开始发表如下言论：\Quote{在我看来，极其奇妙的是，那些容易发现的事，人们却用深奥的思想和言语使之困难；尤其是那些自以为聪明的人，他们想要理解天主的旨意，却把天主当作人，甚至比人还不如：因为一个人若不透露自己的想法，没有人能知道他的意图或心思；也没有人能学会一门专业，除非他长期受老师指导。更何况，没有人能知道无形不可理解的天主的心意或工程，除非祂亲自派遣一位先知来宣布祂的旨意，并阐述祂创造的方式，直到人可合法学习的程度！因此，我认为可笑的是，人们用自然的方式判断天主的权能，认为这对天主可能，那不可能，或这个更大，那个更小，而对于一切都无知；他们是不义的人，却判断正义的天主；无知的，判断设计者；腐败的，判断不朽的；受造物，判断造物主。}\par

% structure: p0118 source=h2 target=subsection reason=relative-heading-rank; base=h2

\subsection{第59章 真先知}

但我不要你们以为，我说这话是剥夺了判断事物的能力；而是劝告人，不要行走在弯曲之处，陷入无尽的错误。因此，我不仅劝告有智慧的人，也劝告所有渴望知道什么对他们有益的人，要寻求真先知；因为只有祂知道一切，知道每个人在寻求什么、如何寻求。因为祂在我们每个人的心中，但在那些没有渴望认识天主及其正义的人身上，祂是不活动的；然而，祂在那些寻求有益于自己灵魂之事的人身上工作，并在他们心中点燃知识之光。因此，先寻求祂；如果你们找不到祂，就不要指望能从别人那里学到什么。但对于那些因爱真理而殷勤寻求祂、且灵魂没有被邪恶占据的人，祂很快就会被找到。因为祂与那些以纯洁精神渴望祂的人同在，他们耐心忍受，因爱真理而从心底发出叹息；但祂离开恶毒的心灵，\BibleRef{智 1:4}因为祂作为先知知道每个人的思念。因此，不要让人以为自己能凭自己的智慧找到祂，除非，正如我们所说，他清空心中一切的邪恶，并怀有纯洁而忠信的渴望去认识祂。因为当一个人这样预备自己时，祂自己作为先知，看见一颗为祂预备好的心，就主动将自己显示给他。\par

% structure: p0120 source=h2 target=subsection reason=relative-heading-rank; base=h2

\subsection{第六十章 祂的裁决不容置疑}

因此，如果有人希望学习一切，他不能通过逐一讨论来做到；因为作为必死的人，他将无法追寻天主的计谋，也无法审视无限本身。但是，如果正如我们所说，他渴望学习一切，就让他寻求真先知；当他找到祂时，不要以问题、辩论和论据来与祂应对；但如果祂作出了任何回应或宣判，就不能怀疑这是确定的。因此，在一切事之前，应寻求真先知，并抓住祂的话语。关于这些，每个人都应讨论这一件事，以便使自己确信这是否真是祂的先知性话语；也就是说，它们是否包含着对未来之事不容置疑的信德，是否标明了确定的时间，是否保持了事物的秩序，是否没有将后来的事说成最先的，也没有将最先的事说成后来的，是否不包含任何诡辩、任何用魔术手法编造来欺骗的内容，或者是否没有将启示给别人的事据为己有，并将其与谎言混杂。当通过正确的判断讨论了这一切，并确认它们是先知的言语时，就应当立刻相信它们关于一切已说过和回应过的事情。\par

% structure: p0122 source=h2 target=subsection reason=relative-heading-rank; base=h2

\subsection{第六十一章 哲学家的无知}

让我们仔细考虑天主眷顾的工作。因为哲学家们引入了一些微妙而难懂的词语，以致连他们演讲中使用的术语都不能为众人所知和理解，天主却表明，那些自以为巧言善辩的人在真理的知识方面完全无知。因为由真先知传授的事物知识是简单、明了、简洁的；而那些行走在弯曲之处、经过言语的荆棘困难的人，对此全然无知。因此，对于谦逊而纯朴的心灵，当他们看到预言的事情成就时，就足够了，甚至远远足够，使他们能从最确定的预知中获得最确定的知识；而对于其余的事，他们可以安息，因为已经获得了明显的真理知识。因为其他一切事物都是由意见处理的，其中不能有稳固的东西。有什么言论是不能反驳的呢？有什么论证是不能被另一个论证推翻的呢？因此，通过这种辩论，人们永远无法达到知识与学问的终点，反而会在问题结束之前找到自己生命的终点。\par

% structure: p0124 source=h2 target=subsection reason=relative-heading-rank; base=h2

\subsection{第六十二章 会议的结束}

\Quote{因此，既然在这些\Emph{哲学家们}中间有不确定的事，我们就必须来到真先知这里。天主圣父愿意祂被众人所爱，因此祂已完全消灭了那些源于人、毫无确定性的意见，好使祂\Emph{真先知}更被寻求，并使那曾被他们遮蔽的祂向人显示真理之路。因为为此天主也创造了世界，并由祂充满世界；因此，即使在最遥远的地极，祂也亲近那些寻求祂的人。但是，如果有人不纯洁、不圣善、不忠信地寻求祂，祂确实在他内，因为祂无处不在，并在所有人的心灵中被找到；但是，正如我们以前所说，祂对不信者来说是沉睡的，并被那些不信祂存在的人视为不在。}当伯多禄说了这话以及更多关于真先知的同样内容后，他遣散了人群；当他非常恳切地请那老人留下来与我们同住时，他未能奏效；但老人也离开了，约定次日回来。此后，我们和伯多禄一起回到住处，享用了惯常的饮食和休息。\par

\SourceRecord{Translated by Thomas Smith.；Ante-Nicene Fathers；Vol. 8.；Edited by Alexander Roberts, James Donaldson, and A. Cleveland Coxe.；Buffalo, NY: Christian Literature Publishing Co.,；1886.；<http://www.newadvent.org/fathers/080408.htm>.}

% structure: p0001 source=h1 target=section reason=page-title

\section{\WorkTitle{本课(卷九)}}

% structure: p0002 source=h2 target=subsection reason=relative-heading-rank; base=h2

\subsection{第一章 解释}

次日，\Person{伯多禄}和我们一起，早早赶到昨天举行讨论的地方。他看到一大群人聚集在那里听讲，又看到那位老人也在其中，便对他说：\Quote{老人，昨天我们约定，今天你应与\Person{克勒孟}讨论。你要么证明没有什么是脱离\Emph{\OriginalTerm{宿命}{genesis}}而发生的，要么\Person{克勒孟}应证明根本没有所谓\Emph{宿命}，而是我们的行为出於我们自己的权能。}老人回答说：\Quote{我记得约定的内容，也记得你在约定之后所说的话。你教导说，人若不从真先知学习，就不可能知道任何事。}于是\Person{伯多禄}说：你不明白我的意思，我现在向你解释。我所讲的是天主在创世之前所怀有的旨意和计划。凭这个计划，祂创造了世界、安排了时间、颁布了法律、应许来世给义人以奖赏他们的善行，并按公义的判决给不义之人定下惩罚。我说，天主的这个计谋和旨意是人无法探知的，因为没有人能靠推测和意见揣度天主的思想，除非天主所派遣的先知宣告它。所以我不是说，任何教义或学问没有先知就不能被探知或认识；因为我知道，各种技艺和学问是人可以知晓并实践的，它们不是从真先知那里学来的，而是从人的导师那里学来的。\par

% structure: p0004 source=h2 target=subsection reason=relative-heading-rank; base=h2

\subsection{第二章 预备事项}

\Quote{既然你自称熟悉星辰的位置和天体运行，并可以从这些说服\Person{克勒孟}，说一切事物都受\Emph{宿命}支配，或者你将从他那里学到万事都由天主眷顾治理，并且我们的行为有些出於我们自己的权能，那么现在正是你们两人着手进行此事的时候了。}老人回答说：\Quote{确实，如果我们能向真先知学习，并听到确定的命题说任何事情取决于我们和我们的自由意志，就没有必要提出这类问题了。因为你昨天的讲论深深影响了我，你争论了先知之能。因此我也同意并确认你的判断：人不能确定无疑地知道任何事，因为他只有短暂的生命，微弱的呼吸以此维持生命。但是，既然我在听到先知之能前已向\Person{克勒孟}承诺，要证明一切事物都受\Emph{宿命}支配，或者从他那里学到我们自身有一些权能，那么就请让他先开始，提出并解释可能遭到反对的问题：因为从你那里听到关于预言之能的几句话后，我承认自己已经困惑，考虑到预知如此伟大；我不认为从推测和意见中收集的任何东西应该被接受。}\par

% structure: p0006 source=h2 target=subsection reason=relative-heading-rank; base=h2

\subsection{第三章 讨论的开始}

老人说完这番话后，我\Person{克勒孟}开始如下讲论：天主借祂的圣子创造了世界，如同一个双层的房屋，被这穹苍（称为天）隔开。祂指派天使之力居住在高层，而让众多人生于这可见的世界，从中祂可以为祂的圣子选择朋友，与祂一同喜乐，并为祂预备如同新娘为新郎一样。但直到婚姻之时（即来世显现），祂指派了一种权能，从生于这个世界的人中拣选并看顾善人，为他们保存给祂的圣子，将他们分别安置在世界一个没有罪的地方；那里已经有了一些人，他们正在预备，如同新娘装饰好等候新郎的来临。因为这世界和今世的王子如同淫夫，败坏和侵犯人的心智，引诱他们离开对新郎的纯爱，去爱奇怪的情人。\par

% structure: p0008 source=h2 target=subsection reason=relative-heading-rank; base=h2

\subsection{第四章 为何创造了恶王子}

但有人会说：那么为什么必须创造那个王子，他要使人心远离真君王？因为天主，正如我所说，愿意为祂的圣子预备朋友，祂不希望他们是因本性的必然而不能成为其他的人，而是希望他们凭自己的选择和意志渴望成为善人；因为不可欲的事物不值得称赞，非有意追求的事物不被判为善。因为如果你的本性必然不允许你改变，那么做那个就没有功德。因此，天主的眷顾愿意许多人诞生在这世界上，以便那些选择善生的人能从众人中被拣选出来。因为祂预见到今世只能由多样和不平等组成，祂按今世事物的多样性赋予每个心智以行动的自由，并指派这个王子，通过他对那些相反事物的暗示，使对更好事物的选择依赖于修德的操练。\par

% structure: p0010 source=h2 target=subsection reason=relative-heading-rank; base=h2

\subsection{第五章 不平等的必然性}

但为了更清楚地表达我们的意思，我们将通过具体例子解释。例如，世上所有的人是否都适合做国王、王子、主人、教师、律师、几何学家、金匠、面包师、铁匠、文法家、富人、农夫、香料商、渔夫或穷人？显然，不可能所有人都成为这些。然而人类生活需要所有这些职业，甚至更多，没有它们就无法度过；因此不平等在此世是必然的。因为不可能有国王，除非他有臣民可以统治和管辖；也不可能有主人，除非他有可以支配的人；其余也如此。\par

% structure: p0012 source=h2 target=subsection reason=relative-heading-rank; base=h2

\subsection{第六章 世界为修德而有的安排}

因此，造物主知道，由于人逃避劳作，没有人会自愿前来参与竞赛——即我们提到的那些职业的实践，通过这些职业可以显明每个人的正义或仁慈——便为人造了一个会饥饿、会口渴、会寒冷的身躯，好使人为了维持身体的缘故，被迫因生计的需要而从事我们提到的所有职业。因为我们为了吃喝穿戴而学习操练每一种技艺。其中便显明了每个人心意的目的：他是要借着偷窃、谋杀、伪证及其他诸如此类的罪行来满足饥渴和寒冷的需求，还是以正义、仁慈和节制，借着职业的实践和双手的劳作，来履行迫切的必需之服务。若他以正义、虔诚和仁慈供应身体所需，他便在所置于他的竞赛中胜出，并被选为圣子的朋友。但若他以欺诈、不义和罪行服侍肉欲，他便成为此世的首领和所有邪魔的朋友；他也被他们教导，将自己恶行的错误归咎于星辰的运行，尽管他是故意且甘愿选择这些恶行的。因为正如我们所说，在饮食欲望的驱使下，人们学习并实践技艺；这欲望当真理的知识到来时便减弱，简朴取而代之。那些只用水和面包、仅仰望天主的人，会有何耗费呢？\par

% structure: p0014 source=h2 target=subsection reason=relative-heading-rank; base=h2

\subsection{第七章 旧的出生与新的出生}

因此，正如我们所说，在世界的安排中存在着某种必要的不平等。因为并非所有人能知晓所有事物、完成所有工作，然而所有人都几乎需要所有事物的使用和服务。因此，有必要一个人劳作，另一个人为他的劳作付酬；一个人为仆，另一个人为主；一个人臣服，另一个人为王。但这种为了人类生活所必要的不平等，天主眷顾已将其转化为正义、仁慈和人道的时机：好让这些事在人与人之间进行时，每个人有机会对他支付工资的人行正义，对或许因疾病或贫穷而无法还债的人行仁慈，对那些因受造而似乎臣服于他的人行仁道，以及对臣属保持柔和，并按照天主的法律行一切事。因为祂赐下了法律，以此辅助人的心智，使他们更容易明白该如何对待每件事，如何逃避邪恶，如何趋向未来的福乐；以及如何在水里重生，借着善工熄灭他们旧出生的火焰。因为我们的第一次出生经由欲望之火而降下，因此，按天主的安排，这第二次出生由水引入，可以熄灭火的本质；并且灵魂被天上的圣神光照，能弃绝第一次出生的恐惧——只要它如此生活直到将来，全然不寻求此世的任何享乐，而是像一个朝圣者和异乡人，以及另一个城邦的公民。\par

% structure: p0016 source=h2 target=subsection reason=relative-heading-rank; base=h2

\subsection{第八章 恶的用途}

但或许你会说，在那些因自然的需要而要求技艺和劳作服务的事上，任何人或许都有能力保持正义，并对他自己的欲望或行为施加他所愿的约束。但对降临到人身上的疾病和虚弱，一些人被邪魔、热病、寒战所折磨，另一些人被疯狂或失去理智所攻击，以及所有那些用无数不幸压倒人类种族的事，我们又要说什么呢？对此我们说，若有人思考整个奥迹的缘由，他会称这些事比我们已经解释的那些更公义。因为天主赐给人一个本性，借此他们可以受教认识善，抵抗恶；也就是说，他们可以学习技艺，抵抗享乐，并在一切事上将天主的法律置于眼前。为此，祂允许某些敌对权势在世上飘荡，与我们争战，理由如前所述，以便借着与它们争战，胜利的棕榈枝和赏报的功绩归于义人。\par

% structure: p0018 source=h2 target=subsection reason=relative-heading-rank; base=h2

\subsection{第九章 \Quote{在罪中受孕}}

因此，有时会发生这样的事：倘若有人行为不节制，且不愿抵抗而宁愿屈服，将这些邪魔容纳在自己里面，那么因它们有害的气息，便会生出不节制、病态且有疾病的后裔。因为当肉欲被完全满足，而在交合中毫无谨慎时，无疑，软弱的后裔会受那些邪魔的缺陷和软弱影响，因为这一切是受它们的煽动而做。因此，父母对他们子女这类缺陷负有责任，因为他们没有遵守交合的法律。尽管还有更隐秘的原因使灵魂服从这些邪恶（现在陈述它们并非我们的目的），但每人仍应承认天主的法律，好从中学习生育的规矩，避免不洁的原因，使所生的是洁净的。因为，在栽种灌木和播种庄稼时，尚需寻找适宜的季节，清理土地，并适当预备一切，免得播下的种子受损而亡；那么，在超越这一切的人身上，难道在播撒其种子时竟不应有丝毫注意或谨慎吗？\par

% structure: p0020 source=h2 target=subsection reason=relative-heading-rank; base=h2

\subsection{第十章 涂上沥青的麻絮}

但有人会说：有些人在幼年时没有任何身体缺陷，然而随着时间推移却陷入那些恶行，以至于有些人甚至被暴力推向死亡，这又如何解释？关于这些，解释也近在眼前，几乎是相同的：因为我们所说那敌对人类的势力，藉着许多不同的贪欲，以某种方式被邀请进入每个人的心中，并找到了进入的途径；它们在其中具有这样的影响力和能力，只能鼓励和煽动，却不能强迫或完成。因此，如果有人同意它们，以致做出他所邪恶欲求的事，他的同意和行为将得到毁灭和最坏死亡的报应。但如果他想到将来的审判，被恐惧所约束，并控制自己，以致没有将他邪恶思想中所构思的付诸行动，他不仅会逃脱现时的毁灭，也会逃脱未来的惩罚。因为每一种罪的根由似乎都像涂了沥青的麻絮，一遇到火焰的热力就立刻燃烧起来；这火的点燃被理解为魔鬼的工作。因此，如果有人被发现涂满了罪恶和贪欲如同涂了沥青，火就容易掌控他。但如果麻絮没有浸透罪的沥青，而是浸透净化和重生的水，魔鬼的火就不能在其中被点燃。\par

% structure: p0022 source=h2 target=subsection reason=relative-heading-rank; base=h2

\subsection{第十一章 恐惧。}

但有人会说：我们现在该怎么办，因为我们已经被罪恶如同沥青般涂抹了？我回答：什么都不用做，但要赶紧被洗净，以便藉着呼求圣名，从你们身上清除那火的燃料，并且将来你们要以对将来审判的恐惧抑制你们的贪欲，以坚定的决心击退那敌对势力，每当它们接近你们的感官。但你又说：如果有人陷入爱情，他怎么能自制，即使他眼前看到那所谓的\Quote{火河}\OriginalTerm{火河}{Pyriphlegethon}？这是那些不愿悔改之人的借口。但现在我不希望你谈论火河。将人类的惩罚放在你面前，看看恐惧有何威力。当有人因爱的罪行被带到刑罚处，被绑在火刑柱上焚烧时，他那时还能对他所爱之人产生任何欲望，或将她的形象放在眼前吗？你一定会说：绝对不能。你看，现时的恐惧能断绝不义的欲望。但那些信天主、承认将来审判和永火惩罚的人——如果他们不远离罪恶，那肯定他们不是以十足的信德相信：因为如果信德是确定的，恐惧也成为确定的；但如果信德有任何缺欠，恐惧也就减弱，然后相反的势力便找到进入的机会。当他们同意这些势力的劝说时，他们必然也屈服于它们的权力，并在它们的煽动下被推向罪恶的悬崖。\par

% structure: p0024 source=h2 target=subsection reason=relative-heading-rank; base=h2

\subsection{第十二章 占星家。}

因此，占星家对这些奥秘一无所知，认为这些事情是由天体的运行所导致；因此，在回答那些前来咨询他们未来之事的人时，他们在许多事情上受骗。这并不奇怪，因为他们不是先知；但通过长期实践，错误的创立者在他们所欺骗的事物中找到某种避难所，并引入某些\OriginalTerm{危险期}{Climacteric Periods}，以便假装对不确定之事有所认识。因为他们把这些危险期描绘为危险时期，在其中一个人有时被毁灭，有时不被毁灭，却不知道支配这些的不是星宿的运行，而是魔鬼的作为；那些魔鬼急于证实占星术的错误，用数学计算欺骗人犯罪，以便当他们因罪受罚时，或因天主的允许，或因法律判决，占星家似乎说出了真相。然而，即使在这件事上也受骗；因为如果人迅速转向悔改，并记念和畏惧将来的审判，死亡的惩罚就会被赦免给那些藉着圣洗的恩宠而皈依天主的人。\par

% structure: p0026 source=h2 target=subsection reason=relative-heading-rank; base=h2

\subsection{第十三章 现世或来世的报应。}

但有人会说：许多人甚至犯了谋杀、奸淫及其他罪行，却没有遭受任何灾祸。这在人类中确实很少发生，但对于那些不明白天主计划的人，似乎常常发生。但天主知晓一切，他知道犯罪的人如何及为何犯罪，以及什么原因导致每个人犯罪。然而，一般要注意的是：如果有人是邪恶的，与其说是心思，不如说是行为，并且不是出于故意煽动而犯罪，惩罚就更快地施加在他们身上，更在今世；因为无论何时何地，天主都按每人的行为报应各人，按祂看为合宜的。但那些存心作恶的人，甚至有时对恩待他们的人发怒，且不顾悔改——他们的惩罚被推迟到将来。因为这些人不像我们前面所说的那些人，不配在今世终结他们罪行的惩罚；而是被允许按照自己的意愿占用现世，因为他们的改正不需要暂罚，而是要求地狱永火的惩罚；在那里，他们的灵魂将寻求悔改，却找不到。\par

% structure: p0028 source=h2 target=subsection reason=relative-heading-rank; base=h2

\subsection{第十四章 知识抑制情欲。}

然而，若他们在此生中就将那时将要受的惩罚摆在眼前，他们必定会节制自己的情欲，绝不会陷入罪恶。因为灵魂中的理智大有能力斩断一切欲望，尤其是当它获得了对天上事物的认识，借此接受了真理之光，便会远离一切恶行的黑暗。因为正如太阳以其光芒的明亮遮蔽并掩盖所有星辰，同样，心智借着知识的之光，使灵魂的一切情欲都变得无效和不活跃，将未来审判的思想如光芒般投射于其上，以致它们再不能在灵魂中出现。\par

% structure: p0030 source=h2 target=subsection reason=relative-heading-rank; base=h2

\subsection{第十五章 对人的畏惧与对天主的畏惧}

\Quote{但作为证据，证明对天主的畏惧在抑制情欲上大有功效，请以对人的畏惧为例。在众人中，谁不贪恋邻人的财物？然而他们因惧怕法律所规定的惩罚而克制自己，行事正直。因畏惧，万民服从君王，军队手持武器而听命。奴隶虽比主人强壮，却因畏惧而服从主人的管辖。甚至野兽也因畏惧而被驯服；最强壮的公牛低头套上轭，巨大的大象因畏惧而听从主人。但我们为何用人间的例子，岂不也有神性的例子？大地本身岂不存于诫命的畏惧之下，以其震动和颤抖为证？大海守住指定的界限；天使保持和平；星辰遵守秩序，江河保持河道：恶魔也因畏惧而逃遁，这是确定的。且不必以太多细节冗长论述，请看对天主的畏惧如何约束一切，使万事保持适当的和谐与固定的秩序。那么，你们岂不更可确信，那在你们心中升起的恶魔的情欲，能因对天主之畏惧的劝诫而被熄灭并完全消除，既然连情欲的煽动者本身也因畏惧的影响力而逃遁？你们知道这些事确实如此；但若你们有何答复，请讲。}\par

% structure: p0032 source=h2 target=subsection reason=relative-heading-rank; base=h2

\subsection{第十六章 不完全的认信}

于是老人说：我的儿子克莱孟明智地构筑了他的论证，以致他使我们对此无话可说；但他关于人性所做的一切论述都有一个主旨，即在于人的意志自由的同时，也有某种外在的恶因，人由此确实受到各种情欲的诱惑，却并不被迫犯罪；并且为此，他说，因为畏惧远比它们更有力，它抵抗并抑制欲望的猛烈，因此，即使自然的情感生起，罪恶仍可不犯，那些煽动并点燃这些情感的恶魔也被驱散。但这些事并不能说服我；因为我意识到某些事，由此我深知，借着天体的安排，人成为杀人犯或奸淫者，并犯下其他恶事；同样，尊贵而贞洁的妇女也被迫行善。\par

% structure: p0034 source=h2 target=subsection reason=relative-heading-rank; base=h2

\subsection{第十七章 占星学识}

\Quote{简言之，当火星居本位中心，与土星成四分相，水星向中心，满月降临其上，在每日的命宫中，它产生杀人犯、将死于刀剑者、嗜血者、醉汉、好色者、魔鬼般的人、探究隐秘者、作恶者、亵圣者，以及诸如此类的人；尤其当没有一颗吉星注视之时。但火星自身与金星成四分相，向中心方向，而无吉星注视时，则产生奸淫者与乱伦者。金星与月亮在土星的界与宫内，若它与土星相合，且火星注视，则产生具有男子气的女子，乐于农耕、建筑及一切男子工作，随意与人通奸而不被丈夫发现，不用香料、香膏、妇女衣饰鞋履，而效男子生活方式。但凶险的金星则使男子如同女子，在一切方面不似男子行事，若它与火星在白羊座；反之，若它在摩羯座或宝瓶座，则使女子生女。}\par

% structure: p0036 source=h2 target=subsection reason=relative-heading-rank; base=h2

\subsection{第十八章 答复}

当老人洋洋洒洒地讲述这个话题，列举了每一种天象图以及天体的位置，借此想要证明畏惧不足以抑制情欲时，我再次回答：诚然，父亲，您论证得最为博学与技巧；理性本身也邀请我对您的言论略作答复，因为我确实通晓算术科学，并乐于与如此博学的人交流。因此请听我回答您所说的，以便您清楚知道命宫绝不来自星辰，并且那些投靠天主的人能够抵抗恶魔的侵袭；而且，如我先前所说，不仅对天主的畏惧能克制自然的情欲，就连对人的畏惧也能，这我们即将向您说明。\par

% structure: p0038 source=h2 target=subsection reason=relative-heading-rank; base=h2

\subsection{第十九章 驳斥占星术}

在每个国家或王国中，都有人订立的律法，或以文字，或仅凭习俗而持久，无人轻易违犯。简言之，最初的赛里斯人，居于世界之初，订有律法：不知杀人、奸淫、淫乱，不犯偷盗，不拜偶像；在那极其广大的整个地区，没有庙宇，没有偶像，没有娼妓，没有淫妇，也没有盗贼受审。并且那里从未有人被杀；也没有人的自由意志，按你们的学说，被火星的炽星强迫动刀杀人；金星与火星相合也不强迫人犯奸淫，尽管当然在他们那里，火星每日占据中天。但在赛里斯人那里，对法律的畏惧比命宫的力量更强大。\par

% structure: p0040 source=h2 target=subsection reason=relative-heading-rank; base=h2

\subsection{第二十章 婆罗门}

同样，在印度地区的\Place{巴克特里亚}人中，有不可胜数的\Person{婆罗门}，他们亦因祖先的传统以及和平的习俗与法律，既不杀人，也不犯奸淫，不崇拜偶像，不食荤腥，从不醉酒，从不作恶，而是常怀敬畏天主之心。他们确实如此行，尽管其余印度人杀人、犯奸淫、崇拜偶像、醉酒并施行此类恶事。是的，在印度本土西部有一地，陌生人进入该地被捕杀并吃掉；既无吉星阻止这些人行此恶事、食用可咒之食，也无凶星迫使\Person{婆罗门}行恶。再者，波斯人中有一习俗，与母亲、姐妹及女儿结婚。在整个地区，波斯人缔结乱伦婚姻。\par

% structure: p0042 source=h2 target=subsection reason=relative-heading-rank; base=h2

\subsection{第二十一章 天域}

为使那些研究星象者无法使用以下托词，即他们声称天上有某些区域被赋予特有之事物，那些波斯民族中有人前往外国，他们被称为\OriginalTerm{玛哥斯人}{Magusæi}，其中有些至今仍在\Place{玛代}，有些在\Place{帕提亚}，有些也在\Place{埃及}，并在\Place{迦拉达}和\Place{夫黎基雅}也有相当数量，所有这些人都毫无变异地持守这乱伦传统的形式，并将其传给后代遵守，即使他们已改变了天域；金星与月亮在土星的界和宫内，并有土星和火星注视，并未强迫他们在其他人中拥有\OriginalTerm{出生星位}{Genesis}。\par

% structure: p0044 source=h2 target=subsection reason=relative-heading-rank; base=h2

\subsection{第二十二章 格洛内人的习俗}

在\Place{格洛内人}中也有一种习俗：妇女耕种田地、建造房屋，并做一切男子之事；她们也被允许随意与人交合，不被丈夫责难，也不被称为奸妇；因为她们到处随意交合，尤其与陌生人；她们不使用香料，不穿染色衣裳，也不穿鞋。另一方面，格洛内男子则装饰、梳妆、穿着柔软色彩缤纷的衣服，佩戴黄金，涂抹香料，这并非缺乏男子气概，因为他们极其好战，且是热忱的猎人。然而，所有格洛内妇女出生时金星并不在摩羯座或水瓶座的不利位置；所有男子也并非金星与火星在白羊座相合，迦勒底星术认为这种配置会使人出生时柔媚放荡。\par

% structure: p0046 source=h2 target=subsection reason=relative-heading-rank; base=h2

\subsection{第二十三章 苏萨人的风俗}

此外，在\Place{苏萨}，妇女使用香料，且是上等品，饰以装饰品和宝石；她们外出时有婢女扶持，比男子更为炫耀。然而她们并不培养谦逊，而是随意与任何人（奴隶或客人）交合，其丈夫允许这种自由；她们不仅不因此受责备，反而统治丈夫。然而，所有苏萨妇女的\OriginalTerm{出生星位}{Genesis}并非金星与木星和火星在天空正中木星的宫内。在遥远的东方，若男孩遭受非自然对待，一旦发现，便被其兄弟、父母或任何亲属杀死，且不得埋葬。再者，高卢人中有一古老法律公开允许男孩如此受辱，并不视为耻辱。难道所有在高卢人中受此卑劣对待的人，都曾经在土星的宫室和火星的疆界中有\OriginalTerm{晓星}{Lucifer}与水星同在吗？\par

% structure: p0048 source=h2 target=subsection reason=relative-heading-rank; base=h2

\subsection{第二十四章 各国的不同习俗}

在\Place{不列颠}地区，多个男子共有一妻；在\Place{帕提亚}，多个女子共有一夫；世界各地都固守自己的风俗习惯。\Person{阿玛宗人}没有丈夫，而是像动物一样，每年春分时离开自己的领土一次，与邻近民族的男子共同生活，同时举行某种仪式，怀孕后返回；若生男婴，便抛弃，只抚养女性。既然所有人的出生都在同一季节，便无需假设在男性出生时火星与土星同时均分，而在女性出生时则从未如此；也无需假设他们不曾有金星与水星在其本宫相合，以产生画家、雕刻家或兑换银钱者；或在金星的宫中，以产生调香师、歌手或诗人。在撒拉逊人、上利比亚人、摩尔人、大洋河口居民、德国偏远地区、萨尔马提亚人、斯基泰人、本都海岸地区所有民族以及克律塞岛，从未有过兑换银钱者、雕刻家、画家、建筑师、几何学家、悲剧诗人或诗人。因此水星和金星的影响在他们中间必然缺失。\par

% structure: p0050 source=h2 target=subsection reason=relative-heading-rank; base=h2

\subsection{第二十五章 非由出生星位，而是自由意志}

在全世界中，唯有玛待人极其谨慎地将尚在呼吸的人扔给狗吃；然而，在他们每日的\OriginalTerm{星命}{Genesis}中，火星与月亮并非都处于巨蟹座。印度人焚烧死者，死者的妻子自愿献身，与丈夫一同被焚。但并非所有被活活烧死的印度妇女，在夜间的\OriginalTerm{星命}{Genesis}中，太阳都位于地下，且火星位于火星的区域。许多日耳曼人用绳索结束自己的生命；但并非所有日耳曼人，因此都有月亮与\OriginalTerm{Hora}{Hora}被土星和火星环绕。由此可见，对法律的恐惧在各国占据主导地位，而由圣神植入人内的自由意志遵循法律；\OriginalTerm{星命}{Genesis}既不能强迫塞里斯人杀人，也不能强迫婆罗门吃肉，不能强迫波斯人避免乱伦，不能强迫印度人停止焚烧，不能强迫玛待人避免被狗吃，不能强迫帕提亚人多妻，不能强迫美索不达米亚妇女保持贞洁，不能强迫希腊人进行体育训练，不能强迫高卢男孩受虐待；它也不能迫使野蛮民族学习希腊人的学问；反而，如我们所说，每个民族根据自由意志遵守自己的法律，并以法律的严苛废除了\OriginalTerm{星命}{Genesis}的定数。\par

% structure: p0052 source=h2 target=subsection reason=relative-heading-rank; base=h2

\subsection{第二十六章　\OriginalTerm{气候带}{Climates}}

但某个精通数学科学的人会说，\OriginalTerm{星命}{Genesis}被分为七个部分，他们称之为\OriginalTerm{气候带}{climates}，每一个气候带由七个天体之一统治；我们所提及的那些不同的法律，并非由人制定，而是由那些主宰星辰按其意志颁布的，星辰所喜悦的，人便作为法律遵守。对此，我们首先回答：世界并非分为七个部分；其次，即使如此，我们在一个部分、一个国家内也会发现许多不同的法律；因此，既没有与天体数目相应的七种\Emph{法律}，也没有与星座数目相应的十二种，更没有与十度分区数目相应的三十六种；而是无法计数的。\par

% structure: p0054 source=h2 target=subsection reason=relative-heading-rank; base=h2

\subsection{第二十七章　\Quote{\OriginalTerm{气候带}{Climates}}学说站不住脚}

此外，我们应当记住前述之事：在同一个印度国家中，既有人吃人肉，也有人连羊肉、鸟肉和一切活物都不吃；\OriginalTerm{Magusæi}{Magusæi}不仅在波斯娶母和女为妻，而且在任何他们居住的民族中，都坚持其乱伦的习俗。然后，此外，我们还提到了无数民族，他们完全不认识文学学问，而且有些智者在若干地方修改了法律本身；有些法律因无法遵守或因其可耻而被主动废弃。当然，我们可以轻易查明，有多少统治者改变了被征服民族的法律和习俗，并将其置于自己的法律之下。罗马人明显做了这事，他们几乎将全世界以及所有先前生活在各自不同法律和习俗下的民族都纳入了罗马法和民事法令。因此，被罗马人征服的民族的星辰，便失去了它们的气候带和份额。\par

% structure: p0056 source=h2 target=subsection reason=relative-heading-rank; base=h2

\subsection{第二十八章　犹太人的习俗}

我还要加添一件事，即使最怀疑的人也能满意。所有生活在梅瑟法律下的犹太人，都在第八天毫无例外地给自己的儿子行割损礼，使幼婴流血。但从来没有一个外邦人在第八天接受这个；另一方面，也从来没有一个犹太人省略过它。那么，既然犹太人遍布世界各地，与外邦人混杂，且在第八天忍受割损肢体，\OriginalTerm{星命}{Genesis}的说法又如何成立呢？并且，正如我所说，除了他们自己，没有一个外邦人这样做，这不是受任何星辰的强迫，也不是因血流的驱使，而是出于他们宗教的法律；无论他们身处世界何地，这个标记对他们都是熟悉的。此外，他们所有人，无论身在何处，都有同一名字，这也是通过\OriginalTerm{星命}{Genesis}而来的吗？而且，他们中出生的孩子从未被遗弃，并且每第七天，无论在哪里，他们都安息，不行路，不用火。\BibleRef{出 35:3}那么，为什么没有一个犹太人被\OriginalTerm{星命}{Genesis}强迫在那天行路、建造、买卖任何东西呢？\par

% structure: p0058 source=h2 target=subsection reason=relative-heading-rank; base=h2

\subsection{第二十九章　福音比\Quote{\OriginalTerm{星命}{Genesis}}更有力}

但我要就此事提供一个更强的证据。因为，看哪，自从正义而真正的先知来临以来，才过了不到七年；在这期间，各国的人来到犹太，因所见到的征兆和奇迹，以及祂教义的伟大而感动，接受了祂的信德；然后回到自己的国家，便弃绝了外邦人无法无天的仪式和乱伦的婚姻。简言之，在帕提亚人中——正如在他们中间宣讲福音的多默给我们写信所说的——现在很多人不再有多妻的恶习；在玛待人中间，很少人再把死人扔给狗吃；波斯人也不以与母亲或女儿乱伦为乐；苏西亚妇女也不再行她们被允许的奸淫；宗教的教导约束了那些\OriginalTerm{星命}{Genesis}无法强迫其犯罪的人。\par

% structure: p0060 source=h2 target=subsection reason=relative-heading-rank; base=h2

\subsection{第三十章　\Quote{\OriginalTerm{星命}{Genesis}}与天主的公义不符}

看，就从我们现在讨论的这件事本身来推断，从我们目前所处的环境来得出结论：仅仅因为一个传言传到人的耳中，说有一位先知在犹太地出现，用神迹和奇事教导人敬拜独一天主，所有人就都带着预备好且热切的心期待着——甚至在我主伯多禄到来之前——有人会向他们宣布那个出现者所教导的内容。但为避免我显得列举过多，我将告诉你应从整体得出什么结论。既然天主是公义的，且祂亲自创造了人的本性，又怎会设置\OriginalTerm{星盘}{genesis}来与我们对抗，强迫我们犯罪，然后当我们犯罪时又惩罚我们？因此，可以肯定的是，天主无论是今生还是来世都不惩罚任何罪人，除非因为祂知道那人本可以得胜，却忽视了胜利。因为即使在今世，祂也向人施行报复，就像对待那些在洪水中灭亡的人一样，他们全部在一日之内，甚至在一小时之内被毁灭，尽管可以肯定他们并非都按\OriginalTerm{星盘}{genesis}的秩序在同一小时出生。但如果说，若没有先前的罪，我们遭受灾祸是天性使然，那就极其荒谬了。\par

% structure: p0062 source=h2 target=subsection reason=relative-heading-rank; base=h2

\subsection{第三十一章　知识的价值}

因此，如果我们渴望救恩，就应当首先追求知识，确信：如果我们的心智留在无知中，我们不仅要承受\OriginalTerm{星盘}{genesis}带来的灾祸，还会承受魔鬼随己意加给我们的任何外在灾祸——除非对法律和将来审判的恐惧抵制我们的一切欲望，并遏止犯罪的暴力。因为连人的恐惧也做了许多善事，也做了许多恶事，这是\OriginalTerm{星盘}{genesis}所不知的，正如我们上文所表明的。因此，我们的心智在三个方面容易出错：来自恶习所带给我们的；来自身体自然地在我们内激起的私欲；或来自敌对力量所强迫我们的。但心智的本性有能力对抗和争战这些，当真理的知识照耀它时，这种知识便赋予了对将来审判的恐惧——这是心智的合适治理者，能将它从私欲的悬崖边挽回。因此，这些事在我们的能力之内，已经充分说明了。\par

% structure: p0064 source=h2 target=subsection reason=relative-heading-rank; base=h2

\subsection{第三十二章　顽固的事实}

\Quote{现在，老人家，如果你还有什么要回答这些事的，请说吧。}然后老人说：\Quote{我的孩子，你论述得最为充分；但我，正如一开始所说的，被自己的意识所阻止，不能同意你这无与伦比的陈述。因为我知道我自己的\OriginalTerm{星盘}{genesis}和我妻子的\OriginalTerm{星盘}{genesis}，而且我知道那些事已经发生了，正如我们的\OriginalTerm{星盘}{genesis}为我们每个人所预定的。我现在不能因言语而背离那些我通过事实和行为所确定的事。简言之，既然我看出你在这方面学识极为精湛，请听我妻子的\OriginalTerm{星盘}{horoscope}，你便会发现那已经发生的结果的格局。因为她的金星与火星合于中天以上，月亮落于火星的宫位和土星的界内。这种格局导致女人成为淫妇，爱上自己的奴隶，并在异乡旅行和水中结束生命。这事果然发生了。因为她爱上了她的奴隶，同时害怕危险和耻辱，便与他私奔，到了国外，在那里满足了她的爱欲，并死于海中。}\par

% structure: p0066 source=h2 target=subsection reason=relative-heading-rank; base=h2

\subsection{第三十三章　临近的相认}

然后我回答：\Quote{你怎么知道她在国外与奴隶同居并死在他身边？}老人说：\Quote{我确知无疑；我并不知道她嫁给了那个奴隶，实际上我甚至没有发现她爱他。但当她走后，我的兄弟告诉了我整个故事，说起初她爱上了他自己；但他作为兄弟是正派的，不肯用乱伦的污点玷污兄弟的床。但她既害怕我，又承受不住不幸的责难（尽管她不该为\OriginalTerm{星盘}{genesis}所迫的事受责备），便假装做了一个梦，对我说：\InnerQuote{有一个人出现在我的异象中，命令我立即带着我的两个双胞胎离开这城。}我听到这话，为她和我的儿子们担忧，便立刻送走了她和孩子们，自己留下了那个年幼的。因为她说，那在她睡眠中警告她的人允许她这样做。}\par

% structure: p0068 source=h2 target=subsection reason=relative-heading-rank; base=h2

\subsection{第三十四章　故事的另一方面}

然后我克莱孟，明白他或许就是我的父亲，眼泪汪汪；我的兄弟们也准备冲上前去揭露此事；但伯多禄制止了他们，说：\Quote{安静，等我允许。}于是伯多禄回答老人说：\Quote{你幼子的名字是什么？}他说：\Quote{克莱孟。}伯多禄说：\Quote{如果我今天把你最贞洁的妻子和你的三个儿子归还给你，你会相信一颗谦逊的心能克服不合理的冲动，并且我们所说的一切都是真实的，\OriginalTerm{星盘}{genesis}算不得什么吗？}老人说：\Quote{正如你不可能实现你所承诺的，同样，任何事也不可能脱离\OriginalTerm{星盘}{genesis}而发生。}伯多禄说：\Quote{我希望所有在场的人作证，我今天将把生活得极为贞洁的妻子和你的三个儿子交还给你。现在，请从这一点取得这些事的证据：我知道整个故事比你准确得多；我将按次序叙述全部经过，好让你知道，也让在场的人了解。}\par

% structure: p0070 source=h2 target=subsection reason=relative-heading-rank; base=h2

\subsection{第三十五章　启示}

他这样说后，便转向群众，开始说道：\Quote{诸位，你们所看见的这个身着褴褛衣衫的人，是罗马城的公民，出身于凯撒的家族。他名叫福斯蒂尼安。他娶了一位最高贵的女子为妻，名叫玛提狄。他与她生了三个儿子，其中两个是双胞胎；那个较年幼的，名叫克莱孟，就是这个人！}他说这话时，用手指着我。\Quote{他的双胞胎儿子就是这些人：尼刻塔和阿奎拉，其中一人从前叫福斯蒂努斯，另一人叫福斯图斯。}但伯多禄刚一说出我们的名字，老人的四肢便瘫软了，昏厥过去。我们这些作儿子的赶紧跑过去，拥抱并亲吻他，担心无法唤回他的灵魂。当这些事情发生时，众人因惊奇而困惑不已。\par

% structure: p0072 source=h2 target=subsection reason=relative-heading-rank; base=h2

\subsection{第三十六章 新的启示。}

但伯多禄命令我们从拥抱父亲中起来，免得我们害死他；他亲自抓住老人的手，将他从深睡中扶起，并逐渐使他苏醒，便开始向他叙述真实发生的全部经过：他的兄弟如何爱上了玛提狄，而她因极其端庄，不愿将兄弟的无耻爱情告诉丈夫，以免在兄弟间挑起敌意，使家族蒙羞；她如何明智地编造了一个梦，藉此被命令带着双胞胎儿子离开城市，将年幼的儿子留给父亲；他们在航海中如何因风暴的猛烈而遭遇海难；当他们被抛到一个叫安塔拉杜斯的岛上时，玛提狄被浪涛抛到岩石上，而她的双胞胎孩子被海盗掳走，带到凯撒勒雅，卖给了一位虔诚的妇人，她把他们当作儿子抚养，让他们接受教育，成为有教养的人；海盗如何改了他们的名字，一个叫尼刻塔，另一个叫阿奎拉；后来他们如何通过\Emph{共同}的学业和交往，依附了西满；当他们看出他是术士和骗子时，便离开了他，来到匝加略处；随后他们又如何与自己结盟；克莱孟如何为了寻求真理而从城市出发，通过结识巴尔纳伯来到凯撒勒雅，与自己相识，并依附于他，从他那里学习了信仰的道理；他又如何找到并在安塔拉杜斯认出了行乞的母亲，全岛如何因他的认亲而喜乐；还有她如何与最贞洁的女主人同住，以及自己为她所行的治愈；克莱孟如何慷慨对待那些善待他母亲的人；后来尼刻塔和阿奎拉问起那陌生妇人是谁，听了克莱孟的整个故事后，他们喊出自己是她的双胞胎儿子福斯蒂努斯和福斯图斯；他们如何讲述了所发生的一切；后来，在伯多禄本人的劝说下，他们如何谨慎地被带到母亲面前，以免她因突如其来的喜乐而昏厥。\par

% structure: p0074 source=h2 target=subsection reason=relative-heading-rank; base=h2

\subsection{第三十七章 另一次认亲。}

但伯多禄正将这些细节说给老人听时，这段叙述令众人十分愉悦，听众因事件的神奇和对人类苦难的同情而流泪，我的母亲（不知怎的）听说父亲被认出，气喘吁吁地冲进我们中间，喊着说：\Quote{我的丈夫，我的主人福斯蒂尼安在哪里？他长久受苦，从一个城市流浪到另一个城市寻找我？}她像疯了一样喊叫，环顾四周，老人跑上前来，流着泪拥抱并紧紧搂住她。这些事发生时，伯多禄请求群众散去，说再待下去不合适，必须给他们私下相见的机会。\Quote{但明天，}他说，\Quote{如果有人愿意，让他们来听道理。}\par

% structure: p0076 source=h2 target=subsection reason=relative-heading-rank; base=h2

\subsection{第三十八章 \Quote{无意中接待了天使}。}

当伯多禄说完这话，群众便散去了。我们也打算回住处去，这时家主对我们说：\Quote{让这样伟大的人物住旅店，实属卑劣且邪恶，而我家中几乎整栋楼空着，床铺很多，一切必需品都备齐了。}伯多禄却推辞。家主妻子便带着子女俯伏在他面前，恳求说：\InnerQuote{我恳求您，请与我们同住。}但即便如此，伯多禄仍未同意，直到那些求他的人的女儿——她被不洁之魔折磨已久，被锁链捆住，关在一间小屋里——魔鬼从她身上被赶出，小屋的门也开了，她带着锁链前来，俯伏在伯多禄足前，说：\InnerQuote{主啊，今天您理当在我这里庆祝我的获救之宴，不要使我或我的父母悲伤。}伯多禄问及她锁链与话语的含义，她父母因女儿出乎意料的康复而喜出望外，竟惊愕得哑口无言；但侍候的仆人们说：\InnerQuote{这女孩从七岁起就被魔鬼附身，常撕咬、啃啮，甚至撕碎任何试图接近她的人，二十年来从未停息，直至今日。无人能治愈她，甚至无人能接近她，因为她使许多人残废，甚至毁了一些人；她比任何男人都强壮，显然是因魔鬼之力而增强。但如今，如您所见，那魔鬼已因您的出现而逃遁，那用最大力气关上的门也敞开了，她本人正以健全的心智站在您面前，请求您使她和她的父母因她康复而快乐，并与他们同住。}一个仆人说完这话，锁链自动从她手脚松开。伯多禄确知那女孩因他而恢复健康，便同意与他们同住。他又吩咐那些留在旅店的人，包括他的妻子，也过来。我们每个人都分到一间单独的卧室，便住了下来。照常进食后，赞颂了天主，便在各自的房间安歇了。\par

\SourceRecord{Translated by Thomas Smith.；Ante-Nicene Fathers；Vol. 8.；Edited by Alexander Roberts, James Donaldson, and A. Cleveland Coxe.；Buffalo, NY: Christian Literature Publishing Co.,；1886.；<http://www.newadvent.org/fathers/080409.htm>.}

% structure: p0001 source=h1 target=section reason=page-title

\section{\WorkTitle{辩驳书（卷十）}}

% structure: p0002 source=h2 target=subsection reason=relative-heading-rank; base=h2

\subsection{第一章 试炼}

但到了早晨，日出之后，我克莱孟、尼塞塔、阿桂拉，同伯多禄一起，来到我父母睡觉的房间；见他们还在熟睡，我们便在门前坐下。伯多禄这样对我们说：\Quote{最亲爱的同仆们，请听我说：我知道你们对你们的父亲怀有深厚的感情；因此，我担心你们会过早地催促他负起信仰的重轭，而他现在尚未准备好；他或许会因为对你们的爱而同意此事。然而，这并不可靠；因为为人的缘故所做的事不值得赞许，而且很快就会瓦解。所以，在我看来，你们应当允许他按照自己的判断生活一年；在那段时间里，他可以与我们同行，当我们教导他人时，他可以单纯地聆听；当他聆听时，如果他真有认明真理的意向，他自会请求负起信仰的重轭；或者，如果他不愿接受，他仍可保持朋友关系。因为那些并非全心接受的人，当他们开始无法承受时，不仅会抛弃他们所接受的，而且仿佛为他们的软弱寻找借口，开始诋毁信仰之道，并恶言攻击那些他们无法追随或效法的人。}\par

% structure: p0004 source=h2 target=subsection reason=relative-heading-rank; base=h2

\subsection{第二章 一个难题}

对此，尼塞塔回答说：\Quote{我的主伯多禄，我绝不反对你正确而良好的建议；但我想说一件事，好从中了解一些我所不知道的事。如果我的父亲在您建议推迟的那一年内去世了，那该怎么办？他会无助地下到地狱，永远受折磨。}然后伯多禄说：\Quote{我理解你对你父亲的好意，并原谅你在你无知之事上的态度。因为难道你以为，如果有人被认为生活正直，他就会立刻得救吗？难道你不认为，他必须由那洞察人心的主审判，看他如何生活正直：是否依照外邦人的规范，服从他们的制度和法律？或为了人的友谊？或仅仅出于习惯，或任何其他原因？或出于必要，而非为了义德本身并为了天主？因为那些唯独为了天主和他的义德而生活正直的人，必进入永恒的安息，并承受天上王国的永恒。因为救恩不是靠强力获得的，而是靠自由；不是靠人的恩惠，而是靠天主的信德。此外，你还应当考虑：天主是预知的，他知道这人是否是祂的人。但如果他知道他不是，我们又该如何对待那些从起初就由祂所定的事呢？但在可能的范围内，我给出建议：等他醒了，我们坐在一起时，你们可以装作想学些什么，提出一些适合他学习的问题；当我们彼此交谈时，他会受到教诲。然而，先看看他是否自己会提问；如果他这样做，谈话的机会就会更合适。但如果他不问任何问题，就让我们轮流互相提问，如我所说，想学些东西。这是我的判断，请你们说说你们的看法。}\par

% structure: p0006 source=h2 target=subsection reason=relative-heading-rank; base=h2

\subsection{第三章 一个建议}

我们赞扬了他正确的建议后，我克莱孟说：\Quote{在一切事上，结局多半回顾开端，事物的终局与它们的开始相似。因此，我盼望，就我们的父亲而言，既然天主藉着您给了我们一个好的开始，祂也会赐予一个与开始相称、与祂相称的结局。不过，我提出这个建议：如果，如您所说，我们在我父亲面前开始说话，仿佛是为了讨论某个问题或互相学习，那么，我的主伯多禄，您不应扮演一个需要学习任何东西的角色；因为如果他看到这一点，反而会感到不悦。他深信您完全知道一切，正如您确实知道一切一样。那么，如果他看到您假裝无知，结果会怎样呢？正如我所说，这反而会伤害他，因他不知您的计划。但如果我们在交谈中，我们兄弟之间有任何疑问，希望您能给我们一个适当的解答。因为如果他看到连您都犹豫不决、心存疑虑，那么他必然会认为没有人知道真理。}\par

% structure: p0008 source=h2 target=subsection reason=relative-heading-rank; base=h2

\subsection{第四章 自由探讨}

对此，伯多禄回答说：\Quote{我们不必为此担忧；如果确实他应当进入生命之门，天主自会提供适当的时机；并且将有一个来自天主的开始，而非来自人。因此，如我所说，让他与我们同行，聆听我们的讨论；但因我看到你们急切，所以我曾说必须寻找机会；当天主赐下机会时，你们要听从我的建议行事。}我们正在这样交谈时，一个男孩来告诉我们，我们的父亲已经醒了；我们正要进去看他时，他亲自来到我们面前，以亲吻问候我们。我们重新坐下后，他说：\Quote{如果一个人想问问题，是否允许？还是像毕达哥拉斯学派那样强制保持沉默？}伯多禄说：\Quote{我们并不强迫那些来到我们这里的人要么一直沉默，要么提问；而是让他们自由行事，因为我们知道，那对自己救恩焦虑的人，如果他的灵魂某处感到痛苦，是不会让它沉默的。但那忽视自己救恩的人，即使强迫他提问，对他也没有益处，除了这一点：他可能会显得热心和勤勉。因此，如果你想了解什么，尽管问吧。}\par

% structure: p0010 source=h2 target=subsection reason=relative-heading-rank; base=h2

\subsection{第五章 善与恶}

那年长者说：\Quote{在希腊哲学家中有一个非常流行的说法，即认为在人的生命中实际上既没有善也没有恶；而是人们依照生活的习俗和惯例，凭自己的偏见把事物称为善或恶。因为甚至连谋杀也不是真正的恶，因为它使灵魂脱离肉体的束缚。此外，他们说，即使正义的法官也会处决犯罪的人；但如果他们知道杀人是恶，正义的人就不会那样做。他们也不说奸淫是恶；因为如果丈夫不知道或不关心，他们说，其中并无恶。他们说，盗窃也不是恶；因为它从拥有的人那里拿走他所没有的东西。而且，确实，它本应自由而公开地取走；但因为是秘密做的，这反而揭露了那个被秘密拿走东西的人的残忍。因为世界上所有的人本应共同使用所有的事物；但由于不义，一个人说这个是自己的，另一个人说那个是他的，于是人与人之间就产生了分裂。总之，有一位最睿智的希腊人知道事情是这样的，他说朋友们应该拥有一切共同之物。当然，在\Emph{一切事物}中无可争议地包括了妻子。他还说，正如空气和阳光不能被分割，世上被赐予所有人共同拥有的其他事物也不应被分割，而应如此拥有。但我想说这些，是因为我渴望转向行善，而除非我先知道什么是善，否则我不能行善；如果我能理解善恶，我就会由此认识到什么是恶，即与善相对。}\par

% structure: p0012 source=h2 target=subsection reason=relative-heading-rank; base=h2

\subsection{第六章 伯多禄的权威}

\Quote{但我希望你们中的一位，而不是\Person{伯多禄}，来回答我所说的话；因为从他那里接受话语和教导时不宜提出问题；而当他就某一主题做出裁决时，那应被持守而不应再反驳。因此，让我们把他作为仲裁者；这样，如果我们的讨论在任何时候没有得出结果，他可以宣布他认为正确的东西，从而为可疑的问题提供一个无可置疑的结束。因此，现在我可以相信，只凭他一个人的意见就满足，如果他表达了任何意见的话；而这将是我最终要做的。然而，我首先想看看是否有可能通过讨论找到所寻求的东西。因此，我希望\Person{克莱孟}首先开始，并表明在实体或行为中是否有任何善或恶。}\par

% structure: p0014 source=h2 target=subsection reason=relative-heading-rank; base=h2

\subsection{第七章 克莱孟的论证}

我回答说：\Quote{既然你确实希望从我这里知道，在自然或行为中是否有任何善或恶，或者是否并非人们由于习俗的偏见，认为某些事物是善的，另一些是恶的，因为他们已经把共同的事物（如你所说是应当像空气和阳光一样共同使用的）作了划分；我认为，我不应该从其他任何地方向你们引述任何陈述，而应从你们所精通并支持的那些研究中引述，以便我所说的你们会毫不迟疑地接受。你们给所有元素和天体指定了某些界限，你们说，这些界限在有些情况下无害地交汇，如婚姻；但在另一些情况下则有害地结合，如奸淫。你们说有些事物是所有人都普遍的，但另一些事物不属于所有人，也不是普遍的。但为了不进行冗长的讨论，我将简要地谈谈这个问题。干燥的大地需要添加和混合水，以便能够产生果实，没有果实人无法生存：因此这是一种合法的结合。相反，如果霜的寒冷与大地混合，或热与水混合，这种结合产生腐败；而在这类事物中，这就是奸淫。}\par

% structure: p0016 source=h2 target=subsection reason=relative-heading-rank; base=h2

\subsection{第八章 公认的恶}

然后我的父亲回答说：\Quote{但是正如不和谐的星辰或元素的结合的有害性立即暴露出来，奸淫也应当立即被表明是一种恶。}然后我说：\Quote{首先告诉我这一点，你是否像你自己所承认的那样，从不协调和不和谐的结合中产生恶？然后在那之后我们再调查另一件事。}然后我的父亲说：\Quote{事物之性正如你所说，我的儿子。}然后我回答说：\Quote{既然你希望学习这些东西，看看有多少事物没有人怀疑是恶。你认为热病、火灾、叛乱、房屋倒塌、谋杀、监禁、拷问、痛苦、哀悼以及诸如此类的事物是恶吗？}然后我的父亲说：\Quote{不错，我的儿子，这些事物是恶，而且是很大的恶；无论如何，谁否认它们是恶，就让他去承受它们吧！}\par

% structure: p0018 source=h2 target=subsection reason=relative-heading-rank; base=h2

\subsection{第九章 根据占星学原理恶的存在}

于是我回答说：\Quote{既然我要与精通占星学的人打交道，我将根据那门学问来与你讨论这件事，使你从我以你所熟悉的事物为出发点的方法中，更能欣然接受。现在请听：你承认我们提到的那类事情是恶事，例如热病、火灾等。按照你的说法，这些是由凶星造成的，如湿润的土星和炽热的火星；而与之相反的事物则由吉星造成，如温和的木星和湿润的金星。难道不是这样吗？}我的父亲回答说：\Quote{是的，我的儿子；不可能不是这样。}然后我说：\Quote{既然你说好事由吉星——例如木星和金星——造成，那么让我们看看当任何一颗凶星与吉星混合时会产生什么，并让我们明白那是恶的。因为你规定金星负责婚姻，如果它与木星成一定相位，它便使婚姻贞洁；但如果木星不在观看，而火星在场，那么你就宣称婚姻因奸淫而败坏。}我的父亲说：\Quote{正是如此。}我回答说：\Quote{因此奸淫是恶的，因为它是由凶星的混合而产生的；简而言之，所有你所说的吉星因与凶星混合而遭受的事情，无疑都应被宣称为恶的。因此，那些我们说过产生热病、火灾及其他类似恶事的星星——按照你的说法——也产生谋杀、奸淫、偷盗，并产生傲慢和愚钝的人。}\par

% structure: p0020 source=h2 target=subsection reason=relative-heading-rank; base=h2

\subsection{第十章 如何进步}

我的父亲说：\Quote{你确实简洁而无比地表明了在行为上有恶的存在；但我仍希望知道这一点：如果\InnerQuote{\OriginalTerm{命宫}{Genesis}}迫使人犯罪，天主如何公正地审判那些犯罪的人？}我回答说：\Quote{我不敢对您说什么，父亲，因为我理应尊敬您；否则我本有一个答案给您，如果合适的话。}我的父亲说：\Quote{我的儿子，请说出你想到的；因为不是你，而是探究的方法在犯错误，就像一个贞洁的妇人对一个不节制的男子，如果她为自己的贞洁和荣誉而愤慨。}我回答说：\Quote{如果我们不坚持我们已经承认和宣信的原则，而是那些已经被界定的东西总被遗忘所松动，我们就会像是在编织\Person{佩涅洛佩}的织物，拆解我们已经做成的东西。因此，在我们还没有勤勉地审视所提出的教理之前，我们不应该轻易地赞同；或者一旦赞同，并且命题已被同意，那么我们就应该遵守已经决定的东西，以便我们可以继续对其他事情进行探究。}我的父亲说：\Quote{你说得好，我的儿子；我知道你为何说这话：因为在昨天关于自然原因的讨论中，你表明某种邪恶的力量，将自己转移到星宿的秩序中，激发人的欲望，以各种方式挑动他们犯罪，却并不强迫或产生犯罪。}我对此回答说：\Quote{您记得很好；然而，虽然您记得，您却陷入了错误。}我的父亲说：\Quote{请原谅我，我的儿子；因为我对这些事情还不够熟练：事实上，您昨天的讲论因着真理迫使我同意您；但在我的意识中，仍有几分热病的残余，这些残余使我暂时远离信德，就如远离健康一样。因为我感到困扰，因为我知道许多事情，甚至几乎所有的事情，都是按照\OriginalTerm{命宫}{Genesis}发生在我的身上。}\par

% structure: p0022 source=h2 target=subsection reason=relative-heading-rank; base=h2

\subsection{第十一章 占星学的试验}

于是我回答说：因此我要告诉您，父亲，数学占星的特性是什么，请您按我所说的去做。您去一位占数学家那里，首先告诉他，在某个时候，如此这般的祸事降临到您身上，并且您想从他那里知道，这些祸事从何处、如何、或透过哪些星辰降临到您身上。他无疑会回答您，凶恶的火星或土星统治了您的那段时期，或者其中某颗星是周期性的，或者某颗星以直径相对、合相或中心的方式注视您；或者其他类似的回答，他会加上说，所有这些之中，某颗星与那颗凶星不和谐，或者不可见，或者位于图形中，或者超出分界，或者被遮蔽，或者没有接触，或者处于昏暗的星辰之中；他会根据自己的理由给出许多其他类似的说法，并详述具体细节。之后，您去另一位占数学家那里，告诉他相反的情况，即在那时候有如此这般的好事发生在您身上，提到同一时间，并问他，这些好事从您命宫的哪些部分而来，并且注意，正如我所说，这些时间应与您询问祸事的时间相同。当您在这些时间上欺骗了他之后，看他将为您编造出什么图形，以表明在那同一时刻本应有好事降临到您身上。因为对于处理人命宫的人来说，不可能不在天体的每一个所谓\InnerQuote{\OriginalTerm{象限}{quadrant}}中找到一些星辰位置有利，一些位置不利；因为根据数学占星，圆环在每个部分都是同样完整的，允许各种不同的原因，他们可以借此机会说出他们喜欢的任何说法。\par

% structure: p0024 source=h2 target=subsection reason=relative-heading-rank; base=h2

\subsection{第十二章 占星学被自由意志挫败}

因为，正如通常发生的那样，当人们看见不吉利的梦，却无法从中得出任何确切结论时，一旦某事发生，他们就把梦中所见与所发生之事联系起来；占星术也是如此。因为在任何事情发生之前，没有任何确切的宣告；但在某事发生后，他们就收集事件的原因。因此，常常当他们出错，事情的结果与预测不符时，他们便归咎于自己，说是某颗星相克，而他们没有看到；他们不知道，自己的错误并非源于对技艺的不娴熟，而是源于整个体系的不一致性。因为他们不知道那些我们确实渴望去做，但在其中我们并不放纵自己欲望的事情是什么。但我们这些学习了这奥秘之理的人，知道其原因，因为我们拥有\OriginalTerm{自由意志}{freedom of will}，有时抵制我们的欲望，有时又向它们屈服。因此，人类行为的结果是不确定的，因为它取决于自由意志。因为占星家确实可以指出恶意力量所产生的欲望；但是，这个欲望的行动或结果是否会实现，在事情完成之前无人知晓，因为它取决于自由意志。这就是为何无知的占星家为自己发明了关于\OriginalTerm{关键期}{climacterics}的说法，作为在不确定中的避难所，正如我们昨天已经充分展示的那样。\par

% structure: p0026 source=h2 target=subsection reason=relative-heading-rank; base=h2

\subsection{第十三章　接纳众人}

\Quote{如果你有什么要对此说的，请说吧。}然后我的父亲说：\Quote{没有比你所陈述的更真实的了，我的儿子。}当我们彼此这样说着的时候，有人告诉我们，外面站着一大群人，他们聚集起来是为了听讲。于是伯多禄命令让他们进来，因为那地方宽敞而便利。他们进来后，伯多禄对我们说：\Quote{如果你们中有人愿意，让他对众人讲话，讨论关于\OriginalTerm{偶像崇拜}{idolatry}的事。}我克莱孟回答说：\Quote{您极大的仁慈、温和以及对众人的耐心鼓励了我们，使我们敢在您面前说话，并询问我们喜欢的事；因此，正如我所说，您性情的温和吸引并鼓励所有人接受救恩教义的训导。这是我以前从未在任何其他人身上见过的，只有在您身上才有，您既无嫉妒也无愤慨。您认为如何？}\par

% structure: p0028 source=h2 target=subsection reason=relative-heading-rank; base=h2

\subsection{第十四章　无人具有普世知识}

然后伯多禄说：\Quote{这些事不仅来自嫉妒或愤慨；有时在某些人身上还有一种羞怯，唯恐他们可能无法完全回答可能提出的问题，于是他们避免暴露自己的不熟练。但没有人应当为此感到羞耻，因为没有人应该宣称自己知道一切；因为只有一位知道一切，就是那创造万有的。如果我们自己的主尚且宣称不知道那日子和时辰，尽管祂预言了它的征兆，并将一切完全交托给圣父，我们怎能承认对某些事无知是丢脸的呢？因为我们在这件事上有主的榜样？但我们只宣称知道我们从\OriginalTerm{真先知}{true Prophet}所学到的事；而这些事是真先知所交付给我们的，祂判断这些对于人类知识是足够的。}\par

% structure: p0030 source=h2 target=subsection reason=relative-heading-rank; base=h2

\subsection{第十五章　克莱孟的揭示}

于是我克莱孟接着这样说道：\Quote{在\Place{特里波利斯}，当您与外邦人辩论时，我的主伯多禄，我对您大为惊叹，尽管您按照希伯来人的方式和您自己法律的规矩受教于您的父亲，从未受过希腊学问研究的玷污，却如此卓越且无与伦比地辩论；您甚至提及了关于诸神历史的一些事情，这些通常在剧场中被公开宣讲。但我看出这些神话和亵渎的话您并不十分熟悉，我将从最初开始，在您面前讲述它们，如果您愿意的话。}然后伯多禄说：\Quote{请说吧；你做得好，协助我的宣讲。}于是我说：\Quote{因此，我将说话，因为您命令我，不是以教导您的方式，而是公开外邦人对诸神所持有的愚蠢意见。}\par

% structure: p0032 source=h2 target=subsection reason=relative-heading-rank; base=h2

\subsection{第十六章　\Quote{惟愿天主的全体人民都是先知}}

但当我正要说话时，\Person{尼塞塔}咬着嘴唇，示意我安静。伯多禄看见他，就说：\Quote{你为什么想要压制他慷慨的性情和高贵的本性，要他为了我的荣誉（这算不得什么）而保持沉默？难道你不知道，如果所有民族听了我传讲的真理并相信后，都投身于教导，那就会为我赢得更大的荣耀，如果你真以为我贪求荣耀的话？因为有什么比预备门徒归于基督更荣耀呢？这些门徒不是保持沉默、只为自己得救，而是讲说所学之事，造福他人？我确实希望你们俩，你尼塞塔和你亲爱的\Person{阿奎拉}，能帮助我宣讲天主的圣言，特别是因为外邦人所犯的错误是你们所熟知的；不仅如此，所有听我讲的人，正如我所说，我希望他们如此听和学，以致也能教导：因为世界需要许多帮助者，通过他们，人们可以从错误中被唤醒。}他说完这话后，对我说：\Quote{那么，克莱孟，继续你已经开始的话题吧。}\par

% structure: p0034 source=h2 target=subsection reason=relative-heading-rank; base=h2

\subsection{第十七章　外邦人的宇宙起源论}

\Quote{我立刻回答说：当你在\Place{特里波利斯}辩论时，如我所说，你关于外邦人的神祇谈了很多有益且令人信服的内容，我愿在你面前陈述关于他们起源的可笑传说，既使你不至于对这种虚妄迷信的谎言一无所知，也使在场的听众知道他们错误的可耻性质。那么，外邦人中的智者说，最先存在的是混沌；混沌经过很长一段时间使其外部凝固，为自己设定了界限和某种基础，仿佛是凝聚成一个巨大蛋的形状和样式；在蛋内，经过长时间，如同在蛋壳内，孕育并活化了某种动物；后来，那个巨大的球体破裂，产生了一种双性的人，他们称为\OriginalTerm{阴阳同体}{masculo-feminine}。他们称此为\OriginalTerm{法涅塔斯}{Phanetas}，意为显现，因为当它显现时，据说光也照耀出来。他们说，由此产生了实体、审慎、运动与交合，由此又产生了天与地。他们又说，从天产生了六个男性，称为提坦；同样，从地产生了六个女性，称为提坦尼得斯。这些从天而生的男性的名字是：\Person{俄刻阿诺斯}、\Person{科俄斯}、\Person{克瑞俄斯}、\Person{许珀里翁}、\Person{伊阿珀托斯}、\Person{克洛诺斯}（在我们中称为\Person{萨图尔努斯}）。同样，从地而生的女性的名字是：\Person{忒亚}、\Person{瑞亚}、\Person{忒弥斯}、\Person{谟涅摩绪涅}、\Person{泰西斯}、\Person{赫柏}。}\par

% structure: p0036 source=h2 target=subsection reason=relative-heading-rank; base=h2

\subsection{第十八章 萨图尔努斯的家族}

\Quote{在这些当中，天的长子娶了地的长女；次子娶了次女，其余以此类推。因此，娶了第一个女性的第一个男性因她的缘故而下降；但第二个女性因所嫁之人的缘故而上升；如此各按顺序，停留在婚配抽签所分到的位置。他们说，从他们的交合中产生了无数其他者。但在六个男性中，那个叫萨图尔努斯的娶了瑞亚，并受到一个神谕的警告：凡他所生的儿子将比他更强大，并将他逐出王国，于是他决定吞吃所有将生的儿子。首先，他生了一个儿子叫阿得斯，在我们中称为奥尔库斯；为了上述原因，他抓住并吞吃了。之后他生了第二个儿子，叫涅普顿；同样吞吃了。最后，他生了他们称为朱庇特的；但他的母亲瑞亚怜悯他，在父亲将要吞吃他时，用计谋将他从父亲身边带走。首先，为了不让人注意到孩子的哭声，她让某些库瑞忒斯敲铙钹和鼓，以使震耳的声音掩盖婴儿的哭声。}\par

% structure: p0038 source=h2 target=subsection reason=relative-heading-rank; base=h2

\subsection{第十九章 他们的命运}

\Quote{但当他从瑞亚腹部的减小得知孩子已经出生时，他要求孩子以便吞吃；于是瑞亚给了他一块大石头，告诉他这就是她所生的。他接过并吞了下去；石头被吞下后，推挤并驱赶出他先前吞下的那些儿子。因此，奥尔库斯首先出来，下降并占据了下界，即冥界。第二个在他之上——他们称为涅普顿——被推向水面。第三个，借其母亲瑞亚的计谋得以幸存，她将他放在一只母山羊上，送上了天堂。}\par

% structure: p0040 source=h2 target=subsection reason=relative-heading-rank; base=h2

\subsection{第二十章 朱庇特的所作所为}

\Quote{但这些老妇的寓言和外邦人的族谱已经够了；因为如果我要详述所有他们称为神的那些人的世代及其邪恶行为，那将没完没了。但作为例子，我略过其他，只详述他们视为最伟大和首要的那一位的恶行，他们称他为朱庇特。因为据说他拥有天堂，作为高于其他者；他长大后，娶了自己的妹妹，他们称为朱诺，在此他马上变得像野兽。朱诺生了伏尔甘；但据他们所说，朱庇特并非他的父亲。然而，朱庇特本人使她成为美狄亚的母亲；朱庇特得到神谕：凡她所生之子将比他更强大，并将他逐出王国，于是他抓住并吞吃了她。之后，朱庇特从他的头脑中生出密涅瓦，从他的大腿生出巴克斯。此后，当他爱上忒提斯时，据说普罗米修斯告诉他，如果他与她同寝，她所生之子将比父亲更强大；出于对此的恐惧，他把她嫁给了珀琉斯。随后，他与珀耳塞福涅同寝，她是朱庇特与刻瑞斯所生的女儿，与她生了狄俄尼索斯，后者被提坦撕碎。但据说，他想起也许他自己的父亲萨图尔努斯可能生另一个儿子，比他更强大，并可能将他逐出王国，于是他带着他的兄弟提坦与父亲开战；击败他们后，他最终将父亲投入监狱，割下他的生殖器，扔进海里。但从伤口流出的血，与海浪混合，并被不断搅动变成泡沫，产生了他们称为阿佛洛狄忒的那位，在我们中称为维纳斯。从他与这位自己的妹妹交合，据说这个朱庇特生了库普里斯，他们说库普里斯是丘比特的母亲。}\par

% structure: p0042 source=h2 target=subsection reason=relative-heading-rank; base=h2

\subsection{第二十一章 黑暗清单}

此乃其乱伦之行；今将述其奸淫之恶。彼污穢\Person{Oceanus}之妻\Person{Europa}，所生\Person{Dodonæus}；\Person{Pandion}之妻\Person{Helen}，所生\Person{Musæus}；\Person{Asopus}之妻\Person{Eurynome}，所生\Person{Ogygias}；\Person{Oceanus}之妻\Person{Hermione}，所生\OriginalTerm{美惠三女神}{Graces}：\Person{Thalia}、\Person{Euphrosyne}、\Person{Aglaia}；其妹\Person{Themis}，所生\OriginalTerm{时序三女神}{Hours}：\Person{Eurynomia}、\Person{Dice}、\Person{Irene}；\Person{Inachus}之女\Person{Themisto}，所生\Person{Arcas}；\Person{Minos}之女\Person{Idæa}，所生\Person{Asterion}；\Person{Alphion}之女\Person{Phœnissa}，所生\Person{Endymion}；\Person{Inachus}之女\Person{Io}，所生\Person{Epaphus}；\Person{Danaus}之女\Person{Hippodamia}与\Person{Isione}，其中\Person{Hippodamia}为\Person{Olenus}之妻，\Person{Isione}为\Person{Orchomenus}或\Person{Chryses}之妻；\Person{Phœnix}之女\Person{Carme}，所生\Person{Britomartis}，此乃\Person{Diana}之侍女；\Person{Lycaon}之女\Person{Callisto}，所生\Person{Orcas}；\Person{Munantius}之女\Person{Lybee}，所生\Person{Belus}；\Person{Latona}，所生\Person{Apollo}与\Person{Diana}；\Person{Eurymedon}之女\Person{Leandia}，所生\Person{Coron}；\Person{Evenus}之女\Person{Lysithea}，所生\Person{Helenus}；\Person{Bellerophon}之女\Person{Hippodamia}，所生\Person{Sarpedon}；\Person{Macarius}之女\Person{Megaclite}，所生\Person{Thebe}与\Person{Locrus}；\Person{Phoroneus}之女\Person{Niobe}，所生\Person{Argus}与\Person{Pelasgus}；\Person{Neoptolemus}之女\Person{Olympias}，所生\Person{Alexander}；\Person{Prometheus}之女\Person{Pyrrha}，所生\Person{Helmetheus}；\Person{Deucalion}之女\Person{Protogenia}与\Person{Pandora}，彼生\Person{Æthelius}、\Person{Dorus}、\Person{Melera}与\Person{Pandorus}；\Person{Proteus}之女\Person{Thaicrucia}，所生\Person{Nympheus}；\Person{Asopus}之女\Person{Salamis}，所生\Person{Saracon}；\Person{Atlas}之女\Person{Taygete}、\Person{Electra}、\Person{Maia}、\Person{Plutide}，彼分别生\Person{Lacedæmon}、\Person{Dardanus}、\Person{Mercury}、\Person{Tantalus}；\Person{Phoroneus}之女\Person{Phthia}，彼生\Person{Achæus}；\Person{Aramnus}之女\Person{Chonia}，彼生\Person{Lacon}；\Person{Chalcea}（一宁芙），所生\Person{Olympus}；\Person{Charidia}（一宁芙），所生\Person{Alcanus}；\Person{Chloris}（\Person{Ampycus}之妻），所生\Person{Mopsus}；\Person{Lesbus}之女\Person{Cotonia}，所生\Person{Polymedes}；\Person{Anicetus}之女\Person{Hippodamia}；\Person{Peneus}之女\Person{Chrysogenia}，所生\Person{Thissæus}。\par

% structure: p0044 source=h2 target=subsection reason=relative-heading-rank; base=h2

\subsection{第二十二章 朱庇特之卑劣变形}

彼尚有无数奸淫之事而未产子嗣，不一而足。然于前述诸人之中，彼有时变形如魔术师而行淫。简言之，彼化为萨梯尔诱惑\Person{Nycteus}之女\Person{Antiope}，所生\Person{Amphion}与\Person{Zethus}；化为其夫\Person{Amphitryon}诱\Person{Alcmene}，所生\Person{Hercules}；化为鹰诱惑\Person{Asopus}之女\Person{Ægina}，所生\Person{Æacus}。亦化为鹰奸污\Person{Dardanus}之子\Person{Ganymede}；化为熊奸污\Person{Phocus}之女\Person{Manthea}，所生\Person{Arctos}；化为黄金奸污\Person{Acrisius}之女\Person{Danæ}，所生\Person{Perseus}；化为公牛奸污\Person{Phœnix}之女\Person{Europa}，所生\Person{Minos}、\Person{Rhadamanthus}与\Person{Sarpedon}；化为蚂蚁奸污\Person{Achelaus}之女\Person{Eurymedusa}，所生\Person{Myrmidon}；化为秃鹫奸污宁芙\Person{Thalia}，所生\Person{Palisci}（在西西里）；化为阵雨奸污\Person{Geneanus}之女\Person{Imandra}（在罗得岛）；化为其夫\Person{Phœnix}奸污\Person{Cassiopeia}，所生\Person{Anchinos}；化为天鹅奸污\Person{Thestius}之女\Person{Leda}，所生\Person{Helen}；又化为星辰奸污同一人，所生\Person{Castor}与\Person{Pollux}；化为田凫奸污\Person{Lamia}；化为牧羊人奸污\Person{Mnemosyne}，所生九位缪斯；化为鹅奸污\Person{Nemesis}；化为火奸污卡德摩斯之女\Person{Semele}，所生\Person{Dionysius}。与其女\Person{Ceres}生\Person{Persephone}，后又化为龙奸污\Person{Persephone}。\par

% structure: p0046 source=h2 target=subsection reason=relative-heading-rank; base=h2

\subsection{第二十三章 何以称神？}

彼亦与其叔\Person{Oceanus}之妻\Person{Europa}及其妹\Person{Eurynome}通奸，并惩罚其父；与其子\Person{Atlas}之女\Person{Plute}通奸，并罚其所生\Person{Tantalus}；由\Person{Orchomenus}之女\Person{Larisse}生\Person{Tityon}，亦加惩罚；夺其子\Person{Ixion}之妻\Person{Dia}，并加永久惩罚；其奸生之子几尽遭暴死；其墓冢多为世人所知。此弑亲者自毁其叔父、污其妻，奸其姊妹，多形之魔术师，其墓竟显于克里特人中；彼等虽知且承认其可怖乱伦之行，且传述于众，然不以认其为神为耻。故我以为奇哉，超然奇哉：彼于恶行罪愆超越众人，竟获至圣至善之名，超乎万名之上，被称为诸神与人之父；若非那喜人恶者说服不幸之灵魂，使其将尊荣独归于彼见其于恶行超越众人者，以诱众人仿效其恶行，更有何由？\par

% structure: p0048 source=h2 target=subsection reason=relative-heading-rank; base=h2

\subsection{第二十四章 多神之愚}

然其子（被这些\Emph{外邦人}视为神者）之墓亦公然指出，一处一处各别：\Person{Mercury}之墓在\Place{Hermopolis}；塞浦路斯的\Person{Venus}之墓在\Place{塞浦路斯}；\Person{Mars}之墓在\Place{色雷斯}；\Person{Bacchus}之墓在\Place{底比斯}（传说彼处被撕碎）；\Person{Hercules}之墓在\Place{推罗}（彼处被火焚）；\Person{Æsculapius}之墓在\Place{Epidaurus}。凡此诸神，非仅被视为已死之人，且被视为因罪受罚之恶人；然愚人仍拜之为神。\par

% structure: p0050 source=h2 target=subsection reason=relative-heading-rank; base=h2

\subsection{第二十五章 死者封神}

但若他们争辩，并断言这些地方与其说是他们埋葬或死亡之地，不如说是他们的出生地，那么从近在眼前且依然新近的事迹，便可揭穿远古的事迹，因为我们已证明他们所敬拜的正是他们自己所承认曾为凡人、曾死去、甚至曾受惩罚者：如叙利亚人敬拜阿多尼斯，埃及人敬拜奥西里斯；特洛伊人敬拜赫克托耳；在琉科涅苏斯敬拜阿喀琉斯，在本都敬拜帕特罗克洛斯，在罗得岛敬拜马其顿的亚历山大；还有许多其他的人，在不同地方被敬拜，而他们毫不怀疑这些人都是死人。由此可见，他们的前人亦陷入同样错误，将死人尊为神明，这些死人或许曾拥有某种能力或技艺，尤其是曾以魔法幻象迷惑愚钝之人。\par

% structure: p0052 source=h2 target=subsection reason=relative-heading-rank; base=h2

\subsection{第二十六章　变形记}

\Quote{因此，如今还有一点：诗人也用辞藻的美饰来粉饰错误的虚妄，以言辞的甘甜说服人相信凡人变为了不朽；不仅如此，他们还宣称人变成了星辰、树木、动物、花卉、飞鸟、泉源与河流。除非显得浪费言辞，我几乎能一一列举他们断言由人变成的所有星辰、树木、泉源与河流；然而，为举例起见，我至少每类提及一个。他们说刻甫斯的女儿安德洛墨达变成了星辰；拉冬河神的女儿达佛涅变成了树木；阿波罗所爱的雅辛托斯变成了花朵；卡利斯托变成了他们称为大熊座的星群；普罗克涅、菲洛墨拉与忒柔斯变成了飞鸟；奇里乞亚的提斯柏化为了泉源，在同一地方的皮拉摩斯化为河流。他们断言几乎所有星辰、树木、泉源、河流、花卉、动物与飞鸟，都曾一度是人类。}\par

% structure: p0054 source=h2 target=subsection reason=relative-heading-rank; base=h2

\subsection{第二十七章　多神论者的矛盾}

但伯多禄听了这话后说：\Quote{按照他们所说，那么在人类变成星辰以及你所提及的其他事物之前，天上就没有星辰，地上就没有树木和动物；也没有泉源、河流和飞鸟。而没有这些，后来变成那些事物的人自己又如何生活呢？显然，没有这些东西，人不可能在地上生活。}于是我答道：但他们甚至不能始终如一地遵守对自己神祇的敬拜；因为他们所敬拜的每一位，都有某样东西献给他，敬拜者应当禁戒那些东西：如他们说橄榄献给了米涅瓦，母山羊献给了朱庇特，种子献给了刻瑞斯，酒献给了巴克斯，水献给了奥西里斯，公羊献给了哈蒙，牝鹿献给了狄安娜，鱼和鸽子献给了叙利亚人的魔鬼，火献给了武尔坎；对每一位，正如我所说，都有某样特别祝圣之物，敬拜者为尊敬那些祝圣的对象，必须禁戒之。但若一个人禁戒这一样，另一个人禁戒那一样，因尊崇一位神而招致所有其他神的愤怒；因此，若想平息所有神，就必须为尊敬所有神而禁戒一切事物，这样，在审判日之前便会被公义的判决自我定罪，因饥饿而惨死。\par

% structure: p0056 source=h2 target=subsection reason=relative-heading-rank; base=h2

\subsection{第二十八章　外教信仰的支柱}

\Quote{但让我们回到正题。有什么理由——更确切地说，是什么疯狂占据了人的心智——使他们敬拜并尊崇一位明知其为不敬、邪恶、亵渎——我指朱庇特——乱伦、弑亲、通奸的人为神，甚至在剧院的歌曲中公然宣扬这些事？而若因这些行为他配得上成为神，那么当他们听说任何杀人犯、通奸者、弑亲者、乱伦者时，也应当将他们敬奉为神。但我无法理解，为何他们在别人身上憎恶的事，却在他身上尊崇。}于是伯多禄答道：\Quote{既然你说你不能理解，那就从我这里学习他们为何在他身上尊崇邪恶。首先，是为了当他们自己行同样的事时，知道他们必蒙他悦纳，因为他们不过是在他的邪恶上效仿了他。其次，因为古人巧妙地将这些事写进著作，典雅地编织进诗句。如今，借助青年教育，因这些知识附着于他们柔嫩而单纯的心智，便难以将其剥离并抛弃。}\par

% structure: p0058 source=h2 target=subsection reason=relative-heading-rank; base=h2

\subsection{第二十九章　寓意解释}

伯多禄说完这些，尼刻塔答道：\Quote{不要以为，我的主伯多禄，外邦的博学之士没有某些似乎有理的论据，用以支持那些看似可责与可耻之事。我说这话，并非希望证实他们的错误（我绝不会有此念头）；然而我知道，在他们较聪明的人中，有某些辩护之词，他们惯于用它们来支持并粉饰那些看似荒谬之事。若您乐意，我就陈述其中一些——因我对此略知一二——我会遵命而行。}伯多禄允许后，尼刻塔便继续如下。\par

% structure: p0060 source=h2 target=subsection reason=relative-heading-rank; base=h2

\subsection{第三十章　俄耳甫斯的宇宙起源论}

在希腊人当中，所有关于远古起源的文献都基于多种权威，但主要是两位：\Person{俄耳甫斯}和\Person{赫西俄德}。他们的著作在意义上分为两部分——即字面义和寓意义；普通大众趋附字面义，但哲学家和学者们的全部口才却倾注于对寓意义的赞叹。因此，\Person{俄耳甫斯}说，起初有\OriginalTerm{混沌}{chaos}，永恒、无限、非受造，万物由它而生。他说这混沌既不是黑暗也不是光明，既非湿也非干，既非热也非冷，而是万物的混合，始终是一团未成形的质料；然而，终于，仿佛一个巨大的卵，它生出并产生了某种双重形式，这种形式经过漫长的时代才被塑造而成，他们称之为\OriginalTerm{阴阳合体}{masculo-feminine}，是由如此多样性相反混合而成的具体形式；这就是万物的本原，出自纯粹的物质，它在出现时实现了四元素的分离，并由前两种元素\Emph{火和气}构成了天，由另两种元素\Emph{土和水}构成了地；他们说，现在万物都是通过这些元素的相互参与而产生和生成的。这便是\Person{俄耳甫斯}所言。\par

% structure: p0062 source=h2 target=subsection reason=relative-heading-rank; base=h2

\subsection{第三十一章 赫西俄德的宇宙起源论。}

但\Person{赫西俄德}对此补充道，在混沌之后，天地立刻被造，由此他说产生了那十一位（有时他也说十二位），其中他列为六男五女。以下是男性之名：\Person{俄刻阿诺斯}、\Person{科俄斯}、\Person{克吕奥斯}、\Person{许珀里翁}、\Person{伊阿珀托斯}、\Person{克罗诺斯}（亦称\Person{萨图尔努斯}）。女性之名：\Person{忒亚}、\Person{瑞亚}、\Person{忒弥斯}、\Person{谟涅摩绪涅}、\Person{忒提斯}。他们用寓意解释这些名字。他们说数目是十一或十二：第一是自然本身，他们称之为\Person{瑞亚}，意即流淌；其余十个是它的属性，也称为性质；然而他们又加上第十二位，即\Person{克罗诺斯}，在我们这里称为\Person{萨图尔努斯}，他们视之为时间。因此，他们断言\Person{萨图尔努斯}和\Person{瑞亚}是时间和物质；当它们与湿和干、热和冷混合时，便产生了万物。\par

% structure: p0064 source=h2 target=subsection reason=relative-heading-rank; base=h2

\subsection{第三十二章 寓意诠释。}

因此，据说她（即\Person{瑞亚}，或自然）产生了一个气泡，这气泡已经聚集了很长时间；它从水中的精神逐渐聚集而膨胀，一度在物质表面被驱动（它从物质中出来如同出自子宫），因寒冷之严酷而硬化，并因冰块增加而不断增大，最后脱落而沉入深渊，被自身重量牵引而下到冥界；因为它变得不可见而被称为\Person{哈得斯}，也称为\Person{俄耳库斯}或\Person{普路托}。由于它从顶部沉到底部，便为湿元素凝聚流到一起腾出了空间；粗重部分即大地，因水退去而裸露。因此，他们说，水的这种自由（先前因气泡的存在而受制）在气泡达到最低处后被称为\Person{尼普顿}。此后，当冷元素因冰冷气泡的凝结而被吸入低层区域，干湿元素已经分离，此时再无阻碍，暖元素因自身的力和轻飘而冲向上层空气区域，靠风和风暴托起。因此，这场风暴——在希腊文中称为\OriginalTerm{}{\GreekText{καταιγίς}}——他们称之为\Emph{埃吉斯}，即母山羊；而上升到上层的火焰他们称之为\Person{朱庇特}；因此，他们说\Person{朱庇特}骑着一只母山羊登上了\Place{奥林匹斯山}。\par

% structure: p0066 source=h2 target=subsection reason=relative-heading-rank; base=h2

\subsection{第三十三章 朱庇特等寓意的诠释。}

希腊人之所以称\Person{朱庇特}，是因为他活着或赋予生命，而我们的人则因为他提供救助。因此，他们说，这就是活着的实体，置于上层区域，借着热的影响如同脑的旋卷将一切吸引到自身，并通过某种调和的适度来安排它们，据说从他的头部生出了智慧，他们称之为\Person{密涅瓦}，希腊人因其永生而称她为\OriginalTerm{}{\GreekText{᾿Αθήνη}}；因为万有之父借他的智慧创造了万物，故说她从头顶产生，出自万有首要之地，并且被描绘成通过元素的调节混合而塑造和装饰了整个世界。因此，印在物质上以便世界得以形成的形状，因受热力约束，据说由\Person{朱庇特}的能量维系。既然这些形状已经足够，无需添加新的东西，每一种事物都靠自身种子的产物来修复，故\Person{萨图尔努斯}的手被\Person{朱庇特}捆绑；因为，如我所说，时间不再从物质中产生任何新东西；而种子的温暖按种类恢复一切；\Person{瑞亚}（即流动物质的增长）不再有新的诞生。因此，他们将元素的首次分割称为\Person{萨图尔努斯}的阉割，因为他不能再产生世界了。\par

% structure: p0068 source=h2 target=subsection reason=relative-heading-rank; base=h2

\subsection{第三十四章 其他寓意诠释。}

\Quote{关于维纳斯，他们提出如下的寓意：他们说，当海洋被置于空气之下，天空的光辉从水面反射而更加愉悦地照耀时，事物的可爱——因水面而显得更加美丽——被称为维纳斯；\Emph{她}与空气结合，如同与\Emph{自己的}兄弟结合，以产生作为欲望对象的美，据说便生下了丘比特。因此，正如我们所说，他们教导说，克洛诺斯（即萨图尔努斯）寓意着时间；瑞亚是物质；哈得斯（即奥尔库斯）是地府的深渊；涅普顿是水；朱庇特是空气——即热元素；维纳斯是事物的可爱；丘比特是欲望，存在于万物之中，藉以繁衍后代，甚或是事物的理性，当它被智慧地审视时便带来愉悦。赫拉（即朱诺）据说是从天上降到地上的中间空气。他们称狄安娜为普罗瑟皮娜，并将下方的空气交给她。他们说阿波罗就是太阳本身，环绕天空运行；墨丘利是言语，万物的理由由此给出；马尔斯是无节制的火，吞噬一切。但为了不因逐一列举而耽搁你们，那些对这类事物有更深奥理解的人认为，通过将这种寓意应用于他们所崇拜的每一个对象，他们给出了公正合理的解释。}\par

% structure: p0070 source=h2 target=subsection reason=relative-heading-rank; base=h2

\subsection{第三十五章　这些寓言之无用}

尼塞塔这样说完后，阿奎拉回答道：无论谁是这些事物的作者和发明者，在我看来他都是非常不虔敬的，因为他掩盖了那些似乎愉快而体面的事物，使他的迷信仪式建立在卑劣可耻的习俗之上，因为按字面所写的那些事显然是不体面且卑劣的；他们宗教的全部遵守就在于此，即通过这些罪行和不虔敬，他们教导人们模仿他们所崇拜的神祇。因为这些寓言对他们有什么益处呢？尽管它们被构造成体面的，但对敬拜或道德改良毫无用处。\par

% structure: p0072 source=h2 target=subsection reason=relative-heading-rank; base=h2

\subsection{第三十六章　寓言乃事后之思}

因此，更为明显的是，谨慎的人们见到普通的迷信如此可耻、如此卑劣，却尚未学会任何修正之道或任何知识，便竭尽所能用论据和解释，以体面的言辞掩盖不体面的事，而不是像他们所说的，以不体面的寓言来掩盖体面的理由。因为倘若情况如此，他们的雕像和图画就绝不会\Emph{描绘}他们的恶习和罪行。那只与勒达通奸的天鹅不会被描绘出来，与欧罗巴通奸的公牛也不会；他们也不会把那位他们认为超越一切的神祇变成千种怪物形状。而且，的确，如果他们中间的伟大智慧之人知道这一切都是虚构而非真理，他们岂不会指控那些为损害神灵而展示此类图画或雕刻此类形象的人为不虔敬和亵渎吗？简言之，让他们把他们当代的国王绘成公牛、鹅、蚂蚁或秃鹫的样子，写上国王的名字，并将这样的雕像或形象立在公共场所，他们很快就会感受到自己行为的错误及其惩罚的严重性。\par

% structure: p0074 source=h2 target=subsection reason=relative-heading-rank; base=h2

\subsection{第三十七章　神如其人，敬拜者亦如是}

然而，既然那些由公共丑行所证实的事才是真实的，而掩饰只是由谨慎之人寻求和捏造出来，以便用体面的言辞为之辩解，因此它们不仅不被禁止，甚至就在神秘仪式中，也出现萨图尔努斯吞噬其子嗣的形象，以及被科律班忒斯的钹鼓声掩盖的孩童；至于萨图尔努斯被阉割之事，还有什么比他的敬拜者为了尊崇他们的神而遭受类似的悲惨命运——被阉割——更能证明其真实性呢？既然这些事是显而易见的，谁会如此缺乏理智，甚至如此迟钝，以致看不出那些关于不幸之神的事是真的，而更为不幸的敬拜者以身体的伤残阉割为其作证呢？\par

% structure: p0076 source=h2 target=subsection reason=relative-heading-rank; base=h2

\subsection{第三十八章　诗人的作品}

\Quote{但是，如果正如他们所说，这些如此值得称道且虔敬的事却由如此不光彩且不虔敬的仪式来施行，那么，无论是最初传下这些事的人，还是在这些事不幸被传下后仍坚持履行它们的人，无疑是亵渎者。我们又如何看待诗人的书卷呢？如果它们以卑劣的寓言玷污了神祇们光荣而虔敬的事迹，难道不该立刻将它们丢弃并投入火中，以免它们引诱尚处柔嫩之龄的男孩，以为众神之首朱庇特本人弑父弑母、与姐妹和女儿乱伦，甚至对男童行淫；维纳斯和马尔斯是通奸者，以及上面所提到的一切？我主伯多禄，您对此有何看法？}\par

% structure: p0078 source=h2 target=subsection reason=relative-heading-rank; base=h2

\subsection{第三十九章　一切都是最好的安排}

\Quote{亲爱的阿奎拉，你要确信，万事皆由天主眷顾安排，使得那与真理对立的原因不仅软弱无力，而且卑劣。因为如果错误的论断更强大、更像真理，任何被其欺骗的人就不易返回真理之道。即便如今，当如此多邪恶可耻的事被述及外教人的神祇时，几乎没有人离弃那卑劣的错误，更何况其中若有什么体面而像真理的事呢？因为心灵很难从幼年时浸染的那些事物中转移；正是因此，如我所说，是天主眷顾促成了这事，使错误的实质既软弱又卑劣。但其他一切事物，天主眷顾也适当而有益地安排，尽管神圣眷顾的方式——因善且最佳——对不明事物原因的我们并不清晰。}\par

% structure: p0080 source=h2 target=subsection reason=relative-heading-rank; base=h2

\subsection{第四十章 进一步探求知识}

伯多禄这样说时，我克莱孟请求尼塞塔为了教导的缘故，将他曾仔细研究过的关于外邦人寓言的某些事情解释给我们听；我说：\Quote{因为与外邦人辩论时，对这些事有所了解是有益的。}于是尼塞塔说：\Quote{如果我的主伯多禄允许，我可以按你的请求去做。}伯多禄说：\Quote{今天我已准许你发言反驳外邦人，如你所知。}尼塞塔说：\Quote{那么，克莱孟，告诉我你要我说什么。}我对他说：\Quote{请告诉我们，外邦人是如何讲述关于诸神的飨宴之事，即他们在珀琉斯与忒提斯婚礼上所举行的盛宴。他们如何看待牧羊人帕里斯，以及三位女神朱诺、弥涅耳瓦和维纳斯之间他担任裁判的事情？墨丘利又代表什么？苹果呢？以及其他按顺序出现的事物？}\par

% structure: p0082 source=h2 target=subsection reason=relative-heading-rank; base=h2

\subsection{第四十一章 神话的解释}

尼塞塔说：\Quote{诸神飨宴之事解释如下。他们说宴席代表世界，众神就座的次序代表天体的位置。赫西俄德称之为天地最初所生的六男六女，他们把这十二位对应于环绕世界的十二宫。他们说宴席上的菜肴是事物的理由和原因，甘美而令人向往，它们以星座位置和星辰运行推论的形式，解释世界如何被统治和管理。然而他们说这些事物以宴席的自由方式存在，因为每个人的心智都有选择是否品尝此类知识，或是否戒除的自由；如同在宴席上无人被迫，人人随意饮食，同样，哲学化的方式取决于意志的选择。他们说纷争是肉体的欲望，它反抗心智的意图，阻碍对哲学化的渴望；因此他们说那段时间便是举行婚礼的时刻。这样，他们使珀琉斯与仙女忒提斯代表干燥与潮湿的元素，二者的混合构成了身体的实体。他们认为墨丘利是言语，借以将教导传达给心智；朱诺是贞洁，弥涅耳瓦是勇气，维纳斯是情欲，帕里斯是理解力。因此，他们说，如果一个人有着粗野、未开化的理解力，不懂得正确的判断，他就会藐视贞洁和勇气，将奖品即苹果交给情欲；由此，不仅他本人，而且他的同胞和整个种族都将招致毁灭与败亡。这些事，他们可以随心所欲地从任何题材中编造；然而它们可以适用于每一个人；因为如果有人有着牧羊人般的、粗野而未受教的理解力，且不愿受教导，当身体的欲火向他暗示情欲的享乐时，他立刻藐视学问的德行和知识的福乐，转而追求肉体的快乐。由此，无法平息的战争兴起，城池被毁，国土沦丧，正如帕里斯通过诱拐海伦，武装了希腊人和野蛮人，使他们互相毁灭。}\par

% structure: p0084 source=h2 target=subsection reason=relative-heading-rank; base=h2

\subsection{第四十二章 圣经的解释}

伯多禄称赞了他的陈述，说：\Quote{据我看来，聪明的人从他们阅读的事物中提取了许多貌似真实的道理；因此必须十分谨慎，当诵读天主的法律时，不应按我们自己心智的理解去诵读。因为神圣的圣经中有许多话可以被引向各人预先设想的含义；这是不应该做的。因为你不可寻求一种外来的、从外部带来的、你可用圣经权威加以证实的含义，而应从圣经本身获取真理的含义；因此，你应当从那个按照从祖先传下来的真理保存圣经含义的人学习圣经的意义，以便他能够权威地宣布他所正确领受的。但当一个人从圣经接受了完整而稳固的真理准则后，如果他从普通教育和自由学科中，也许是他少年时曾投入其中的，为确立真道而贡献一些东西，也并非不当；不过，当他认识了真理，就应弃绝虚谎和伪装。}\par

% structure: p0086 source=h2 target=subsection reason=relative-heading-rank; base=h2

\subsection{第四十三章 劝勉之言}

他说完这话，就望着我们的父亲说：\Quote{因此，老人家，如果你实在关心你灵魂的得救，希望当你愿意与身体分离时，因这短暂的归化而找到永远的安息，就请问你任何想了解的事，并寻求建议，以便你能够摆脱心中尚存的任何疑惑。因为即使是年轻人，生命的期限也不确定；但对老年人来说，甚至不确定也说不上，因为毫无疑问他们余下的时间不多了。因此，无论老少都应当非常认真地对待自己的归化和悔改，并致力于用最宝贵的装饰来修饰他们的灵魂，以迎接未来，比如真理的教义、贞洁的恩宠、正义的光辉、虔敬的美丽，以及一切一个有理智的心灵应当装饰的事物。此外，他们还应断绝与不当和不忠的同伴来往，与信徒为伴，经常参加那些探讨贞洁、正义和虔敬主题的集会；衷心地常向天主祈祷，求祂赐予那些应向天主祈求的事物；感谢祂；真诚地为他们过去的行为忏悔；在某种程度上，如果可能，也通过向穷人施行仁慈之举来帮助他们的忏悔：因为通过这些方式，赦免将更易获得，仁慈将更快地施予怜悯人的人。}\par

% structure: p0088 source=h2 target=subsection reason=relative-heading-rank; base=h2

\subsection{第四十四章 热忱}

但如果前来悔改的人年事已高，他更应当感谢天主，因为在认识了真理之后，肉欲的全部暴力已被打破，他不再需要战斗来压制身体反抗心灵的快感。因此，他应当操练于真理的学习和仁慈的事工，结出与悔改相称的果实；并且不要以为皈依的证明在于时间的长短，而在于虔诚与决心的强度。因为心灵在天主面前是显明的；祂不计算时间，而重视内心。祂赞许那些一听到真理的宣讲，不迟延、不虚度光阴，而是立刻——可以说就在同一时刻——憎恶过去，开始向往未来，并以天国的爱燃烧的人。\par

% structure: p0090 source=h2 target=subsection reason=relative-heading-rank; base=h2

\subsection{四十五、众人都当悔改。}

所以，你们中谁也不可再拖延或回顾，而应自愿地趋近天国的福音。不要让穷人说：\Quote{等我富足了，那时我便悔改。}天主并不向你索求金钱，而是要求一颗仁慈的心和一个虔诚的心灵。也不让富人因世俗的挂虑而延误他的悔改，不要思想如何处置丰盛的果实；内心不要说：\Quote{我该怎么办？我把果实放在哪里？}也不要对他的灵魂说：\Quote{你存有许多财物，够多年之用；吃喝快乐吧！}因为有人要对他说：\Quote{愚昧人！今夜就要索回你的灵魂，你所预备的一切将归谁呢？}因此，每个年龄、每个性别、每个身份的人都应速速悔改，以获永生。青年应当感恩，因为他们在情欲最猛烈之时，就已将颈项置于纪律的轭下。老年人自己也值得称赞，因为他们因敬畏天主而改变了长久不幸陷于其中的积习。\par

% structure: p0092 source=h2 target=subsection reason=relative-heading-rank; base=h2

\subsection{四十六、确凿的预言}

因此，谁也不可拖延，谁也不可迟延。行善有什么理由可以耽搁呢？或者你害怕，当你行善之后，得不到你所期望的赏报吗？即使没有赏报，你行善又有什么损失呢？光凭良心不就够了吗？但如果你如愿以偿，岂非以小博大，以暂换永？我说这话是为不信的人。因为我们所宣讲的，正如我们所宣讲的那样；因为不可能另有别样，因为那是藉着预言的话所应许的。\par

% structure: p0094 source=h2 target=subsection reason=relative-heading-rank; base=h2

\subsection{四十七、\Quote{这话是确实的，值得完全接纳。}}

如果有人想精确地了解我们宣讲的真理，就让他来听，让他认清什么是真先知；然后，一切疑惑都将离他而去，除非他以顽固的心态抗拒他所发现为真实的事。因为有些人的唯一目的就是无论如何要取得胜利，并为此而寻求赞美，而非寻求自己的救恩。这些人连一句话也不应当对他们说，惟恐崇高的话语受损害，并使损害它的人被判永远死亡。因为在哪一点上有人可以反对我们的宣讲呢？或者，我们的宣讲在何处与真实和可敬的信念相悖呢？它宣称：应当恭敬天主圣父，万有的创造者，以及祂的圣子——只有祂认识天父和祂的旨意，只有祂值得相信祂所命令的一切。因为只有祂是法律、立法者和公义的审判者；祂的法律吩咐要以清醒、贞洁、公义和仁慈的生活恭敬万有的主天主，并将一切希望唯独寄托于祂。\par

% structure: p0096 source=h2 target=subsection reason=relative-heading-rank; base=h2

\subsection{四十八、哲学家的错误}

但有人会说，哲学家们也给出了这类训诫。并非如此：因为他们确实制定了关于公义和节制的诫命，但他们不知道天主是善恶之事的赏报者；因此，他们的法律和训诫只能逃避公开的控告者，却不能净化良心。因为一个人既然不知道有隐秘之事的见证者和审判者，为何要害怕在暗中犯罪呢？此外，哲学家在他们的训诫中甚至加上应当恭敬众神——即魔鬼——即便在其他方面似乎值得赞许，仅此一点就足以定他们犯有最可怕的亵渎之罪，并以自己的判断定他们的罪，因为他们虽然宣称只有一位天主，却命令崇拜许多神，以此迁就人类的错误。而且，哲学家们说天主不发怒，却不知他们在说什么。因为忿怒若是扰乱心灵、使其失去正确判断，便是邪恶的。但惩罚恶人的忿怒并不给心灵带来纷乱；可以说，它是同一个情感，将赏报分配给善者，将惩罚分配给恶者；因为如果祂赐福给善人和恶人，将同等的赏报赋予虔诚者和不虔者，那么祂看起来不是良善的，而是不公义的。\par

% structure: p0098 source=h2 target=subsection reason=relative-heading-rank; base=h2

\subsection{四十九、天主的忍耐}

但你说：\Quote{天主也不应该作恶。}你说得对，祂也不作恶。但那些被祂所造的人，因为不相信自己将受审判，放纵享乐，已从虔诚和正义中堕落。但你会说：\Quote{如果惩罚恶人是正确的，那么他们一作恶就应该立即受惩罚。}你急切地想要这样，固然很好；但那位永恒者、没有任何事能向祂隐藏的，因祂是无终的，同样祂的忍耐也是无限的；祂不看重报复的快慢，而看重救恩的缘由。因为祂喜悦罪人悔改，胜过喜悦他死亡。\BibleRef{则 18:33} 因此，简言之，祂赐予人类圣洗圣事，如果有人速速接受它，并从此保持无污点，那么他在无知期间所犯的一切罪都将被涂抹。\par

% structure: p0100 source=h2 target=subsection reason=relative-heading-rank; base=h2

\subsection{五十、哲学家并非人类的恩人}

哲学家对人类生活有何贡献呢？他们说天主不对人发怒，这不过是教导人不要畏惧任何惩罚或审判，从而消除对罪人的一切约束。那些说没有天主，一切皆由偶然和意外发生的人，对人类有何益处呢？无非是使人听闻后，以为没有审判者，没有万物的守护者，便毫无畏惧地冲向一切由愤怒、贪婪或情欲所驱使的行为。确实，那些说没有什么是可以脱离\OriginalTerm{命数}{Genesis}而成就的人，对人类生活大有裨益——也就是说，每个人都可将自己犯罪的缘由归于命数，从而在犯罪中宣称自己无罪，不以悔改洗刷罪过，反而将责任推给命运，使罪过加倍。还有那些主张应当敬拜神灵，且是像不久前提及的那类神灵的哲学家，我该如何说呢？这难道不是规定要敬拜恶习、罪行和卑鄙之事吗？我为你们感到羞耻，也怜悯你们，如果你们尚未发现这些事不值得相信，且是不敬虔、可憎恶的；或者，你们虽已发现并确认它们是邪恶的，却仍视它们为善，甚至为至善而加以敬拜。\par

% structure: p0102 source=h2 target=subsection reason=relative-heading-rank; base=h2

\subsection{第五十一章 基督，真先知}

\Quote{然后，此外，有些哲学家竟敢谈论关于天主的事，尽管他们是有死的，只能凭意见谈论不可见的事，或世界的起源（因为他们不在场），或世界的终结，或灵魂在地狱的待遇和审判，忘记了认识现在和可见的事固然属于有理性的人，但唯有先知性的预知才能知道过去、未来和不可见的事。因此，这些事不应从猜测和意见中收集（人在其中大受欺骗），而应从对先知真理的信德中获取，正如我们的道理一样。因为我们不说自己的话，也不宣布由人的判断所收集的事；因为那样做就是欺骗听众。但我们宣讲那由真先知交托并启示给我们的事。关于祂的先知预知和能力，如果有人（如我所说）愿意接受明确的证明，就让他立刻来，留心听，我们将提供明显的证据，使他不仅用耳朵听到先知预知的能力，甚至用眼睛看见，用手触摸；当他对此建立了确信的信德后，他将毫不费力地担负起公义与虔敬的轭；见：\BibleRef{玛 11:30} 并且他将从中体验到如此大的甘饴，以致他不仅不会抱怨其中有任何劳苦，反而渴望再加添一些东西在他身上。}\par

% structure: p0104 source=h2 target=subsection reason=relative-heading-rank; base=h2

\subsection{第五十二章 阿皮翁与阿努比翁}

然后我的父亲听见这事，就欢喜，对彼得说：\Quote{如果你允许，我想去问候阿皮翁和阿努比翁，因为他们是我的好友；也许我能说服阿努比翁与克肋孟就\OriginalTerm{命数}{Genesis}的问题进行辩论。} 于是彼得说：\Quote{我同意；我赞扬你，因为你尊重你的朋友。但请想想，所有事情如何按你的愿望藉着天主的眷顾发生；因为，看，不仅\Emph{正当情感的对象}已由天主的安排归还给你，连你朋友的到来也为你安排了。} 然后我的父亲说：\Quote{我确实认为正如你所说。} 当他这样说后，他就往阿努比翁那里去了。\par

% structure: p0106 source=h2 target=subsection reason=relative-heading-rank; base=h2

\subsection{第五十三章 变形}

但我们整夜与彼得坐在一起，提问并向他学习许多课题，因他的教导和言语的甘饴而十分喜悦，一直保持清醒；到天亮时，彼得看着我及我的兄弟们，说：\Quote{我想知道你们父亲遭遇了什么。} 他正说话时，我的父亲进来了，发现彼得正在跟我们谈论他。他致意后，便开始道歉，并解释在外过夜的原因。但我们看着他，感到惊骇；因为我们看到他脸上是西蒙的面容，却听见我们父亲的声音。我们躲避他并诅咒他时，父亲对我们如此粗暴和野蛮的对待感到惊讶。然而彼得是唯一看到他本来面貌的人；他对我们说：\Quote{你们为什么诅咒你们的父亲？} 我们和我们的母亲一同回答他说：\Quote{他看起来是西蒙，虽然有我们父亲的声音。} 于是彼得说：\Quote{你们确实只认识他的声音，这声音没有被巫术改变；但对我来说，他的面容（在别人看来被西蒙的技艺改变了）我也认出是你们父亲福斯提尼亚努斯的。} 他看着我父亲，说：\Quote{你妻子和儿子们惊慌的原因是这样的——你的面容看起来不再是从前的样子，而是可憎恶的西蒙的脸显现在你身上。}\par

% structure: p0108 source=h2 target=subsection reason=relative-heading-rank; base=h2

\subsection{第五十四章 安提约基雅的骚动}

当他这样说时，先前前往安提约基雅的人中有一位回来了，对伯多禄说：\Quote{我主伯多禄，我愿您知道，西满在安提约基雅公开行了许多征兆和奇事，向百姓灌输的尽是激起对您的仇恨，称您为魔术师、巫士、凶手；他如此激起了对您的仇恨，以至于他们若在任何地方找到您，甚至渴望吞噬您的肉。因此，我们这些先前被派去的人，见全城对您大为骚动，便私下聚在一起，商议应当怎样行。}\par

% structure: p0110 source=h2 target=subsection reason=relative-heading-rank; base=h2

\subsection{第五十五章　计谋}

我们见无法摆脱困境，这时有百夫长科尔乃略，奉凯撒之命因公事前往凯撒勒雅总督处。我们单独请了他来，告诉他我们忧愁的缘故，恳求他若能为力，就帮助我们。他立刻十分情愿地答应，只要我们协助他的计划，他就能使西满立即逃走。我们承诺必积极行事，他便说：\InnerQuote{凯撒已下令在罗马城及各省份搜捕并消灭巫士，其中许多人已被处死。因此我将通过我的朋友放出风声，说我是来捉拿那个巫士的，奉凯撒之命为此而来，好让他与其余同伙一同受罚。所以你们那些在他身边伪装的人，要暗中告知他，仿佛从某处听说我是来捉拿他的；他一听见，必定会逃走。或者你们若想到更好的办法，请告诉我。}何须多说？我们的那些人当时在他身边，本是为监视他而伪装的，就这样做了。西满得知此事临到，便将这消息视为他们的大恩，于是逃走。他因此离开安提约基雅，据我们所知，与雅典诺多鲁斯一同来到这里。\par

% structure: p0112 source=h2 target=subsection reason=relative-heading-rank; base=h2

\subsection{第五十六章　西满改变容貌的计谋}

\Quote{我们这些先于您前往的人，都认为您暂时不应上安提约基雅，直到我们看见他因离开而在百姓中所播下的对您的仇恨是否有所缓和。}那位从安提约基雅来的人通报完毕后，伯多禄看着我们的父亲说：\Quote{福斯蒂尼亚努斯，你的面容已被西满魔术师改变了，这是明显的；因为他以为凯撒正在追捕他，便惊恐逃走，将自己的面容放在你身上，若你被当作他而捉拿处死，好使他能给你的儿子们带来悲伤。}我的父亲听见这话，喊着流泪说：\Quote{你判断得对，伯多禄！因为与我非常友好的阿努比翁，曾以一种神秘的方式开始告知我他的阴谋；但不幸的是，我没有相信他，因为我从未伤害过他。}\par

% structure: p0114 source=h2 target=subsection reason=relative-heading-rank; base=h2

\subsection{第五十七章　极大的悲痛}

当我们所有人，连同我的父亲，都因悲伤和哭泣而激动时，其间阿努比翁来到我们这里，告知我们西满已在夜间逃走，往犹太去了。但看见我们的父亲哀叹自怜，说：\Quote{我真不幸，听到他是巫士时竟不相信！我这可怜的人遭遇了什么？在一天之内，被妻子和儿子认出，却不能与他们一同欢乐，反而被抛回从前漂泊时所受的苦难！}——我的母亲则披散着头发，哭得更加悲切——我们也被父亲容貌的改变所困惑，仿佛被雷击而神魂超拔，不明白究竟发生了什么。但阿努比翁见我们如此痛苦，竟哑口无言。于是伯多禄看着我们这些儿子们说：\Quote{你们要相信我，这确实是你们的父亲；因此我命你们尊敬他如父。因为天主必会赐下一个机会，使他能脱去西满的容貌，恢复你们父亲——即他自己的——明显形像。}\par

% structure: p0116 source=h2 target=subsection reason=relative-heading-rank; base=h2

\subsection{第五十八章　事情如何发生}

然后，他转向我的父亲说：\Quote{我准许你向阿皮翁和阿努比翁致意，你说他们是你自幼的朋友，但并未准许你与西满说话。}我的父亲说：\Quote{我承认我犯了罪。}阿努比翁说：我也与他一同恳求您宽恕这位老人——他本是良善高尚的人。他不幸被那巫士诱惑欺骗；因为我来告诉您事情的经过。当他来向我们致意时，恰好那时我们正围站在西满身边，听他告诉我们他打算当晚逃走，因他听说有人甚至已来到这拉奥迪雅城要奉皇帝之命捉拿他，但他想将他们所有的怒气转嫁到这新近来到的福斯蒂尼亚努斯身上。他对我们说：\InnerQuote{只要你们让他与我们共进晚餐，我便调制一种油膏，他饭后以此涂面，此后必在众人眼中看来有我相同的容貌。但你们先用一种草汁涂抹面部，以免被他容貌的变化所欺，这样在你们以外的所有人眼中，他都像是西满。}\par

% structure: p0118 source=h2 target=subsection reason=relative-heading-rank; base=h2

\subsection{第五十九章　哀悼场景}

\Quote{当他说这话时，我对他说：\InnerQuote{你做这事能得到什么好处？}然后\Person{西满}说：\InnerQuote{首先，那些寻找我的人可以抓住他，从而放弃搜寻我。但如果他被凯撒惩罚，他的儿子们——他们离弃了我，逃到\Person{伯多禄}那里，如今成了他的助手——将会非常悲伤。}现在，\Person{伯多禄}，我向你承认实话。那时我不敢告诉\Person{福斯提尼亚诺}；但\Person{西满}也没有给我们私下和他谈话的机会，向他完全揭露\Person{西满}的计谋。与此同时，大约半夜时分，\Person{西满}已经逃往\Place{犹太}。\Person{阿特诺多罗}和\Person{阿皮翁}去护送他；但我假装身体不适，以便留在家里，让他尽快回到你那里，希望他能在你那里隐藏起来，免得被搜寻\Person{西满}的人抓住，带到凯撒面前，无辜丧命。现在，我因担心他，来看望他，并在那些去护送\Person{西满}的人回来之前返回。} \Person{阿努比翁}转向我们说：\Quote{我，\Person{阿努比翁}，确实看到了你父亲的真实面容，因为我之前曾被\Person{西满}本人傅油，正如我告诉你们的，好让\Person{福斯提尼亚诺}的真实面容显现在我眼前；因此我惊讶并惊叹于\Person{西满}术士的技艺，因为你们站在这里却不认识自己的父亲。} 当我的父亲和母亲，以及我们所有人，为所遭遇的事哭泣时，\Person{阿努比翁}动了怜悯之心，也哭了。\par

% structure: p0120 source=h2 target=subsection reason=relative-heading-rank; base=h2

\subsection{第六十章 反间计}

于是\Person{伯多禄}动了怜悯之心，答应恢复我们父亲的面容，对他说：\Quote{听着，\Person{福斯提尼亚诺}：一旦你被改变的面容的错误为我们带来一些好处，并有助于我们目前的计划，那时我就会恢复你面容的真实形貌；但有一个条件：你必须首先执行我命令你的事。} 我父亲答应他会尽全力完成一切可能托付给他的事，只要他能恢复自己的面容，\Person{伯多禄}就这样开始说：你亲耳听到，先前被派去的人中有一个从\Place{安提约基雅}回来，告诉我们\Person{西满}在那里如何煽动群众反对我，如何煽动全城仇恨我，宣称我是术士、谋杀犯和骗子，以至于他们急切地想要见到我，甚至吃我的肉。因此，照我说的去做：把\Person{克肋孟}留给我，你和你的妻子以及你的儿子\Person{福斯都}和\Person{福斯提诺}先到\Place{安提约基雅}去。我还会派其他我认为合适的人与你们同行，他们会遵守我所命令的一切。\par

% structure: p0122 source=h2 target=subsection reason=relative-heading-rank; base=h2

\subsection{第六十一章 布设陷阱}

\Quote{因此，当你们和他们一起到了\Place{安提约基雅}时，由于你们会被认为是\Person{西满}，要站在公共场所，宣布你的悔改，说：\InnerQuote{我\Person{西满}向你们宣告并承认，我所有关于\Person{伯多禄}的话都是假的：因为他既不是诱惑者，也不是术士，也不是谋杀犯，也不是我攻击他时所说的任何东西；但我是在疯狂驱使下说了所有这些话。因此，我本人——正是此前给你们造成仇恨他的理由的人——恳求你们，不要对他存有任何这样的想法。但要放下你们的仇恨，停止你们的愤怒，因为他确实是天主为世界的救恩而派遣的——是真先知的门徒和宗徒。因此，我劝告、敦促并嘱咐你们，当你们向他宣讲真理时，要听从他并相信他，免得你们轻视他，你们的城市突然毁灭。但我要告诉你们，为什么我现在向你们做这个承认。今夜，一位天主的使者因我的邪恶而责备我，并严厉地鞭打我，因为我是真理宣扬者的敌人。因此，我恳求你们，即使我自己以后再来到你们这里，试图说任何反对\Person{伯多禄}的话，你们也不要接待或相信我。因为，我向你们承认，我曾是术士、诱惑者、骗子；但我悔改，因为通过悔改可以涂抹先前的恶行。}}\par

% structure: p0124 source=h2 target=subsection reason=relative-heading-rank; base=h2

\subsection{第六十二章 良心问题}

当\Person{伯多禄}向我的父亲提出这个建议时，他回答说：\Quote{我知道你的意思；不必再麻烦了：因为我明白并且知道当我到了那里时，我该承担什么。} \Person{伯多禄}进一步指示他说：\Quote{因此，当你到了那里，看到百姓因你的讲话而转变，放下仇恨，重新渴望见到我时，就派人告诉我，我会立刻到来；当我到来时，我会立即将你从这奇怪的面容中释放出来，恢复你那为所有朋友所熟知的本来面目。} 说了这话，他命令我的兄弟们和他一起去，同时还有我们的母亲\Person{玛提迪亚}和一些朋友。但我的母亲拒绝和他一起去，说：\Quote{如果我与\Person{西满}的面容来往，似乎我会成为通奸者；但如果我被迫和他一起走，至少不可能和他同床共枕；但我不知道我是否能同意甚至和他一起走。} 当她坚决拒绝时，\Person{阿努比翁}开始劝她说：\Quote{请相信我和\Person{伯多禄}。难道他的声音不能说服你，他是你的丈夫\Person{福斯提尼亚诺}吗？我确实爱你一样地爱他。简而言之，我自己也要和你们一起去。} \Person{阿努比翁}说了这话后，我母亲答应会和他一起去。\par

% structure: p0126 source=h2 target=subsection reason=relative-heading-rank; base=h2

\subsection{第六十三章 虔诚的诡计}

然后我说：\Quote{天主按我们的意愿安排我们的事务；因为我们同有一位名叫阿努比翁的占星家在一起，若我们到了安提约基雅，我们将与他以最诚挚的态度就星命之学进行辩论。}当我们的父亲在午夜后与伯多禄所派陪同他的人以及阿努比翁一同出发后，早上在伯多禄前往讨论之前，那些护送西满的人，即阿丕翁和雅典诺多，回来了，并到我们这里询问我父亲的下落。但伯多禄得知他们到来后，吩咐他们进来。他们就座后，问道：\Quote{法乌斯蒂尼安在哪里？}伯多禄回答说：\Quote{我们不知道；因为自从他晚间到你们那里去以后，他的任何朋友都没有见过他。但昨天早上西满来询问他的下落；因我们没有回答他，我不知道他是什么意思，但他自称是法乌斯蒂尼安。但没有人相信他，他就去哀哭，并威胁说要自杀；之后他向海边走去。}\par

% structure: p0128 source=h2 target=subsection reason=relative-heading-rank; base=h2

\subsection{第六十四章　谎言的比拼}

阿丕翁和他同来的人听到这话，大声哀号说：\Quote{你们为什么这样做？你们为什么不接待他？}当雅典诺多正要告诉我那就是我的父亲法乌斯蒂尼安本人时，阿丕翁阻止了他，并说：\Quote{我们从某人那里得知他同西满一起走了，而且是出于法乌斯蒂尼安自己的请求，因他不愿见他的儿子们，因为他们都是犹太人。因此我们听到这事，就来此询问他的下落；既然他不在这里，那告诉我们他已同西满一起走了的人必定说了真话。这就是我们要告诉你们的事。}但我克莱孟明白伯多禄的计谋：他想让他们以为老人会从他们那里被追索，从而害怕和逃走，于是我开始协助他的计谋，对阿丕翁说：\Quote{亲爱的阿丕翁，请听！我们认为好的东西，也愿意传授给我们的父亲；但如果他不愿接受，反而如你所说，因厌恶我们而逃走——这话或许说得有些严厉——我们就不管他了。}我这样说了之后，他们就离开了，诅咒我的残忍，并跟随西满的踪迹，正如我们次日所获悉的。\par

% structure: p0130 source=h2 target=subsection reason=relative-heading-rank; base=h2

\subsection{第六十五章　计谋成功}

在此期间，伯多禄每日照常教导百姓，行许多奇迹和医治。十天后，我们的一位从安提约基雅来的人，由我父亲所派，向我们报告我父亲如何在公开场合控告西满——他的面容似乎就是西满的——并以无与伦比的赞美颂扬伯多禄，把他推荐给所有百姓，使众人向往他，以致所有人都因他的言论而改变，渴望见他；而且许多人变得非常爱慕伯多禄，以致他们对我父亲（在他扮演西满的角色时）发怒，并想要对他动手，因为他竟对伯多禄行了如此不义！\Quote{因此，}他说，\Quote{要赶紧，免得他可能被杀害；因为他极其惧怕，派我急速到你们这里，请求你们立即前来，好使你们能找到他活着，并且当这城市对你们的情感正在增长之际，你们能出现在这有利的时机。}他还告诉我们，当我父亲一进城后，全体百姓都聚集到他那里，以为他是西满；他便开始向众人公开认错，按照人民复原所需的要求——因为凡来的人，无论贵族平民、富户穷人，都希望他照常行些神迹奇事——他这样对他们讲道：——\par

% structure: p0132 source=h2 target=subsection reason=relative-heading-rank; base=h2

\subsection{第六十六章　说谎的唇讲真话}

\Quote{天主的忍耐长久包容我这最不幸的西满；因为你们在我身上所惊讶的一切，并非藉真理而成，而是靠魔鬼的谎言和诡计，为要颠覆你们的信仰并定自己的罪。我承认我关于伯多禄所说的一切都是谎言；他从未是术士或凶手，而是天主为拯救你们众人所派遣的；若从此刻起你们以为他可被轻视，要确信你们的城市可能会突然被毁灭。但\Emph{你们会问}，我为何自愿向你们作此承认？今夜我受到天主的一位天使的严厉责斥和鞭打，因为我曾是祂的仇敌。因此我恳求你们，即使从此刻起连我自己若再开口攻击他，你们也要把我从你们眼前赶走；因为那讨厌的魔鬼，那与人类得救为敌的，藉我的口攻击他，为使你们不能藉他获得生命。因为魔术能藉我向你们显示什么奇迹？我使铜狗吠叫，使雕像走动，使人改变容貌并突然从人眼前消失；为这些事你们本应诅咒那魔术，它用魔鬼的锁链捆绑你们的灵魂，为使我向你们显示虚妄的奇迹，使你们不信伯多禄，他奉那派遣他者的名医治病人、驱逐魔鬼、使瞎子看见、使瘫痪者恢复健康、并使死人复活。}\par

% structure: p0134 source=h2 target=subsection reason=relative-heading-rank; base=h2

\subsection{第六十七章　法乌斯蒂尼安恢复原貌}

当他作了这些和类似的陈述时，百姓开始咒骂他，并哭泣哀叹，因为他们曾得罪了伯多禄，相信他是个术士或恶人。但当天晚上，法乌斯蒂尼安恢复了自己的面容，西满魔士的形貌离开了他。西满听说他的面容在法乌斯蒂尼安身上促成了伯多禄的荣耀，便急忙赶来要抢在伯多禄前面，打算用他的法术让他的形貌从法乌斯蒂尼安身上离开——然而基督已按照祂宗徒的话成就了这事。但尼凯塔和阿奎拉看见他们父亲的面容在必要的宣告之后恢复原状，便感谢天主，并不再容他向百姓讲话。\par

% structure: p0136 source=h2 target=subsection reason=relative-heading-rank; base=h2

\subsection{第六十八章　伯多禄进入安提约基雅}

但西满开始秘密地前往他的朋友和熟人中间，比以前更厉害地诋毁伯多禄。于是众人往他脸上吐唾沫，将他赶出城，说：\Quote{你若再敢来此地反对伯多禄，你自己的死将由你自己负责。} 这些事在\Emph{\Place{劳狄刻雅}}传开，伯多禄便吩咐民众次日集会；他祝圣了一位跟随他的人为他们作主教，另一些人为长老，并为多人授洗，治愈了所有患病或被邪魔折磨的人；他在那里又住了三天；一切安排妥当后，他辞别他们，从劳狄刻雅出发，安提约基雅的民众非常想念他。通过尼塞塔和阿奎拉，全城开始听说伯多禄即将到来。于是安提约基雅全城的民众听说伯多禄抵达，都去迎接他，几乎所有老人和贵族都头上撒灰前来，以此表示忏悔，因为他们曾听从术士西满，反对伯多禄的宣讲。\par

% structure: p0138 source=h2 target=subsection reason=relative-heading-rank; base=h2

\subsection{第六十九章 伯多禄的感恩}

他们叙述这些和类似的事，将那些患病、被邪魔折磨、瘫痪以及遭受各种危难的人带到他们面前；聚集的病人不计其数。伯多禄见他们不仅因西满的影响而悔改了对他的恶念，而且对天主表现出如此全然的信德，相信他能治愈所有患各种疾病的人，便向天伸开双手，流着泪倾心祈祷，感谢天主说：\Quote{当受一切赞美的圣父，我颂扬祢，因为祢曾屈尊成就了祢圣子的一切话语和应许，使一切受造物知道，惟有祢是天主，在天上及地上。}\par

% structure: p0140 source=h2 target=subsection reason=relative-heading-rank; base=h2

\subsection{第七十章 奇迹}

说了这些话后，他登上一处高地，吩咐所有病人都排列在他面前，对他们所有人说：\Quote{既然你们看见我如同你们一样是人，不要以为你们能从我这里恢复健康，而是通过那从天降下、向信祂的人显示身体和灵魂完美良药的那一位。因此，让全体民众为你们的宣告作证：你们全心相信主耶稣基督，好叫他们知道，他们自己也能因祂而得救。} 所有病人异口同声呼喊，说伯多禄所宣讲的祂是真天主；忽然，天主恩宠的不可抗拒之光显现在民众中间；瘫痪者得治愈，开始跑到伯多禄脚前；盲人因复明而欢呼；跛子因重获行走能力而感谢；病人因恢复健康而喜乐；甚至有些几乎已死、失去知觉、不能说话的人，也被扶起；所有癫狂者和附魔者都得到释放。\par

% structure: p0142 source=h2 target=subsection reason=relative-heading-rank; base=h2

\subsection{第七十一章 成功}

圣神在那一天显示了祂大能的如此大恩宠，以致全体从最小到最大，异口同声承认主；长话短说，七天之内，超过一万人相信天主，受洗并藉圣化被祝圣；因此，德敖斐罗，那位在该城所有权贵中最为显赫的人，以极大的渴望将他宫殿的大厅祝圣为教堂，全体民众在那里为宗徒伯多禄安置了一个座位；众人每日聚集听道，相信那由治愈效力所证实的健康教义。\par

% structure: p0144 source=h2 target=subsection reason=relative-heading-rank; base=h2

\subsection{第七十二章 圆满结局}

于是，我克莱孟和我的兄弟们以及我们的母亲对父亲说话，问他心中是否还有不信的残余。他说：\Quote{来，你们将在伯多禄面前看到，我的信德增长了多大。} 于是福斯提尼亚努斯上前，俯伏在伯多禄脚前说：\Quote{祢言语的种子，我心田已领受，现已发芽，并如此进展到丰硕的成熟，只缺祢用那属灵的镰刀将我分离出糠秕，安置在主的仓库里，使我分沾神圣的筵席。} 伯多禄极欢欣地握住他的手，将他交给我克莱孟和我的兄弟们，说：\Quote{正如天主将你们的儿子归还给你们父亲，同样，你们的儿子将他们的父亲归还给天主。} 他向全体民众宣布禁食，在下一个主日为他授洗；在民众当中，他藉他的皈依为契机，叙述了他的一切经历，以致全城待他如同天使，待他不亚于对待宗徒。这些事传开，伯多禄便吩咐民众次日集会；他祝圣了一位跟随他的人为主教，另一些人为长老，又为大批人授洗，并治愈了所有患病受苦的人。\par

\SourceRecord{Translated by Thomas Smith.；Ante-Nicene Fathers；Vol. 8.；Edited by Alexander Roberts, James Donaldson, and A. Cleveland Coxe.；Buffalo, NY: Christian Literature Publishing Co.,；1886.；<http://www.newadvent.org/fathers/080410.htm>.}
